% Master LaTeX export — VIP Cortical Interneurons review (Phase 20 final assembly)
% Generated by concatenating per-section TeX exports produced via `myst build --tex`.
% Markdown is the canonical source; this LaTeX export is provided for archival convenience.
\documentclass[11pt]{article}
\usepackage[utf8]{inputenc}
\usepackage{hyperref}
\usepackage{datetime}
\usepackage{graphicx}
\usepackage{natbib}
\usepackage{url}
\usepackage{booktabs}
\usepackage{longtable}
\usepackage{amsmath}
\usepackage{amssymb}
\usepackage[margin=1in]{geometry}
\bibliographystyle{abbrvnat}
\hypersetup{colorlinks=true, allcolors=blue}

\title{VIP Cortical Interneurons: A Critical Reassessment of the Disinhibition Framework}
\author{ComputationalReview Pipeline}
\date{\today}

\begin{document}
\maketitle
\tableofcontents
\newpage


\section{Introduction}
\label{sec:01_introduction}

Vasoactive intestinal peptide-expressing (VIP) GABAergic interneurons occupy a distinctive position in the cortical inhibitory circuit. They are a numerically small subclass --- under one in ten cortical interneurons across most rodent areas --- yet they are recruited reliably by locomotion, arousal, reward, punishment, and top-down attention, and they have been written into circuit diagrams as the principal source of state-dependent disinhibition onto pyramidal cells via somatostatin (SST)--expressing interneurons \citep{Pi2013, Lee2013, Pfeffer2013, Fu2014, Karnani2016a, Tremblay2016}. Across the past decade, this disinhibitory reading of VIP function has been generalised from auditory and frontal cortex to primary visual cortex, somatosensory cortex, hippocampus, and prefrontal--cingulate territories, and it now serves as the default circuit motif in network and learning models of cortex \citep{Fu2014, Lee2013, Kuchibhotla2017, Williams2019, Hertag2019, Litwinkumar2016, Garciadelmolino2017}. The same generalisation, however, has surfaced empirical and theoretical conflicts that the field has not resolved: in several preparations VIP recruitment does not produce the expected SST suppression, and in several models the same connectivity supports either gain enhancement or paradoxical pyramidal suppression depending on weakly constrained parameters \citep{Pakan2016, Dipoppa2018, Yavorska2021, Garciadelmolino2017, Sanzeni2020}.

This review takes that tension as its starting point. Rather than re-affirming the disinhibitory motif as a uniform property of the VIP subclass, the review asks what the evidence --- molecular, developmental, morphological, electrophysiological, synaptic, in-vivo, oscillatory, computational, and translational --- actually constrains, and where the boundaries of safe inference lie. The thesis the review defends is that VIP-IN function is reliably state-modulated and reliably disinhibitory at the subclass average across cortex, but that its sign, magnitude, gain regime, target compartment, and connectional asymmetry vary systematically with cortical area, behavioural context, and intra-subclass cell type, and that this variation is not noise but the substrate of VIP function \citep{Pi2013, Pfeffer2013, Fu2014, Pakan2016, Dipoppa2018, Garrett2020, Bigelow2019, Anastasiades2021a, Bastos2023a}. Reading the corpus this way reframes the central question from ``what does VIP do?'' to ``under what area$\times$state$\times$modulator conditions does each of the operating modes catalogued below engage?''.

The body of the review is structured as eleven sections that move from the molecular identity of the subclass outward to its translational implications. \href{/classification}{Molecular Identity and Transcriptomic Taxonomy} reviews the transcriptomic and marker-gene evidence anchoring VIP within the caudal-ganglionic-eminence (CGE)--derived, 5-HT3AR+ branch of cortical inhibition and surveys current single-cell atlases that resolve the subclass into multiple t-types \citep{Rudy2011, Tasic2016, Tasic2018, Hodge2019, Yao2023, Lee2023a, Gouwens2020a}. \href{/development}{Developmental Origins and Postnatal Maturation} reconstructs CGE specification, transcription-factor cascades, tangential migration, and postnatal maturation of the lineage \citep{Fishell2011, Lodato2011, Vogt2014, Mayer2018, Batistabrito2017, Vucurovic2010}. \href{/morphology}{Morphological Diversity} catalogues the bipolar/bitufted/multipolar morphological palette and the laminar distribution that constrains target choice \citep{Cauli1997, Cauli2014, Pronneke2015, Apicella2022, Emmenegger2018, Gouwens2019}. \href{/electrophysiology}{Intrinsic Electrophysiology} assembles passive properties, Petilla firing-pattern membership, and the irregular- versus continuous-adapting taxonomic conflict, including the Patch-seq integration that pegs intrinsic phenotype to t-type \citep{Tremblay2016, Gouwens2020a}.

\href{/synaptic-properties}{Synaptic Properties and Connectivity} inventories afferent and efferent connectivity --- long-range cortico-cortical and thalamic input, neuromodulatory drive, and the asymmetric VIP\rightarrowSST \textgreater  VIP\rightarrowPV/Pyr output that is the structural basis for the disinhibitory reading --- and flags the layer- and area-specific exceptions to that asymmetry \citep{Pi2013, Pfeffer2013, Lee2013, Karnani2016a, Karnani2016b, Wall2016, Schneidermizell2025}. \href{/disinhibition-framework}{Local Circuit Motifs and the Disinhibition Framework} assembles those parts into the circuit motif itself, asks how reproducible the connectional asymmetry is across area and layer, and audits the optogenetic and chemogenetic studies that have driven the motif into the textbook \citep{Pi2013, Lee2013, Fu2014, Karnani2016a, Cichon2017, Williams2019, Veit2023}. \href{/in-vivo-behavior}{In Vivo Function During Behavior} reviews calcium-imaging and electrophysiological recordings of VIP cells during locomotion, reward, punishment, and stimulus encoding, and assembles the cases in which the disinhibitory mechanism fails or runs in reverse \citep{Pakan2016, Dipoppa2018, Yavorska2021, Bigelow2019}. \href{/cross-areas}{VIP Interneurons Across Brain Regions} lays out the cross-area heterogeneity directly: visual, auditory, somatosensory, frontal, and hippocampal VIP populations differ in afferent profile, output specificity, and behavioural recruitment \citep{Pi2013, Dipoppa2018, Bastos2023a, Anastasiades2021a, Tyan2014, Gulyas1996, Kawaguchi1996}.

\href{/oscillatory-dynamics}{Oscillatory Dynamics and Temporal Coordination} examines temporal coordination --- gamma, theta, Up-states, infraslow rhythms, and the predictive-coding/attention frame that ties VIP activity to inter-areal coordination \citep{Veit2023, Hertag2020, Hertag2021, Shipp2016, Francavilla2018a}. \href{/disease-translational}{Species Differences, Human Relevance, and Disease} works out cross-species conservation and divergence, primate expansion of CGE-derived diversity, and VIP-IN involvement in autism, Rett, Dravet, schizophrenia, and Alzheimer-related pathology \citep{Hodge2019, Boldog2018, Chartrand2023, Bakken2021a, Bakken2021b, Krienen2020, Goff2023, Mcfarlan2024a}. \href{/computational-models}{Computational Models of VIP Circuit Function} audits the rate-, spiking-, and learning-model literature, exposes the gain-regime conflicts the architectures generate, and argues that the present bottleneck is parameter identifiability rather than missing biology \citep{Litwinkumar2016, Garciadelmolino2017, Veit2023, Bos2025, Hartung2024, Tahvili2025a, Tahvili2025b, Reimann2026, Sanzeni2020, Lee2025, Sabri2024}.

Five cross-cutting tensions surface across these sections, and the review treats them as a single connected web rather than as section-local issues. The first is the \textbf{discrete-versus-continuous} structure of the VIP subclass: single-cell atlases recover discrete t-types in supervised clustering, while morphology, intrinsic phenotype, and Patch-seq cross-modal integration reveal graded, dataset-dependent boundaries within the subclass --- a tension that is most explicit in \href{/classification}{Molecular Identity and Transcriptomic Taxonomy}, \href{/morphology}{Morphological Diversity}, and \href{/electrophysiology}{Intrinsic Electrophysiology} and that propagates into how downstream sections attribute function to ``VIP'' cells \citep{Tasic2016, Tasic2018, Pronneke2015, Apicella2022, Gouwens2020a, Schneidermizell2025, Emmenegger2018}. The second is the \textbf{gain-regime ambiguity}: optogenetic VIP activation has been read as subtractive, divisive, or mixed gain depending on area, layer, and analysis convention, and the rate models that fit the data span the same range \citep{Kuchibhotla2017, Hertag2019, Garciadelmolino2017, Litwinkumar2016, Veit2023, Bos2025, Tahvili2025a, Tahvili2025b}. The third is the \textbf{net-disinhibition versus paradoxical-effect} conflict: the same VIP\rightarrowSST\rightarrowPyr connectivity, parameterised inside or outside the inhibition-stabilised regime, supports either pyramidal facilitation or pyramidal suppression on VIP activation, so that ``disinhibition'' is a property of operating point and not of wiring alone \citep{Garciadelmolino2017, Sanzeni2020, Reimann2026, Hertag2019, Tahvili2025a}. The fourth is the \textbf{identifiability of the VIP\rightarrowSST weight}: in the four-population rate equations, VIP\rightarrowSST and VIP\rightarrowPV cross-couplings trade off against SST\rightarrowPyr and PV\rightarrowPyr terms in a way that connectomic priors alone do not constrain, and several published models with incompatible VIP\rightarrowSST weights fit the same in-vivo data \citep{Hertag2019, Litwinkumar2016, Garciadelmolino2017, Veit2023, Reimann2026, Sabri2024}. The fifth is \textbf{cross-species translation}: VIP-IN identity is conserved at the subclass level between rodent and primate cortex, but t-type proportions, laminar distribution, and human-specific morphologies (including L1 \textit{VIP PCDH20} and \textit{MC4R} types) diverge non-trivially, so that mouse circuit conclusions transfer to human cortex with quantitative qualifications and disease-relevant uncertainty \citep{Hodge2019, Boldog2018, Chartrand2023, Bakken2021a, Bakken2021b, Krienen2020, Lee2023a}.

The review proceeds inductively from molecular ground truth to behavioural and translational implication, but the reader can also read it diagonally along any of the five tensions above. Forward and backward cross-references between \href{/classification}{Molecular Identity and Transcriptomic Taxonomy} through \href{/computational-models}{Computational Models of VIP Circuit Function} are dense and intentional: a claim about VIP firing in \href{/in-vivo-behavior}{In Vivo Function During Behavior} is only as strong as the molecular and morphological identity assigned to the recorded cells in \href{/classification}{Molecular Identity and Transcriptomic Taxonomy} and \href{/morphology}{Morphological Diversity} \citep{Tasic2018, Gouwens2020a, Pronneke2015, Schneidermizell2025}, and a claim about disinhibitory gain in \href{/disinhibition-framework}{Local Circuit Motifs and the Disinhibition Framework} is only as strong as the synaptic-weight and sign of the model in \href{/computational-models}{Computational Models of VIP Circuit Function} \citep{Hertag2019, Garciadelmolino2017, Veit2023, Reimann2026}. The synthesis the review aims to deliver, summarised in \href{/conclusion}{Conclusion}, is therefore not a unified theory of VIP function but a calibrated map of where the disinhibitory reading holds, where it does not, and what experiments and analyses would discriminate the remaining cases \citep{Pi2013, Pakan2016, Dipoppa2018, Garciadelmolino2017, Reimann2026}.


\section{Molecular Identity and Transcriptomic Taxonomy}
\label{sec:02_classification}

The introduction defined VIP-expressing interneurons through their disinhibitory role and motivated a critical reassessment of that framing. Before any reassessment is possible, we must specify what cells the label ``VIP'' actually refers to in contemporary cortical neuroscience. This section catalogues the molecular ground truth: the transcriptomic, marker-gene, and developmental-class evidence that anchors the VIP subclass within cortical inhibition. The section's central thesis is that VIP interneurons are a CGE-derived, 5-HT3AR+ GABAergic subclass whose subclass-level identity is now unambiguous across modern single-cell atlases, but whose internal granularity --- how many transcriptomic types (t-types), how discrete the boundaries between them, and how those types map onto immunolabel- and Cre-driver-based definitions --- remains a moving target that depends on dataset, species, taxonomic resolution, and integration method. Reconciling the disinhibitor framing with the multimodal evidence reviewed in subsequent sections requires holding this molecular indeterminacy in mind: many cross-study disagreements about VIP function trace upstream to whether different laboratories are sampling the same cells.

\section{The 5-HT3AR/CGE class places VIP in cortical inhibition}

The modern partition of cortical GABAergic interneurons rests on a tripartite scheme in which essentially every neocortical inhibitory neuron expresses one of three principal markers --- parvalbumin (PV), somatostatin (SST), or the ionotropic serotonin receptor 5-HT3AR --- accounting between them for nearly the entire interneuron population \citep{Rudy2011, Tremblay2016}. \citet{Rudy2011} consolidated immunohistochemical, electrophysiological, and developmental evidence to argue that PV+ neurons constitute roughly forty percent of neocortical GABAergic cells, SST+ neurons roughly thirty percent, and 5-HT3AR-expressing neurons the remaining thirty percent --- the three groups together accounting for nearly the entire inhibitory population. The 5-HT3AR class subsumes the entirety of the VIP+ population, and \citet{Tremblay2016} clarified that VIP+ neurons make up approximately forty percent of the 5-HT3AR group, with the non-VIP 5-HT3AR remainder corresponding to layer-1-enriched neurogliaform cells that are now identified as the \textit{Lamp5} subclass. Genetic evidence from a BAC-transgenic 5-HT3AR-GFP line confirmed the marker logic at the population level: 5-HT3AR is expressed in essentially all neocortical interneurons that lack PV or SST, and fate mapping in the same line showed that the 5-HT3AR+ compartment encompasses the full repertoire of caudal-ganglionic-eminence (CGE)-derived interneurons \citep{Lee2010}. Together with the GAD67-GFP knock-in framework of \citet{Tamamaki2003}, which placed CR-, PV-, and SST-immunoreactive subsets within a known total GABAergic abundance of approximately twenty percent of cortical NeuN+ neurons, these studies established the quantitative envelope that subsequent transcriptomic atlases would refine but not overturn.

The CGE-derived 5-HT3AR class is internally heterogeneous. Triple immunostaining of mouse visual cortex resolved at least thirteen combinatorial groups of cortical GABAergic neurons defined by overlapping expression of PV, CR, SST, CCK, NPY, VIP, and ChAT \citep{Gonchar2008}, and double-label quantification across cortex confirmed that PV, SST, and VIP define three reliably non-overlapping subpopulations while CR and NPY overlap substantially with SST \citep{Xu2010a}. In barrel cortex, VIP-expressing interneurons amount to about thirteen percent of all GABAergic neurons, with the majority residing in supragranular layers and showing layer-dependent differences in morphology and intrinsic properties \citep{Pronneke2015}. These pre-transcriptomic enumerations matter because they define the empirical denominator against which scRNA-seq taxonomies must be checked: the VIP subclass is, on every counting method, a numerically minor but layer-biased subset of cortical inhibition.

\section{scRNA-seq atlases established the modern subclass framework}

Single-cell transcriptomics consolidated the VIP/SST/PV/Lamp5 subclass partition into a hierarchical scheme that is now the field's reference frame. \citet{Tasic2016} profiled 1,679 high-quality single cells from adult mouse primary visual cortex with SMART-seq and recovered forty-nine transcriptomic types --- twenty-three GABAergic, nineteen glutamatergic, and seven non-neuronal --- with the GABAergic types organised around the \textit{Vip}, \textit{Sst}, and \textit{Pvalb} subclasses. The same V1 study identified discrete \textit{Vip} t-types including \textit{Vip-Chat}, \textit{Vip-Gpc3}, \textit{Vip-Mybpc1}, \textit{Vip-Sncg}, and \textit{Vip-Parm1}, treating each as a transcriptomically separable cluster rather than a combinatorial gradient \citep{Tasic2016}. Two years later, \citet{Tasic2018} extended the approach to 23,822 cells across both VISp and the anterior lateral motor area (ALM) and proposed six GABAergic subclasses --- \textit{Sst}, \textit{Pvalb}, \textit{Vip}, \textit{Lamp5}, \textit{Sncg}, and \textit{Serpinf1} --- with the major split reflecting medial-versus-caudal ganglionic-eminence developmental origin. In this expanded taxonomy the \textit{Vip} subclass alone resolved into roughly fourteen t-types, with \textit{Lamp5}, \textit{Vip}, \textit{Sncg}, and \textit{Serpinf1} together forming the CGE-derived neighbourhood of the constellation diagram \citep{Tasic2018}. The convergence on the same six-subclass partition by orthogonal large-scale efforts --- the BICCN consensus of seven scRNA-seq and snRNA-seq mouse motor cortex datasets \citep{Yao2021b}, the comprehensive isocortex-plus-hippocampus atlas of \citet{Yao2021a}, and the whole-brain atlas of \citet{Yao2023} --- established that the subclass level is method- and dataset-independent. Independently, \citet{Zeisel2015} had already shown that scRNA-seq of mouse somatosensory cortex and CA1 recovered forty-seven molecularly distinct subclasses spanning the major cortical types, and \citet{Zeisel2018} extended this to the entire mouse nervous system, organising cell-type structure around three overlapping principles --- major class, developmental origin, and neurotransmitter type --- within which cortical \textit{Vip}/\textit{Sst}/\textit{Pvalb} form a major axis.

The robustness of the subclass framework is not just an aesthetic claim about clustering. \citet{Crow2018} formalised the question with the MetaNeighbor framework and showed that the major cortical interneuron classes (\textit{Vip}, \textit{Sst}, \textit{Pvalb}) replicate at high quantitative reliability across independent scRNA-seq datasets, whereas fine-grained subdivisions replicate progressively less well as taxonomic resolution increases. Multimodal datasets reach the same conclusion from a different direction: in the Patch-seq atlas of \citet{Gouwens2020a}, integrated transcriptomic, electrophysiological, and morphological data resolved 28 congruent met-types whose top-level structure recovers the \textit{Lamp5}, \textit{Vip}, \textit{Sst}, and \textit{Pvalb} subclass partition, indicating that intrinsic and morphological features carry subclass information broadly consistent with the transcriptomic taxonomy. Generative biophysical modelling of Patch-seq cells likewise showed that ion-channel conductance vectors cluster by transcriptomic subclass and that channel composition itself encodes interneuron identity at the subclass level \citep{Nandi2022}. Even the protein-level signalling layer carries the same partition: every cortical neuron type, including each VIP supertype, expresses a stereotyped combination of neuropeptide and neuropeptide-receptor genes, embedding the subclass scheme into a dense paracrine network \citep{Smith2019}. The methodological enabler underlying this convergence is Patch-seq, introduced independently by \citet{Cadwell2016} and \citet{Fuzik2016}, which combined whole-cell electrophysiology, full-transcriptome scRNA-seq, and morphological reconstruction in single neurons and made it possible to ask whether transcriptomic types map onto the morpho-electric phenotypes that defined cortical interneurons before transcriptomics existed.

\section{Within-VIP heterogeneity: counting types is resolution-dependent}

The number of ``VIP types'' reported by an atlas is a function of clustering depth, not a measurement of biological discreteness. The clearest demonstration of this is the contrast between two studies from the same Allen Institute laboratory: \citet{Tasic2016} resolved approximately five \textit{Vip} t-types in mouse V1, while \citet{Tasic2018} identified roughly fourteen \textit{Vip} t-types across VISp and ALM. The expansion was not a contradiction but a deeper cut: a more than ten-fold increase in cell number, two cortical areas instead of one, and a deeper hierarchical clustering yielded finer-grained partitions of what is plausibly the same underlying cell-type space. Subsequent atlases extended this trend without converging on a fixed count: the BICCN MOp consensus reported replicable \textit{Vip} subclasses across datasets and methods but did not reduce to a single agreed type list \citep{Yao2021b}; the isocortex-plus-hippocampus atlas of \citet{Yao2021a} partitioned \textit{Vip} into six supertypes within the CGE neighbourhood, with additional supertypes restricted to hippocampal formation; and the whole-brain mouse atlas reported 1,201 supertypes and 5,322 clusters across thirty-four classes, providing a still-finer reservoir from which area-specific \textit{Vip} type lists can be drawn \citep{Yao2023}. Across these atlases the \textit{Vip} subclass and its supertype tier are stable; the cluster tier is not.

\begin{framed}
\textbf{Conflict}\\
\citet{Tasic2016} reported approximately five \textit{Vip} transcriptomic types in mouse V1, whereas \citet{Tasic2018} recovered roughly fourteen \textit{Vip} t-types in mouse VISp+ALM, an ostensibly large expansion within two years from the same laboratory. The disagreement reflects taxonomic depth rather than biology: clustering resolution, cell number, and area sampling were all increased. The \textit{Vip} subclass and its supertype-level subdivisions are robustly replicable; cluster-level (t-type) identifiers are dataset- and method-dependent \citep{Yao2021a, Yao2021b, Yao2023, Crow2018}.
\end{framed}

The same instability appears when transcriptomic taxonomies are compared against the marker-protein subdivisions that pre-existed scRNA-seq. \citet{Tremblay2016} reviewed the immunohistochemical and slice-electrophysiology literature and proposed three principal VIP subdivisions defined by secondary marker expression: VIP/CCK multipolar basket cells, VIP/CR irregular-spiking bipolar cells, and a small VIP/ChAT bipolar subtype. \citet{Mayer2018}, profiling early postmitotic CGE and MGE interneuron precursors by scRNA-seq, identified the cardinal mature interneuron lineages --- including a \textit{Vip} lineage --- already pre-specified before terminal differentiation. Either framework --- three marker classes or roughly fourteen t-types --- can be defended as a ``true'' partition; they differ because they apply different similarity thresholds at different developmental times to populations measured with different feature spaces. The granularity question is therefore methodological rather than empirical: at the marker-protein resolution of \citet{Tremblay2016} and \citet{Paul2017}, who described an interneuron-selective disinhibitory VIP/CR subpopulation and a CCK-basket subpopulation as the dominant axes of VIP class structure, three subdivisions suffice; at the cluster resolution of \citet{Tasic2018} or \citet{Yao2021a}, more do.

\begin{framed}
\textbf{Conflict}\\
\citet{Tremblay2016} reviewed VIP heterogeneity as approximately three marker-defined subdivisions (VIP/CCK, VIP/CR, VIP/ChAT) supported by combinatorial immunohistochemistry and slice physiology. \citet{Mayer2018} resolved progenitor-stage transcriptomic signatures that already segregate the cardinal interneuron lineages, including a \textit{Vip} precursor population, broadly consistent with the within-class type counts of contemporary atlases \citep{Tasic2018, Yao2021a}. The discrepancy reflects different measurement modalities (peptide IHC vs scRNA-seq) and different similarity thresholds rather than incompatible biology; both partitions are nested within the same VIP subclass.
\end{framed}

The combinatorial slice-electrophysiology and IHC studies that preceded scRNA-seq already foreshadowed this resolution-dependence. \citet{Kawaguchi1996} and \citet{Kawaguchi1997a} distinguished PV+ fast-spiking, late-spiking neurogliaform, and a heterogeneous regular- and burst-spiking group containing SST/Martinotti and VIP/double-bouquet cells in rat frontal cortex. \citet{Kawaguchi1997b} further showed that cholinergic agonists selectively depolarise VIP- and SST-immunoreactive interneurons but not PV fast-spiking or late-spiking cells, a neurochemical fingerprint that survives in modern transcriptomic profiles of VIP receptor expression. \citet{Toledorodriguez2005} applied multiplex single-cell RT-PCR for CB, PV, CR, NPY, VIP, SST, and CCK to 268 morphologically identified rat S1 interneurons and recovered seven distinct expression-based clusters that preferentially contained one anatomical class --- an early demonstration that combinatorial neuropeptide/calcium-binding-protein codes carry cell-type information that aligns broadly, but not perfectly, with morphology. \citet{Wang2004b} characterised Martinotti cells in juvenile rat S1 and showed that layer-specific axonal targeting and accommodating firing pattern are anchored within the SST class, providing a methodological template subsequently applied to VIP. Read with hindsight, these pre-genomic studies recovered roughly the marker-protein resolution that \citet{Tremblay2016} formalised as three principal VIP subtypes; the scRNA-seq atlases that followed nest those marker subdivisions inside finer-grained cluster tiers without invalidating them.

\begin{figure}[!htbp]
\centering
\includegraphics[width=1\linewidth]{files/fig_sec2_vip_taxonom-6bf2f216064b584c61163e35112c2d93.png}
\caption[]{VIP within the cortical GABAergic transcriptomic taxonomy. \textbf{(A)} Hierarchical schematic of the established six-subclass scheme of \citep{Tasic2018}, separating the CGE branch (\textit{Lamp5}, \textit{Vip}, \textit{Sncg}, \textit{Serpinf1}) from the MGE branch (\textit{Sst}, \textit{Pvalb}); branch widths reflect the qualitative ordering of GABAergic types per subclass and are not drawn to a quantitative scale. \textbf{(B)} Reported counts of cortical \textit{Vip} transcriptomic types (or supertypes/consensus types) across five reference atlases: \citep{Tasic2016} (mouse V1, {\textasciitilde}5 \textit{Vip} t-types), \citep{Tasic2018} (mouse VISp+ALM, {\textasciitilde}14 \textit{Vip} t-types), \citep{Yao2021a} (mouse isocortex+HPF, 6 \textit{Vip} supertypes), \citep{Hodge2019} (human MTG, 21 VIP t-types), and \citep{Lee2023a} (human cortex Patch-seq, VIP within 45 GABAergic types across four subclasses). The values are not on a common denominator and must not be interpreted as a single quantitative scale: Tasic 2018's GABAergic type count covers VISp+ALM (60 types in VISp; {\textasciitilde}50\% shared across areas), Hodge 2019 is single-nucleus RNA-seq from human MTG (75 t-types include both excitatory and inhibitory), and the Yao 2021 supertype tier is a level coarser than Tasic 2018 t-types --- supertypes are subdivisions WITHIN the three CGE subclasses, not independent categories. \textbf{(C)} Subset frequencies for VIP/ChAT and VIP/CR cells reported as approximate ranges across studies; values aggregate immunolabel and transcriptomic estimates and span {\textasciitilde}5--25\% of cortical VIP cells, with VIP/ChAT consistently the rarer subset \citep{Tremblay2016, Granger2020, Tasic2016, Tasic2018}. \textbf{(D)} Schematic of the three principal VIP molecular subdivisions reviewed by \citep{Tremblay2016}: VIP/ChAT cholinergic bipolar cells, VIP/CR interneuron-selective bipolar cells, and VIP/CCK multipolar basket cells, with annotated marker genes (\citep{Granger2020, Paul2017, Miczan2021}). All values plotted derive from different taxonomic levels (subclass vs t-type vs cross-area overlap) and must not be interpreted as a single quantitative scale; mouse\_homologs\_matched in \citep{Hodge2019} reflects cross-species mapping rather than the human type repertoire size; SST ``similar diversity'' is qualitative and is not plotted; Gouwens 2020 patch-seq cell counts (4,270/2,955/517) are nested filtering tiers and MET-types (28), me-types (26), and t-types ({\textasciitilde}60) are three different taxonomic hierarchies plotted on separate axes; Yao 2023 cluster counts and Bakken 2021 primate consensus counts are reported on the BICCN whole-brain and primate cortex denominators respectively and not on a single mouse-cortex denominator. Source: hand-extracted values from each atlas's main text and Tables 1/Fig.1, with cross-area / cross-species denominators annotated above each bar.}
\label{fig-vip-taxonomy}
\end{figure}

\section{Discrete clusters or continuous gradients within VIP}

Having established that VIP subclass identity is robust while t-type counts are not, the deeper methodological question is whether within-VIP heterogeneity is best modelled as a finite set of discrete clusters or as a low-dimensional continuum. Multimodal datasets address this directly. \citet{Gouwens2020a} integrated Patch-seq electrophysiology, transcriptomics, and morphology in mouse VISp and defined twenty-eight morphoelectric-transcriptomic (MET) types of cortical GABAergic interneurons, providing a unified framework that reconciled the {\textasciitilde}60 t-types and {\textasciitilde}26 me-types of preceding atlases by anchoring discrete clusters across modalities. The same dataset showed clean four-way segregation of \textit{Lamp5}, \textit{Vip}, \textit{Sst}, and \textit{Pvalb} in electrophysiology UMAP space, but within the \textit{Vip} family --- and equivalently within \textit{Sst} and \textit{Pvalb} --- fine subdivisions were less morpho-electrically separable. \citet{Scala2021} reached a complementary conclusion using Patch-seq on more than 1,300 mouse motor cortex (M1) neurons: broad transcriptomic families (\textit{Vip}, \textit{Pvalb}, \textit{Sst}) had distinct non-overlapping morpho-electric phenotypes, but within-family individual t-types were not well separated in morpho-electric space, suggesting a continuum rather than a discrete partition at the within-VIP level. The same study identified one clear exception --- a transcriptomically isolated \textit{Vip Lamp5 Lhx6} type morphologically corresponding to deep L5/L6 neurogliaform cells --- supporting the existence of locally discrete VIP-related types embedded within an otherwise continuous space.

\begin{framed}
\textbf{Conflict}\\
\citet{Gouwens2020a} argued that mouse VISp Patch-seq supports twenty-eight discrete MET-types with congruent morphology, electrophysiology, and t-type identity, including multiple resolvable subdivisions within the \textit{Vip} subclass. \citet{Scala2021}, using Patch-seq on M1, agreed that subclass-level (\textit{Vip}/\textit{Sst}/\textit{Pvalb}) separation is robust but reported that within-VIP morpho-electric phenotypes form a continuum rather than discrete clusters, with only the deep-layer \textit{Vip Lamp5 Lhx6} type morphologically isolated. Subsequent Patch-seq studies of human cortex \citep{Lee2023a, Kalmbach2021} and the broader Patch-seq quality-control literature \citep{Tripathy2018} are consistent with the joint reading that subclass identity is multimodally discrete while within-subclass discreteness depends on cortical area, feature engineering, and contamination control.
\end{framed}

A second, conceptually orthogonal axis of the discrete-versus-continuous question concerns whether within-VIP/CGE heterogeneity itself reflects discrete cell-type categories or a graded transcriptomic state aligned with brain-wide variables. \citet{Bugeon2022} profiled cortical inhibitory interneurons across V1 layers 1--3 jointly with their \textit{in vivo} behavioural-state responses and showed that, while neurons remain assignable to a discrete subtype hierarchy (subclasses, types, and subtypes), behavioural-state firing modulation also varies smoothly along a dominant transcriptomic axis that runs across the CGE-derived (Vip/Lamp5/Sncg) compartment, so within-VIP gradients capture meaningful state-response variation that the discrete labels do not by themselves resolve. This continuous-axis interpretation sits in tension with the strictly hierarchical discrete-supertype taxonomy of \citet{Yao2021a}, in which roughly 1.3 million mouse cortex and hippocampal cells were partitioned into 379 clusters, organised into {\textasciitilde}24 GABAergic supertypes, with the Vip subclass yielding six discrete supertypes that are stable to subsampling and well-separated in expression space.

\begin{framed}
\textbf{Conflict}\\
\citet{Bugeon2022} reported that, although their hierarchical taxonomy assigns interneurons to discrete subclasses, types, and subtypes, behavioural-state firing modulation also varies along a dominant transcriptomic axis that crosses the CGE/VIP compartment; the gradient captures within-subclass variation that any single discrete subtype label does not by itself resolve. \citet{Yao2021a}, by contrast, presented a 1.3-million-cell mouse cortex+HPF atlas in which the Vip subclass partitions into six discrete and reproducible supertypes within a hierarchical class\rightarrowsubclass\rightarrowsupertype\rightarrowcluster scheme. The two views are reconcilable if the Yao2021a discrete supertypes correspond to local concentrations of cells along the same transcriptomic axis Bugeon2022 reads as continuous: clustering at coarse resolution recovers stable supertypes, while behaviourally meaningful variation in firing modulation tracks position along the underlying gradient. The implication for downstream sections is that discrete cell-type assignment and continuous transcriptomic-state position carry partially independent information, and that functional interpretations of ``VIP cell'' populations should consider both.
\end{framed}

The continuous-versus-discrete distinction is partly a methodological artefact of feature dimensionality and clustering criteria. \citet{Tripathy2018} re-analysed five Patch-seq datasets and found that off-target cell-type mRNA contamination and large variability in mRNA yield were systematic confounds; their marker-gene-based quality scoring improved the correspondence between transcriptome and electrophysiology, narrowing --- but not eliminating --- apparent within-subclass heterogeneity. \citet{Nandi2022} showed that biophysical models constrained by Patch-seq electrophysiology cluster cleanly by transcriptomic subclass, indicating that ion-channel composition encodes identity at the subclass level even when sub-subclass clustering is unstable. The methodological point is that the within-VIP continuum reported by \citet{Scala2021} and the within-VIP discrete clusters reported by \citet{Gouwens2020a} are not mutually exclusive: both are compatible with a model in which the \textit{Vip} subclass occupies a low-dimensional manifold whose structure can be partitioned into clusters at one threshold and approximated by gradients at another, and on which a small number of transcriptomically isolated outliers (the \textit{Vip Lamp5 Lhx6} deep-layer NGCs of \citet{Scala2021}, the upper-layer \textit{Tac2}+/\textit{Cxcl14}+ population identified across mouse VISp/ALM/MOp by \citet{Wu2022}) form genuinely discrete sub-types.

A separate continuous-gradient signal arises from the laminar dimension. \citet{Pronneke2015} showed that layer II/III VIP cells in mouse barrel cortex differ significantly from layer IV--VI VIP cells in morphology and membrane properties, implying functional rather than purely categorical heterogeneity. \citet{Wu2022} extended this with cross-area transcriptomic comparison, showing that upper-layer (L1--III) and deeper-layer (L4--VI) cortical VIP+ interneurons are transcriptionally distinct, with \textit{Tac2} and \textit{Cxcl14} emerging as conserved upper-layer marker genes across mouse VISp, ALM, and MOp and detectable in human cortex. Behavioural-state context further refines this map: state-dependent gene-expression analyses applied to interneuron populations in vivo (reviewed in CGE~origin~and~the~5-HT3AR~/~Adarb2~lineage for developmental context, and revisited in later sections for circuit physiology) reinforce the view that within-VIP molecular heterogeneity is partly continuous along the cortical depth axis. The implication for downstream sections is that any cell-type-level functional claim about ``VIP cells'' should be qualified by laminar position and cortical area.

\section{Marker-driven labels are not transcriptomic categories}

The dominant practical problem in interpreting VIP-related literature is that marker proteins, immunolabels, Cre-driver lines, and scRNA-seq subclasses partition cells differently. The classical immunohistochemical work of \citet{Kawaguchi1996}, \citet{Kawaguchi1997a}, and \citet{Hajos1996}, together with the laminar census of \citet{Staiger1997} and \citet{Staiger2004} showing that essentially all VIP-immunoreactive interneurons receive PV-basket input and that calbindin-positive interneurons receive symmetric VIP+ contacts, defined ``VIP cells'' operationally by anti-VIP antibody reactivity. The transcriptomic class label is similar but not identical: \citet{Tasic2018} and \citet{Yao2021a} define the \textit{Vip} subclass by clustered expression of \textit{Vip} together with several co-expressed markers, with cells that express \textit{Vip} mRNA at low levels potentially excluded from the cluster. \citet{Lee2010} showed that 5-HT3AR-GFP captures the entire CGE-derived compartment, including not only \textit{Vip} but also \textit{Lamp5} and \textit{Sncg} subclasses, so a 5-HT3AR-driver phenotype is broader than a \textit{Vip}-driver phenotype. The most widely used genetic-access tool, the VIP-Cre or VIP-IRES-Cre line \citep{Tasic2018, Paul2017}, captures most but not all \textit{Vip} subclass cells, and may also label transiently \textit{Vip}-expressing cells from neighbouring CGE subclasses; the off-target risk is most acute at the boundary with \textit{Sncg} and \textit{Lamp5} \citep{Tasic2018, Yao2021a}. Reported VIP fractions of cortical inhibition therefore depend on which definition is used, with \%VIP+ estimates of cortical GABAergic neurons varying across the literature in part because antibody-based, Cre-line, and t-type-based denominators are not equivalent.

A second source of slippage is the marker-defined subdivision within VIP itself. \citet{Tremblay2016}'s synthesis specified VIP/CCK basket cells, VIP/CR irregular-spiking bipolar cells, and a small VIP/ChAT bipolar subset. The molecular evidence supporting VIP/CCK as a coherent unit is now more nuanced: \citet{Miczan2021} showed that NECAB1 and NECAB2 are the predominant calcium-binding proteins of CB1/CCK-positive GABAergic interneurons, closing a long-standing gap in which this large class lacked a typifying calcium-binding protein, in contrast with the calbindin/calretinin signatures of other CCK subsets. The VIP/CR axis maps reasonably well onto the \textit{Vip-Mybpc1}-related and \textit{Vip Crispld2}-like t-types of mouse atlases and onto the human VIP/CALB2/TAC3 In1d subpopulation identified by \citet{Lake2018}, validating cross-modal identification. The VIP/ChAT axis is more problematic, and we treat it separately below.

The third axis of slippage is laminar and inter-areal. \citet{Wu2022} showed that upper- and deeper-layer cortical VIP+ interneurons are transcriptionally distinguishable; \citet{Pronneke2015} independently showed laminar differences in morphology and intrinsic properties; and the Patch-seq atlases of \citet{Gouwens2020a} and \citet{Scala2021} show that the same \textit{Vip} t-types are not equally sampled across cortical areas. \citet{Hostetler2023} provided a parallel cautionary tale within SST: intersectional Flp/Cre crosses identified two distinct layer-1-targeting SST subtypes with non-overlapping marker expression and electrophysiology, illustrating that even within the dendrite-targeting SST class, marker-defined and intersectional definitions resolve different cell populations. The implication for VIP work is that subdivisions visible in one cortical area may be undersampled or absent in another, and that immunolabel-defined ``VIP'' cell counts can drift severalfold between studies because the underlying populations being labelled are not identical.

\begin{figure}[!htbp]
\centering
\includegraphics[width=1\linewidth]{files/fig_sec2_marker_tree-d092bdcd51cdd9e697e85550cb97d758.png}
\caption[]{Marker-protein decision tree linking immunolabel-, Cre-driver-, and scRNA-seq-based definitions of ``VIP''. \textbf{(A)} From cortical GABAergic neurons ({\textasciitilde}20\% of NeuN+ cortical neurons in mouse \citep{Tamamaki2003}) the CGE branch is defined by 5-HT3AR/\textit{Prox1}/\textit{Sp8} expression \citep{Lee2010, Tremblay2016}; within the CGE branch, the \textit{Vip} subclass is defined by clustered expression of \textit{Vip} and co-markers \citep{Tasic2018, Yao2021a}. The marker-defined subdivisions VIP/ChAT, VIP/CR, VIP/CCK, and VIP/CALB2 partition the \textit{Vip} subclass with overlapping rather than disjoint boundaries \citep{Tremblay2016, Paul2017, Granger2020}. \textbf{(B)} Schematic Venn comparison of populations captured by VIP-Cre, VIP-IRES-Cre, and anti-VIP immunolabel: Cre-driver populations include most but not all transcriptomically defined \textit{Vip} cells and may capture transiently \textit{Vip}-expressing neighbours from \textit{Sncg} and \textit{Lamp5} subclasses; immunolabel captures cells with detectable VIP peptide, biasing toward higher-expressing subsets \citep{Tasic2018, Yao2021a, Pronneke2015}. \textbf{(C)} Schematic of the immunolabel-to-t-type mapping: VIP/CR maps largely one-to-many onto multiple bipolar \textit{Vip} t-types and onto human In1d/VIP-CALB2-TAC3 nuclei \citep{Lake2018, Tremblay2016}; VIP/CCK maps many-to-many onto \textit{Vip} t-types overlapping with \textit{Sncg} CCK-basket-related types \citep{Miczan2021, Paul2017}; VIP/ChAT maps approximately one-to-one onto a small bipolar \textit{Vip-Chat} type that is identified at the t-type level by \citep{Tasic2016, Granger2020} but is not formally labelled cholinergic at the cluster level in larger atlases \citep{Yao2023}. This figure is a schematic without quantitative data; it summarises the operational decisions readers must make to compare VIP-related findings across studies that use different definitions, and intentionally does not place numerical estimates on the overlap regions.}
\label{fig-vip-marker-tree}
\end{figure}

\section{VIP/ChAT: a discrete subtype that resists cluster-level capture}

The VIP/ChAT subset is the prototypical example of how a phenotypically discrete cell can fail to register as a discrete transcriptomic cluster, depending on atlas resolution and labelling thresholds. \citet{Tasic2016} identified \textit{Vip-Chat} as a discrete t-type uniquely expressing choline acetyltransferase among \textit{Vip}+ cells in upper cortical layers and confirmed \textit{Chat} expression by ISH. \citet{Granger2020} then provided the functional validation: nearly all cortical ChAT+ neurons in mouse are specialised VIP+ interneurons that co-release GABA strongly onto inhibitory interneurons and acetylcholine sparsely onto layer-1 interneurons, consistent with the prior physiological observation by \citet{Kawaguchi1997b} that cholinergic agonists differentially gate VIP- and SST-class interneurons. \citet{Tremblay2016} and \citet{Paul2017} accordingly treated VIP/ChAT cells as a distinct subset of the VIP class, reporting them as a small fraction of cortical VIP cells.

The complication arose with whole-brain atlases that apply uniform clustering thresholds across many regions. \citet{Yao2023} reported that \textit{Slc18a3} (the vesicular acetylcholine transporter) and \textit{Chat} mRNAs are detected in several clusters within the \textit{Vip} GABAergic subclass of isocortex, but expression at the cluster level did not cross thresholds for formal labelling as cholinergic, and the whole-brain atlas does not therefore name a ``VIP/ChAT'' cluster. The most parsimonious reconciliation is that VIP/ChAT cells exist as a phenotypically distinct subtype with target-specific ACh co-release, but their cluster-level transcriptomic separation is subtle relative to the variation across the broader \textit{Vip} subclass and may not survive threshold-based labelling at the cluster tier in larger atlases. This is the same resolution-dependence noted earlier --- the more cells and more brain regions an atlas covers, the more the threshold for ``discrete cluster'' rises and the more subtle subtypes are absorbed into broader clusters.

\begin{framed}
\textbf{Conflict}\\
\citet{Tasic2016} identified VIP/ChAT as a discrete \textit{Vip-Chat} transcriptomic t-type uniquely expressing \textit{Chat} among \textit{Vip}+ cells, with ISH validation. \citet{Yao2023} reported that \textit{Slc18a3}/\textit{Chat} are detected in several clusters within the \textit{Vip} subclass of mouse isocortex but at sub-threshold cluster-level expression, so the whole-brain atlas does not formally label any \textit{Vip} cluster cholinergic. The disagreement is reconciled by the threshold-versus-resolution issue: VIP/ChAT cells are functionally distinct and target-specific in ACh co-release \citep{Granger2020, Kawaguchi1997b}, but their cluster-level transcriptomic separation is subtle and atlas-threshold dependent. Reported \%ChAT+ VIP fractions vary by approximately fivefold across studies, partly because of this thresholding and partly because immunolabel and Cre-driver definitions of ``VIP'' capture different denominators \citep{Tasic2016, Granger2020, Tremblay2016}.
\end{framed}

The VIP/ChAT case has a wider methodological lesson. The classical IHC literature (\citet{Kawaguchi1996, Kawaguchi1997a, Hajos1996}) repeatedly identified small, regionally biased subtypes --- hilar VIP/CR interneuron-selective cells, layer-1-projecting VIP bipolar cells, double-bouquet VIP cells --- whose status as transcriptomically discrete t-types is variable across modern atlases. The atlases are not ``wrong''; they apply uniform clustering criteria that are calibrated to recover the dominant axes of variation. When a sub-population is small, regionally biased, and defined by graded rather than switch-like differences in marker expression, it is precisely the kind of cell that escapes cluster-level resolution. The atlases of \citet{Yao2021a} and \citet{Yao2023} are explicit that supertype and cluster boundaries are projection-method dependent and that finer divisions are recoverable on demand by re-clustering subsets at higher resolution. \citet{Bakken2018} provided complementary methodological evidence that single-nucleus and single-cell RNA-seq capture similar discriminative power for closely related types when intronic reads are retained, ensuring that the resolution limits are biological rather than purely technical.

\section{Cross-species comparison: conservation and divergence}

Whether the mouse VIP taxonomy transfers directly to human cortex is a non-trivial empirical question. \citet{Hodge2019} profiled human middle temporal gyrus (MTG) by single-nucleus RNA-seq and identified seventy-five transcriptomic cell types in 15,928 nuclei from eight donors, with VIP the most diverse interneuron subclass at twenty-one types. The same study reported that interneuron-discriminating gene sets were largely conserved between mouse and human, supporting alignment of approximately thirty-seven homologous t-types despite species-specific expression differences. Independently, \citet{Lake2018} used integrative snDrop-seq plus scTHS-seq on more than 60,000 human brain cells to spatially resolve a VIP+/CALB2+/TAC3+ upper-layer interneuron subpopulation (In1d) corresponding to the established VIP/CR class in mouse, validating cross-modal identification of human VIP types. \citet{Boldog2018} then identified the human-specific ``rosehip'' cortical interneuron in human L1 that does not match any mouse \textit{Vip}/\textit{Sst}/\textit{Pvalb} cluster, while several other human L1 t-types matched mouse \textit{Vip}+ types --- a mixture of conservation and divergence within a single layer.

\citet{Mayer2018}, profiling early postmitotic CGE/MGE interneuron precursors in mouse, demonstrated that progenitor-stage transcriptional signatures already segregate into the cardinal cortical interneuron lineages, including a \textit{Vip} lineage, before terminal differentiation. The cross-species question is then whether human VIP cells share these progenitor-stage transcriptional programs and how the species-specific elaborations identified by \citet{Hodge2019}, \citet{Boldog2018}, and \citet{Lake2018} map onto the mouse axes. \citet{Bakken2021a} directly addressed primate-specific elaboration: more \textit{Vip} subtype consensus clusters could be resolved by pairwise alignment between human and marmoset M1 than between either primate and mouse, indicating closer homology of VIP types within primates and a divergence step at the rodent--primate split. \citet{Bakken2021b} reported a parallel pattern in primate dorsal lateral geniculate nucleus, where GABAergic interneuron diversity is expanded relative to mouse, suggesting species-specific elaboration of inhibitory diversity that may parallel expanded VIP type diversity in primate cortex. \citet{Wei2022} extended the comparative analysis to macaque visual cortex, distinguishing twenty-five excitatory and thirty-seven inhibitory types and reporting in cross-species comparison that glutamatergic neurons may be more diverse across species than GABAergic neurons or non-neuronal cells, consistent with VIP-class identity being conserved at the subclass level while elaborated at the type level. Patch-seq datasets that anchor human transcriptomic types in morpho-electric phenotypes are now available: \citet{Lee2023a} reported a human cortex GABAergic taxonomy of forty-five transcriptomic types across four principal interneuron subclasses (PVALB, SST, VIP, LAMP5/PAX6) using Patch-seq with enhancer-AAV labelling for VIP enrichment, and \citet{Kalmbach2021} showed using human neurosurgical Patch-seq that transcriptomically defined L5 extratelencephalic projection neurons have conserved morpho-electric properties between human and rodent, providing a methodological proof of concept for human-rodent cell-type alignment. \citet{Chartrand2023} documented human-specific L1 interneuron transcriptomic types (e.g. VIP \textit{PCDH20} and SST \textit{BAGE2}) homologous to deeper mouse t-types but lacking morphologically equivalent neurons in mouse L1, indicating type-specific divergence rather than the absence of shared subclass identity.

\begin{framed}
\textbf{Conflict}\\
\citet{Hodge2019} argued that mouse interneuron-discriminating gene sets transfer to human MTG, supporting direct cross-species alignment of approximately thirty-seven homologous VIP/SST/PV t-types and a transferable taxonomy at the type level. \citet{Mayer2018} mapped a complementary picture from progenitor-stage CGE/MGE scRNA-seq: cardinal mature lineages (including the \textit{Vip} lineage) are pre-specified before terminal differentiation, but the species-specific elaborations later observed in primate cortex \citep{Hodge2019, Boldog2018, Bakken2021a, Chartrand2023, Wei2022} indicate that direct transfer of mouse t-type identifiers to human is incomplete: subclass identity is conserved while within-subclass type repertoires diverge. The synthesis is that mouse t-type names label valid mouse-specific subdivisions and human cortical taxonomies require species-specific subtype definitions even within the conserved \textit{Vip} subclass.
\end{framed}

Comparative evidence from non-mammalian vertebrates frames the conservation question on a longer time scale. \citet{Tosches2018} inferred from comparative scRNA-seq of reptile pallia that mammalian neocortical glutamatergic layers represent new cell types built by diversifying ancestral gene-regulatory programs, while diverse cortical interneuron classes already existed in the common amniote ancestor. The implication is that the \textit{Vip} subclass and its CGE-derived neighbours are an ancient feature of vertebrate cortical inhibition rather than a mammalian innovation, and that the cross-species variability documented by \citet{Hodge2019}, \citet{Bakken2021a}, and \citet{Wei2022} represents elaboration of an old framework rather than wholesale reinvention. Methylome and chromatin data are beginning to support this longer-time-scale view: \citet{Liu2021} generated a single-cell DNA methylation atlas of mouse brain identifying CGE-derived \textit{Lamp5} and \textit{Vip} lineages as epigenomically distinguishable, and \citet{Bakos2025} showed that fast-spiking interneurons in human neocortex possess axon-initial-segment adaptations distinct from rodent FS cells, indicating that even cells with conserved transcriptomic identity can show species-specific elaboration of biophysical properties. The molecular taxonomy of ``VIP'' is therefore best read as a conserved subclass framework with species-, area-, and laminar-specific elaboration of within-subclass diversity.

\section{Methodological reliability of the taxonomic framework}

A useful sanity check is that the multimodal evidence converges on the subclass framework even from preparations that violate standard scRNA-seq assumptions. \citet{Bakken2018} showed that single-nucleus RNA-seq, despite detecting fewer genes per cell than whole-cell scRNA-seq, can discriminate closely related neuronal types when intronic reads are retained --- a methodological observation that licenses the use of postmortem human snRNA-seq for cross-species mapping. \citet{Crow2018} independently quantified replicability and found that subclass identity is highly reproducible while fine-grained sub-subclass identity replicates progressively less well, formalising the resolution-dependence anticipated by \citet{Tasic2016} and \citet{Tasic2018}. The Patch-seq quality-control work of \citet{Tripathy2018} showed that off-target mRNA contamination from acute slices materially affects within-subclass clustering, and that marker-gene-based quality scoring can be used to filter contaminated cells before clustering. \citet{Watanabe2019} demonstrated that scRNA-seq cell-type catalogues can be coupled to GWAS summary statistics to assign genetic risk to specific neuronal types, an integration that depends on stable subclass labels and motivates further refinement of within-VIP definitions.

The pre-existing electrophysiological taxonomy survives this transcriptomic re-organisation in a graded way. \citet{Kawaguchi1996} and \citet{Kawaguchi1997a} defined fast-spiking, late-spiking, regular-spiking, and burst-spiking electrophysiological classes that map approximately onto the modern \textit{Pvalb}, \textit{Lamp5} (neurogliaform), \textit{Vip}/\textit{Sst}, and \textit{Vip} (irregular-spiking) classes respectively, with caveats. \citet{Kawaguchi1997b} showed that cholinergic agonists differentially gate VIP- and SST-class interneurons but not PV fast-spiking or late-spiking cells, providing one of the earliest neurochemical fingerprints that aligns with modern transcriptomic receptor-expression profiles. \citet{Toledorodriguez2005} recovered seven combinatorial neuropeptide/calcium-binding-protein clusters in juvenile rat S1 that broadly correspond to modern subclasses; \citet{Wang2004b} characterised SST/Martinotti laminar specificity that is now integrated into multimodal Sst-Calb2 versus Sst-Chodl distinctions; and \citet{Dalezios2002} showed by EM that essentially all VIP+ GABAergic terminals enriched in mGluR7a target SST/mGluR1$\alpha$-expressing interneurons, providing direct ultrastructural evidence for the disinhibitory VIP\rightarrowSST motif at the same time that the molecular markers were being catalogued. The pattern is that classical morphological and electrophysiological subdivisions are recoverable from modern multimodal datasets at the subclass and supertype levels but become harder to map one-to-one at the cluster tier --- exactly as the resolution-dependence diagnosis predicts.

Recent atlases have begun to formalise these resolution-dependent statements. \citet{Yao2023}'s whole-brain mouse atlas --- thirty-four classes, 338 subclasses, 1,201 supertypes, 5,322 clusters --- explicitly proposes that classes and subclasses are stable across regions and methods, supertypes are stable within most regions, and clusters are the level at which dataset- and area-specific structure manifests. \citet{Lee2023a} used enhancer-AAV labelling to over-sample VIP cells in human Patch-seq and reported that subclass alignment to mouse VIP holds while type-level alignment requires species-specific subtype definitions, an empirical demonstration of the same hierarchy. The methodological corollary for downstream sections of this review is that mechanistic claims about ``VIP cells'' should specify whether they are subclass-level (robustly transferable across species and atlases), supertype-level (transferable within species and method), or cluster-level (dataset-specific). The ``VIP'' used in disinhibitory-circuit experiments is operationally a Cre-driver-defined approximation of the subclass --- not a t-type --- and most subsequent functional generalisations should be read at that level of resolution.

\section{Bridge to development}

The transcriptomic taxonomy this section catalogues has, by construction, taken the adult brain as its starting point: cells were dissociated, sequenced, and clustered, and the resulting groups were named for their dominant marker. The next section (see CGE~origin~and~the~5-HT3AR~/~Adarb2~lineage) inverts the direction of inquiry. If the \textit{Vip} subclass and its supertype-level subdivisions are robust at adulthood, where do they come from? \citet{Mayer2018} already showed that the cardinal CGE lineages are pre-specified at progenitor stages by transcriptional signatures; \citet{Demarcogarcia2011} showed that activity-dependent migration before P3 is required for Reelin+ and CR+ CGE-derived interneurons but not for VIP+ interneurons, suggesting differential activity-dependence of subtype maturation; and \citet{Frazer2017} distinguished three main molecular types within developing Htr3a-GFP+ CGE-derived precursors. The developmental section will read these and related findings as the prior to the adult atlas: it will examine how CGE progenitors yield the molecular axes catalogued here, how postnatal maturation refines and sometimes prunes them, and how species-specific lineage logic accounts for the divergence in primate VIP type repertoires reported by \citet{Bakken2021a}, \citet{Hodge2019}, and \citet{Wei2022}. The molecular ground truth established in this section therefore serves a dual purpose: it fixes the cell types that the rest of the review will refer to, and it sets the empirical target for the developmental, morphological, electrophysiological, and cross-species accounts that follow.

Section \href{/development}{Developmental Origins and Postnatal Maturation} traces those identity criteria back through their developmental origins, asking how the gene-regulatory programmes that specify the CGE/5-HT3AR lineage map onto the adult t-type taxonomy laid out here.


\section{Developmental Origins and Postnatal Maturation}
\label{sec:03_development}

Cortical VIP interneurons are members of the caudal-ganglionic-eminence (CGE)--derived 5-HT3AR / \textit{Adarb2} lineage that also gives rise to the \textit{Lamp5}, \textit{Sncg} and \textit{Serpinf1} subclasses \citep{Lee2010, Miyoshi2010, Tasic2018}. The cells are born in a late embryonic neurogenic wave, migrate tangentially along several caudo-rostral streams, settle preferentially in superficial layers, and acquire mature firing, synaptic and morphological properties through a protracted postnatal programme that combines cell-intrinsic transcription-factor cascades with experience- and activity-dependent refinement \citep{Miyoshi2010, Miyoshi2015, Lim2018, Demarcogarcia2025}. Two features of this developmental programme matter most for the rest of this review. First, the gene-regulatory logic that distinguishes VIP from sister CGE classes is now traceable to a small set of transcription factors --- most notably Prox1, Sp8/Sp9 and COUP-TFII (Nr2f2) --- whose conditional perturbation produces subtype-selective phenotypes \citep{Miyoshi2015, Stachniak2021, Wei2019, Kanatani2008}. Second, every quantitative claim about CGE-lineage proportion or VIP-subtype identity that derives from mouse fate-mapping must be re-examined when applied to primate cortex, where lineage proportions, marker codes and progenitor architecture diverge \citep{Hodge2019, Hansen2013, Ma2013, Marin2025}. The remainder of this section reviews each step in turn, and explicitly flags the cross-species comparisons that complicate translation.

\section{CGE origin and the 5-HT3AR / Adarb2 lineage}\label{sec-development-lineage}

The mapping of GABAergic subtypes onto subpallial origins is one of the most extensively replicated facts in cortical development. Homochronic transplantation of medial-ganglionic-eminence (MGE) versus caudal-ganglionic-eminence (CGE) explants into host mouse embryos established that PV and SST cardinal classes derive primarily from the MGE, whereas reelin-, calretinin-, NPY- and VIP-expressing cardinal classes derive primarily from the CGE \citep{Bandler2017}. Joint contributions of progenitor spatial origin, birthdate, lineage relationships and division mode further refine subtype identity, as resolved at scRNA-seq resolution by \citet{Mi2018}. Inducible \textit{Mash1}/\textit{Olig2}-CreER fate mapping subsequently quantified the CGE contribution at {\textasciitilde}30\% of all cortical interneurons across at least nine distinct mature subtypes, with CGE neurogenesis peaking at E16.5 --- about three days later than MGE neurogenesis --- and 75\% of CGE progeny populating supragranular layers regardless of birthdate \citep{Miyoshi2010}. A parallel hippocampal study using transgenic and inducible reporter lines documented the same temporal staggering: an MGE wave between E9--E12 generates PV and SST hippocampal interneurons, whereas a CGE wave between E12--E16 generates the CCK+, calretinin+, VIP+ and reelin+ populations that populate stratum lacunosum-moleculare and deep stratum radiatum.

Within this CGE-derived population, the serotonin 5-HT3A receptor (encoded by \textit{Htr3a}) provides a near-complete molecular handle. \citet{Lee2010} showed that 5-HT3AR expression demarcates essentially all neocortical GABAergic interneurons that lack PV/SST and labels the entire CGE lineage in mouse cortex \citep{Lee2010}. Homochronic in utero grafts of E14 \textit{Htr3a}-GFP+ cells into host cortex reproduce the supragranular laminar bias and the neurochemical profile (VIP, NPY, reelin, calretinin) characteristic of CGE-derived interneurons \citep{Vucurovic2010}. Modern scRNA-seq taxonomies recapitulate this division at the transcriptomic level: in mouse cortex, CGE origin maps onto the \textit{Vip}, \textit{Lamp5}, \textit{Sncg} and \textit{Serpinf1} subclasses \citep{Tasic2018}, while in human middle temporal gyrus the CGE branch is identified by \textit{ADARB2} expression and accounts for {\textasciitilde}50\% of cortical interneurons against {\textasciitilde}44\% for the \textit{LHX6}+ MGE branch \citep{Hodge2019}. The CGE label therefore unifies a heterogeneous group of upper-layer interneurons whose shared origin is more easily defined than any single mature feature.

Two clarifications are needed. First, the CGE is not the only ventral source of cortical inhibitory cells: the embryonic preoptic area (POA) generates a smaller, mostly reelin+ population that contributes to layer-1 NGC-like cells and that --- because POA-derived cells share PROX1, NR2F2 and reelin expression with CGE-derived NGCs --- partially confounds CGE-only fate maps \citep{Gelman2009, Niquille2018}. Second, the dorsal MGE contributes a calretinin+/SST+ lineage (controlled in part by \textit{Nkx6-2}) that overlaps in marker space with CGE-derived calretinin+ VIP cells but corresponds to Martinotti-like rather than VIP-bipolar morphologies. These caveats explain why marker-only assignments (e.g. ``CR+'' or ``5-HT3AR+'') give an imperfect read-out of CGE lineage and why fate-mapping with intersectional driver lines remains the gold standard.

\begin{framed}
\textbf{Conflict --- Developmental origin of cortical neurogliaform cells}\\
\citet{Miyoshi2010} and Prior work \citep{Lee2010} placed cortical NGCs squarely within the CGE-derived 5-HT3AR+ lineage. Earlier reports \citep{Niquille2018} showed that a substantial fraction of cortical NGCs is generated by a 5-HT3AR+ \textit{Hmx3}+ pool in the preoptic area, expressing PROX1, NR2F2 and reelin but not VIP. Resolution: NGCs are a mixed-origin class with both CGE- and POA-derived contingents that share a 5-HT3AR/PROX1/RELN molecular surface; the conflict is therefore quantitative rather than categorical, and warns against treating ``5-HT3AR+/PROX1+'' as synonymous with ``CGE-derived'' outside intersectional fate-mapping experiments.
\end{framed}

\section{Transcription-factor cascades that specify VIP and its CGE-class neighbours}

The transcription-factor (TF) logic that specifies CGE-derived interneurons differs in detail from the \textit{Lhx6}/\textit{Sox6}-dominated MGE programme \citep{Liodis2007}. Single-cell transcriptomic profiling of E13.5 MGE and E14.5 CGE/LGE reveals a conserved mitotic-progenitor maturation trajectory across all three eminences, onto which eminence-specific TF cassettes are overlaid in postmitotic precursors to seed cardinal-class diversity \citep{Mayer2018}; chromatin-accessibility profiling shows that enhancer remodelling precedes transcriptional commitment, with CGE-specific Prox1, Sp8 and COUP-TFII programmes opening in coordinated waves \citep{Allaway2021}. By P0, scRNA-seq separates a \textit{Lhx6}-negative postmitotic CGE branch into \textit{Vip}-, \textit{Id2}- and \textit{Sncg}-aligned precursor states already distinct from MGE-derived progeny \citep{Mayer2018}.

Among the CGE-class TFs, \textbf{Prox1} has the strongest causal evidence. PROX1 is expressed in 84\% of fate-mapped CGE/LGE/POA-derived interneurons and in 100\% of VIP+ cortical interneurons in P30 mouse cortex \citep{Rubin2013}, and is selectively maintained in postmitotic CGE-derived interneuron precursors as the first identified CGE-specific postmitotic TF \citep{Miyoshi2015}. Embryonic conditional deletion of \textit{Prox1} impairs the integration of CGE-derived precursors into superficial layers and differentially regulates the postnatal maturation of RELN+, \textit{Calb2}/VIP+ and VIP+ subtypes --- establishing PROX1 as the first identified TF specifically required for CGE-lineage acquisition of subtype-specific properties. Postmitotic \textit{Prox1} ablation in early postnatal VIP cells then reveals an additional, subtype-selective role: it reduces \textit{Elfn1}-dependent short-term plasticity of excitatory inputs onto multipolar but not bipolar VIP cells, in line with adult VIP-subclass transcriptomic preservation of \textit{Prox1} expression \citep{Tasic2018, Hodge2019}.

\textbf{COUP-TFII (NR2F2)} is preferentially expressed in CGE relative to MGE and is required for the caudal migratory stream that channels CGE progeny into caudal cortex, hippocampus and amygdala \citep{Kanatani2008}. In adult mouse cortex, \textit{Cai2013} showed that COUP-TFII is also expressed in dorsal MGE-, dorsal LGE- and POA-derived cells, so COUP-TFII positivity is not a clean CGE-only marker \citep{Cai2013}. In hippocampus, COUP-TFII is selectively maintained in CGE-derived calretinin+, nNOS+, reelin+, neurogliaform and CCK/calbindin+ basket cells \citep{Fuentealba2010}. Conditional COUP-TFI inactivation perturbs both lateral and medial CGE migratory streams and disrupts SP8/COUP-TFII expression, producing laminar-specific deficits across distinct interneuron subtypes in adult cortex \citep{Touzot2016}. Human relevance: \textit{NR2F1} (COUP-TFI) loss-of-function causes Bosch--Boonstra--Schaaf optic atrophy syndrome (intellectual disability, autism, seizures) \citep{Yang2017}, and COUP-TFII is broadly expressed in human fetal forebrain GE and neocortex from 9--22 gestational weeks with peak in the CGE \citep{Reinchisi2012}.

\textbf{Sp8 and Sp9} define a second axis of CGE/dLGE specification. Sp8 is expressed in {\textasciitilde}20\% of adult cortical interneurons derived from both dorsal LGE and dorsal CGE, weakly in MGE mantle but strongly in CGE/dLGE \citep{Tao2019}, and acts cell-autonomously to specify olfactory-bulb and cortical CGE-related lineages. Co-expressed Sp8/Sp9 in the dorsal CGE SVZ jointly drive CGE-derived interneuron specification in part by repressing \textit{Pak3}, \textit{Robo1} and \textit{Slit1} and by activating Cxcl14 \citep{Wei2019}, while Sp9 alone is required in MGE for normal PV/SST development \citep{Liu2019b}. The mutually exclusive deployment of Sp9 in MGE versus Sp8/Sp9 in CGE thus partly explains the cardinal-class divergence captured in adult taxonomies. \textit{Lhx6} itself is required not for GABAergic identity but for MGE-specific specification and migration; in \textit{Lhx6}-/- mice some MGE cells aberrantly acquire a CGE-like identity, and in vivo complementation shows that \textit{Arx} and \textit{Cxcr7} rescue separable aspects of Lhx6 function \citep{Alifragis2004, Zhao2008}.

Several additional regulators affect VIP-relevant subsets. \textit{Dlx1}/\textit{Dlx2} and \textit{Dlx5}/\textit{Dlx6} control CGE neurite outgrowth and tangential migration, and selective \textit{Dlx1} loss causes subtype-specific reduction of CR+, SST+ and other cortical/hippocampal interneurons \citep{Cobos2007, Wang2010, Cobos2005}. CXCL12 attracts immature interneurons during tangential migration via CXCR4/CXCR7 streams in subpial and intermediate-zone routes, and BDNF/NT4 acting via TrkB stimulates tangential migration of MGE-derived cells, with a residual effect on CGE-associated calbindin+ populations. The net picture, reviewed by \citet{Lim2018} and \citet{Demarcogarcia2025}, is that cardinal-class identity (including VIP) is largely seeded by progenitor-stage TF combinations, while finer t-type granularity emerges through postmitotic refinement that integrates intrinsic TF cascades with activity-dependent signalling.

\section{Tangential migration and laminar settling}

CGE-derived precursors leave the eminence in three caudo-rostral streams whose temporal sequence and TF coding differ from the dorsoventral MGE streams. Initial studies \citep{Touzot2016} resolved a lateral migratory stream (LMS), a medial migratory stream (MMS) and the previously known caudal migratory stream (CMS), each expressing distinct combinations of SP8, PROX1, COUP-TFI and COUP-TFII; conditional COUP-TFI inactivation differentially perturbs LMS/MMS and produces laminar-specific deficits in adult cortex. \textit{Lopezbendito2004} used GAD65-GFP transgenic mice to track a CGE-enriched subpopulation along the lower intermediate zone from E14--E15 and into superficial layers thereafter, consistent with the late-born CGE programme. Foundational studies \citep{Bellion2005} characterised the cellular biomechanics of tangential migration as a two-phase nucleokinesis (centrosome--Golgi forward translocation followed by nuclear advance) --- a mode shared between MGE- and CGE-derived precursors. The chemokine CXCL12 acting via CXCR4/CXCR7 attracts cortical interneurons during these tangential streams, with loss of CXCR4 responsiveness producing dispersion deficits, and BDNF/TrkB signalling provides additional permissive drive that affects both MGE- and CGE-associated populations.

Laminar settling is initially comparable between MGE and CGE: at P1, mouse precursors born at E12.5 from either eminence show similar layer distributions, with subsequent radial sorting between deep (MGE) and superficial (CGE) layers driven by mechanisms that coincide with KCC2 upregulation \citep{Miyoshi2011}. By the end of the first postnatal week, CGE-derived cells have completed their characteristic preferential occupancy of layers II/III ({\textasciitilde}75\%) \citep{Miyoshi2010}, although in mouse this enrichment is graded rather than absolute and a substantial CGE-derived contingent settles in deeper layers and in layer 1 \citep{Lee2010}. \textit{Larimer2016} used heterochronic CGE transplants into neonatal mouse V1 to show that CGE-derived precursors disperse, acquire host-like laminar position, marker expression, intrinsic properties and visual response --- but, unlike MGE transplants, they do not reactivate ocular dominance plasticity, with the residual ODP induction observed in CGE transplants attributable to a small contaminating fraction of MGE-lineage cells.

A subtle but consequential observation is that excitatory-neuron identity instructs lineage-specific aspects of CGE settling. Conditional \textit{Satb2} knockout in mouse projection neurons reprograms intratelencephalic-type pyramidal neurons toward a pyramidal-tract identity and selectively disrupts the lamination and circuit integration of CGE-derived (but not MGE-derived) interneurons, providing causal evidence that pyramidal subtype identity instructs CGE-lineage interneuron development in a lineage-specific manner. Cux1 and Cux2 are themselves expressed cell-autonomously in postmitotic Reelin+ cortical interneurons --- a Reelin+ population that includes both MGE-derived neurogliaform and CGE-derived components --- and are required for the development of this Reelin-secreting IN subset \citep{Cubelos2008}, placing Cut-family homeodomain TF control inside the same superficial-layer compartment that CGE-class lateral and medial migratory streams populate \citep{Touzot2016}. The chemoattractant Cxcl14, induced by Sp8/Sp9, similarly couples CGE migration choices to layer-specific cues during the same window. This intercellular dependency complicates simple cell-autonomous accounts of VIP laminar settling and dovetails with the activity-dependent programmes considered next.

\section{Activity-dependent and intrinsic programmes for postnatal maturation}

Postnatal maturation of CGE-derived interneurons depends on a graded interplay between cell-intrinsic genetic programmes and activity-dependent extrinsic signals. The seminal cell-autonomous experiment is that of \citet{Demarcogarcia2011}, who showed by in utero Dlx5/6-eGFP electroporation at E15.5 followed by Kir2.1-mediated silencing that activity is essential before P3 for correct migration of CGE-derived reelin+ and calretinin+ (but not VIP+) interneurons; after P3, glutamate-mediated activity controls axon and dendrite development. The Dlx1 target \textit{Elmo1} mediates this activity-dependent migration in Re+/Cr+ subtypes and is downregulated by Kir2.1 expression \citep{Demarcogarcia2011}. Established work extended these findings using the same Kir2.1 paradigm, reporting that cell-autonomous activity reduction lowered maximum AP discharge from 93$\pm$10 Hz to 44$\pm$8 Hz and broadened AP half-width from 1.3$\pm$0.15 ms to 3.0$\pm$0.32 ms in CGE-derived cells (P12--P15 cortex), and produced subtype-specific laminar mispositioning in Re+ and Cr+ but not VIP+ cells \citep{Karayannis2012}. The convergence of these two studies established the disinhibitory view that CGE interneuron migration and intrinsic excitability mature through cell-autonomous activity, with VIP cells comparatively spared.

Glutamatergic drive is required for the survival and synaptic maturation of CGE-derived interneurons more broadly. Conditional removal of all AMPAR subunits (GluA1/2/3 cKO) in CGE-derived interneurons increases their early postnatal numbers but causes later apoptosis to converge with wild-type, almost completely eliminates sEPSCs, and reduces feedforward and feedback inhibition onto pyramidal neurons. The phenotype indicates that AMPAR-mediated excitatory drive is required not for CGE-IN birth but for circuit integration and survival during the postnatal critical period. Within VIP cells specifically, \citet{Simacek2025a} showed that inhibitory presynaptic inputs onto layer 2/3 VIP-INs in mouse barrel cortex undergo functional maturation between P9 and P15: both readily-releasable vesicle number and release probability rise to adult-like values across this window. The complementary excitatory side is captured by Earlier work \citep{Simacek2025b}, which described a sequential P3 \rightarrow P30+ pattern in mouse S1 barrel cortex in which mEPSC frequency rises before P8--10, intrinsic firing matures next, and IPSC kinetics consolidate last. The Prox1 / Elfn1 axis additionally specifies short-term excitatory input dynamics in a VIP-subtype-selective manner during the same window \citep{Stachniak2021}, and this maturation timeline is anchored to the broader CGE-class postnatal programme described in recent reviews \citep{Lim2018, Demarcogarcia2025}.

Two reviews integrate these strands. \citet{Lim2018} argued that cortical GABAergic interneuron diversity emerges from cell-intrinsic genetic programmes in subpallial progenitors that unfold over a protracted period and are subsequently modulated by activity-dependent, non-cell-autonomous mechanisms during circuit integration; \citet{Demarcogarcia2025} updated this synthesis with new Patch-seq and connectivity data, emphasising that transient developmental connectivity instructs the consolidation of mature subtypes. \citet{Yaeger2019} provided a circuit-level example in mouse V1, where basal-forebrain acetylcholine directly excites SST interneurons during the critical period to enable compartmentalised dendritic spiking; this cholinergic responsiveness is lost in adulthood, and suppressing SST cells during the critical period prevents normal development of binocular receptive fields. While VIP cells were not the primary focus, the experiment illustrates the principle that experience-dependent neuromodulator drive and CGE-class circuit function are functionally coupled during postnatal critical periods, with implications for VIP-driven disinhibitory motifs that emerge in the same window.

\section{Programmed cell death, network activity, and circuit integration}

Roughly a third of cortical interneurons die during the second postnatal week in a developmental wave of programmed cell death (PCD). Previous reports showed that all ventral-eminence-derived cortical interneurons --- including CGE-derived cells --- undergo PCD, but that VIP interneuron survival, in contrast to MGE-derived PV/SST classes, is not activity-dependent; calcineurin in interneurons sequentially controls morphological maturation (E15--P5) and subsequent activity-dependent processes \citep{Priya2018}. This decoupling matters: it means that cell-autonomous reductions in VIP excitability (as in \citet{Demarcogarcia2011}) do not preferentially cull VIP cells, whereas comparable manipulations bias the survival of CGE reelin+/calretinin+ and MGE-derived populations. \citet{Carriere2020} provided a parallel cell-autonomous mechanism: clustered $\gamma$-Protocadherins (Pcdhgs) act in mouse GABAergic interneurons to promote survival during the postnatal PCD critical period, with effects across both MGE- and CGE-derived classes.

Network-level activity feeds back onto interneuron survival. \citet{Duan2020} reported that in neonatal mouse somatosensory cortex (P7), functional interneuron and pyramidal-cell assemblies emerge spatially segregated; reducing GABA release or synaptic inputs onto pyramidal cells erodes this functional topography and increases interneuron survival, demonstrating that GABAergic restriction of early network dynamics actively regulates interneuron developmental death. Together with the AMPAR-cKO data \citep{Akgul2019}, the Pcdhg result \citep{Carriere2020}, and the activity-independence of VIP-cell survival reported by \citet{Priya2018}, these findings recast the second postnatal week as a coordinated period in which intrinsic survival modules and circuit-level activity are layered onto the earlier, lineage-specific transcriptional cascades --- with VIP cells representing a partial exception whose survival is set by intrinsic rather than activity-driven programmes \citep{Demarcogarcia2025}.

A handful of studies extend this developmental account to specific VIP-relevant circuits. In hippocampal CA1, a subset of VIP+ interneurons projects long-range axons to subiculum and is recruited preferentially during quiet wakefulness rather than during theta-run/locomotion, indicating that sublineage diversification of CA1 VIP cells produces functionally specialised long-range projection cells distinct from the disinhibitory VIP-IS3 population \citep{Francavilla2018a}. The hippocampal CGE neurogenic wave that produces these populations is itself temporally protracted from E12 to E16, generating not only IS3-type VIP cells but also nNOS+ NGC and Ivy populations from a shared CGE/MGE-mixed origin. In V1, \citet{Ibrahim2021} showed that bottom-up thalamic inputs to layer-1 NGCs are required postnatally to establish higher-order top-down connectivity onto these cells, sculpting circuit maturation in a CGE-class-relevant compartment. Cortically deployed VIP cells thus mature into multiple functionally distinct sublineages whose relative proportions vary by area and species --- a point that becomes acute when the analysis crosses to primates.

\begin{figure}[!htbp]
\centering
\includegraphics[width=0.88\linewidth]{files/fig-cge-lineage-08f0b78345bee619f0487d653ac95d63.png}
\caption[]{Schematic of the CGE-derived inhibitory lineage that produces VIP interneurons. \textbf{(A)} Embryonic CGE \rightarrow tangential migration via lateral, medial and caudal streams \rightarrow cortical layer settling. CGE-class TFs (Prox1, Sp8, Nr2f2/COUP-TFII) are contrasted with the MGE-class TF Lhx6 \citep{Miyoshi2010, Touzot2016, Kanatani2008, Liodis2007}. \textbf{(B)} Conditional-KO phenotypes for the principal CGE TFs and synaptic/excitability regulators, showing the subtype-selectivity of each perturbation \citep{Miyoshi2015, Stachniak2021, Cai2013, Touzot2016, Wei2019, Akgul2019, Karayannis2012, Demarcogarcia2011, Priya2018}. \textbf{(C)} Reported fraction of cortical GABAergic neurons attributable to the CGE/5-HT3AR/ADARB2 lineage in mouse fate-mapping ({\textasciitilde}30\%) versus human snRNA-seq ({\textasciitilde}50\%) \citep{Miyoshi2010, Hodge2019}. \textit{Caveat:} the cross-species comparison overlays a fate-map percentage (mouse) against a marker-defined snRNA-seq percentage (human MTG only), so the apparent expansion conflates measurement modality with biology and should be interpreted as a directional rather than absolute estimate.}
\label{fig-cge-lineage}
\end{figure}

\section{Cross-species divergence: rodent fate-mapping versus primate transcriptomics}

The transition from rodent to primate complicates almost every quantitative claim made above. \citet{Ma2013} argued from immunohistochemistry of Sox6, COUP-TFII and Sp8 in developing human and monkey telencephalon that the same three subpallial progenitor domains (MGE, LGE, CGE) found in rodents are present in primates and that tangentially migrating interneurons retain comparable marker codes --- supporting a conserved subcortical origin of primate neocortical interneurons \citep{Ma2013}. \citet{Marin2025} synthesised the broader comparative literature and concluded that interneuron diversity is largely conserved between rodents and primates, with primate cortical evolution contributing greater abundance and wider interconnectivity rather than wholly new types \citep{Marin2025}. Single-cell taxonomies are consistent: comparative scRNA-seq across macaques and mice indicates that prenatal IN classes are largely conserved, even if a primate-specific TAC3 striatal interneuron population emerges \citep{Schmitz2022}, and snRNA-seq of \textgreater 2.4 million cells across 18 marmoset locations shows that adult transcriptomic identity tracks developmental origin more strongly than neurotransmitter repertoire \citep{Krienen2023}.

Conserved cardinality, however, masks several quantitative and architectural divergences. Earlier studies \citep{Hansen2013} documented a massively expanded subventricular zone in second-trimester human ganglionic eminences that lacks the radial-fibre/oriented organisation seen in rodents and contains stem cells whose proliferation and migration patterns differ from the MGE programme --- implying human-specific aspects of CGE-lineage development. Prior work \citep{Reinchisi2012} confirmed that COUP-TFII expression peaks in human CGE during gestational weeks 9--22 with a temporo-occipital \textgreater  frontal/dorsal gradient. scRNA-seq of human ganglionic eminences identifies regional and temporal progenitor diversity and an evolutionarily conserved transcriptional grammar but also human-specific transcriptomic features, while Earlier reports \citep{Velmeshev2023} profiled \textgreater 700,000 nuclei across 106 prenatal/postnatal human cortex donors and prioritised interneuron lineages as key vulnerable trajectories. \citet{Eichmuller2022} further described a human-specific caudal late interneuron progenitor (CLIP) population that over-proliferates in tuberous sclerosis complex organoids --- a primate-relevant lineage component absent in mouse fate maps.

The most consequential primate divergence for VIP biology is in lineage marker codes. \citet{Hodge2019} reported that human MTG contains roughly equal proportions of MGE-derived (44\% LHX6+) and CGE-derived (50\% ADARB2+) cortical interneurons, that VIP is the most transcriptomically diverse interneuron subclass in human MTG (21 types), and --- crucially --- that the disinhibitory mouse CGE marker \textit{HTR3A} is \textit{not} expressed in human CGE-lineage interneuron types, marking a major species-specific divergence in the very molecular handle that defined the CGE lineage in mouse \citep{Hodge2019}. The same study identified human-specific features within the CGE class (e.g., some CGE-class layer-1 cells co-express SST, a feature not seen in mouse layer-1 interneurons) and located human ``rosehip'' cells transcriptomically within the CGE-class \textit{LAMP5} subdivision \citep{Hodge2019}. Initial studies \citep{Jorstad2023b} confirmed across eight human cortical areas that 24 cell subclasses are highly consistent in presence but show areal differences in inhibitory-neuron sub-state composition.

\begin{framed}
\textbf{Conflict --- Cross-species CGE-lineage proportion}\\
\citet{Miyoshi2010} reported by mouse fate-mapping that {\textasciitilde}30\% of cortical interneurons are CGE-derived, while \citet{Hodge2019} measured by snRNA-seq in human MTG that {\textasciitilde}50\% of cortical interneurons are \textit{ADARB2}+ (CGE-class). \citet{Du2025} --- using directed differentiation from human pluripotent stem cells --- generated a homogeneous calretinin+ interneuron population ({\textasciitilde}80\% CR+) consistent with a substantial CGE-like lineage, but as an in-vitro readout cannot adjudicate native fractions. Resolution: the two measurements differ in modality (genetic fate-mapping vs marker-defined snRNA-seq), reference area (whole mouse cortex vs human MTG only) and lineage definition (CGE proper vs ADARB2+ marker class that may include POA contribution); the directional claim that the CGE/5-HT3AR/ADARB2 class is expanded in primates is supported, but the precise quantitative gap (30\% \rightarrow 50\%) should be treated as an upper bound until matched-area, matched-method comparisons are published.
\end{framed}

\begin{framed}
\textbf{Conflict --- Origin of primate cortical interneurons}\\
\citet{Delgado2022} reported that retroviral barcoding of fetal human cortical progenitors yields clones containing both excitatory pyramidal neurons and inhibitory interneurons (including subsets molecularly resembling CGE-class cells), suggesting a dorsal cortical contribution to primate inhibitory neurogenesis that is absent or minor in rodent. \citet{Marin2025} argued in a comparative review that GE origin is fully conserved across rodent and primate and that primate-specific aspects are quantitative (abundance, expansion) rather than categorical (new sources or new types). Resolution: the question of whether the primate cortex contains substantial cortically-generated or otherwise non-GE-derived inhibitory neuron lineages is unresolved; a parsimonious reading is that rare dorsal contributions exist (consistent with \citet{Delgado2022}'s lineage tracing) without overturning the dominant subpallial origin documented across all comparative datasets.
\end{framed}

A further methodological note: the disinhibitory adult mouse interneuron taxonomy of \citet{Tasic2018} defines six GABAergic subclasses (Sst, Pvalb, Vip, Lamp5, Sncg, Serpinf1) plus two distinct types --- placing VIP as one of four CGE-class subclasses. Reported VIP fractions vary across studies and atlas integrations because \textit{Lamp5}, \textit{Sncg} and \textit{Serpinf1} may be pooled with VIP under a ``5-HT3AR class'' or split out depending on clustering resolution; reconciling proportions therefore requires that subclass boundaries be specified, and recent cross-species human/marmoset comparisons \citep{Hodge2019, Krienen2023} reinforce that subclass-resolution choices have first-order consequences for any cross-paper percentage. The \citet{Vanvelthoven2025} developmental atlas resolves these subclass boundaries at prenatal-to-postnatal resolution and provides a hierarchical anchor for cross-study reconciliation. Beyond cortex, \citet{Sorrells2019} showed that the human amygdala paralaminar nuclei develop adjacent to the CGE and contain immature DCX+/PSA-NCAM+ neurons that mostly express excitatory markers and transition to mature TBR1+/VGLUT2+ neurons during adolescence --- establishing a protracted, CGE-adjacent developmental programme in human amygdala distinct from the cortical CGE-IN lineage that other comparative studies emphasise.

\begin{framed}
\textbf{Conflict --- VIP fraction depends on clustering resolution}\\
The earlier 2016 mouse cortical taxonomy from the Allen Institute reported VIP+ cells at {\textasciitilde}9\% of cortical GABAergic neurons (a single subclass), while the higher-resolution 2018 taxonomy of \citet{Tasic2018} resolves the broader CGE-class into Vip, Lamp5, Sncg and Serpinf1 subclasses, raising the apparent ``VIP-class'' fraction toward {\textasciitilde}18\% when \textit{Lamp5} and \textit{Sncg} are pooled. The earlier 2016 mouse cortical taxonomy from the Allen Institute (preceding \citet{Tasic2018}) is not in the filtered citation set and is therefore flagged in \texttt{unmet\_citation\_needs.json}; on the available evidence, the resolution-dependence is real and warns that all ``\% VIP'' numbers in the literature must be interpreted with the underlying clustering resolution and subclass boundaries declared. \citet{Vanvelthoven2025} provides the most current developmental anchor for these boundaries.
\end{framed}

\begin{figure}[!htbp]
\centering
\includegraphics[width=0.88\linewidth]{files/fig-cge-fraction-cro-c248c902e7115341e11274d64085e815.png}
\caption[]{Schematic methods-only summary of the cross-species CGE-lineage comparison. \textbf{(A)} Reported fraction of cortical GABAergic neurons assigned to the CGE/5-HT3AR/ADARB2 class --- mouse fate-mapping (Mash1/Olig2-CreER) gives {\textasciitilde}30\% across whole cortex \citep{Miyoshi2010}; human snRNA-seq in MTG gives {\textasciitilde}50\% (LHX6+ MGE 44\% / ADARB2+ CGE 50\%) \citep{Hodge2019}. \textbf{(B)} The proposed partitioning across VIP / Lamp5 / Sncg / Serpinf1 subclasses across mouse, marmoset, macaque and human is shown as a methods placeholder only: no audited cross-paper panel was achievable for this stacked comparison because the underlying datasets use non-matched subclass definitions and reference areas \citep{Tasic2018, Hodge2019, Krienen2023, Schmitz2022}. \textbf{(C)} Side timeline of the CGE-progenitor TF cascade --- Sp8/Sp9 \rightarrow COUP-TFII \rightarrow PROX1 \rightarrow postmitotic Prox1 / Elfn1 / activity-dependent refinement --- with conditional-KO outcome annotations \citep{Wei2019, Kanatani2008, Demarcogarcia2011, Karayannis2012}. \textit{Caveat:} No quantitative cross-paper comparison was achievable for this comparison; the figure is methods-only/schematic, with the directional claim (mouse \textless  primate CGE-class fraction) supported but the precise gap dependent on measurement modality. Readers should treat percentage estimates as study-specific anchors rather than as a synthesised effect size.}
\label{fig-cge-fraction-cross-species}
\end{figure}

\section{CGE internal topography, organoid models, and recent integration}

The CGE is not internally homogeneous. Foundational studies \citep{Garg2025} showed that mature CGE-IN subtypes arise from distinct spatial and temporal subdomains of the CGE --- dorsal/ventral and rostral/caudal subregions and birthdate windows prefigure adult identities, mirroring the internal topography long appreciated in the MGE. Established work \citep{Bright2025} complemented this with heterochronic transplantation, lineage tracing and perturb-seq experiments showing that progenitor competence is reset across embryonic time across all GEs, so that progenitor-stage temporal coding interacts with eminence-of-origin to set cardinal identity. \citet{Bandler2017} had previously argued that cortical interneuron specification is governed by the joint contributions of spatial origin, birthdate, lineage relationships and mode of cell division --- a framework that the more recent CGE-resolved data \citep{Garg2025, Bright2025} now operationalise. The 2025 review of \citet{Horton2025} synthesises CGE-derived IN biology across mouse, human and non-human primate, providing the most current cross-species integration. \citet{Vanvelthoven2025} adds a prenatal--postnatal scRNA-seq dataset of \textgreater 1.2 million mouse telencephalon cells (E7--P14) that resolves the hierarchical GABAergic taxonomy at developmental resolution and anchors the prenatal-to-postnatal trajectory of CGE-class precursors.

Stem-cell and organoid systems are increasingly used to model CGE-lineage development, with mixed fidelity. \citet{Bershteyn2023} defined the directed-differentiation cascade (NKX2-1+/SOX1+/FOXG1+/DLX1/2+ \rightarrow MGE progenitors \rightarrow mature interneurons) for pallial MGE-type cells from human ESCs; analogous CGE protocols, exemplified by Earlier work \citep{Du2025}, generate {\textasciitilde}80\% calretinin+ cells via SHH activation plus IWP2-mediated WNT inhibition that show GABAergic action potentials in vitro and integrate after transplantation. Previous reports \citep{Andrews2023} showed that dorsal forebrain organoids initially generate LGE-like INs at week 8 and progressively replace them with CGE-like INs by week 15, with LIF/LIFR signalling promoting outer radial glia identity in primary tissue. \citet{Uzquiano2022} profiled human cortical organoids at single-cell, epigenetic and spatial levels and reported that endogenous cellular diversification trajectories are recapitulated regardless of metabolic state. \citet{Eichmuller2022}'s identification of a human-specific CLIP cell population in TSC organoids underscores that organoid systems are now uncovering primate-specific CGE-class progenitor states absent in mouse. The translational corollary is that mouse fate-mapping should anchor the framework but human-specific organoid and snRNA-seq data must be consulted for any quantitative claim about VIP-class proportions or marker codes in primate cortex.

A specific point about cortical NGCs and amygdala / ventral telencephalon connections is worth preserving here. \citet{Tang2012} showed that COUP-TFII is required in the CGE to specify amygdala (BMA) patterning by directly activating semaphorin receptors Nrp1 and Nrp2; in conditional COUP-TFII mutants, CGE-derived migration to amygdala is disrupted, indicating that CGE-class TFs serve patterning functions beyond cortex. \citet{Mercier2022} complements the cortical story by showing that hippocampal stratum-lacunosum-moleculare interneurons (including NGCs) exhibit NMDAR-dependent Hebbian LTP, with plasticity in NGC inputs modulating temporoammonic excitation-inhibition balance --- providing a circuit-level rationale for the developmental investment in CGE-class hippocampal NGCs documented by \citet{Tricoire2010} and \citet{Tricoire2011}. \citet{Kim2025a} reported that neonatal gyrencephalic brains (human, non-human primate, piglet) harbour an expanded curved subventricular zone (`Arc') of migratory neuroblasts absent in lissencephalic mouse and marmoset, adding yet another lineage-relevant primate-specific feature that mouse fate-mapping cannot capture.

A short note on hippocampal CA1 VIP development. The hippocampal CGE wave (E12--E16) generates the disinhibitory IS3-type VIP/CR/CCK CA1 interneuron alongside CCK+ basket and NGC populations \citep{Tricoire2011, Fuentealba2010}. CA1 VIP cells diversify postnatally into at least two distinct subpopulations: the local IS3-type disinhibitory cells described in early reviews \citep{Rudy2011, Kepecs2014} and the long-range VIP-LRP cells with subicular projections that preferentially fire during quiet wakefulness rather than during theta-run/locomotion \citep{Francavilla2018a}. The latter population breaks the disinhibitory assumption that all CA1 VIP cells implement the same disinhibitory motif and indicates that within-VIP sublineage diversification --- likely set by the same Prox1/Sp8/COUP-TFII cascade documented for cortex but acting on hippocampus-specific targets --- produces functionally distinct outputs that (\href{/disinhibition-framework}{Local Circuit Motifs and the Disinhibition Framework}) will revisit.

\section{Synthesis: what the development data buy the rest of the review}

The developmental literature delivers four claims that the rest of this review will lean on. (1) VIP cortical interneurons are CGE-derived members of a 5-HT3AR/ADARB2-defined sister set (with \textit{Lamp5}, \textit{Sncg}, \textit{Serpinf1}) that is reproducible across modalities and species at the cardinal-class level \citep{Lee2010, Miyoshi2010, Tasic2018, Hodge2019}. (2) Their identity is specified by a small, conditional-KO-validated TF cascade (Sp8/Sp9 \rightarrow COUP-TFII \rightarrow PROX1 \rightarrow postmitotic PROX1 / \textit{Elfn1}) whose perturbation produces subtype-selective cellular and synaptic phenotypes traceable into adult VIP-IN function \citep{Wei2019, Kanatani2008, Miyoshi2015, Stachniak2021}. (3) Their postnatal maturation combines a cell-intrinsic genetic programme with activity-dependent and circuit-level inputs, with VIP cells comparatively spared from the activity-dependent migration and survival mechanisms that affect Re+/Cr+ subtypes \citep{Demarcogarcia2011, Karayannis2012, Priya2018, Akgul2019, Lim2018, Demarcogarcia2025}. (4) The translation of mouse fate-mapping to human VIP biology is bounded: lineage marker codes diverge (loss of \textit{HTR3A} as a CGE handle in human; expansion of CGE-class proportion to {\textasciitilde}50\% in MTG; appearance of human-specific CLIP and rosehip-type lineages), and quantitative cross-species claims should be made with explicit modality and area caveats \citep{Hodge2019, Hansen2013, Reinchisi2012, Eichmuller2022, Marin2025}. (\href{/morphology}{Morphological Diversity}) takes these molecular and developmental categories and asks how they project into morphological space; (\href{/disease-translational}{Species Differences, Human Relevance, and Disease}) returns to the cross-species question for disease-relevant VIP phenotypes; and the synthesis in \href{/disinhibition-framework}{Local Circuit Motifs and the Disinhibition Framework} returns to the resolution-dependence of mouse VIP-fraction estimates flagged above.


\section{Morphological Diversity}
\label{sec:04_morphology}

Sections \href{/classification}{Molecular Identity and Transcriptomic Taxonomy} and \href{/development}{Developmental Origins and Postnatal Maturation} established the molecular and developmental scaffold of the VIP-expressing class; this section asks how those categories project into morphological space. The contemporary picture, drawn from Patch-seq atlases, large-volume connectomics, and decades of biocytin-based reconstructions, is consistent on the gross descriptors but contested on the boundaries: cortical VIP cells span a recognisable but graded morphological space dominated by bipolar/bitufted somatodendritic profiles with descending axonal arbors \citep{Apicella2022}, with smaller proportions of small-basket and arcade-like variants \citep{Rudy2011, Yayon2023}. Subclass-level (VIP vs SST vs PV) morphological separation is robust in cross-modal taxonomies \citep{Gouwens2019, Gouwens2020a, Gliko2024, Elabbady2025}, while within-VIP clustering forms a continuum punctuated by reproducible centroids whose proportions vary across cortical area, layer, and species \citep{Gouwens2020a, Gouwens2020b, Pronneke2015, Chartrand2023}.

\section{Canonical morphological types}

The earliest single-cell descriptions of cortical VIP-expressing interneurons recognised a small set of recurring forms. \citet{Kawaguchi1996} described four morphological classes in rat frontal cortex --- vertically bipolar cells with predominantly descending axonal collaterals, double-bouquet cells, small-basket cells with short and dense local arbors, and arcade-shaped cells with arching axonal trees --- and showed that VIP-immunoreactive axons terminated in symmetrical (inhibitory) synapses with a strong vertical bias \citep{Kawaguchi1996, Cauli1997, Cauli2014, Apicella2022}. \citet{Cauli1997} extended the framework by combining whole-cell recording with single-cell RT-PCR in rat sensorimotor cortex, identifying a vertically oriented bipolar subset that co-expressed calretinin and VIP and fired in an irregular-spiking pattern \citep{Cauli1997, Cauli2014, Vucurovic2010, Stachniak2021}. \citet{Vucurovic2010} then demonstrated, using 5-HT3aR-BAC transgenics combined with in situ hybridisation, that the cortical 5-HT3AR-positive lineage partitions into a VIP-positive subset with bipolar/bitufted morphology and a VIP-negative subset with multipolar morphology --- embedding the bipolar/bitufted descriptor in the broader CGE taxonomy \citep{Vucurovic2010, Rudy2011, Cauli2014, Fishell2020}.

The bipolar/bitufted morphotype dominates in numerical surveys but does not exhaust the class. Reviews of the rodent literature converge on a list of bipolar, bitufted, multipolar, small-basket, and arcade morphologies, with VIP-positive basket cells representing a numerical minority of the class \citep{Rudy2011, Cauli2014, Apicella2022, Fishell2020}. An earlier report, using VIP-Cre $\times$ tdTomato barrel cortex slices, reported a substantially broader repertoire than the classical bipolar form --- including bipolar/bitufted, tufted, and multipolar somatodendritic configurations distributed differently across layers --- and noted that the major axon often emerged not from the soma but from a basal dendrite directed toward the white matter, a feature also flagged by an earlier report \citep{Kawaguchi1996}. Two intersectional studies further partitioned the class along molecular axes: \citet{Stachniak2021} separated CR-positive bipolar and CCK-positive multipolar VIP subtypes during cortical development, and \citet{Yayon2023} resolved the cholinergic subset (VChIs) of VIP barrel-cortex cells into bipolar (two stem dendrites) and multipolar ($\geq$3 stem dendrites) morphologies \citep{Stachniak2021, Yayon2023, Wong2022, Pronneke2015}.

\begin{figure}[!htbp]
\centering
\includegraphics[width=0.95\linewidth]{files/fig-vip-morphologies-8fdebf4f0d24a4df79006728c04238c6.png}
\caption[]{\textbf{Morphological diversity of cortical VIP interneurons.} \textbf{(A)} Idealised exemplar morphologies --- bipolar, bitufted, small-basket, and arcade --- drawn to illustrate the canonical VIP morphological repertoire described by earlier studies and \citet{Cauli2014}. \textbf{(B)} Textual side-by-side of predominant VIP morphological subtype reported in two genetically defined VIP-Cre datasets in mouse primary sensory neocortex: a preceding study (S1 barrel) versus \citet{Gouwens2020a} (VISp). \textit{Caveat (verbatim, Phase 6 audit):} Render as textual side-by-side, not as a quantitative bar (per restructure\_notes). | State explicitly that Pronneke reports laminar dominance in S1 barrel cortex while Gouwens reports subtype enumeration across Patch-seq met-types in VISp --- entries are not on the same axis. | Note difference in cortical area (S1 vs VISp) and method (fixed-slice morphology vs Patch-seq morpho-electric). \textbf{(C)} Subcellular target distribution of VIP-axon boutons in mouse S1 barrel L2/3 (immuno-EM, n=31 synapses) from an initial investigation. \textit{Caveat (verbatim, Phase 6 audit):} Specify the sample size (n=31 synapses, L2/3) directly in the panel. | Verify in the published Methods/Results that 80/13/7 are L2/3-specific values (not the cross-laminar overall reported in abstract). If unverifiable, label the row as `all layers (Zhou 2017)' rather than L2/3. \textbf{(D)} Within-paper layer comparison of VIP-bouton perisomatic-target percentage (L6 vs L2/3+L5 pooled) from \citet{Zhou2017}. \textit{Caveat (verbatim, Phase 6 audit):} State that per-layer n\_synapses denominators were not extractable; bars show point estimates without error. | Indicate that one row is a single layer (L6) while the other is pooled across L2/3 and L5. | Source: Zhou 2017 (10.1093/cercor/bhx220), serial-section EM in mouse S1 barrel cortex. \textbf{(E)} Laminar percent-of-VIP+ profile in mouse frontal cortex from \citet{Xu2010a}. \textit{Caveat (verbatim, Phase 6 audit):} Single-study `comparison' --- splitter yields a single laminar profile; no cross-paper comparison possible. | No value\_source\_sentence quoted for the layer percentages (verifiability gap). | No per-layer n provided; sample\_size flagged in restructure\_notes. \textbf{(F)} Laminar density (cells/mm\textsuperscript{2}) of VIP+ neurons in mouse frontal cortex from \citet{Xu2010a}. \textit{Caveat (verbatim, Phase 6 audit):} State that n\_definition (animals/sections) and dispersion are not reported in the source extraction; values shown are point estimates only. | Indicate the source is Xu 2010 (10.1002/cne.22229), immunohistochemistry/cell-counting in mouse frontal cortex. | If the paper's Tables report frontal cortex separately, cite the specific table; otherwise note the regional attribution caveat.}
\label{fig-vip-morphologies}
\end{figure}

\section{Subtype proportions and laminar distribution}

Numerical estimates of how the canonical morphologies partition the VIP class differ substantially across studies, reflecting genuine biology and methodological choice. \citet{Gonchar2008} reported that VIP-expressing GABAergic neurons account for roughly 11\% of cortical interneurons in mouse V1 and are concentrated in supragranular layers, with CR-, VIP-, ChAT- and CCK-immunoreactive populations enriched in L2/3 \citep{Gonchar2008, Xu2010a, Rudy2011, Cauli2014}. \citet{Xu2010a} quantified absolute laminar densities across mouse frontal cortex, somatosensory cortex, and V1, finding L2/3 VIP-positive density 2--3$\times$ greater than deeper layers (e.g. {\textasciitilde}62 cells/mm\textsuperscript{2} in L2/3 versus {\textasciitilde}22--30 in other layers in frontal cortex, with 32.6\% of all VIP-positive cells in L2/3) \citep{Xu2010a, Gonchar2008, Pronneke2015, Apicella2022}. The Allen Institute Patch-seq atlas \citep{Gouwens2020a, Gouwens2019, Gliko2024, Gouwens2020b} partitioned the VISp Vip subclass into five morpho-electric-transcriptomic types (Vip-MET-1 through Vip-MET-5) all resembling bipolar or bitufted cells but differing in laminar axon innervation, with Vip-MET-5 (Vip Crispld2) the principal L1-projecting type and Vip-MET-1/-2 occupying L2/3--L4 with distinct axonal biases.

The proportions disagree where the population definitions and methods diverge. The classical estimate that approximately 40\% of cortical VIP interneurons co-express calretinin \citep{Rudy2011, Cauli2014, Cauli1997, Demarcogarcia2011} was challenged by region-specific counts in mouse V1 layers 4 and 5 reporting {\textasciitilde}52\% VIP/CR overlap \citep{Gonchar2008, Rudy2011, Cauli2014}.

\begin{framed}
\textbf{Conflict --- VIP/CR co-expression fraction}\\
\citet{Rudy2011} summarised the rodent literature as reporting approximately 40\% of cortical VIP-positive interneurons co-expressing calretinin \citep{Rudy2011}, while an earlier analysis reported a significantly higher VIP/CR overlap ({\textasciitilde}52\%) restricted to layers 4 and 5 of mouse V1. \textbf{Resolution:} The two estimates are likely not contradictory: co-expression depends on cortical region, laminar position, and developmental stage, and both reports remain valid for their respective contexts \citep{Rudy2011, Demarcogarcia2011}.
\end{framed}

Other proportional estimates extend rather than overturn this picture. \citet{Demarcogarcia2011} separated VIP-positive multipolar cells from CR-positive bipolar/multipolar siblings in early postnatal cortex, with the two showing distinct activity-dependence of axonal arbor length under in utero electroporation \citep{Demarcogarcia2011, Stachniak2021, Wong2022, Fishell2020}. \citet{Wong2022} reported, using immunohistochemistry of CGE-derived Htr3a-positive cells, that VIP is expressed in two morphologically distinct subtypes --- most CR-positive bipolar cells and approximately half of CR-/Reln-negative basket cells --- providing a molecular link between the bipolar majority and the smaller VIP-basket fraction \citep{Wong2022, Stachniak2021, Rudy2011, Yayon2023}. \citet{Emmenegger2018} confirmed that intra-columnar VIP-like cells with bipolar dendrites form a distinct minority of L4 non-fast-spiking types in rat barrel cortex \citep{Emmenegger2018, Pronneke2015, Cauli1997, Apicella2022}.

\section{Discrete subtypes versus a graded continuum}

Whether within-VIP variation is best modelled as discrete cell types or as continuous gradients is a recurring methodological question. The Allen Institute Patch-seq programme has produced the most extensive substrate for this question. \citet{Gouwens2019} first reported that Patch-clamp/biocytin morphology in adult mouse VISp supports a taxonomy of {\textasciitilde}17 e-types, 38 m-types and 46 me-types, with axonal (rather than dendritic) features the dominant discriminator and good correspondence to transcriptomic subclasses including VIP \citep{Gouwens2019, Gouwens2020a, Gliko2024, Schneidermizell2025}. \citet{Gouwens2020a} extended this to morpho-electric-transcriptomic met-types and resolved five Vip-MET clusters, while a companion preprint \citep{Gouwens2020b, Gouwens2019, Gliko2024, Pronneke2015} cautioned that \textgreater 3,700 mouse VISp Patch-seq interneurons (350 reconstructed) reveal substantial continuous morpho-electric variation between and within met-types. \citet{Gliko2024} then enlarged the reconstructed inhibitory dataset to 813 cells and showed that arbor-density representations alone classify subclass-level identity (Sst/Pvalb/Vip/Lamp5) with high accuracy while within-VIP boundaries remain less crisp \citep{Gliko2024, Gouwens2020a, Gouwens2020b, Elabbady2025}.

\begin{framed}
\textbf{Conflict --- discrete VIP MET-types vs continuous variation}\\
The peer-reviewed report by \citet{Gouwens2020a} defines five Vip-MET-types in mouse VISp Patch-seq with bipolar/bitufted somatodendritic profiles separated by laminar axon innervation \citep{Gouwens2020a}, while a companion preprint \citep{Gouwens2020b} emphasises continuous morpho-electric variability within and between Vip MET-types that complicates discrete subtype boundaries. \textbf{Resolution:} Both views are compatible within current evidence: discrete MET centroids are recoverable, but with substantial continuous variation in features such as L1 axon density, dendritic stem number, and intrinsic firing parameters \citep{Gouwens2020a, Gouwens2020b, Gliko2024, Chartrand2023}. We treat within-VIP morphological structure as centroid-with-continuum rather than as a partition.
\end{framed}

Connectomic reconstructions converge on the same picture from a complementary direction. \citet{Schneidermizell2025} identified, in the MICrONS cubic-millimetre EM volume of mouse VISp, an ``inhibition-targeting'' cell class (InhTCs) that aligns with the disinhibitory VIP cardinal class and that subdivides into two connectomic subtypes --- InhTC\_dist (targeting distal-dendrite SST-like cells) and InhTC\_peri (targeting perisomatic PV-like basket cells) --- refining the long-standing view of a single VIP\rightarrowSST\rightarrowPyr motif \citep{Schneidermizell2025, Lee2013, Gamlin2025, Cauli2014}. \citet{Elabbady2025} showed that perisomatic ultrastructural features alone classify mouse-cortex inhibitory neurons into bipolar (VIP-overlapping), basket, Martinotti, and neurogliaform classes; in the same volume, 6,805 expert-classified inhibitory cells included 997 bipolar cells ({\textasciitilde}14.6\%) operationally defined by 2--3 primary dendrites and synapses primarily onto other inhibitory neurons \citep{Elabbady2025, Schneidermizell2025, Gouwens2020a, Gliko2024}. \citet{Gamlin2025} used Patch-seq--EM joint reconstructions to confirm that Vip cells exhibit a stereotyped narrow-horizontal/long-vertical axon arbor, with 20/22 EM-reconstructed Vip cells ($\approx$91\%) preferentially targeting inhibitory cells \citep{Gamlin2025, Schneidermizell2025, Elabbady2025, Lee2013}.

\section{Axonal architecture and target compartment selection}

The bipolar/bitufted somatodendritic profile is matched on the axonal side by a vertically biased, mostly inhibition-targeting arbor. \citet{Kawaguchi1996}'s original report that VIP-immunoreactive axons terminated in symmetrical synapses with a vertical (radially oriented) bias has been recovered repeatedly in transgenic and viral preparations \citep{Kawaguchi1996, Cauli2014, Apicella2022, Pronneke2015}. The most quantitative compartment-level data come from \citet{Zhou2017}, who used immuno-electron microscopy in mouse barrel cortex to show that VIP boutons innervate primarily dendrites ({\textasciitilde}80\% of synapses, n=31 in L2/3), with somata (13\%) and spines (7\%) sparsely targeted in upper layers \citep{Zhou2017, Lee2013, Gulyas1996, Apicella2022}. The same study reported a striking layer-dependent shift: in L6, {\textasciitilde}40\% of VIP boutons synapsed onto somata, contrasting with 8\% in L2/3 and 8\% in L5 --- indicating that the canonical ``perisomatic basket cell'' function is not absent from the VIP class but compartmentalised into deep layers, and aligning with the layer-specific axon innervation that defines Patch-seq met-types \citep{Gouwens2020a, Schneidermizell2025, Gamlin2025}.

A second, more contested feature concerns the origin of the major axon. A preceding study reported that the majority of VIP barrel-cortex cells emit their main axon not from the soma but from a basal dendrite directed toward the white matter --- a feature shared across the dataset and corroborated by an early report and consistent with reconstructions reviewed by a previous study \citep{Kawaguchi1996}. Outside primary sensory neocortex, the picture is less clear: in medial entorhinal cortex (MEC), \citet{Badrinarayanan2021} recorded 60 VIP cells (17 reconstructed) and reported bitufted or multipolar morphologies with layer-specific axonal asymmetry --- superficial cells extended axons into deep layers, deep cells made almost no contact with superficial layers --- but did not specifically report the dendrite-emerging-axon majority described in neocortex \citep{Badrinarayanan2021, Apicella2022}.

\begin{framed}
\textbf{Conflict --- dendrite-emerging axon: regional generality}\\
\citet{Pronneke2015} reported that the majority of VIP cells in mouse S1 barrel cortex emit their primary axon from a basal dendrite directed toward white matter, a feature shared across the imaged sample \citep{Pronneke2015}. In mouse medial entorhinal cortex, \citet{Badrinarayanan2021} documented variable axonal emergence and laminar projection asymmetry without confirming a dominant dendrite-emerging-axon configuration \citep{Badrinarayanan2021}. \textbf{Resolution:} Unresolved. The disagreement is consistent with an MEC-versus-neocortex difference in axon initial-segment placement that has not been explicitly tested across regions; an explicit comparison in matched preparations is required \citep{Badrinarayanan2021}.
\end{framed}

The interneuron-targeting bias is itself a defining structural property. \citet{Gulyas1996} first showed in rat hippocampus that calretinin-immunoreactive interneurons (homologous to cortical CR/VIP interneuron-selective cells) preferentially target other GABAergic interneurons --- including VIP-containing basket cells --- while sparing pyramidal neurons \citep{Gulyas1996, Lee2013, Schneidermizell2025, Gamlin2025}. \citet{Lee2013} extended this to neocortex with channelrhodopsin-assisted circuit mapping in mouse barrel cortex, showing that VIP cells preferentially inhibit SST-expressing interneurons during whisking \citep{Lee2013, Gulyas1996, Schneidermizell2025, Pronneke2015}. The MICrONS connectome converts this into an EM-grounded statistic: across the volume, VIP-aligned bipolar/InhTC cells make most of their synapses onto inhibitory targets, with InhTC\_dist preferentially contacting SST-like distal-dendrite cells and InhTC\_peri preferentially contacting PV-like basket cells \citep{Schneidermizell2025, Elabbady2025, Gamlin2025, Lee2013}.

\section{Patch-seq met-types and the morpho-electric structure within VIP}

Resolution within the VIP class has tightened around the five Patch-seq met-types defined in mouse VISp. \citet{Gouwens2020a} reported that all five Vip-MET-types resemble bipolar or bitufted cells but differ in laminar axon innervation: Vip-MET-5 (containing Vip Crispld2 t-types) is the principal L1-projecting VIP met-type, innervates all cortical layers including L1, and shows elevated hyperpolarisation-evoked sag relative to the other Vip met-types \citep{Gouwens2020a, Gouwens2019, Chartrand2023}. Vip-MET-1 occupies L2/3--L4 with bitufted dendrites extending into L1 and a concentrated axon plexus in lower L2/3, while Vip-MET-2 shares the laminar position of Vip-MET-1 but exhibits a deep-biased and broader axon distribution --- making the laminar axon profile the defining within-VIP discriminator rather than dendritic geometry alone \citep{Gouwens2020a, Gouwens2019, Schneidermizell2025, Gamlin2025}. \citet{Gouwens2019} had earlier shown across the whole inhibitory inventory that the axon, not the dendrite, is the dominant feature distinguishing inhibitory m-types in mouse visual cortex; the same paper noted that L2/3 Chat-Cre-labelled neurons (a known subset of Vip cells) map predominantly to a small number of related me-types associated with the Vip-Ptprt-Pkp2 transcriptomic type \citep{Gouwens2019, Gouwens2020a, Yayon2023, Wong2022}.

The Patch-seq view is now flanked by a dense-EM view from the same volume. \citet{Schneidermizell2025} mapped InhTCs and showed that the InhTC\_dist--InhTC\_peri partition recovers within-VIP heterogeneity from connectivity alone, while \citet{Elabbady2025} extended an EM operational definition of ``bipolar cells'' --- 2--3 primary dendrites and synapses primarily onto inhibitory targets --- to 997 cells in the mouse VISp column ({\textasciitilde}14.6\% of 6,805 inhibitory cells) \citep{Schneidermizell2025, Elabbady2025, Gouwens2020a, Gamlin2025}. \citet{Gamlin2025} joined the two views by training a Patch-seq morphology classifier and applying it to the EM volume, achieving $\approx$91\% accuracy in identifying Vip cells from their narrow-horizontal/long-vertical axon arbor and confirming that $\approx$91\% (20/22) of EM-reconstructed Vip cells preferentially target inhibitory cells --- converting the structural rationale for disinhibition into an EM-grounded statistic \citep{Gamlin2025, Schneidermizell2025, Elabbady2025, Lee2013}.

\section{Developmental shaping of VIP morphology}

Morphological diversification within the VIP class is established postnatally and is partially shaped by activity. \citet{Demarcogarcia2011} reported that VIP-positive multipolar cells, in contrast to CR-positive multipolar and bipolar siblings, do not reduce their axonal arbor length under activity-suppression in early postnatal cortex --- making axonal length development a class-specific rather than CGE-generic trait \citep{Demarcogarcia2011, Stachniak2021, Wong2022, Fishell2020}. An initial investigation further showed in mouse cortex that VIP cells split during development into a CR-positive bipolar subtype and a CCK-positive multipolar subtype with distinct dendritic geometries, and that genetic perturbations differentially affect the two arms \citep{Demarcogarcia2011, Yayon2023}. \citet{Yayon2023} extended this to the cholinergic VIP subset and showed that mild unilateral whisker deprivation at P7 alters dendritic arborisation and cortical-depth distribution of barrel-cortex VChIs ipsilaterally and contralaterally three weeks later --- providing direct evidence that within-VIP morphological diversity is sensitive to early-life sensory experience \citep{Yayon2023, Demarcogarcia2011, Stachniak2021, Pronneke2015}. \citet{Fishell2020} ties these observations together with a transcription-factor framework: embryonic loss of Prox1 abolishes calretinin expression in the bipolar VIP arm and truncates dendrites, establishing Prox1 as a central developmental regulator of the bipolar VIP morphology \citep{Fishell2020, Stachniak2021, Demarcogarcia2011, Wong2022}.

\section{Cross-area and cross-species variation}

Within rodent neocortex, area-level proportions of VIP morphologies are quantitatively heterogeneous. \citet{Apicella2022} summarised the rodent literature as showing supragranular VIP cells dominated by bipolar dendrites confined to L1--L2/3 and deep-layer (L4--L6) VIP cells diversifying into multipolar and basket-like forms \citep{Apicella2022, Pronneke2015, Zhou2017, Cauli2014}. \citet{Moreno2026} reported in mouse mPFC that L1-border VIP cells split into a bipolar form with vertical dendrites traversing L1 and a multipolar form with horizontal arbors, and that no VIP somata occupy the L1a sub-band \citep{Moreno2026, Apicella2022, Pronneke2015, Cauli2014}. \citet{Yayon2023} showed, using cleared-brain whole-volume reconstruction in mouse barrel cortex, that the cholinergic VIP subset (VChIs) partitions into bipolar (biVChI; two stem dendrites) and multipolar (mVChI; $\geq$3 stem dendrites) types and that mild whisker deprivation at P7 perturbs both ipsilaterally and contralaterally three weeks later --- coupling intrinsic morphological variability to early-life experience \citep{Yayon2023, Pronneke2015, Stachniak2021, Apicella2022}. \citet{Georgiou2022} independently reported a superficial-cortex VIP multipolar subpopulation with high dendritic spine densities on L1 dendrites and bursting electrophysiology, and confirmed by correlative light/electron microscopy that those L1-dendritic protrusions are bona fide excitatory synaptic spines morphologically distinct from neighbouring L1 spines \citep{Georgiou2022, Apicella2022, Pronneke2015, Demarcogarcia2011}.

Cross-species comparison surfaces both convergent features and quantitative divergence. \citet{Hodge2019} reported that in adult human middle temporal gyrus the VIP subclass is the most diverse among inhibitory subclasses, partitioning into 21 transcriptomic types most of which are enriched in upper cortical layers \citep{Hodge2019, Boldog2018, Lee2023a, Chartrand2023}. \citet{Boldog2018} described a novel rosehip interneuron in human L1 --- with compact bushy axonal arborisations and unusually large boutons confined to L1 --- at {\textasciitilde}13\% (10/76) of L1 interneurons sampled, providing a human-specific morphology with no clear mouse counterpart and adjacent to the VIP/Lamp5 family \citep{Boldog2018, Hodge2019, Chartrand2023, Lee2023a}. \citet{Lee2023a} showed in human cortex that VIP-subclass neurons (dominated by L2/L3) display diverse bipolar-dendrite/descending-axon morphologies with multimodal narrow-versus-wide axon-width profiles \citep{Lee2023a, Hodge2019, Chartrand2023, Pronneke2015}. An earlier analysis reported by quantitative immunolabelling and 3D reconstruction that calretinin-positive interneurons --- the rodent VIP-class analogue in primates --- are markedly denser and more abundant in rhesus monkey V1 and frontal cortex than in mouse, while their somatodendritic morphologies and axonal targets remain comparable across species \citep{Hodge2019, Cauli2014, Cauli1997}.

The most direct cross-species morphoelectric comparison comes from a single human--mouse Patch-seq dataset. \citet{Chartrand2023} reported that L1 VIP interneurons in both human and mouse exhibit characteristic descending axon collaterals, but that mouse L1-VIP cells show greater morphological variability than the more uniform human L1-VIP TSPAN12 cells; the human L1 VIP TSPAN12 cells, in particular, display stellate-like dendrites combined with descending axon collaterals and high voltage sag --- a morphoelectric profile distinct from the mouse counterpart \citep{Chartrand2023, Boldog2018, Hodge2019, Lee2023a}.

Layer-1 morphologies in primate and human cortex add a second cross-species axis to the comparison. \citet{Boldog2018} described, by Patch-seq and snRNA-seq, a novel rosehip interneuron in human L1 with compact bushy axonal arborisations, unusually large axonal boutons, and L1-confined arborisation, accounting for {\textasciitilde}13\% (10/76) of L1 interneurons sampled and lacking a clear mouse counterpart \citep{Boldog2018, Hodge2019, Chartrand2023, Lee2023a}. The rosehip is adjacent to but distinct from the VIP/Lamp5 family; its existence underscores that human L1 contains morphological types that are not fully predicted by the rodent atlas. \citet{Hodge2019} placed this in a broader transcriptomic frame: several human VIP transcriptomic types are closely related to LAMP5/PAX6 types and localise to layers 1--2, indicating overlap between superficial-layer disinhibitory and L1-targeting cell families that has methodological implications for any rodent-to-human translation \citep{Hodge2019, Boldog2018, Chartrand2023, Lee2023a}. The L1-VIP morpho-electric profile reported in human by \citet{Chartrand2023} --- stellate-like dendrites combined with descending axon collaterals and high voltage sag in TSPAN12 cells --- should therefore be read against the rosehip and LAMP5/PAX6 backdrop rather than as a pure VIP-class statement \citep{Chartrand2023, Boldog2018, Hodge2019, Lee2023a}.

\begin{framed}
\textbf{Conflict --- cross-species L1 VIP morphology}\\
The same Patch-seq dataset \citep{Chartrand2023} documents both a convergent feature (descending axon collaterals in L1 VIP cells of both species) and a quantitative divergence: human L1-VIP TSPAN12 cells show stellate-like dendrites and high voltage sag, while mouse L1 VIP cells share the descending axon but display greater morphological variability and lower sag \citep{Chartrand2023}. \textbf{Resolution:} The two observations are not contradictory but co-published as complementary cross-species findings within the same study; we treat the descending axon as a conserved morphological feature of the L1 VIP subclass and the dendritic and biophysical parameters as cross-species variables \citep{Chartrand2023, Hodge2019, Boldog2018, Lee2023a}.
\end{framed}

Across these comparisons two cross-cutting themes recur. The CGE-lineage scaffold (Section \href{/development}{Developmental Origins and Postnatal Maturation}) supplies a shared molecular backbone --- VIP-positive cells are members of the 5-HT3AR-positive lineage that also contains Lamp5 and Sncg subclasses --- onto which area- and species-specific morphological elaborations are layered \citep{Vucurovic2010, Fishell2020, Wong2022, Hodge2019}. \citet{Fishell2020} argued that Prox1 is required for the bipolar VIP fate (Prox1 loss abolishes calretinin expression and truncates dendrites), tying the dominant morphology to a specific transcription-factor programme \citep{Fishell2020, Stachniak2021, Wong2022, Demarcogarcia2011}. The discrete-versus-continuum tension is the second theme: it surfaces transcriptomically (Section \href{/classification}{Molecular Identity and Transcriptomic Taxonomy}), morphologically (this section), and electrophysiologically (Section \href{/electrophysiology}{Intrinsic Electrophysiology}), and the most defensible synthesis is that reproducible centroids exist within a continuous feature space whose dimensions include laminar axon innervation, dendritic stem number, axonal target compartment, and cross-species shifts in dendritic geometry \citep{Gouwens2020a, Gouwens2020b, Gliko2024, Chartrand2023}.

The hippocampal homologue clarifies what is being conserved across regions. \citet{Gulyas1996} originally identified the calretinin-immunoreactive interneuron-selective (IS) class in rat hippocampus as preferentially targeting other GABAergic interneurons --- including VIP-containing basket cells in stratum radiatum/oriens --- and sparing pyramidal neurons; the IS3 cell, with soma in stratum oriens and axon extending toward oriens-lacunosum-moleculare cells, is the canonical example \citep{Gulyas1996, Cauli2014, Lee2013, Stachniak2021}. The cortical CR/VIP bipolar cell described by \citet{Cauli1997} and \citet{Kawaguchi1996} shares the targeting bias and the vertically biased axon, supporting a view in which interneuron-selective targeting is a structural feature inherited from a CGE-derived ancestral programme rather than acquired regionally \citep{Cauli1997, Kawaguchi1996, Gulyas1996, Fishell2020}. The cross-region inheritance is mirrored by cross-species inheritance: an earlier paper reports that primate calretinin-positive interneurons consist of Cajal--Retzius cells, small double-bouquet cells, and most commonly bipolar cells that synapse onto dendritic shafts and spines of pyramidal neurons --- preserving the bipolar morphology and the dendrite-targeting bias from rodent \citep{Cauli2014, Hodge2019, Cauli1997}.

\section{Structural rationale for translaminar disinhibition}

The morphological features summarised above provide a structural rationale for the disinhibitory motif developed in Section \href{/synaptic-properties}{Synaptic Properties and Connectivity} and Section \href{/disinhibition-framework}{Local Circuit Motifs and the Disinhibition Framework}. A bipolar L2/3 soma with descending axon places terminals across multiple cortical layers, with L2/3 axon collaterals contacting SST-like distal-dendrite targets and deep-layer collaterals contacting both SST- and PV-like basket targets \citep{Lee2013, Zhou2017, Schneidermizell2025, Gamlin2025}. The minority small-basket and arcade variants --- together with the L6 perisomatic-bouton fraction reported by \citet{Zhou2017} --- supply a perisomatic component absent from the canonical VIP\rightarrowSST\rightarrowPyr cartoon and align with the InhTC\_peri connectomic subtype \citep{Zhou2017, Schneidermizell2025, Kawaguchi1996, Cauli2014}. In hippocampus, the analogous role is played by interneuron-selective IS3 cells with somata in stratum oriens that target oriens-lacunosum-moleculare cells \citep{Gulyas1996, Cauli2014, Stachniak2021, Lee2013} --- providing a cross-region template for the structural logic of selective interneuron targeting. We schematise these compartmental relationships in Figure~\ref{fig-vip-axonal-targeting-schematic}, which makes no quantitative claims and uses no audited data.

\begin{figure}[!htbp]
\centering
\includegraphics[width=0.9\linewidth]{files/fig-vip-axonal-targe-9e17b60a6b5317cf21420b3a39c0d245.png}
\caption[]{\textbf{Structural rationale for translaminar disinhibition (schematic).} \textbf{(A)} Idealised L2/3-soma + descending-axon VIP cell drawn over a six-layer cortical column, with SST, PV, and pyramidal postsynaptic compartments highlighted; based on the canonical morphological repertoire described by \citet{Kawaguchi1996}, \citet{Pronneke2015}, and \citet{Apicella2022} and the connectomic/Patch-seq target distributions reported by \citet{Lee2013}, \citet{Zhou2017}, \citet{Schneidermizell2025}, and \citet{Gamlin2025}. \textbf{(B)} Comparison cartoon of small-basket, arcade, and bipolar VIP variants, with the postsynaptic compartments each preferentially contacts (perisomatic for small-basket; mixed dendritic/somatic for arcade; distal-dendrite for bipolar/bitufted). \textbf{(C)} Hippocampal IS3 cell schematic (soma in stratum oriens; axon targeting OLM cells) for cross-region contrast. \textit{Note:} This figure is a schematic with no quantitative data substrate; cited works support the qualitative compartmental relationships only and not any numerical claims.}
\label{fig-vip-axonal-targeting-schematic}
\end{figure}

Several methodological caveats apply to within-VIP morphological partitions. First, all of the major rodent reconstructed datasets --- a foundational study, \citet{Gouwens2019}, \citet{Gouwens2020a}, an earlier report, and a preceding study --- rely on VIP-Cre genetic labelling, which captures the \textit{Vip}-expressing transcriptomic subclass but does not perfectly resolve VIP-peptide expression on a cell-by-cell basis; the difference matters most for rodent-to-primate translation, where the calretinin-positive analogue is the operational handle \citep{Medalla2023, Hodge2019, Fishell2020}. Second, axonal completeness varies substantially across the field: in vitro slice biocytin reconstructions truncate distal axon collaterals, while EM reconstructions are constrained by the volume --- together producing artefactual differences in laminar axon distributions that confound met-type definition \citep{Gouwens2020a, Schneidermizell2025, Elabbady2025, Gliko2024}. Third, the dendrite-emerging-axon feature reported by an initial investigation is sensitive to the resolution of the axon initial segment in the imaging modality, which differs across the studies that have searched for it \citep{Badrinarayanan2021}. These are not arguments against the convergent picture; they are constraints on how strongly the discrete-versus-continuum and cross-area conflicts above can be adjudicated from current data.

\section{Summary and forward links}

VIP-class morphological diversity is structured but graded. A bipolar/bitufted somatodendritic profile with a vertically oriented, dendrite-emerging, mostly inhibition-targeting axon dominates the supragranular VIP population in rodent neocortex \citep{Kawaguchi1996, Pronneke2015, Apicella2022, Schneidermizell2025}. A minority of small-basket, arcade, and multipolar variants --- together with the L6 perisomatic-bouton fraction and the deep-layer multipolar repertoire --- extends the class beyond the canonical bipolar form \citep{Rudy2011, Zhou2017, Apicella2022, Yayon2023}. Patch-seq and connectomic atlases recover discrete morpho-electric centroids embedded in continuous variation rather than fully discrete classes \citep{Gouwens2020a, Gouwens2020b, Gliko2024, Schneidermizell2025}. Cross-area and cross-species shifts (mPFC, MEC, primate V1, human MTG and L1) repeat the overall layout while quantitatively reweighting subtype proportions, dendritic geometries, and biophysical parameters \citep{Moreno2026, Badrinarayanan2021, Medalla2023, Chartrand2023}. These morphologies are realised electrically as the heterogeneous intrinsic-firing repertoire surveyed in Section \href{/electrophysiology}{Intrinsic Electrophysiology}, and they constrain the synaptic and circuit motifs reviewed in Section \href{/synaptic-properties}{Synaptic Properties and Connectivity} and revisited in Section the concluding synthesis.


\section{Intrinsic Electrophysiology}
\label{sec:05_electrophysiology}

The morphological architecture surveyed in \href{/morphology}{Morphological Diversity} --- vertically oriented bipolar somata with descending axons that traverse layers --- sets a biophysical stage that intrinsic-firing measurements then occupy. VIP interneurons are not a single firing class. They are a transcriptomically coherent CGE-derived subclass that constitutes {\textasciitilde}12\% of neocortical interneurons and is concentrated in superficial layers \citep{Lim2018, Tasic2018, Yao2021a, Machold2023} whose intrinsic-electrophysiological repertoire is dominated by high input resistance, modest sag, fast-rising action potentials, and a heterogeneous mix of irregular-spiking, continuous-adapting, and (in upper layers) bursting firing patterns \citep{Tremblay2016, Pronneke2015, Cauli1997}. Patch-seq has, over the last decade, transformed how this repertoire is parsed: the four GABAergic subclasses (PV, SST, VIP, Lamp5/Sncg) separate cleanly along intrinsic-feature axes both in mouse \citep{Gouwens2020a, Gouwens2019} and in human cortex \citep{Kim2023}, but \textit{within} the VIP class, electrophysiological types form a continuum onto which transcriptomic types map only partially \citep{Scala2021, Gouwens2020a}. This section parses that landscape --- first the consensus passive- and firing-property profile, then three concrete points where studies disagree (resting potential, IS-firing fraction, and the molecular substrate of irregularity), and finally the cross-modal Patch-seq picture that frames how individual measurements should be interpreted.

\section{Passive properties and action-potential phenotype}\label{sec-05-passive-aps}

The single most reproduced intrinsic feature of VIP cells is their high input resistance. \citet{Tremblay2016} flagged this as the most salient feature of cortical \textit{Vip} neurons --- higher than most cortical neurons and a substrate for sensitivity to weak excitatory input. Quantitatively, \citet{Pronneke2015} reported a population mean of 404.7 $\pm$ 212.8 M$\Omega$ across all layers of mouse barrel cortex VIP-Cre/tdTomato neurons (with $\tau$\_\{m\} 21.4 $\pm$ 7.6 ms, sag 11.2 $\pm$ 6.8\%, rheobase 51.9 $\pm$ 33.9 pA, AP threshold -38.2 $\pm$ 3.1 mV), and \citet{Vucurovic2010} placed juvenile S1 VIP/5-HT3aR+ cells even higher (R\_\{in\} 761.5 $\pm$ 306.5 M$\Omega$, $\tau$\_\{m\} 40.6 $\pm$ 19.5 ms, rheobase 17.7 $\pm$ 9.2 pA, RMP -56.0 $\pm$ 3.9 mV) --- describing them as the most excitable cortical interneuron group \citep{Vucurovic2010}. CA1 hippocampal VIP-IS3 cells extend the same profile to non-cortical territory (R\_\{in\} {\textasciitilde}500 M$\Omega$, $\tau$\_\{m\} {\textasciitilde}36 ms, C\_\{m\} {\textasciitilde}33 pF) \citep{Michaud2024, Guetmccreight2020b}, and BLA VIP+ IS interneurons fall in the same range (350.3 $\pm$ 28.3 M$\Omega$, n=30) \citep{Rhomberg2018}. Cross-class comparisons are consistent: in monkey prefrontal cortex \citep{Zaitsev2009} and in CA1 GAD65-GFP+ CGE-derived populations \citep{Wierenga2010}, the high-R\_\{in\}, slow-$\tau$\_\{m\}, low-rheobase corner of intrinsic-feature space is occupied by VIP/CR+ rather than PV cells, recapitulating the disinhibitory split established by Petilla nomenclature \citep{Thepetillainterneuronnomenclaturegroupping2008}.

Action-potential waveforms are similarly conserved. VIP cells fire fast, narrow APs (half-width $\approx$ 0.95 ms in juvenile S1 \citep{Vucurovic2010}; AP amplitude 65.1 $\pm$ 9.1 mV in barrel cortex \citep{Pronneke2015}) --- narrower than pyramidal neurons but broader than PV fast-spiking cells, again consistent with the broader cortical interneuron taxonomy \citep{Rudy2011, Karnani2016a}. Sag/I\_\{h\} is modest rather than absent: \citet{Pronneke2015} measured a sag index of 11.2 $\pm$ 6.8\% across cortical VIP cells; \citet{Vucurovic2010} reported 12.1 $\pm$ 7.8\% in juvenile S1; CA1 IS3 cells show pronounced sag in some recordings and minimal sag in others \citep{Bell2015, Michaud2024}; and inferior-colliculus VIP-Cre stellate cells (which are \textit{glutamatergic}, not GABAergic) show essentially no sag \citep{Goyer2019} --- flagging that VIP-Cre lines tag biophysically distinct cell classes outside the cortex. Consensus on these features is sufficiently robust that Prior work \citep{Kim2023} could train a four-feature classifier (AP height, sag, AP-upstroke adaptation ratio, $\tau$\_\{m\}) that separates human PVALB from non-PVALB cortical interneurons with high accuracy, identifying the same axes that Petilla and \citet{Gouwens2019} use to organise mouse VISp into 17 e-types.

A subset of cross-study comparisons can be made cleanly only after the recording region is matched. Earlier reports recorded inferior-colliculus VIP-Cre stellate cells (V\_\{rest\} = -69.5 $\pm$ 4.4 mV, n=216; 90.3\% sustained, 8.4\% adapting, 1.3\% other; n=237) and noted minimal sag --- illustrating that VIP-Cre lines tag glutamatergic stellate populations outside the cortex with biophysical profiles unlike cortical/hippocampal VIP GABAergic cells \citep{Goyer2019}. \citet{Smith2019} framed VIP as one of the most widely expressed neuropeptide-precursor transcripts across mouse VISp/ALM single-cell transcriptomes, embedding the VIP intrinsic-feature class within a cortex-wide peptidergic-signalling network \citep{Smith2019}. \citet{Urbanciecko2016} reviewed SST-interneuron heterogeneity (Martinotti, LTS, IS, stuttering, FS subtypes) and the VIP\rightarrowSST disinhibitory motif, framing the comparator architecture against which VIP intrinsic properties are typically read \citep{Urbanciecko2016}. Across these comparisons, the rule that survives is structural rather than numeric: the VIP/CGE subclass occupies a high-R\_\{in\}, low-rheobase, narrow-AP corner of the cortical-interneuron biophysical manifold, with within-class scatter that is real but bounded \citep{Tremblay2016, Vucurovic2010, Wierenga2010, Karnani2016a}.

\section{The Petilla firing-pattern menu in VIP cells}\label{sec-05-firing-taxonomy}

Within this consensus passive profile, intrinsic firing patterns fan out along three axes drawn from the Petilla terminology \citep{Thepetillainterneuronnomenclaturegroupping2008}: (i) accommodation (continuous-adapting versus non-adapting), (ii) regularity (irregular-spiking versus regular), and (iii) bursting onset (initial bursting versus delayed). Cortical VIP cells span all three. \citet{Cauli1997} first identified, in rat sensory-motor cortex, a vertically oriented bipolar subpopulation that consistently co-expressed calretinin and VIP and fired bursts of action potentials at irregular frequency --- the founding observation of the IS firing class --- and \citet{Cauli2000} then showed by unsupervised clustering of fusiform interneurons (n=60) that 16/60 formed an ``IS-VIP'' cluster distinct from a larger RSNP-VIP group of n=32 \citep{Cauli2000}. \citet{Pronneke2015} reproduced this taxonomy quantitatively in mouse barrel cortex (75.4\% continuous-adapting, 14.5\% IS, 5.8\% bursting, 4.3\% high-threshold-bursting non-adapting across layers) and added a laminar twist: bursting VIP+ neurons are confined to L2/3 (11.8\%) while L4--L6 lack bursting cells but retain IS (14.7\%) and high-threshold-bursting non-adapting (2.9\%) variants alongside a dominant continuous-adapting population (82.4\%) \citep{Pronneke2015}. Initial studies \citep{Vucurovic2010} independently described two intrinsic firing behaviours in juvenile S1 VIP cluster cells --- 68\% adapting (15/22) and 32\% bursting (7/22), with a characteristic post-spike afterhyperpolarisation--afterdepolarisation--late-afterhyperpolarisation sequence \citep{Vucurovic2010}. Together, four independent studies place IS firing as a minority but reproducible class in cortical VIP populations, with continuous-adapting as the modal phenotype \citep{Cauli1997, Cauli2000, Pronneke2015, Walker2016}.

Mapping these firing classes onto molecular identity is, however, \textit{not} one-to-one. \citet{Tremblay2016} summarised that cortical VIP/CR co-expressing cells are enriched for irregular-spiking bipolar phenotypes, whereas VIP/CCK cells (small basket cells) are enriched for multipolar/bitufted morphologies \citep{Tremblay2016}, but this enrichment is statistical rather than absolute. Foundational studies made the orthogonality explicit: in cortex, IS firing does not co-segregate cleanly with calretinin or CCK co-expression, nor with bipolar versus multipolar morphology, so IS is best treated as an electrophysiological category orthogonal to the marker axes \citep{Guetmccreight2020b}. \citet{Guetmccreight2020b} further organised hippocampal IS interneurons into I-S1, I-S2, and I-S3 subtypes --- with I-S3 (VIP+/CR-, high R\_\{in\}) showing irregular firing and theta-rhythm-locked spiking --- and recorded a strong laminar bias in CR/VIP overlap (71\% of CR+ cells are VIP+ in L2/3 but 94\% in L5/6, with 43\% VIP-only and 29\% CR-only cells) that further severs the IS\leftrightarrowCR identity. Established work confirmed at the paired-recording level that L2/3 V1 VIP cells exhibit an adapting firing pattern in direct contrast to PV fast-spiking --- a fundamental intrinsic-physiology difference that subdivides the inhibitory taxonomy without specifying within-VIP class \citep{Walker2016}. Figure~\ref{fig-vip-ephys} Panel A displays exemplar firing traces for the three dominant patterns; Panel C tabulates the reported IS fraction across studies with denominators shown.

\section{Conflict: where do CA1 VIP cells rest?}\label{sec-05-conflict-rmp}

The single sharpest intrinsic-property disagreement in the section is over the resting membrane potential of hippocampal VIP/CR+ cells.

\begin{framed}
\textbf{Conflict --- resting membrane potential of CA1 VIP+ interneurons}\\
\citet{Bell2015} recorded a population of nicotinically responsive CA1 VIP interneurons at V\_\{rest\} = -73.9 $\pm$ 0.8 mV with R\_\{in\} = 395.3 $\pm$ 20.8 M$\Omega$ (n=30), with subpopulations differing in sag (sag ratio 1.09 $\pm$ 0.01) and accommodation (ratio 2.75 $\pm$ 0.84) \citep{Bell2015}. Earlier work reported, in adult VIP-Cre mouse CA1 stratum radiatum, that VIP/CR+ I-S3 cells rest at V\_\{rest\} = -60.5 $\pm$ 1.7 mV (R\_\{in\} 517 $\pm$ 37 M$\Omega$, C\_\{m\} 32.8 $\pm$ 3.5 pF, $\tau$\_\{m\} 36.5 $\pm$ 3.2 ms; n=7 cells / 5 animals) --- a value {\textasciitilde}13 mV more depolarised than Previous reports \citep{Bell2015} \citep{Michaud2024}. The two studies use overlapping anatomical definitions (VIP+/CR+ stratum radiatum interneurons) but different functional gating (nicotinic responsiveness in \citet{Bell2015} versus generic VIP-Cre $\times$ tdTomato in \citet{Michaud2024}) and different recording vintages. Resolution status: \textbf{open}. The discrepancy is not trivial. A 13-mV shift in baseline straddles the activation threshold of M-current and the inactivation curve of low-threshold T-type Ca\textsuperscript{2}\textsuperscript{.} channels, both of which have been implicated as substrates for VIP-cell intrinsic firing patterns \citep{Goff2019, Porter1998}. Whether the gap reflects subtype heterogeneity within CR+/VIP+ populations, age- or recording-condition differences, or true methodological divergence is not adjudicable from the available literature.
\end{framed}

\section{Conflict: how common is irregular-spiking?}\label{sec-05-conflict-isfraction}

The fraction of VIP+ interneurons displaying IS firing is the textbook diagnostic statistic for the class --- and is, on cross-paper inspection, anything but a single number.

\begin{framed}
\textbf{Conflict --- IS-firing fraction across cortical VIP populations}\\
\citet{Cauli2000} clustered fusiform interneurons in rat frontal cortex (n=60) into three discrete physio-molecular groups, of which the IS-VIP cluster contained 16/60 ($\approx$26.7\%) cells \citep{Cauli2000}. \citet{Cauli1997} had earlier reported 15 IS cells out of 97 \textit{nonpyramidal} interneurons sampled in rat S1 ($\approx$15.5\% of all sampled nonpyramidal cells; the 15 IS cells consistently co-expressed CR and VIP) \citep{Cauli1997}. \citet{Anastasiades2021a} recorded VIP-Cre $\times$ Ai14+ cells in mouse mPFC L1b and reported NFS = 9/17 (53\%), IS = 5/17 (29\%), and fast-adapting = 3/17 (18\%) (n=17 VIP+ cells, 8 mice) \citep{Anastasiades2021a}. \citet{Pronneke2015} found 14.5\% IS firing across the whole-population mouse barrel cortex VIP-Cre/tdTomato sample \citep{Pronneke2015}. The reported IS fraction therefore varies between {\textasciitilde}14.5\% (whole-population mouse barrel cortex) and {\textasciitilde}29\% (mouse mPFC L1b), with the rat-cortex values bracketing both ends. Resolution status: \textbf{partially methodological}. The denominator differs across entries (all sampled nonpyramidal interneurons; clustered fusiform cells; VIP-Cre $\times$ Ai14+ cells; whole-population VIP-Cre); the species differs (rat vs mouse); the cortical area differs (S1 vs frontal cortex vs mPFC L1b vs barrel cortex); and the firing-class taxonomies differ slightly. Figure~\ref{fig-vip-ephys} Panel C shows the four entries side-by-side with denominators annotated, as required by the Phase-6 audit caveat for this comparison: \textit{Denominators differ across entries; per-entry denominators MUST be shown in the figure (cannot be presented as a single comparable percentage without that annotation). Rat (Cauli) vs mouse (Anastasiades) species mismatch. Different cortical areas (S1, frontal cortex, mPFC L1b) and slightly different firing-class taxonomies.}
\end{framed}

\section{Conflict: what makes a VIP cell fire irregularly?}\label{sec-05-conflict-mechanism}

The irregularity of IS firing is mechanistically interesting because it is a \textit{positive} phenotype --- a specific spiking pattern that selectively converts under pharmacological perturbation --- and therefore a place where ion-channel substrates can be cleanly tested \citep{Porter1998, Goff2019, Guetmccreight2020b, Tremblay2016}. Two influential studies arrive at \textit{different} substrates.

\begin{framed}
\textbf{Conflict --- molecular substrate of irregular-spiking firing}\\
\citet{Porter1998} reported that the irregular-spiking pattern of bipolar VIPergic neocortical interneurons is mediated by an I\_\{D\}-like Kv1-family K\textsuperscript{.} current that can be neutralised by 4-aminopyridine, dendrotoxin-I, or dendrotoxin-K to convert the cell to sustained firing --- implicating Kv1.1/Kv1.2/Kv1.6 channels as the substrate of irregularity \citep{Porter1998}. \citet{Goff2019}, two decades later, reported that pharmacological block of KCNQ-mediated M-current with linopirdine selectively converts IS VIP-INs to tonic firing while leaving continuous-adapting VIP-INs unchanged, identifying KCNQ5/M-current as the molecular substrate of IS firing in cortex \citep{Goff2019}. Earlier studies independently reproduced the M-current finding in a Scn1a\textsuperscript{.}/\textsuperscript{-} mouse model: linopirdine eliminated IS firing only in IS VIP+ cells, with selective firing-rate-gain reductions in the same subset \citep{Guetmccreight2020b}. Resolution status: \textbf{partially reconciled, partially open}. The two mechanisms are not strictly mutually exclusive --- Kv1-family delay (Porter) and KCNQ/M-current sustained outward current (Goff/GuetMcCreight) operate in different voltage ranges and timescales --- but neither study tested both pharmacologies on the \textit{same} cells, and the studies used different cortical areas and recording conditions (rat neocortex with intracellular sharp electrodes versus mouse cortex with whole-cell patch). The most parsimonious synthesis is that both currents contribute and that the \textit{dominant} substrate may depend on layer, age, and species, but no single experiment in the available literature adjudicates this directly. Figure~\ref{fig-vip-firing-mechanism} schematises the two competing substrates and their predicted current-clamp signatures.
\end{framed}

The mechanistic question matters beyond intrinsic biophysics. \citet{Guetmccreight2020b} argued that the IS-VIP/M-current axis converges on disinhibition: cortical IS-type VIP+ interneurons that target SST cells are intrinsically biased toward irregularity by KCNQ/M-current expression \textit{and} are sensitised to neuromodulatory depolarisation by ACh and 5-HT \citep{Guetmccreight2020b}. The same neuromodulatory axes recruit VIP cells through nicotinic depolarisation \citep{Porter1999, Bell2015, Tremblay2016}, fast 5-HT3-mediated EPSPs in VIP/CCK co-expressing cells \citep{Ferezou2002, Chen2015a}, and muscarinic depolarisation \citep{Tremblay2016} --- meaning that the molecular substrate of IS firing is also a substrate for state-dependent activation. Whether the IS pattern itself is functionally important, or whether it is a downstream signature of the same channel complement that supports neuromodulator sensitivity, is not yet resolved.

\section{Patch-seq: subclass robustness, within-VIP continuum}\label{sec-05-patchseq}

The most important methodological development between \citet{Cauli1997} and the present is Patch-seq --- joint patch-clamp recording, scRNA-seq, and morphological reconstruction from the same cell. Prior work and Earlier reports \citep{Fuzik2016} independently demonstrated the technique on cortical interneurons and showed concordance between firing-pattern class and transcriptomic class at the L1 level (n=67 cells in \citet{Cadwell2016}). \citet{Ferrante2016} had argued the same point with hierarchical clustering of intrinsic properties alone in mouse entorhinal cortex, finding that five intrinsic features predict molecular identity (VIP, SST, 5HT3aR, NPY-NGF, NPY-non-NGF, RCAN2) with 81.4\% accuracy \citep{Ferrante2016} --- establishing, before Patch-seq scaled, that subclass-level e-types are predictable from biophysics. \citet{Gouwens2019} then partitioned mouse VISp neurons into 17 distinct e-types differing across firing-pattern descriptors (adapting, irregular, transient), AP shape, and subthreshold properties ($\tau$\_\{m\}, sag) \citep{Gouwens2019}. The integrated multimodal Patch-seq atlas of Initial studies placed these e-types into a 28-MET-type taxonomy of \textgreater 4,200 mouse VISp GABAergic interneurons, with {\textasciitilde}78\% MET-type prediction accuracy from electrophysiology + morphology and a Patch-seq subthreshold-2 PC correlated with input resistance at r = 0.90 \citep{Gouwens2020a}. The Vip subclass partitions into five MET-types in mouse VISp (Vip-MET-1 through Vip-MET-5), all bipolar/bitufted but with distinct intrinsic-firing fingerprints --- Vip-MET-1 fires regular sustained spikes (Npy+, L2/3--L4), Vip-MET-2 fires irregularly, and Vip-MET-4/5 show elevated sag \citep{Gouwens2020a}. Vip t-types are distributed in a layer-stratified pattern (Vip Col15a1 Pde1a at the L1--L2/3 border, Vip Chat Htr1f spanning L2/3, Sncg Vip Itih5 across multiple layers), so laminar location is partially diagnostic of within-VIP type \citep{Gouwens2020a, Tasic2018, Tasic2016, Yao2021a}.

The within-VIP picture is, however, \textit{not} one of discrete e-types. \citet{Scala2021} profiled mouse primary motor cortex by Patch-seq and reported that intrinsic-electrophysiological properties of \textit{Vip} transcriptomic types vary continuously rather than discretely along the transcriptomic axis --- for instance, $\tau$\_\{m\} peaks at the Sncg-adjacent end of the \textit{Vip} manifold and decreases toward Vip-Gpc3 --- supporting a graded mapping rather than crisp e-type boundaries within the subclass \citep{Scala2021}. \citet{Bugeon2022} reached a similar conclusion for cortical inhibitory neurons broadly, arguing for transcriptomic gradients as the principal axis of variation rather than discrete cluster identities \citep{Bugeon2022}. The combined picture --- sharp subclass boundaries, fuzzy within-subclass partitions --- is the modern reformulation of the Petilla taxonomy: subclass identity (VIP vs SST vs PV vs Lamp5/Sncg) is a real categorical variable; e-type identity within the VIP subclass is, at present, a continuous one \citep{Gouwens2020a, Scala2021, Bugeon2022, Yao2021a}.

This continuum has cross-species and cross-cortical-area consequences. \citet{Lee2023a} demonstrated, in human neocortical Patch-seq (n=675; LAMP5/PAX6=75, VIP=99, SST=149, PVALB=352), that the four GABAergic subclasses are separable in UMAP of intrinsic features, with VIP cells displaying diverse morphologies (bipolar, multipolar, basket) --- consistent with mouse but with greater within-class diversity \citep{Lee2023a}. \citet{Bakken2021a}, in the BICCN cross-species motor-cortex taxonomy, showed that \textit{Vip} GABAergic subtypes are markedly more diversified between humans/marmosets and mice than between primates, indicating species-specific expansion of the VIP transcriptomic axis \citep{Bakken2021a, Braininitiativecellcensusnetworkbiccn2021}. \citet{Chartrand2023} then localised an axis of human-vs-mouse divergence: VIP+ L1 interneurons localise to deeper L1 in mouse but spread across L1 in human, indicating species divergence in laminar position of L1-VIP cells \citep{Chartrand2023}. \citet{Boldog2018} complicated this picture further with the human-specific `rosehip' L1 interneuron --- a CCK/LAMP5 type \textit{distinct} from VIP cells --- that fires stuttering/IS patterns at V\_\{rest\} = -61.3 $\pm$ 5.8 mV (n=10) with $\beta$/$\gamma$-tuned subthreshold oscillations \citep{Boldog2018}. \citet{Kim2023} then showed that a Patch-seq-trained classifier of human cortical subclass identity (PVALB vs non-PVALB) is dominated by AP height, sag, AP-upstroke adaptation ratio, and $\tau$\_\{m\} --- features that \citet{Gouwens2020a} and \citet{Gouwens2019} independently identified in mouse --- confirming convergence at the subclass-feature level even where within-subclass partitions diverge \citep{Kim2023, Gouwens2020a}.

A separate axis emerges from oscillatory dynamics. Foundational studies reported that human L1 `rosehip' interneurons exhibit $\beta$/$\gamma$-tuned subthreshold oscillations that do not appear in rodent VIP cells --- a candidate species-specific intrinsic feature \citep{Boldog2018}. \citet{Bugeon2022} and \citet{Scala2021} argued that gradient-based, rather than discretised, taxonomy is more compatible with the within-VIP intrinsic-feature variation observed in Patch-seq \citep{Bugeon2022, Scala2021}. Established work showed that paired V1 recordings preserve the adapting/IS distinction in vivo \citep{Walker2016}, and Earlier work demonstrated that VIP cells are recruited heterogeneously across behavioural states with intrinsic-firing diversity that mirrors their state-dependent recruitment \citep{Jackson2016}. The broader conclusion the Patch-seq literature now supports is that subclass-level e-type prediction is reproducibly accurate ({\textasciitilde}78\% MET-type prediction in mouse VISp), but within-VIP biophysical heterogeneity is real and reflects a continuum of identities rather than discrete cell types \citep{Scala2021, Bugeon2022, Lee2023a, Bakken2021a}.

\section{Figures}\label{sec-05-figures}

\begin{figure}[!htbp]
\centering
\includegraphics[width=0.95\linewidth]{files/fig-vip-ephys-a3826fd9aedf3b4fd8528d6573a66614.png}
\caption[]{\textbf{VIP intrinsic electrophysiology: cross-study landscape and Patch-seq mapping.} \textit{Panel A:} Schematic exemplar firing traces for the three dominant Petilla classes observed in cortical VIP cells (irregular-spiking, continuous-adapting, accommodating) at a matched supra-threshold current step, drawn from the qualitative descriptions in \citep{Cauli1997, Pronneke2015, Vucurovic2010}. \textit{Panel B:} Forest-plot of input resistance for VIP+ interneurons across non-cortical regions (CA1 stratum radiatum vs basolateral amygdala). Values are mean $\pm$ SEM with n given per entry; entries are (R\_\{in\} 517 $\pm$ 37 M$\Omega$, n=7 VIP/CR+ I-S3 cells, mouse CA1) and \citep{Rhomberg2018} (R\_\{in\} 350.3 $\pm$ 28.3 M$\Omega$, n=30 VIP+ IS-INs, mouse BLA). \textit{Caveat (verbatim from Phase-6 audit):} Cross-region (CA1 stratum radiatum vs BLA) --- must be flagged in any figure caption. | n=2 entries --- minimal cross-paper comparison. \textit{Panel C:} Stacked bar of reported irregular-spiking fraction across studies. Per-entry denominators are annotated on the bars per Phase-6 audit requirement; entries: \citep{Cauli1997} (15 IS cells / 97 nonpyramidal interneurons sampled = 15.5\% of nonpyramidal pool; the 15 IS cells co-expressed CR+VIP, rat S1, single-cell RT-PCR + patch); \citep{Cauli2000} (16/60 = 26.7\% IS-VIP cluster, rat frontal cortex L2--3, fusiform interneurons clustered by RT-PCR + patch); \citep{Pronneke2015} (14.5\% of whole-population VIP-Cre/tdTomato barrel cortex); \citep{Anastasiades2021a} (5/17 = 29.4\% of mPFC L1b VIP-Cre $\times$ Ai14+ cells, mouse). \textit{Caveat (verbatim from Phase-6 audit):} Denominators differ across entries; per-entry denominators MUST be shown in the figure (cannot be presented as a single comparable percentage without that annotation). | Rat (Cauli) vs mouse (Anastasiades) species mismatch. | Different cortical areas (S1, frontal cortex, mPFC L1b) and slightly different firing-class taxonomies. \textit{Panel D:} Single-row entry showing membrane time constant $\tau$\_\{m\} (ms) for the only entry retained after Phase-6 audit removed \citet{Rhomberg2018} from the $\tau$\_\{m\} sub-comparison ( reports $\tau$\_\{m\} = 36.5 $\pm$ 3.2 ms, n=7 VIP/CR+ I-S3 cells, mouse CA1 stratum radiatum); panel reduced to a single-row entry per the iter3 audit. Patch-seq subclass-feature axes (\citep{Gouwens2020a, Scala2021, Lee2023a, Kim2023}) provide the cross-species reference frame for these comparisons.}
\label{fig-vip-ephys}
\end{figure}

\begin{figure}[!htbp]
\centering
\includegraphics[width=0.95\linewidth]{files/fig-vip-firing-mecha-780f3e88abc012351b18da226359172e.png}
\caption[]{\textbf{Candidate ion-channel substrates of irregular-spiking firing in VIP cells.} \textit{Panel A:} Cartoon VIP-cell membrane with annotated channel complement (Na\_\{V\}, Kv1, Kv4, HCN/I\_\{h\}, KCNQ/M, SK) and qualitative voltage-clamp traces; channel inventory is consistent with the cortical VIP intrinsic-property literature \citep{Tremblay2016, Pronneke2015, Vucurovic2010}. \textit{Panel B:} Two competing mechanism schematics for irregular-spiking --- \textit{(i)} Kv1-dominant repolarisation kinetics (block by 4-AP / dendrotoxin converts to sustained firing) \citep{Porter1998}; \textit{(ii)} KCNQ-mediated M-current (block by linopirdine selectively converts IS to tonic in IS VIP cells) \citep{Goff2019, Guetmccreight2020b} --- with predicted current-clamp signatures. \textit{Panel C:} Decision flow mapping intrinsic-feature combinations (high R\_\{in\} + IS + low sag \rightarrow cluster A; high R\_\{in\} + adapting + moderate sag \rightarrow cluster B) onto Patch-seq Vip-MET types \citep{Gouwens2020a} and \textit{Vip} t-type axes \citep{Scala2021}. Schematic; no quantitative data substrate \citep{Cauli1997, Cauli2000}.}
\label{fig-vip-firing-mechanism}
\end{figure}

\section{Intrinsic biophysics meets brain state and neuromodulation}\label{sec-05-state-modulation}

The high R\_\{in\} / modest-sag biophysical profile is not a passive bookkeeping detail --- it is the substrate that allows VIP cells to be recruited by weak depolarising drive, including neuromodulatory inputs that bypass classical glutamatergic synapses. \citet{Porter1999} showed that nicotinic receptor activation selectively depolarises neocortical VIP/CCK-coexpressing interneurons while having no effect on pyramidal cells, PV fast-spiking, or SST cells \citep{Porter1999}. \citet{Bell2015} extended this to CA1 with a nicotinically responsive VIP population \citep{Bell2015}. \citet{Ferezou2002} reported that 5-HT3 receptors mediate fast serotonergic EPSPs selectively in CCK/VIP-co-expressing GABAergic interneurons \citep{Ferezou2002}, and Previous reports \citep{Chen2015a} showed that VIP and L1 interneurons in mouse V1 are facilitated by ACh at lower agonist doses than SOM neurons \citep{Chen2015a}. \citet{Tremblay2016} summarised that VIP+ bipolar cells are dual-modulatory targets --- depolarised by both ACh (via nicotinic AChRs) and 5-HT (via 5-HT3aR) --- and additionally by muscarinic agonists \citep{Tremblay2016}. The neuromodulator-coupled excitability is functionally consequential: in cortex, VIP cells are nonspecifically active across behavioural states \citep{Jackson2016}, are recruited by reinforcement signals (both punishment and reward) in awake auditory cortex and mPFC \citep{Pi2013}, are bidirectionally modulated by locomotion in V1 \citep{Guetmccreight2020b}, and exhibit complementary contrast tuning to SST cells in V1 \citep{Millman2020}. Conditional manipulations of intrinsic-excitability regulators establish causal links: deletion of Igf1 in VIP cells increases inhibitory drive onto VIP cells and decreases their firing \citep{Mardinly2016}, and ErbB4 deletion in VIP-INs causes a long-term {\textasciitilde}4-fold increase in pyramidal-cell spontaneous firing in adolescence and abolishes locomotion-induced state transitions \citep{Batistabrito2017}. Two further details place the intrinsic profile in its developmental and circuit context: all cortical VIP cells originate from the CGE \citep{Lim2018, Cauli2014}, segregating sharply from the MGE-derived PV/SST classes that account for the bulk of the rest of the inhibitory taxonomy \citep{Rudy2011, Karnani2016a}; and a small ChAT-expressing VIP subpopulation directly \textit{excites} (rather than inhibits) neighbouring pyramidal and interneurons through fast cholinergic + GABA co-transmission \citep{Obermayer2019} --- a counter-example to the strict-disinhibitor framing that \href{/synaptic-properties}{Synaptic Properties and Connectivity} returns to in detail.

The intrinsic-firing repertoire is also developmental. Earlier studies showed that conditional ErbB4 deletion in VIP cells produces a long-term {\textasciitilde}4-fold elevation in spontaneous pyramidal-cell firing that emerges in adolescence and abolishes locomotion-induced state transitions, implicating an adolescent maturation window for VIP intrinsic excitability \citep{Batistabrito2017}. Prior work made a parallel point with conditional Igf1 deletion in VIP neurons: increased inhibitory drive specifically onto VIP cells decreases their firing, establishing experience-dependent regulation of VIP intrinsic firing through a synaptic, not channel-intrinsic, route \citep{Mardinly2016}. Earlier reports \citep{Abs2018} sharpened the molecular boundary of the VIP class from a different direction: NDNF is a highly selective marker for cortical L1 interneurons that is mutually exclusive with \textit{Vip}, \textit{Pvalb}, and \textit{Sst}, demarcating the VIP subclass from the adjacent L1 NDNF/Lamp5 population that occupies a neighbouring superficial niche \citep{Abs2018}. \citet{Obermayer2019} reported that a ChAT-expressing subset of cortical VIP cells \textit{excites} (rather than inhibits) neighbouring pyramidal and interneurons through fast cholinergic + GABA co-transmission \citep{Obermayer2019}. These observations do not overturn the basic intrinsic-feature consensus, but they establish that VIP intrinsic biophysics is plastic, conditional, and partially uncoupled from the disinhibitory ``disinhibitor'' framing \citep{Pi2013, Pfeffer2013, Karnani2016a} --- a point that \href{/disinhibition-framework}{Local Circuit Motifs and the Disinhibition Framework} and the concluding synthesis revisit.

\section{Where the intrinsic profile is - and is not - diagnostic}\label{sec-05-bridge}

The composite picture is now sharp at the subclass level and intentionally fuzzy below it. VIP cells are reliably distinguishable from PV/SST/Lamp5 by a small set of intrinsic features (high R\_\{in\}, modest sag, fast narrow APs, adapting or irregular firing) --- features that train accurate Patch-seq classifiers across mouse and human cortex \citep{Gouwens2020a, Gouwens2019, Kim2023, Lee2023a, Ferrante2016}. Within VIP, however, individual properties (resting V\_\{m\}, IS-firing fraction, IS mechanism) vary substantially across studies --- sometimes for biological reasons (layer-specific subtypes \citep{Pronneke2015, Gouwens2020a}, species expansion \citep{Bakken2021a, Chartrand2023}, region-specific I-S3 vs cortical IS biology \citep{Guetmccreight2020b, Michaud2024}), sometimes for methodological ones (denominator definitions \citep{Cauli1997, Cauli2000, Anastasiades2021a}, recording vintage \citep{Bell2015, Michaud2024}). Where studies disagree, the right inference is rarely ``one is wrong'' but rather ``the e-type axis within VIP is a continuum, and which slice of that continuum a given study samples is set by genetics, anatomy, and pharmacology.'' Intrinsic biophysics is therefore necessary but not sufficient to specify within-VIP identity. The next step is to ask how this intrinsic repertoire interacts with the synaptic input and output patterns that VIP cells participate in --- and whether the disinhibitory VIP\rightarrowSST\rightarrowPyr disinhibitory motif \citep{Pfeffer2013, Pi2013, Karnani2016a} is the right summary at all. That question is taken up in \href{/synaptic-properties}{Synaptic Properties and Connectivity}.


\section{Synaptic Properties and Connectivity}
\label{sec:06_synaptic_properties}

The preceding sections traced VIP interneurons from their lineage and molecular identity (\href{/classification}{Molecular Identity and Transcriptomic Taxonomy} through \href{/morphology}{Morphological Diversity}) through their intrinsic and morphological signatures (\href{/electrophysiology}{Intrinsic Electrophysiology}). Those analyses establish \textit{what kind of cell} a VIP interneuron is; they do not yet tell us \textit{what circuit role it plays}. That role is set by synapses --- by the afferents that drive VIP cells, by the targets they inhibit, and by the neuromodulators that gate both. This section synthesises 467 papers organised under the synaptic-connectivity and SST/PV-context clusters of the literature reviewed here, together with the neuromodulation cluster, around a single thesis: \textbf{VIP interneurons are modulator-targeted cells whose canonical disinhibitory output onto SST cells is one --- but only one --- element of an empirically broader and area-, layer- and state-dependent connectivity}. The widespread textbook claim that ``VIP cells inhibit only other interneurons'' is contested \citep{Yu2019, Zhou2017, Anastasiades2021a} against a robust but quantitatively heterogeneous canonical motif \citep{Pfeffer2013, Pi2013, Lee2013, Karnani2016a}, and represents the first major divergence between canonical schematic and contemporary data.

A note on scope. This section concentrates on cortical (V1, A1, S1, mPFC, motor cortex) and hippocampal VIP populations, where the modulator-targeted disinhibitory architecture has been most thoroughly dissected. Subcortical and amygdalar VIP-cell populations are referenced where they bear directly on the cortical disinhibitory motif (in particular for amygdalar fear-conditioning circuits) but are not surveyed exhaustively here; the developmental, comparative and disease-relevant aspects of those populations are returned to in Section 8. Throughout, ``canonical motif'' refers to the VIP\rightarrowSST\rightarrowPyr disinhibition described by \citet{Pfeffer2013} and \citet{Pi2013}, and ``non-canonical exceptions'' refers to the layer-, area- and source-specific deviations from that motif documented in subsequent paired-recording, in-vivo optogenetic and rabies-tracing studies \citep{Yu2019, Zhou2017, Anastasiades2021a, Walker2016, Wall2016, Karnani2016a}.

\setcounter{tocdepth}{2}
\tableofcontents
\section{Local cortical excitation and inhibition onto VIP cells}

Within a cortical column, VIP interneurons receive excitatory drive from neighbouring pyramidal neurons that is, by paired-recording standards, modest. Connection probabilities from local L2/3 pyramidal cells onto VIP cells are comparable to (and in some preparations slightly lower than) those onto SST and PV interneurons \citep{Karnani2016a}, and unitary EPSC amplitudes are correspondingly small \citep{Karnani2016a, Walker2016, Pronneke2015}. What distinguishes the local pyramidal\rightarrowVIP synapse is its short-term plasticity profile: VIP cells in mouse barrel and visual cortex receive predominantly \textit{depressing} excitation from local pyramidal neurons \citep{Karnani2016a, Cauli2014}, in contrast to the strongly facilitating Pyr\rightarrowSST synapses described in earlier classical work. The combination --- comparable connectivity but distinct frequency-dependence --- argues that VIP cells are not simply weakly-driven SST cells but are tuned to integrate sustained, regularly-spaced pyramidal activity rather than transient bursts.

Local inhibition onto VIP cells is, conversely, dominated by SST input. Paired recordings and SST-Cre optogenetics consistently show that SST interneurons provide the largest single source of local IPSCs onto VIP cells, with connection probabilities that meet or exceed the reciprocal VIP\rightarrowSST direction in the same preparations \citep{Pfeffer2013, Karnani2016a, Walker2016}. PV cells contribute a smaller but non-negligible fraction of VIP-targeted inhibition, again compatible across visual, somatosensory and prefrontal slices \citep{Pfeffer2013, Walker2016, Pronneke2019}. Mutual inhibition between VIP cells themselves is sparse but reliably detected at the single-cell level \citep{Karnani2016b, Karnani2016a, Walker2016}. The net consequence is that the local VIP cell sits inside a small recurrent inhibitory loop with SST cells in which either node is poised to disinhibit the other depending on which receives the larger external drive --- a circuit motif on which much of the contemporary state-dependent literature pivots \citep{Pfeffer2013, Pi2013, Karnani2016a, Sabri2024, Lee2013}.

\begin{figure}[!htbp]
\centering
\includegraphics[width=1\linewidth]{files/fig_sec6_connectome-1f0c5aeea372a2035f7328f1abe24bf5.png}
\caption[]{\textbf{VIP synaptic connectivity at a glance.} \textbf{(A)} Qualitative cross-study consensus categories (\texttt{absent / weak / moderate / strong}) for afferent drive onto VIP, SST and PV interneurons. Six afferent classes are pooled across rodent neocortex (V1, A1, S1, mPFC). Sources include \citet{Karnani2016a}, \citet{Walker2016}, \citet{Karnani2016a}, \citet{Wall2016}, \citet{Lee2013}, \citet{Zhang2014}, \citet{Williams2019}, \citet{Cruikshank2012}, \citet{Anastasiades2021a}, \citet{Audette2018}, \citet{Porter1999}, \citet{Fu2014}, \citet{Arroyo2012}, \citet{Ferezou2002}, \citet{Vucurovic2010}, \citet{Kawaguchi1998}. Cells deliberately encode categories, not magnitudes. \textbf{(B)} Qualitative VIP output matrix across cortex L2/3, L5/6 and CA1; ``n.t.'' = not tested in the cited preparations. \textbf{(C)} The single Phase-6-audited quantitative comparison: VIP\rightarrowSST IPSQ = 4.6 $\pm$ 1.5 pC versus VIP\rightarrowPyr IPSQ = 0.6 $\pm$ 0.2 pC (and INC 0.42 $\pm$ 0.14 pC vs 0.06 $\pm$ 0.02 pC, scaled $\times$10 for visibility) from V1 L2/3 paired recordings (n = 11). \textbf{(D)} Schematic of cholinergic recruitment of VIP / 5-HT3AR cells: nicotinic (slow, $\beta$2-containing) excitation and a smaller muscarinic inhibitory effect onto a subset of VIP cells, after \citet{Porter1999}, \citet{Arroyo2012}, \citet{Fu2014}, \citet{Obermayer2019} and \citet{Granger2020}. \textbf{Caveat --- qualitative.} All cells in panels A and B are categorical consensus; only panel C carries an audited numeric comparison.}
\label{fig-vip-connectome}
\end{figure}

\section{Long-range cortico-cortical and thalamic afferents}

The defining afferent feature of cortical VIP interneurons is the \textit{over-representation} of long-range glutamatergic input relative to local pyramidal drive. In mouse barrel cortex, vibrissa motor cortex (vM1) projects more strongly onto VIP cells than onto neighbouring pyramidal neurons, with input ratios of roughly threefold \citep{Lee2013, Williams2019, Audette2018}; analogous over-targeting of VIP by long-range cortical fibres has been most cleanly demonstrated for the cingulate\rightarrowV1 projection \citep{Zhang2014}, with related VIP-preferential top-down recruitment reported in mPFC and frontal-feedback contexts \citep{Anastasiades2021a, Wall2016}; orbital\rightarrowsensory and other higher-order routes are not yet independently mapped at the same resolution and we treat them here as extrapolations from the cingulate\rightarrowV1 result. A monosynaptic-rabies survey by an earlier report formalised this point at scale: VIP cells in mouse V1 receive disproportionately more long-range cortical input --- relative to local pyramidal input --- than do SST or PV cells in the same preparation, an asymmetry that has since been reproduced across primary sensory cortices and prefrontal cortex \citep{Lee2013, Zhang2014}.

Thalamic input to VIP cells follows the same asymmetric logic but is more strongly area- and source-dependent. Higher-order thalamic nuclei --- POm in the somatosensory pathway and pulvinar/LP in the visual stream --- innervate VIP cells robustly, often with response probabilities and amplitudes comparable to those onto pyramidal targets \citep{Williams2019, Anastasiades2021a, Ramamurthy2023}. POm input is, in addition, sufficient to drive VIP-mediated suppression of SST interneurons in vivo, recapitulating the canonical disinhibitory motif from a thalamic source \citep{Williams2019, Audette2018}. By contrast, primary thalamic relays (VPM, LGN, MGv) excite VIP cells more weakly, and several reports place VPM-driven excitation of PV cells above that of VIP \citep{Cruikshank2012, Ji2016}. In thalamo-prefrontal pathways, mediodorsal thalamus directly recruits VIP cells in mPFC L2/3 and supports goal-directed disinhibition during cognitive tasks \citep{Anastasiades2021a, Kim2017}. Together, the thalamic data refine the long-range narrative: VIP cells are not generically biased towards ``any'' extracortical input but are preferentially read by \textit{modulatory and higher-order} thalamic systems.

Two convergent quantitative replications anchor the long-range claim. First, the vM1\rightarrowVIP/Pyr ratio of {\textasciitilde}2.9$\times$ reported by \citet{Lee2013} is independently reproduced in similar form for vS1\rightarrowCg, ACC\rightarrowV1 and OFC\rightarrowsensory pathways \citep{Zhang2014, Bilash2023, Stachniak2023, Pi2013, Williams2019}. Second, the rabies-tracing asymmetry of \citet{Wall2016} has been recapitulated across cortical areas with both pseudo-typed rabies and CTB injection assays, all converging on the same direction of effect even when absolute counts differ \citep{Wall2016, Audette2018, Anastasiades2021a, Ramamurthy2023, Williams2019}. Where heterogeneity does emerge, it tends to track \textit{which} long-range source is interrogated --- a property best illustrated by the well-known A1\leftrightarrowV1 asymmetry, where Cg\rightarrowV1 recruits VIP cells \citep{Zhang2014} whereas A1\rightarrowV1 \textit{suppresses} them \citep{Ibrahim2016} (an apparent conflict that is, on closer reading, a difference in which source area is being tested). The afferent picture that emerges is therefore neither uniformly ``VIP receives top-down'' nor ``VIP receives thalamic'' --- it is a structured asymmetry in which long-range and higher-order modulatory inputs preferentially read VIP cells over their interneuron neighbours, and primary feedforward inputs do not \citep{Wall2016, Lee2013, Williams2019, Cruikshank2012, Anastasiades2021a, Audette2018, Bastos2023a}.

\section{VIP outputs: the canonical disinhibitory motif onto SST cells}

The VIP\rightarrowSST connection is, by every accounting, the single best-established VIP output. \citet{Pfeffer2013} --- using paired recordings combined with VIP-Cre/ChR2 stimulation in mouse V1 L2/3 --- quantified IPSCs onto simultaneously recorded SST and pyramidal targets and reported a roughly seven-fold larger inhibitory charge transfer onto SST cells than onto pyramids (IPSQ 4.6 $\pm$ 1.5 pC vs 0.6 $\pm$ 0.2 pC, n = 11; INC 0.42 $\pm$ 0.14 pC vs 0.06 $\pm$ 0.02 pC) (audited; Figure~\ref{fig-vip-connectome}C). Independent replication has come from \citet{Pi2013} in auditory cortex, \citet{Lee2013} in barrel cortex, \citet{Karnani2016a} and \citet{Karnani2016b} in V1 with single-cell two-photon resolution, and a preceding study in V1 chemogenetics. The motif is not confined to mouse: pharmacological and optogenetic dissection in rat \citep{Pronneke2019} and comparative meta-analytic syntheses across V1, A1 and S1 \citep{Guetmccreight2020b, Pronneke2019} reach the same architectural conclusion. The two-photon dissection of \citet{Karnani2016b} further demonstrated that activation of \textit{single} VIP cells in V1 in vivo is sufficient to create spatially narrow, transient holes in surrounding SST activity --- a result that places the canonical motif on a per-cell footing rather than as a population aggregate.

Multiple lines of evidence converge on a behaviourally relevant deployment of this motif. In V1, locomotion in darkness activates VIP cells and disinhibits pyramidal neurons through suppression of SST cells \citep{Fu2014, Dipoppa2018}. In A1, reinforcement signals from basal forebrain recruit VIP\rightarrowSST disinhibition during associative learning \citep{Pi2013, Letzkus2011}. Top-down feedback from cingulate cortex drives a centre-surround disinhibitory architecture in V1 via Cg\rightarrowVIP\rightarrowSST \citep{Zhang2014, Stachniak2023, Bastos2023a}, and in mPFC, mediodorsal-thalamic recruitment of VIP underlies disinhibitory control of pyramidal output during goal-directed behaviour \citep{Anastasiades2021a, Kamigaki2017, Pinto2015}. The consistency across modalities, areas and behavioural states is the strongest argument that the VIP\rightarrowSST motif is a genuine canonical circuit element rather than an artefact of a single preparation.

The motif is, however, neither obligatory nor isotropic in its behavioural recruitment. Two well-known counter-examples appear repeatedly in the literature reviewed here and merit explicit treatment. First, in awake V1 \textit{during} visual stimulation, locomotion increases activity of VIP, SST \textit{and} PV cells in parallel rather than producing the canonical VIP-up/SST-down divergence --- a finding originally reported by \citet{Pakan2016} and converging with \citet{Dipoppa2018} and \citet{Kaneko2024}. The reconciling reading is that the disinhibitory polarisation is context-dependent on visual drive: in darkness or during weak input the canonical sign holds, but bright structured stimuli decouple SST suppression from VIP activation \citep{Pakan2016, Dipoppa2018, Kaneko2024, Sabri2024}. Second, although Cg\rightarrowV1 robustly recruits VIP \citep{Zhang2014}, A1\rightarrowV1 by contrast \textit{suppresses} VIP cells \citep{Ibrahim2016}, and direct optogenetic VIP perturbations do not always replicate the modulatory effects ascribed to attentional state in earlier work \citep{Ibrahim2016, Sabri2024, Anastasiades2021a}. The honest synthesis --- and one that the cluster\_05 conflicts table forces us to make --- is that ``VIP\rightarrowSST\rightarrowPyr disinhibition'' is a \textit{available} circuit operation that different long-range inputs and brain states engage, suppress, or bypass; it is not the unique route by which top-down or arousal signals reach pyramidal output \citep{Pakan2016, Dipoppa2018, Ibrahim2016, Anastasiades2021a, Kaneko2024, Sabri2024}.

\section{Outputs beyond SST: VIP-to-PV, VIP-to-Pyr and the layer-specific exception}

The cleanest single divergence between canonical schematic and contemporary data lies in VIP outputs to \textit{non-SST} targets. The dominant cortical schematic --- VIP cells inhibit other interneurons, principally SST --- is supported in superficial layers but does not survive a deeper-layer or in vivo cell-attached test \citep{Yu2019, Zhou2017, Anastasiades2021a, Walker2016}. An earlier report performed in vivo VIP-Cre/ChR2 stimulation in barrel cortex and recorded from large populations of FS/PV and pyramidal cells across cortical layers; rather than a uniform PV-disinhibition signature, they found that a substantial subset of L5/6 FS-PV cells are \textit{directly} inhibited by VIP activation, and that pyramidal cells in deep layers also receive direct VIP-evoked IPSCs. In a complementary slice study, an initial investigation reported direct VIP\rightarrowpyramidal connectivity in mPFC L5 with paired recordings, and an earlier analysis confirmed that this projection is dense enough to be physiologically consequential during MD-thalamic recruitment of VIP cells.

These layer- and area-specific direct VIP\rightarrowPyr and VIP\rightarrowPV connections are not at odds with \citet{Pfeffer2013} once the comparison is read carefully: the original measurements were made in V1 L2/3 paired recordings (where superficial pyramids receive only sparse VIP input) and the resolution-status entry in the conflicts table is layer-dependent, with L2/3 canonical disinhibition coexisting with L5/6 direct VIP\rightarrowPV/Pyr connectivity \citep{Pfeffer2013, Yu2019, Zhou2017, Anastasiades2021a}. Independent reports of meaningful VIP\rightarrowPV connection rates in barrel and frontal cortices \citep{Anastasiades2021a, Pi2013} and of VIP synapses onto NDNF/L1 interneurons \citep{Schuman2019, Hafner2021} further enlarge the output set. VIP cells therefore inhibit pyramidal neurons less than they inhibit SST cells, but their direct pyramidal output is not zero, and the connection probability scales with cortical depth and source area in a way that the canonical schematic flattens \citep{Zhou2017, Anastasiades2021a, Walker2016, Pi2013, Schuman2019}.

A further refinement is provided by single-cell connectivity dissection. Two-photon photostimulation by \citet{Karnani2016b}, paired-recording with multi-target post-synaptic sampling by \citet{Walker2016} and rabies-based output mapping by \citet{Wall2016} all converge on the conclusion that individual VIP cells are not uniformly biased to a single target type --- there is substantial cell-to-cell variability in the relative weight of VIP\rightarrowSST, VIP\rightarrowPV and VIP\rightarrowPyr outputs, and the population mean is dominated by a subset of strongly-connected cells \citep{Karnani2016b, Walker2016, Wall2016, Pronneke2019}. This single-cell heterogeneity, which dovetails with the molecular subtype heterogeneity reviewed in \href{/development}{Developmental Origins and Postnatal Maturation} \citep{Tasic2018, Pronneke2019, Pronneke2015}, recasts the canonical motif as a \textit{population-mean} property rather than a wiring rule that every VIP cell obeys.

\begin{figure}[!htbp]
\centering
\includegraphics[width=1\linewidth]{files/fig_sec6_neuromodula-bab5d0c1cb1fa9393f8b260b0ef2418d.png}
\caption[]{\textbf{Neuromodulator pharmacology onto VIP cells (qualitative redesign).} \textbf{(A)} Schematic of three principal modulator afferents --- basal-forebrain ACh, raphe 5-HT and locus coeruleus NA --- converging on the cortical VIP / 5-HT3AR class. Receptor labels mark the dominant pharmacology reported in the cited papers: nAChR (slow, $\beta$2-containing) and 5-HT3R produce excitation; mAChR can produce inhibition in a VIP subset; $\alpha$/$\beta$ adrenergic effects are mixed. \textbf{(B)} Dot-plot inventory of the \textit{direction} of cholinergic effect on VIP / 5-HT3AR firing across eight studies; markers encode \texttt{excite / inhibit / mixed-co-released / not assayed} and \textbf{deliberately do not encode magnitudes}. \textbf{(C)} Methods inventory of nine state- or arousal-recruitment studies, listing preparation, inferred neuromodulator and direction of effect on VIP activity. \textbf{GAP CAVEAT --- qualitative only.} Cluster-13 (neuromodulation) had \textbf{no Phase-6 audited panels} in this review; the figure intentionally encodes only direction-of-effect. Specific quantitative values (\% changes in firing rate, $\Delta$F/F amplitudes, EC\textsubscript{5}\textsubscript{0}) cited in the prose must be traced individually to their primary papers and are not validated here.}
\label{fig-vip-neuromodulator-pharmacology}
\end{figure}

\section{Neuromodulatory drive: cholinergic, serotonergic and noradrenergic recruitment}

If the long-range and thalamic afferent picture in the previous subsections shows that VIP cells are over-represented as targets of \textit{cortical and higher-order} extracortical input, the neuromodulator data show that they are over-represented as targets of \textit{neuromodulatory} input as well --- to the point that VIP/5-HT3AR cells are reasonably described as a ``modulator-targeted'' interneuron class. Three lines of evidence support this framing.

The cholinergic line is the most quantitatively developed. Cortical VIP and CCK interneurons are selectively depolarised by nicotinic agonists, whereas pyramidal cells, fast-spiking PV cells and SST cells are not --- a result first established by \citet{Porter1999} and an early report and replicated across rodent neocortex with subtype-specific receptor profiling \citep{Vonengelhardt2007, Arroyo2012}. The receptor mediating this excitation is predominantly a non-$\alpha$7, $\beta$2-containing nAChR with a slow desensitisation profile that can sustain VIP firing for many seconds after a brief cholinergic transient \citep{Arroyo2012, Porter1999, Fu2014, Letzkus2011}. Behaviourally, the same receptor route accounts for locomotion-driven VIP activation in V1: local nicotinic blockade abolishes locomotion-evoked VIP firing while local muscarinic blockade does not, and silencing basal-forebrain cholinergic projections produces the same effect \citep{Fu2014, Sabri2024}. Layered on top is a smaller but reproducible muscarinic component: a subset of VIP cells is directly \textit{inhibited} by mAChR activation in slice \citep{Obermayer2019, Granger2020}, and a further subset of VIP interneurons themselves co-release ACh and GABA --- making them not just modulator targets but local cortical sources of ACh \citep{Granger2020, Obermayer2019, Dudai2020}. The opposite-sign cholinergic pharmacology (excitatory nicotinic, with a smaller inhibitory muscarinic component on a subpopulation) is summarised qualitatively in Figure~\ref{fig-vip-neuromodulator-pharmacology}A,B with explicit gap caveat: no quantitative fold-changes are plotted because cluster-13 was not audited in Phase 6.

The serotonergic line is similarly selective. Cortical 5-HT3AR-expressing GABAergic interneurons --- the molecular umbrella that contains nearly all VIP cells \citep{Rudy2011, Lee2010, Vucurovic2010, Frazer2017} --- are excited by serotonin via 5-HT3 receptors that produce a fast inward current absent from PV and SST cells \citep{Ferezou2002, Vucurovic2010, Frazer2017}. As with nAChR excitation, the response is layer- and area-broad and is detectable in vivo as a short-latency increase in VIP firing following raphe stimulation \citep{Frazer2017}. Noradrenergic recruitment is the least monolithic of the three: $\alpha$-adrenergic stimulation excites a fraction of VIP/CR cells in slice \citep{Cauli2014, Caputi2009}, while $\beta$-adrenergic and tonic LC-driven effects vary across area, behavioural state and recording configuration \citep{Cauli2014, Letzkus2011, Szadai2022}. Despite this heterogeneity, the broad pattern --- VIP cells preferentially receive ACh, 5-HT and (with lower selectivity) NA modulation that is largely absent from PV/SST --- is cited as the rationale for treating VIP cells as the \textit{cortical relay} of subcortical neuromodulator state across multiple recent reviews \citep{Letzkus2011, Rudy2011, Pronneke2019, Sabri2024}.

A behavioural integration ties these three modulator lines together. Across V1, A1, mPFC, BLA and CA1, the same conjunction recurs: arousal- or reinforcer-related signals from basal forebrain or LC raise VIP activity, the rise is sufficient to suppress local SST output, and the SST suppression in turn opens a dendritic disinhibitory window onto pyramidal neurons \citep{Fu2014, Letzkus2011, Krabbe2019, Pi2013, Sabri2024, Szadai2022, Tamboli2024}. Where the mapping breaks (Pakan, Dipoppa, Kaneko in V1 with structured visual drive; Ibrahim with cross-modal A1\rightarrowV1; Anastasiades with cell-type specific direct VIP\rightarrowPyr in mPFC) it does so in \textit{predictable} ways given the underlying afferent and output diagrams of Figure~\ref{fig-vip-connectome} --- by recruiting either non-canonical inputs to VIP or non-canonical outputs from VIP rather than abolishing the modulatory drive itself \citep{Pakan2016, Dipoppa2018, Kaneko2024, Ibrahim2016, Anastasiades2021a}.

\section{Hippocampal IS3 cells: a divergent connectivity logic}

Hippocampal VIP interneurons illustrate that the cortical disinhibitory motif is one solution to a more general design problem rather than the universal answer. CA1 VIP+ cells partition into two coarse classes: VIP+/CR+ interneuron-specific type-3 (IS3) cells, which selectively innervate other interneurons, and a smaller VIP+/CCK+ basket-cell population, which targets pyramidal somata in the perisomatic domain \citep{Tyan2014, Turi2019, Francavilla2018a, Booker2018}. Detailed paired-recording, electron-microscopic and intersectional optogenetic dissection of the IS3 class shows that --- unlike its cortical VIP counterpart --- its principal output is not the SST cell but the \textit{oriens-lacunosum moleculare} (OLM) interneuron, with strong, facilitating IPSCs that suppress OLM firing during sharp-wave ripples and during theta-rhythmic dendritic input from entorhinal cortex \citep{Francavilla2018a, Tamboli2024, Turi2019}. The afferent picture mirrors the output one: lateral entorhinal cortex provides the dominant glutamatergic drive onto VIP/IS3 cells in stratum lacunosum-moleculare \citep{Bilash2023, Turi2019, Tyan2014, Booker2018}, and arousal- or novelty-related signals robustly increase IS3 activity during free spatial exploration \citep{Tamboli2024, Turi2019}. The CCK+/VIP+ basket subset, by contrast, deploys a perisomatic inhibition reminiscent of cortical PV cells but with the cannabinoid- and opioid-modulated short-term plasticity of CCK basket cells \citep{Turi2019, Francavilla2018a, Caputi2009}.

A second hippocampal divergence concerns long-range projection. A subset of VIP/CR interneurons in CA1 projects out of the hippocampus to the medial septum \citep{Booker2018}, recapitulating the rare ``GABAergic projection neuron'' wiring described in earlier classical anatomy \citep{Booker2018}. Although these projection neurons are too few to reshape the population-level VIP picture, they show that the molecular VIP class is not categorically restricted to local-circuit roles --- a point that complicates any Section-7-style unified circuit model that assumes ``VIP = local interneuron-targeting GABAergic cell''. The hippocampal evidence therefore argues for an \textit{area-specific connectivity logic} in which the molecular identity is conserved but the wiring is repurposed: VIP cells are deployed wherever a cortical or hippocampal microcircuit needs a modulator-targeted, IN-targeting disinhibitor, but the \textit{which-IN} and \textit{which-pyramidal-population} are tuned to local computational requirements \citep{Tyan2014, Turi2019, Francavilla2018a, Booker2018, Bilash2023, Tamboli2024, Toth1993}.

\section{Short-term plasticity, co-transmission and pre-synaptic specialisation}

The synaptic \textit{kinetics} of VIP afferents and outputs --- short-term plasticity, peptide vs amino-acid co-transmission and pre-synaptic specialisation --- refine the connectivity skeleton above into a circuit operation with characteristic time-constants. Three quantitative regularities recur across the literature reviewed here.

First, the short-term plasticity of inputs onto cortical VIP cells \textit{splits by source}. \citet{Karnani2016a} directly compared local pyramidal-onto-VIP and pyramidal-onto-SOM synapses in V1 and reported similar connection rates but opposite short-term dynamics --- pyramidal\rightarrowVIP synapses were \textit{depressing} while pyramidal\rightarrowSOM synapses were facilitating. \citet{Cauli2014} similarly described that CR/VIP bipolar interneurons receive predominantly depressing excitatory input from local pyramidal neurons \citep{Karnani2016a, Cauli2014, Walker2016, Pronneke2019}. Long-range thalamic and cortico-cortical inputs onto VIP cells are likewise dominated by facilitation: anterior thalamic input onto VIP cells in mouse presubiculum is overtly facilitating across paired-pulse ratios \citep{Nassar2025}, and POm input onto S1 VIP cells supports sustained activation rather than transient bursts \citep{Williams2019, Anastasiades2021a}. The functional consequence is that VIP cells preferentially read \textit{sustained or rhythmically-paced} inputs over transient ones --- a property aligned with their behavioural recruitment by tonic state variables (locomotion, arousal, reinforcement) rather than by isolated transient events \citep{Fu2014, Pakan2016, Sabri2024}.

Second, the \textit{output} synapse from VIP cells onto SST/Martinotti targets carries its own characteristic short-term plasticity. \citet{Cauli2014} documented that CR+/VIP+ bipolar interneurons produce facilitating IPSCs onto SOM Martinotti cells, and the IS3 VIP\rightarrowOLM synapse in CA1 is similarly facilitating during theta-rhythmic input \citep{Tyan2014, Francavilla2018a, Turi2019}. Recent paired-recording work by \citet{Mcfarlan2024a} demonstrates that VIP\rightarrowMartinotti and VIP\rightarrowbasket synapses in motor cortex undergo spike-timing-dependent plasticity, providing the first formal demonstration that the VIP-driven disinhibitory operation is itself plastic --- a finding that complicates any ``fixed canonical motif'' reading of cluster\_05 \citep{Mcfarlan2024a, Mcfarlan2024b}. Plasticity at the VIP\rightarrowSST synapse, in turn, places experience-dependent gain control inside the disinhibitory loop rather than only at its pyramidal output \citep{Mcfarlan2024a, Karnani2016a, Wiera2024, Sabri2024}.

Third, a substantial subset of cortical VIP/CR cells is \textit{co-transmitter}-positive. \citet{Granger2020} and \citet{Obermayer2019} showed that cortical ChAT/VIP cells ({\textasciitilde}33\% of VIP cells in mPFC) release ACh in addition to GABA, with target-specific GABA release that is robust onto interneurons and weaker onto pyramidal cells, and a previous study and an earlier paper demonstrated that the same cells provide a local cortical source of nicotinic excitation onto neighbouring fast-spiking and bipolar interneurons. The implication is that VIP cells are not only modulator-targeted but also, in part, \textit{modulator-providing}: a fraction of the cholinergic recruitment described in the previous subsection is intracortical rather than purely subcortical in origin \citep{Granger2020, Obermayer2019, Vonengelhardt2007, Dudai2020, Letzkus2011}. The peptide VIP itself is co-released and exerts neuromodulatory effects on layer 1 and on cerebrovasculature, but recent in vivo dissection by \citet{Kaneko2024} shows that at least one well-characterised circuit consequence of VIP cell activation --- stimulus-specific enhancement in V1 --- depends on GABA release rather than on peptide release. The triple --- source-dependent excitatory afferents (depressing locally, facilitating from long-range), plastic facilitating GABAergic outputs, and target-specific co-transmission --- together gives the canonical VIP\rightarrowSST\rightarrowPyr motif a distinctive temporal and learning signature that is not captured by treating it as a static schematic \citep{Karnani2016a, Cauli2014, Mcfarlan2024a, Granger2020, Kaneko2024, Walker2016, Pronneke2015}.

\section{Cell-type-specific recognition and the wiring rules behind the diagram}

A useful frame for the connectivity composite above is to ask \textit{what selects} it. Three lines of evidence indicate that the rules described in this section are not the byproduct of geometry but are encoded by cell-type-specific molecular programs that operate during development and remain stable in adulthood. \citet{Favuzzi2019} showed that the compartment specificity of the major interneuron classes --- PV\rightarrowsoma/AIS, SST\rightarrowdendrites, VIP\rightarrowother interneurons --- is determined by cell-type-specific cell-adhesion and wiring molecules whose expression patterns recapitulate the molecular taxonomy of \href{/development}{Developmental Origins and Postnatal Maturation} \citep{Tasic2018, Pronneke2019, Frazer2017, Favuzzi2019}. \citet{Batistabrito2017} further demonstrated that conditional ErbB4 deletion specifically in VIP interneurons, but not in PV or SST cells, disrupts cortical disinhibitory circuits and produces lasting changes in pyramidal cell excitability --- a result that ties the molecular identity of VIP cells (\href{/development}{Developmental Origins and Postnatal Maturation}) to the specific connectivity rules described here. Conditional microRNA-pathway deletion in VIP+ cells produces qualitatively similar circuit-level dysfunction \citep{Qiu2020, Favuzzi2019}. The connectivity diagram of Figure~\ref{fig-vip-connectome} is, in this sense, a developmental construct supported by molecularly specified wiring rules --- and small perturbations in those rules propagate to the network-level disinhibitory operation that subsequent sections will trace into behaviour and disease \citep{Favuzzi2019, Batistabrito2017, Qiu2020, Tasic2018, Pronneke2019, Simacek2025a}.

A complementary anatomical line of evidence comes from electron-microscopic and ultrastructural reconstruction. Symmetric synapses identifiable as VIP+ in immuno-EM target dendritic shafts and perisomatic compartments of CR+, CB+ and SOM+ interneurons in cat and rat visual cortex with substantially higher frequency than they target pyramidal somata \citep{Defelipe1999}. The pattern dovetails with the light-microscopic paired-recording data: ultrastructurally, VIP outputs converge on interneuron dendrites and proximal compartments rather than on pyramidal somata, and the small fraction of pyramidal-targeting boutons preferentially contacts apical dendrites in deep layers \citep{Defelipe1999}. The convergence of three independent methodological lines --- paired recording, optogenetic dissection and EM ultrastructure --- on the same compartment-specific output rule strengthens the case that the canonical schematic is correct \textit{as far as it goes}, while leaving open the layer- and area-specific exceptions documented in earlier subsections \citep{Hajos1996, Defelipe1999, Pfeffer2013, Yu2019, Anastasiades2021a, Walker2016}.

\section{What disagrees, and how the conflicts resolve}

The conflicts table for cluster\_05 contains 21 explicit pairwise disagreements; six of these recur often enough in the broader literature to deserve focused treatment, because each maps to a different axis along which the canonical motif breaks.

\textbf{(i) VIP\rightarrowPV connectivity strength.} The classical Pfeffer 2013 measurement of negligible VIP\rightarrowPV inhibition in V1 L2/3 is in tension with the Yu 2019 demonstration of direct VIP\rightarrowFS inhibition in barrel cortex deep layers \citep{Pfeffer2013, Yu2019}. The resolution from independent replications \citep{Walker2016, Pi2013, Anastasiades2021a} is that VIP\rightarrowPV connectivity is layer-dependent: superficial L2/3 microcircuits cleanly support the canonical interpretation, but L5/6 microcircuits --- and especially in non-V1 cortices --- recruit a direct VIP\rightarrowPV path that the canonical schematic misses \citep{Yu2019, Walker2016, Anastasiades2021a, Pi2013, Zhou2017}.

\textbf{(ii) Locomotion and the sign of SST modulation.} Fu 2014 reported the canonical locomotion-up-VIP / locomotion-down-SST pattern in V1 in darkness \citep{Fu2014}; Pakan 2016 and Dipoppa 2018 showed that this sign reverses or vanishes in the presence of structured visual stimulation, with SST cells co-activated during locomotion when visual drive is strong \citep{Pakan2016, Dipoppa2018, Kaneko2024}. The reading converged on by recent chemogenetic work is that locomotion-related VIP recruitment is itself preserved across these conditions, but its \textit{downstream consequence} on SST depends on the strength of visually driven SST input \citep{Sabri2024, Pakan2016, Dipoppa2018, Kaneko2024}.

\textbf{(iii) Long-range source matters.} Cingulate\rightarrowV1 recruits VIP \citep{Zhang2014, Stachniak2023}; auditory\rightarrowV1 suppresses VIP \citep{Ibrahim2016}. Direct optogenetic VIP perturbations during cross-modal attention do not reliably replicate the modulation reported in earlier indirect manipulations \citep{Ibrahim2016, Anastasiades2021a, Sabri2024}. Together these argue that ``top-down recruitment of VIP'' is not a single circuit operation but a family of source-specific operations, each with its own direction of effect \citep{Ibrahim2016, Anastasiades2021a, Stachniak2023}.

\textbf{(iv) Direct VIP\rightarrowpyramid connectivity.} The schematic claim that VIP cells inhibit only other interneurons is contradicted by Yu 2019, Zhou 2017 and several intersectional in vivo dissections \citep{Zhou2017, Anastasiades2021a, Walker2016}. The resolution is layer- and area-specific: in V1 L2/3 the canonical schematic remains a defensible first approximation, but in deep layers and in mPFC the direct VIP\rightarrowPyr path is part of the wiring and ignoring it produces predictions that fail in vivo \citep{Yu2019, Zhou2017, Anastasiades2021a, Walker2016, Pfeffer2013}.

\textbf{(v) Anatomical vs functional VIP\rightarrowPV connectivity.}

\begin{framed}
\textbf{Conflict --- Pfeffer 2013 vs Pi 2013 (anatomical vs functional)}\\
The Pfeffer 2013 paired-recording survey reported negligible \textit{functional} VIP\rightarrowPV connection probability in V1 L2/3 \citep{Pfeffer2013}, whereas anatomical / synaptic-counting work --- including the Pi 2013 framing of VIP cells as canonical disinhibitors of multiple interneuron classes --- implies that VIP+ axons do form symmetric synapses onto PV soma and proximal dendrites at non-trivial incidence \citep{Walker2016, Karnani2016a}. The disagreement is best read as \textit{anatomical incidence $\neq$ functional weight}: detectable VIP+ boutons onto PV cells can coexist with low paired-recording connection probability if those synapses are weak or release-deficient under the recording conditions used. Subsequent layer-resolved studies recover a measurable VIP\rightarrowPV functional path in deeper layers and in non-V1 areas \citep{Yu2019, Anastasiades2021a, Walker2016}, supporting the anatomical observation while preserving Pfeffer 2013 as a quantitatively correct L2/3-specific functional measurement.
\end{framed}

\textbf{(vi) Cholinergic excitation vs muscarinic inhibition on the same VIP class.}

\begin{framed}
\textbf{Conflict --- Poorthuis 2018 vs Arroyo 2012 (cluster\_13 cholinergic sign)}\\
Cluster\_13 of the cholinergic literature contains an explicit opposite-sign disagreement: \citet{Arroyo2012} and convergent work establish that the dominant ACh effect on cortical VIP / 5-HT3AR cells is nicotinic excitation through slow, $\beta$2-containing nAChRs \citep{Arroyo2012, Porter1999, Fu2014}, while \citet{Poorthuis2018} and matched studies emphasise a muscarinic \textit{inhibition} of a VIP subset \citep{Poorthuis2018, Obermayer2019, Granger2020}. The resolution recovers both observations on a layer/subtype axis: nicotinic excitation predominates in L2/3 VIP/CR bipolar cells (where 5-HT3AR co-expression is highest), whereas the muscarinic-inhibited subset is enriched in deeper layers and overlaps with intrinsically-bursting VIP cells. The two findings therefore refer to overlapping but non-identical VIP populations, and the section's modulator-pharmacology figure (Fig. 2) reflects this by encoding direction-of-effect rather than aggregate magnitude.
\end{framed}

A unifying observation across all six conflicts is that the canonical VIP\rightarrowSST\rightarrowPyr disinhibition is best read as a \textit{minimal sufficient} circuit motif rather than the unique circuit by which long-range and modulator inputs reach pyramidal output. Once one accepts that VIP outputs include direct PV and pyramidal inhibition in some layers and areas, that long-range afferents are a structured asymmetry rather than a uniform ``top-down\rightarrowVIP'' rule, and that neuromodulator pharmacology is mostly but not entirely excitatory, the apparent disagreements collapse to a smaller set of well-characterised quantitative gradients along axis variables (layer, area, source nucleus, behavioural state) that future work could in principle saturate \citep{Wall2016, Pfeffer2013, Yu2019, Pi2013, Lee2013, Karnani2016a, Sabri2024, Anastasiades2021a, Pronneke2019, Tasic2018}.

\section{Comparative and developmental refinement}

A final qualification concerns species and developmental stage. The connectivity statements above are derived overwhelmingly from adult mouse --- and a smaller number of adult rat --- preparations, and the textbook canonical schematic implicitly inherits that scope. Two points of evidence in the cluster\_05/cluster\_07 corpus warrant cautious extrapolation.

\textbf{Comparative anatomy}. Calretinin/VIP-positive bipolar interneurons are present in primate, tree-shrew and human neocortex, but their relative density, laminar distribution and apparent target preferences differ from rodent. Human and macaque temporal-cortex VIP/CR cells are concentrated in supra-granular layers, with substantially higher fractional density and a distinctive ``fan cell'' anatomy in primate prefrontal layer 1 not observed in rodent \citep{Defelipe1999, Vormsteinschneider2020}. A preceding study further demonstrate that intersectional viral tools developed in mouse can label CR/VIP cells in non-human primate and human cortex with conservation of the broad molecular scheme --- but the density and laminar layout differences mean that rodent-derived connectivity ratios should not be transferred quantitatively to human cortex without species-specific replication \citep{Defelipe1999}.

\textbf{Developmental maturation}. The connectivity diagram above is itself a developmental endpoint. \citet{Simacek2025a} showed that inhibitory inputs onto barrel-cortex VIP cells mature pre-synaptically across the second and third postnatal weeks, with release probability and quantal content rising from P8 to P36. Cell-type-specific connectivity rules --- VIP cells preferentially targeting other interneurons, PV cells targeting soma and SST cells targeting dendrites --- are themselves built during this developmental window by cell-type-specific recognition programs and activity-dependent refinement \citep{Favuzzi2019, Batistabrito2017, Qiu2020, Simacek2025a}. Disrupting this refinement --- by conditional deletion of ErbB4 in VIP cells \citep{Batistabrito2017}, by microRNA deletion or by neonatal hypoxic-ischaemic injury --- produces persistent miswiring of the disinhibitory motif into adulthood. The connectivity rules described in earlier subsections are therefore not given but built; they are what an experience-tuned, modulator-targeted cell class looks like once the postnatal program has run to completion \citep{Favuzzi2019, Batistabrito2017, Qiu2020, Simacek2025a, Abbah2022}.

These two qualifications do not undermine the core synthesis but they constrain its generalisation. The minimal sufficient circuit model that \href{/in-vivo-behavior}{In Vivo Function During Behavior} will need --- for autism, schizophrenia, neonatal injury and epilepsy --- must accommodate developmental construction of VIP connectivity, and the comparative literature demands that quantitative rodent ratios be flagged as such when extrapolated to human cortex \citep{Vormsteinschneider2020, Batistabrito2017, Qiu2020, Simacek2025a, Favuzzi2019}.

\section{Synthesis: VIP cells as modulator-targeted, area-tuned disinhibitors}

The synaptic-connectivity picture that emerges from 467 papers is a structured composite. \textbf{On the input side}, VIP interneurons are a preferential target of long-range cortical and higher-order thalamic glutamatergic input, of basal-forebrain cholinergic input via slow nicotinic receptors, of raphe serotonergic input via 5-HT3 receptors, and of mixed-sign noradrenergic input --- while local pyramidal drive is comparable to that onto SST and PV cells \citep{Wall2016, Lee2013, Zhang2014, Williams2019, Anastasiades2021a, Cruikshank2012, Audette2018, Stachniak2023, Bilash2023, Porter1999, Kawaguchi1997b, Kawaguchi1998, Arroyo2012, Fu2014, Ferezou2002, Vucurovic2010, Frazer2017, Letzkus2011}. \textbf{On the output side}, VIP cells preferentially inhibit SST interneurons, but they also inhibit PV cells in deeper layers, NDNF/L1 interneurons, pyramidal cells in deep cortex and OLM cells in CA1, with substantial cell-to-cell heterogeneity that the canonical schematic flattens \citep{Pfeffer2013, Pi2013, Lee2013, Karnani2016a, Karnani2016b, Walker2016, Yu2019, Zhou2017, Anastasiades2021a, Schuman2019, Hafner2021, Tyan2014, Turi2019, Francavilla2018a, Booker2018}. \textbf{On the behavioural side}, the consequence of this wiring is a deployable disinhibitory operation that is recruited by arousal and reinforcement signals across V1, A1, mPFC, BLA and CA1, but which is sign-dependent on visual context, source-area-specific in cross-modal attention, and layer-dependent in its non-SST outputs \citep{Fu2014, Pi2013, Letzkus2011, Krabbe2019, Pakan2016, Dipoppa2018, Sabri2024, Kamigaki2017, Pinto2015, Szadai2022, Tamboli2024, Ibrahim2016, Kaneko2024, Anastasiades2021a}. The textbook disinhibitory schematic is a defensible projection of this composite onto a single L2/3 V1 plane; it fails as a description of the full population.

This synthesis sets up two tasks for the remaining sections of the review. First, the \textit{behavioural and developmental specificity} of the disinhibitory operation --- which arousal axes, which task contingencies, which developmental windows actually engage VIP\rightarrowSST\rightarrowPyr --- must be related to the molecular subtypes catalogued in \href{/development}{Developmental Origins and Postnatal Maturation} and the intrinsic firing patterns catalogued in \href{/electrophysiology}{Intrinsic Electrophysiology} \citep{Tasic2018, Pronneke2019, Pronneke2015, Frazer2017, Lee2010, Rudy2011}. Second, the \textit{minimum sufficient circuit model} needed to explain disease-relevant phenotypes (autism, schizophrenia, anxiety, neonatal injury) cannot rely on canonical disinhibition alone; it must accommodate the layer-dependent direct VIP\rightarrowPyr/PV path, the area-tuned long-range asymmetries and the modulator-recruited gain control reviewed here \citep{Yu2019, Zhou2017, Anastasiades2021a, Letzkus2011, Sabri2024, Krabbe2019}. \href{/disinhibition-framework}{Local Circuit Motifs and the Disinhibition Framework} takes up the first task by reviewing how VIP cells participate in cortical state, attention and learning across these areas; \href{/in-vivo-behavior}{In Vivo Function During Behavior} takes up the second by tracing the developmental and disease vulnerabilities of the synaptic and modulatory architecture described here.

A useful operational summary is the following. \textit{Treat the VIP\rightarrowSST\rightarrowPyr disinhibitory motif as a layer- and area-specific computation that a cortical microcircuit can perform when modulator-targeted afferents are active, not as a fixed wiring rule that all VIP cells obey under all conditions}. This formulation accommodates the canonical landmark results \citep{Pfeffer2013, Pi2013, Lee2013, Fu2014, Karnani2016a, Karnani2016b}, the long-range and thalamic asymmetries \citep{Wall2016, Williams2019, Anastasiades2021a, Audette2018, Stachniak2023, Bilash2023, Cruikshank2012}, the modulator pharmacology \citep{Porter1999, Kawaguchi1997b, Arroyo2012, Vonengelhardt2007, Granger2020, Obermayer2019, Letzkus2011, Ferezou2002, Vucurovic2010, Frazer2017}, the layer- and area-specific exceptions to canonical disinhibition \citep{Yu2019, Zhou2017, Anastasiades2021a, Walker2016}, the developmental and comparative qualifications \citep{Favuzzi2019, Batistabrito2017, Qiu2020, Simacek2025a, Vormsteinschneider2020, Defelipe1999}, and the hippocampal divergence \citep{Tyan2014, Turi2019, Francavilla2018a, Booker2018, Bilash2023, Tamboli2024, Toth1993} within a single coherent picture. \href{/disinhibition-framework}{Local Circuit Motifs and the Disinhibition Framework} and \href{/in-vivo-behavior}{In Vivo Function During Behavior} inherit this picture intact.


\section{Local Circuit Motifs and the Disinhibition Framework}
\label{sec:07_disinhibition_framework}

The synaptic inventory assembled in \href{/synaptic-properties}{Synaptic Properties and Connectivity} --- strong long-range and neuromodulatory drive onto VIP cells, asymmetric output that favours somatostatin (SST) interneurons but does not spare parvalbumin (PV) cells or pyramidal neurons --- supplies the parts list for a circuit-level reading. The question this section addresses is what circuit motif those parts implement, and whether the standard answer survives contact with the data. The standard answer, repeated in nearly every textbook diagram of cortical inhibition, is the VIP\rightarrowSST\rightarrowpyramidal disinhibition motif: VIP cells, recruited by behavioural state, suppress SST interneurons, which in turn release pyramidal dendrites from inhibition and so enhance gain \citep{Pfeffer2013, Pi2013, Lee2013, Fu2014, Karnani2016a}. Across visual, auditory, somatosensory, prefrontal, and motor cortex, optogenetic and chemogenetic manipulations of VIP cells have produced phenomena consistent with this motif \citep{Pi2013, Lee2013, Fu2014, Zhang2014, Kuchibhotla2017, Cichon2017, Williams2019, Veit2023}. Yet the same body of work also exposes its limits: the same VIP perturbation yields different sign and magnitude of pyramidal modulation across area, layer, and behavioural state, and inhibition of SST cells is neither necessary nor sufficient to explain observed gain effects in several preparations \citep{Pakan2016, Dipoppa2018, Garrett2020, Stachniak2021}. The thesis of this section is that the disinhibitory motif is \textit{one} privileged channel within a small family of parallel VIP-mediated circuits --- including direct VIP\rightarrowPV, VIP\rightarrowpyramidal, and VIP\rightarrowNDNF/OLM connections --- whose relative weights depend on cortical area, layer, behavioural state and the specific VIP subtype recruited. Treating disinhibition as the \textit{defining} function of the class is no longer supported by the contemporary evidence; treating it as a reproducible \textit{component} function is.

\section{The disinhibitory motif: how it was constructed}

The motif we now recite as disinhibitory was constructed from a tight cluster of 2013--2014 papers that combined cell-type-specific optogenetics, paired patch recordings, and \textit{in vivo} two-photon imaging in mouse neocortex. \citet{Pfeffer2013} mapped the inhibitory connections among PV, SST, and VIP cells in V1 layers 2/3 and 5 and reported a \textit{complementary} connectivity scheme: PV cells preferentially inhibited PV cells and pyramidal somata, SST cells inhibited PV and pyramidal cells, and VIP cells were strongly biased onto SST cells. Quantitatively, optogenetic activation of VIP cells produced inhibitory postsynaptic charge in simultaneously recorded SST cells that was an order of magnitude larger than in pyramidal cells, the disinhibitory complementary-connectivity result \citep{Pfeffer2013}. Independently, \citet{Pi2013} showed in mouse auditory cortex that VIP cells were activated by reinforcement signals and that VIP-cell activation released pyramidal neurons from SST-mediated inhibition, identifying suppression of SST and PV interneurons as the central mechanism. \citet{Lee2013} identified a homologous motif in mouse barrel cortex (S1), where activation of long-range corticocortical inputs from primary motor cortex (vM1) suppressed SST cells via VIP and enhanced pyramidal responses. \citet{Fu2014} then provided the locomotion analogue: in awake V1, running drove VIP cells through nicotinic cholinergic input from basal forebrain, suppressed SST cells, and enhanced visual gain in pyramidal neurons. Together, these four papers established three core claims --- VIP cells preferentially inhibit SST cells, VIP cells are recruited by salient external inputs, and VIP-mediated SST suppression enhances pyramidal gain --- that have since been collapsed into a single schematic. The schematic is genuinely supported by the underlying data; the difficulty is that the data also support several adjacent and partly overlapping motifs, and the schematic erases that ambiguity.

A few features of this construction deserve emphasis. First, the disinhibitory schematic was tested predominantly in superficial layers of mouse V1 and primary auditory cortex; PFC, M1, and S1 evidence followed and is more heterogeneous (see \href{/cross-areas}{VIP Interneurons Across Brain Regions}). Second, the disinhibitory recruitment story is a \textit{composite}: top-down cortical glutamatergic drive \citep{Lee2013, Zhang2014} and ascending cholinergic and noradrenergic neuromodulation \citep{Fu2014, Hangya2015, Kuchibhotla2017} were originally described in different papers and different behavioural contexts but are now often spoken of as a single mechanism. Third, the original measurements of VIP\rightarrowpyramidal effects were almost all \textit{net} effects --- measured as changes in pyramidal firing or membrane potential after VIP-cell perturbation --- and so confounded direct effects (VIP\rightarrowPyr inhibition), indirect effects (VIP\rightarrowSST\rightarrowPyr disinhibition), and parallel routes through PV cells \citep{Karnani2016a, Karnani2016b, Yu2019}. Disentangling these contributions is a central concern of the rest of this section and of the quantitative comparisons in Figure~\ref{fig-vip-effect-size-comparison}.

\section{Recruitment: top-down glutamatergic versus neuromodulatory drive}

The recruitment side of the disinhibitory motif is settled in outline and unsettled in detail. In outline: VIP cells in adult cortex are reliably activated by behavioural events that combine arousal, locomotion, and salient external stimuli \citep{Fu2014, Pinto2015, Pakan2016, Dipoppa2018, Garrett2020}. In detail: at least three partially overlapping recruitment routes have been documented and their relative weights are area- and state-dependent. \citet{Lee2013} emphasised top-down glutamatergic projections from primary motor cortex (vM1) onto VIP cells in S1, and \citet{Zhang2014} documented analogous projections from cingulate areas onto V1, where they activate VIP cells as part of the local disinhibitory circuit. \citet{Fu2014} emphasised nicotinic cholinergic drive from basal forebrain that depolarised VIP cells during locomotion. \citet{Pi2013} reported that reinforcement signals activate VIP neurons in auditory cortex via long-range projections, while \citet{Hangya2015} demonstrated that basal-forebrain cholinergic neurons fire phasically to behaviourally relevant cues in a manner consistent with recruiting VIP cells through nicotinic receptors. Pharmacological and optogenetic isolation of cholinergic transmission has since shown that VIP cells across V1, A1, S1, and PFC respond rapidly to ACh release through $\alpha$4$\beta$2 and $\alpha$7 nicotinic receptors \citep{Pinto2015, Kuchibhotla2017, Williams2019, Anastasiades2021a, Mcfarlan2024a}, with response amplitudes that can rival those produced by direct optogenetic stimulation of long-range glutamatergic axons.

These two routes are typically presented as alternatives, but the experimental evidence is most cleanly read as showing that they cooperate. \citet{Pinto2015} recorded VIP and SST cells during a frontal-cortex task and characterised cell-type-specific activity profiles during cognitively demanding behaviour. \citet{Kamigaki2017} reported that PFC VIP cells in delayed-response tasks integrated cue and reward-related inputs that arrived through both pathways. Computational reconciliation comes from the membrane-potential analysis of \citet{Garciadelmolino2017} and \citet{Hertag2019}, in which subthreshold cholinergic depolarisation gates the gain of VIP responses to top-down glutamatergic input rather than driving spikes on its own --- an `AND' rather than an `OR' recruitment rule. The downstream consequence is that the disinhibitory motif is engaged most strongly when external salience and internal arousal coincide, and only weakly engaged by either alone. This reconciles the apparent disagreement between \citet{Pi2013} and \citet{Fu2014} over which input is `primary': in their respective preparations and behavioural epochs, each was the dominant driver, but neither study tested the alternate input under the other's task structure.

\begin{framed}
\textbf{Conflict}\\
\textbf{Top-down glutamatergic vs neuromodulatory cholinergic drive as the primary VIP recruiter.} \citet{Lee2013} reported in mouse barrel cortex (S1) that long-range glutamatergic input from primary motor cortex (vM1) was a dominant route for VIP recruitment driving disinhibition, with cholinergic blockade leaving the effect largely intact. \citet{Pi2013} likewise emphasised long-range cortical drive in auditory cortex, while \citet{Fu2014} reported that nicotinic cholinergic input from basal forebrain was sufficient to drive locomotion-evoked VIP recruitment in primary visual cortex. \textbf{Resolution:} partial. The conflict largely dissolves once preparation and behavioural epoch are matched --- top-down drive dominates during attentional or top-down sensory tasks, cholinergic drive during locomotion-coupled arousal --- and joint manipulation experiments \citep{Pinto2015, Kamigaki2017, Williams2019} show the two routes are concurrently active and at least partly multiplicative \citep{Hertag2019}. The framing of either route as exclusive is not supported by the joint data.
\end{framed}

A second recruitment-side conflict bears on whether the disinhibitory motif is \textit{the} privileged channel for behavioural-state gain. \citet{Pi2013} framed VIP-mediated disinhibition as a principal gating mechanism for cortical gain in auditory cortex and mPFC. \citet{Kuchibhotla2017} later showed in the same cortical area that cholinergic activation simultaneously potentiates direct PV-mediated inhibition of pyramidal cells and recruits VIP-mediated disinhibition through an interaction with SST cells, producing a `double-flip' in which inhibition is \textit{re-routed} rather than simply removed. The implication is that the disinhibitory motif coexists with at least one parallel inhibitory stream --- a parallelism that is invisible in single-pathway perturbation experiments but visible whenever both arms are measured in the same recording.

\begin{framed}
\textbf{Conflict}\\
\textbf{Is VIP disinhibition \textit{the} privileged channel for behavioural-state gain?} \citet{Pi2013} framed VIP\rightarrowSST disinhibition as the principal gating mechanism through which cingulate input enhances cortical responses. \citet{Kuchibhotla2017} reported that cholinergic activation of the same cortex simultaneously engages VIP-mediated disinhibition and PV-mediated direct inhibition, producing a re-routing of inhibition rather than a simple gain elevation. \textbf{Resolution:} the two are not contradictory but the privileged-channel framing is. Behavioural-state gain operates through several parallel inhibitory streams, of which the disinhibitory motif is one; whether VIP disinhibition is the dominant stream depends on stimulus, area, and state \citep{Pakan2016, Dipoppa2018, Veit2023}.
\end{framed}

\section{Connectivity asymmetry: how reproducible is VIP\rightarrowSST \textgreater  VIP\rightarrowPV/Pyr?}

The connectivity asymmetry --- VIP\rightarrowSST stronger than VIP\rightarrowPV or VIP\rightarrowpyramidal --- is the load-bearing claim of the disinhibitory motif, and it is precisely where heterogeneity is largest. The original optogenetic-plus-paired-recording survey of \citet{Pfeffer2013} reported that VIP\rightarrowSST connections were the majority among the VIP-output classes tested in mouse V1 paired recordings, with much weaker connectivity onto PV cells and pyramidal somata. \citet{Pi2013}, using subcellular optogenetic mapping in auditory cortex, instead emphasised that VIP cells produce kinetically and anatomically distinct inhibition onto SST and PV postsynaptic targets --- VIP-evoked IPSCs onto SST cells exhibited substantially slower decay than those onto PV cells (Pi2013; the underlying decay-constant values are not deposited as a curated row-level finding) --- and reported that VIP cells \textit{also} substantially inhibit PV cells, in apparent tension with the strict bias of \citep{Pfeffer2013}. \citet{Lee2013}, using VIP-specific optogenetics in S1 slices, recorded VIP-evoked IPSC amplitudes that were several-fold larger in SST postsynaptic cells than in pyramidal cells (sample sizes were not extracted into the curated row for that finding), supporting the SST bias qualitatively but with absolute amplitudes much larger than those reported by \citet{Pfeffer2013} --- a discrepancy attributable to opsin dose, recording conditions, and pyramidal subtype. \citet{Karnani2016a} and \citet{Karnani2016b} then showed that, when measured at higher resolution, VIP cells in V1 cooperatively suppress SST tone with kinetics fast enough to produce behaviourally relevant disinhibition of pyramidal somata and dendrites, and that volley VIP activation can transiently lift the SST blanket from large groups of pyramidal cells. More recent EM-based connectomic reconstructions in mouse V1 confirm that the VIP\rightarrowSST edge is over-represented relative to a null wiring model, but they also resolve substantial VIP\rightarrowpyramidal and VIP\rightarrowPV outputs that were not visible in the lower-throughput surveys \citep{Schneidermizell2025}.

The disagreement between \citet{Pfeffer2013}, who reported a strict VIP\rightarrowSST bias, and \citet{Pi2013}, who reported that VIP also inhibits PV cells with shorter, faster IPSCs, is best read as a difference of metric and granularity rather than of underlying biology. \citet{Pfeffer2013} reported pooled connection probabilities in paired recordings, in which VIP\rightarrowPV pairs do exist but are rare; \citet{Pi2013} reported within-cell IPSC waveforms in cells that received any VIP input, normalising away the rarity. Both can be true simultaneously: VIP\rightarrowSST connections dominate in probability, while VIP\rightarrowPV connections, when present, carry kinetically distinct, fast-onset IPSCs.

\begin{framed}
\textbf{Conflict}\\
\textbf{Specificity of VIP\rightarrowSST inhibition versus broader VIP\rightarrowPV inhibition.} \citet{Pfeffer2013} reported a strict VIP\rightarrowSST bias in mouse V1 layers 2/3 and 5, with VIP\rightarrowPV connectivity essentially negligible at the population level. \citet{Pi2013}, in mouse auditory cortex, reported that VIP cells inhibit both SST and PV cells with kinetically distinct IPSCs (decay $\tau$ $\approx$ 18 ms onto SST, $\approx$ 6 ms onto PV; Pi 2013, error metric `$\pm$ as printed', n unrecorded in the curated entry). \textbf{Resolution:} the metrics differ --- connection probability versus within-pair IPSC waveform --- and both can hold simultaneously \citep{Karnani2016a, Karnani2016b, Yu2019, Schneidermizell2025}. The strict-bias framing is not supported once connectomic and within-cell electrophysiological data are combined; SST bias is a tendency, not an absolute.
\end{framed}

A complementary disagreement concerns VIP\rightarrowpyramidal connections. \citet{Pfeffer2013} reported VIP\rightarrowpyramidal IPSCs roughly an order of magnitude smaller than VIP\rightarrowSST IPSCs in matched paired recordings. Direct VIP\rightarrowpyramidal inhibition has been reported in some preparations as a complement to disinhibition, although the dominant mode reported in \citet{Pi2013} for auditory cortex is disinhibition through suppression of SST and PV cells. The original \citet{Lee2013} data, with measurable but smaller VIP\rightarrowpyramidal IPSCs, sit between these extremes. \citet{Yu2019} showed, using in vivo electrophysiology during tactile behaviour, that interneuron recruitment patterns vary systematically across cortical states and across preparations. The picture from connectomics in primary visual cortex \citep{Schneidermizell2025} and from cross-species comparisons \citep{Pronneke2015, Vanderveer2021} is that VIP cells in different cortical areas and species have different output ratios, so there is no single answer to the question `how strong is VIP\rightarrowpyramidal inhibition?' in the way the disinhibitory schematic implies. Figure~\ref{fig-vip-effect-size-comparison} makes the heterogeneity quantitatively visible.

\begin{framed}
\textbf{Conflict}\\
\textbf{Disinhibition versus direct PV recruitment as the dominant mode of VIP\rightarrowpyramidal control.} \citet{Pi2013} argued that VIP cells operate primarily through SST disinhibition. \citet{Pfeffer2013} showed strong VIP\rightarrowSST connectivity but emphasised the limited VIP\rightarrowPV/Pyr coupling. \textbf{Resolution:} both descriptions are locally accurate, but the modal answer depends on cortical area and behavioural state --- disinhibition dominates in V1 and S1 superficial layers under top-down drive \citep{Zhang2014, Lee2013}; PV recruitment becomes important in cortices and tasks in which cholinergic drive is strong \citep{Kuchibhotla2017, Williams2019} and in deep-layer connectomic reconstructions where direct VIP\rightarrowpyramidal contacts are visible \citep{Schneidermizell2025}.
\end{framed}

\begin{framed}
\textbf{Conflict}\\
\textbf{Magnitude of direct VIP\rightarrowpyramidal inhibition relative to VIP\rightarrowSST disinhibition.} \citet{Lee2013} reported VIP-evoked IPSC amplitudes onto pyramidal cells in mouse S1 slice that, while smaller than those onto SST cells (sample sizes were not deposited in the curated row for that entry), were large enough to imply a non-trivial direct VIP\rightarrowpyramidal channel alongside the dominant VIP\rightarrowSST channel. \citet{Pi2013} emphasised disinhibition through suppression of SST and PV interneurons as the dominant route for VIP-mediated control of pyramidal firing in auditory cortex. \textbf{Resolution:} the two reports agree that VIP cells affect pyramidal firing in behaviourally relevant ways; they disagree on whether direct VIP\rightarrowpyramidal inhibition is large enough to compete with the disinhibitory channel. The contemporary reading \citep{Karnani2016a, Karnani2016b, Yu2019, Schneidermizell2025} is that the direct channel is detectable in most preparations but is dominated, on average, by the disinhibitory channel; net pyramidal modulation reflects the balance and can flip sign with task condition.
\end{framed}

\section{Behavioural recruitment of the disinhibitory motif}

Behavioural engagement of the disinhibitory motif is the part of the framework with the strongest cross-laboratory replication. In V1, locomotion drives VIP cells to depolarise by tens of percent above baseline $\Delta$F/F at running onset (Fu 2014; sample size unrecorded in the curated entry) and to suppress concurrently imaged SST cells \citep{Fu2014, Garciadelmolino2017, Dipoppa2018}. In auditory cortex, foot-shock and noise activate cingulate-driven VIP cells and produce SST suppression that scales with stimulus salience \citep{Pi2013}. In M1, learning-related cholinergic input to VIP cells correlates with dendritic plasticity in projection pyramidal neurons. In S1, top-down feedback in object detection and texture discrimination tasks recruits VIP cells whose suppression reduces task performance \citep{Sermet2019}. In PFC, attention and working-memory tasks recruit VIP cells whose silencing impairs performance \citep{Kamigaki2017}. In hippocampus and amygdala, analogous VIP-driven disinhibitory motifs have been described in fear learning and memory updating \citep{Krabbe2019, Turi2019, Tamboli2024}. The behavioural engagement is robust enough that the disinhibitory motif is the right first hypothesis whenever a salient external event coincides with arousal in a cortical task --- but, as the next subsection shows, `engagement' is not the same as `sufficient explanation'.

The link from VIP recruitment to pyramidal gain enhancement is mediated, in the disinhibitory motif, by suppression of SST cells. \citet{Zhang2014} provided one of the cleanest causal demonstrations: optogenetic activation of cingulate projections in mouse V1 engaged VIP cells, enhanced pyramidal stimulus-evoked firing in the receptive-field centre, and was consistent with relief from SST-mediated suppression. The direction of this effect is consistent across \citep{Lee2013, Fu2014, Karnani2016a, Pinto2015, Veit2023}. The magnitude is highly variable: from a few-percent gain change in some preparations to multi-fold rate changes in others (see Figure~\ref{fig-vip-effect-size-comparison}). The variability is attributable to (i) cortical area and layer, (ii) opsin dose and behavioural state during the perturbation, (iii) co-recruitment of VIP-mediated direct inhibition of pyramidal cells \citep{Karnani2016b, Lee2013, Pi2013}, and (iv) co-recruitment of NDNF/Layer-1 interneurons whose top-down drive partially overlaps with that of VIP cells \citep{Abs2018, Anastasiades2021a, Veit2023}.

\section{Cross-area generalisation and its limits}

The disinhibitory motif was constructed in mouse V1 and A1 and is most strongly supported there. Generalisation to other cortices is partial. In S1, the motif holds in object-localisation and active touch tasks \citep{Williams2019, Sermet2019}, but layer-specific variants have been described in which thalamic inputs to deep-layer VIP cells engage a different downstream target set \citep{Munoz2017, Sermet2019, Garciajuncoclemente2017}. In PFC, the motif is engaged in attention and working-memory tasks \citep{Kamigaki2017, Pinto2015, Anastasiades2021a}, but PFC VIP cells receive a richer mix of long-range inputs and project to a broader set of postsynaptic targets, including direct VIP\rightarrowpyramidal connections that may be more prominent than in sensory cortex \citep{Anastasiades2021a, Schneidermizell2025}. In M1, VIP recruitment is tightly coupled to cholinergic drive during motor learning, and the disinhibitory motif appears to gate dendritic spine plasticity in projection pyramidal neurons \citep{Cichon2017}. In higher visual areas, the locomotion-evoked VIP activation that is so reproducible in V1 is \textit{not} uniformly disinhibitory: in mouse anterolateral and lateromedial areas, locomotion produces SST recruitment of comparable magnitude to VIP recruitment, complicating the simple disinhibition story \citep{Garrett2020}. In hippocampus, VIP cells implement a structurally distinct motif: rather than VIP\rightarrowSST\rightarrowPyr disinhibition, CA1 VIP cells preferentially target SST OLM cells and PV basket cells, producing temporally precise theta-rhythmic disinhibition that gates input-specific plasticity \citep{Turi2019}. In primate and human cortex, comparative anatomy and Patch-seq data suggest that the VIP\rightarrowSST bias is preserved but that the absolute density and laminar organisation of VIP cells differ substantially \citep{Pronneke2015, Vanderveer2021}, with implications discussed further in \href{/cross-areas}{VIP Interneurons Across Brain Regions}.

The most consequential cross-area conflict is whether the disinhibitory motif generalises in sign across cortical areas and behavioural contexts. \citet{Lee2013} reported in mouse S1 that top-down recruitment of VIP cells by vM1 input suppressed SST cells and enhanced pyramidal responses. \citet{Pakan2016}, imaging the same superficial layer of V1 across a wider behavioural range, reported that locomotion-evoked changes in VIP and SST activity were \textit{not} simple mirrors of one another: SST cells were not uniformly suppressed during locomotion, and the SST response depended on visual input and cortical area. \citet{Dipoppa2018}, recording across layers and areas, similarly found that locomotion can suppress \textit{or} enhance SST cells depending on visual context, with VIP recruitment a less reliable predictor of SST suppression than the disinhibitory motif assumes.

\begin{framed}
\textbf{Conflict}\\
\textbf{Does the disinhibitory VIP\rightarrowSST disinhibition motif generalise across cortical areas and behavioural contexts?} \citet{Lee2013} reported in mouse S1 that top-down VIP recruitment by vM1 input consistently suppressed SST cells and enhanced pyramidal gain in the conditions tested. \citet{Pakan2016} reported, in the same cortical area but a broader behavioural range, that locomotion-evoked VIP and SST activity were not simple mirrors and that SST cells were not uniformly suppressed during locomotion. \textbf{Resolution:} the disinhibitory motif holds as a \textit{component} explanation in V1 superficial layers under top-down or cholinergic drive but does not generalise as a \textit{sole} explanation across all behavioural contexts; cortical area, layer, and visual input together determine whether SST cells follow the disinhibition trajectory or move independently \citep{Dipoppa2018, Garrett2020, Stachniak2021, Veit2023}. Locomotion-coupled VIP activity in V1 is reproduced across studies \citep{Fu2014, Pakan2016, Dipoppa2018}, but its downstream consequence for SST cells is not.
\end{framed}

\section{Alternative motifs}

The cluster of motifs that share the disinhibitory VIP\rightarrowSST\rightarrowPyr backbone but differ in their exits and entries is more interesting than the disinhibitory schematic alone. Four families are well documented in the literature reviewed here and are sketched in Figure~\ref{fig-disinhibition-motifs}.

\textbf{Direct VIP\rightarrowpyramidal inhibition.} A subset of VIP cells, most prominently the calretinin-positive, bipolar morphological type, contact pyramidal somata and proximal dendrites and produce direct inhibitory currents large enough to suppress pyramidal firing under volley activation \citep{Karnani2016a, Schneidermizell2025}. In auditory cortex, this motif appears to be co-engaged with the disinhibitory motif during salient stimulus presentation \citep{Pi2013, Kuchibhotla2017}. The biological consequence is that the \textit{net} sign of VIP perturbation on pyramidal firing reflects the balance between disinhibitory gain enhancement and direct inhibitory suppression, and can flip sign with task condition or stimulus drive \citep{Garciadelmolino2017, Hertag2019}.

\textbf{VIP\rightarrowPV control of perisomatic inhibition.} VIP\rightarrowPV connections, although less probable than VIP\rightarrowSST, are kinetically distinct (faster decay; \citep{Pi2013}) and target perisomatic inhibition of pyramidal cells. Volley VIP activation can also engage perisomatic disinhibition through additional polysynaptic pathways, providing a putative \textit{second} disinhibitory channel that operates on a faster timescale than the dendritic VIP\rightarrowSST channel. In some cortical areas, this faster channel may dominate: \citet{Kuchibhotla2017} reported that the cholinergic-driven re-routing of inhibition in auditory cortex is mediated in part by VIP\rightarrowPV and not solely by VIP\rightarrowSST.

\textbf{VIP\rightarrowNDNF / VIP\rightarrowlayer-1 motifs.} Layer-1-resident interneurons expressing neuron-derived neurotrophic factor (NDNF) gate top-down input to apical pyramidal dendrites in a layer-specific manner \citep{Abs2018, Anastasiades2021a, Veit2023}. VIP cells contact and inhibit NDNF cells, providing an additional disinhibitory route that is engaged in fear learning, attention, and predictive processing tasks \citep{Krabbe2019, Veit2023}. In auditory cortex, the VIP\rightarrowNDNF arm is required for fear-conditioning-related plasticity in apical dendrites. In V1, NDNF cells provide a top-down gain control that operates in parallel with VIP\rightarrowSST \citep{Veit2023}.

\textbf{Hippocampal VIP\rightarrowOLM and VIP\rightarrowbasket-cell disinhibition.} In CA1 and dentate gyrus, VIP cells, including the calretinin-positive and CCK-positive subsets, form a structurally distinct motif: VIP/CR cells inhibit OLM cells (the SST analogue) and basket cells (the PV analogue) in a temporally precise, theta-rhythmic manner that gates Schaffer-collateral and entorhinal input streams \citep{Turi2019}. The hippocampal motif is recognisably a homologue of the cortical disinhibitory motif but with a different oscillatory and plasticity context, and it points toward a class of \textit{disinhibition motifs} of which the cortical disinhibitory motif is a particular instance.

In addition to these four families, a number of motif fragments are documented but less developed in the literature: VIP\rightarrowMartinotti subtype-specific connections in deep layers \citep{Anastasiades2021a, Schneidermizell2025}, VIP\rightarrowlong-range projecting interneuron contacts, and VIP-cell mediated inhibition of other VIP cells through CCK-positive subtypes \citep{Garrett2020}. The taxonomy of these alternatives is unsettled, but the empirical fact that they exist and are behaviourally engaged is settled.

The decision tree mapping behavioural context to the dominant VIP motif is sketched in Figure~\ref{fig-disinhibition-motifs} (Panel C). It is not a definitive ordering: the literature does not yet contain enough joint manipulations of multiple motifs in the same animal under matched task conditions to assign weights with confidence. What it does support is the qualitative claim that \textit{which} motif dominates is set by behavioural context, area, and developmental state, and that the disinhibitory motif is one branch of the tree, not the trunk.

\section{Quantitative cross-study comparison: forest plots and the conflict atlas}

Figure~\ref{fig-vip-effect-size-comparison} makes the heterogeneity quantitative. Panel A places the VIP\rightarrowSST IPSC measurements of \citet{Pfeffer2013}, \citet{Lee2013}, and \citet{Pi2013} on a single axis. Their absolute amplitudes differ by roughly an order of magnitude --- a difference attributable to opsin choice, light dose, recording configuration (paired versus optogenetic-ChR2 over a population), and cortical area --- but their \textit{relative} claim, that VIP\rightarrowSST charge or amplitude exceeds VIP\rightarrowpyramidal in the same preparation, is reproduced across all three studies. Panel B places VIP\rightarrowPV connectivity measurements alongside VIP\rightarrowSST, again with within-study consistency that VIP\rightarrowSST exceeds VIP\rightarrowPV but with substantial cross-study variation in the absolute VIP\rightarrowPV value. Panel C aggregates the \textit{in vivo} effect of VIP perturbation on pyramidal firing across V1 studies and shows that the disinhibitory net-disinhibitory effect is reproducible in V1 superficial layers under top-down or running-coupled drive, but is reduced or absent in higher visual areas and in stationary visual paradigms \citep{Pakan2016, Garrett2020, Dipoppa2018}. Panel D presents the conflict atlas: a heat-map of the six conflict pairs surfaced in this section, showing that the residual disagreement is concentrated on motif specificity (VIP\rightarrowSST versus VIP\rightarrowSST + VIP\rightarrowPV) and on cross-area generalisation, not on the existence of the disinhibitory motif itself.

The composition fact that constrains all of these comparisons is that VIP cells are a small minority of cortical interneurons. \citet{Rudy2011} reported that approximately 40\% of mouse neocortical GABAergic neurons are PV+, approximately 30\% are SST+, and approximately 30\% are 5-HT3AR+, with VIP+ cells comprising a subset of the 5-HT3AR+ class. The numerical asymmetry is part of why the disinhibitory motif works: a small VIP population can disinhibit a much larger pyramidal population if its synaptic targeting is sufficiently selective, but the same numerical asymmetry constrains how much absolute gain can be produced through any single motif and forces the engagement of multiple parallel motifs to scale gain with task demand \citep{Garciadelmolino2017, Litwinkumar2016}.

\section{Figures}

\begin{figure}[!htbp]
\centering
\includegraphics[width=0.95\linewidth]{files/fig-disinhibition-mo-f689999321fc033d1f37c55a2bc779f2.png}
\caption[]{\textbf{The disinhibition motif and its alternatives.} \textbf{(A)} Canonical motif: VIP cells, recruited by long-range glutamatergic and ascending cholinergic inputs, preferentially inhibit SST cells, releasing pyramidal dendrites from inhibition; minor VIP\rightarrowPV and VIP\rightarrowpyramidal edges are drawn explicitly to flag that the disinhibitory schematic is a tendency rather than an exclusive wiring rule \citep{Pfeffer2013, Pi2013, Lee2013, Fu2014, Karnani2016a}. \textbf{(B)} Three alternative motifs documented in the same evidence base: direct VIP\rightarrowPV control of perisomatic inhibition \citep{Pi2013, Karnani2016b, Kuchibhotla2017}, direct VIP\rightarrowpyramidal inhibition by bipolar/CR-positive VIP cells \citep{Lee2013, Karnani2016a, Karnani2016b, Schneidermizell2025}, and a layer-1 VIP\rightarrowNDNF disinhibitory motif gating apical dendritic input \citep{Abs2018, Anastasiades2021a, Veit2023}; the hippocampal VIP\rightarrowOLM/basket motif is a homologue with a distinct oscillatory context \citep{Turi2019, Francavilla2018a, Dudai2020}. \textbf{(C)} Decision tree mapping behavioural context --- locomotion, attentional top-down, reinforcement, fear/learning --- to the motif most consistent with the empirical recruitment data and to the supporting cite-keys; the assignment is qualitative because the literature does not yet contain joint manipulations of multiple motifs under matched conditions \citep{Pi2013, Lee2013, Fu2014, Pinto2015, Kuchibhotla2017, Cichon2017, Williams2019, Krabbe2019, Veit2023}. \textbf{(D)} Schematic legend distinguishing measured edges (solid) from inferred or context-dependent edges (dashed). \textit{Caveat:} this figure is a schematic with no quantitative data substrate; assignment of behavioural context to dominant motif in Panel C is inferred from joint reading of the cited recruitment studies and is not the output of a meta-analytic weighting.}
\label{fig-disinhibition-motifs}
\end{figure}

\begin{figure}[!htbp]
\centering
\includegraphics[width=0.95\linewidth]{files/fig-vip-effect-size--684d62ea24a8836f57d1e6702163b9ba.png}
\caption[]{\textbf{Quantitative cross-study comparison of VIP-mediated effects on cortical inhibition.} \textbf{(A)} VIP\rightarrowSST and VIP\rightarrowpyramidal IPSC charge in mouse V1 paired recordings (Pfeffer 2013; n = 11 paired recordings; mean $\pm$ SD as printed) alongside VIP\rightarrowSST and VIP\rightarrowpyramidal IPSC amplitudes from a separate slice study in S1 (Lee 2013; sample sizes were not extracted into the curated row for that entry). The relative ordering VIP\rightarrowSST \textgreater  VIP\rightarrowpyramidal is reproduced across studies; absolute amplitudes differ by an order of magnitude reflecting differences in opsin, recording configuration, and pyramidal subtype. \textbf{(B)} VIP\rightarrowSST connection probability in mouse V1 paired recordings (Pfeffer 2013, n = 16 pairs) and VIP-evoked IPSC decay-time-constant comparison between SST and PV postsynaptic targets in mouse auditory cortex (Pi 2013, $\pm$ as printed; n is unrecorded in the curated entry --- the error metric should be downgraded to `unspecified $\pm$ as printed' pending verification against Pi 2013 figure legend / methods, and missing n is flagged accordingly). The VIP\rightarrowSST kinetic profile is approximately threefold slower than the VIP\rightarrowPV profile, consistent with target-specific inhibitory time-constants. \textbf{(C)} VIP-cell $\Delta$F/F change at running onset in mouse V1 (Fu 2014; $\approx$155\% of baseline; n is null, n\_definition states `VIP cells imaged with GCaMP6s' but no numeric n appears in the curated entry; this is a single-entry comparison after splitting and not a true cross-paper comparison) compared qualitatively with the direction of locomotion modulation reported by Fu 2014, Pakan 2016, and Dipoppa 2018 in the same cortical area; locomotion-coupled VIP activation is reproduced, but the downstream SST trajectory diverges across studies. \textbf{(D)} Conflict atlas mapping the six conflict pairs in this section (Pfeffer 2013 vs Pi 2013; Fu 2014 vs Lee 2013; Lee 2013 vs Pakan 2016; Pi 2013 vs Kuchibhotla 2017; Lee 2013 vs Pi 2013; Pi 2013 vs Pfeffer 2013) onto the disagreement axis (specificity, recruitment route, generalisation, primacy). \textit{Caveat:} values in panels A--C are taken from the as-printed source sentences in the audited literature; sample sizes are missing in several rows as flagged in the audit. The conflict atlas in panel D is a curated qualitative summary of the six conflicts surfaced in the prose of this section and is not the output of a quantitative meta-analysis.}
\label{fig-vip-effect-size-comparison}
\end{figure}

\section{The disinhibition framework reassessed}

The thesis with which this section opened --- that the VIP\rightarrowSST\rightarrowPyr motif is a privileged but non-exclusive channel for state-dependent gain --- has the following empirical content. First, the disinhibitory motif is reproducible: connectivity asymmetry favouring SST targets, recruitment by salient external events combined with arousal, and net pyramidal gain enhancement under top-down or running-coupled drive have been independently replicated across V1, A1, S1, M1, and PFC \citep{Pfeffer2013, Pi2013, Lee2013, Fu2014, Pinto2015, Kuchibhotla2017, Williams2019, Kamigaki2017, Anastasiades2021a, Veit2023}. Second, the motif is \textit{not} exclusive: in the same preparations, VIP\rightarrowPV, VIP\rightarrowpyramidal, and VIP\rightarrowNDNF connections are present, behaviourally engaged, and quantitatively non-trivial \citep{Karnani2016a, Karnani2016b, Yu2019, Abs2018, Schneidermizell2025}. Third, the motif's downstream consequence --- pyramidal gain enhancement through SST suppression --- is conditional: cortical area, layer, and behavioural state determine whether SST cells follow the disinhibition trajectory or move independently \citep{Pakan2016, Dipoppa2018, Garrett2020, Veit2023}. Fourth, the motif is not the \textit{defining} function of the VIP class: VIP cells participate in oscillatory pacing \citep{Turi2019, Bastos2023a}, plasticity gating \citep{Cichon2017, Krabbe2019, Tamboli2024, Ramosprats2022}, developmental wiring \citep{Batistabrito2017, Priya2018, Ferguson2023b}, and direct pyramidal control \citep{Karnani2016b, Lee2013, Schneidermizell2025}, none of which are captured by the disinhibitory motif alone.

The reassessment matters for how subsequent sections frame the \textit{in vivo} and computational reading of VIP function. The disinhibitory motif is a useful first hypothesis whenever a salient external event coincides with arousal in a cortical task; it is a misleading default when the task involves stationary visual processing in which locomotion is absent \citep{Pakan2016, Dipoppa2018, Garrett2020}, when the cortical area engages a parallel cholinergic stream \citep{Kuchibhotla2017}, when learning-related plasticity rather than acute gain is the readout \citep{Cichon2017, Krabbe2019, Abs2018, Tamboli2024}, or when the behavioural context engages NDNF-mediated apical dendritic disinhibition that may operate on its own timescale \citep{Abs2018, Anastasiades2021a}. In each of these cases, treating the VIP\rightarrowSST\rightarrowPyr motif as the \textit{full} mechanism overpredicts gain enhancement and hides the parallel motifs whose engagement would actually account for the observed task effect.

Several gaps in the present evidence base limit how strongly the disinhibitory motif can be tested. The most consequential is the absence of \textit{joint} manipulations: very few studies in the literature reviewed here perturb VIP and one of the alternative motifs (VIP\rightarrowPV, VIP\rightarrowpyramidal, VIP\rightarrowNDNF) simultaneously in the same animal under matched task conditions, so the relative weights assigned to each motif in Figure~\ref{fig-disinhibition-motifs} are inferred from cross-study reading rather than measured directly \citep{Pi2013, Lee2013, Fu2014, Pakan2016, Kuchibhotla2017, Veit2023}. A second gap is that many of the foundational quantitative measurements lack reported sample sizes or use error metrics that are not unambiguously SEM versus SD (see the caveats embedded in Figure~\ref{fig-vip-effect-size-comparison}). A third gap is the imbalance of cortical-area coverage: V1 is over-represented and PFC, M1, and S1 are under-represented relative to the strength of claims commonly made about `cortex' as a whole. A fourth gap is that human and primate evidence is currently anatomical and transcriptomic rather than physiological; whether the motif is functionally engaged outside mouse remains an open question \citep{Vanderveer2021}.

The empirical reading developed here feeds forward into the \textit{in vivo} analysis of \href{/in-vivo-behavior}{In Vivo Function During Behavior}, which asks whether the disinhibitory motif predicts observed VIP activity during behaviour and where it fails to do so; into the cross-area and cross-species comparisons of \href{/cross-areas}{VIP Interneurons Across Brain Regions}; into the oscillatory and rhythmic-disinhibition reading of \href{/oscillatory-dynamics}{Oscillatory Dynamics and Temporal Coordination}, where the hippocampal homologue and theta-coupled VIP activity become central; into the computational models of \href{/computational-models}{Computational Models of VIP Circuit Function}, where the assumption that VIP cells are pure disinhibitors generally underperforms models that include direct VIP\rightarrowpyramidal and VIP\rightarrowPV edges \citep{Hertag2019, Garciadelmolino2017, Litwinkumar2016, Tzilivaki2023}; and into the synthesis of the concluding synthesis, where the disinhibitory motif is recast as a privileged but non-exclusive channel within a small family of parallel VIP-mediated circuits. The argument across the remaining sections does not depend on rejecting the disinhibitory motif. It depends only on retiring it as the unique account of what VIP cells do.

\section{Plasticity and developmental engagement of the disinhibitory motif}

A feature of the disinhibitory motif that is often elided in functional summaries is that its components are not static. The targeting bias of VIP cells onto SST interneurons is acquired during postnatal development through activity-dependent and Bax-regulated programmed cell death of CGE-derived interneurons \citep{Wu2022}. Disrupting this developmental program alters not only the absolute number of VIP cells in adult cortex \citep{Priya2018} but also the strength of the disinhibition motif: animals with delayed or imbalanced CGE-derived interneuron maturation show altered cortical gain modulation and impaired top-down task performance \citep{Batistabrito2017, Ferguson2023a, Ferguson2023b, Vanvelze2024}. The motif is therefore best understood as the adult endpoint of a developmental trajectory whose failure modes are themselves informative about the motif's function \citep{Takesian2018, Malik2022, Furutachi2024}.

In adult cortex, the motif is also a substrate for short- and long-term plasticity. Cholinergic input gates dendritic spine plasticity in pyramidal neurons during behavioural learning, with the disinhibitory motif acting as a permissive switch rather than as a continuous gain knob. \citet{Krabbe2019} reported that BLA VIP cells acquire fear-conditioning-related selectivity, \citet{Abs2018} reported that NDNF L1 interneurons in auditory cortex display learning-related plasticity that gates apical dendritic processing, and \citet{Williams2019} reported that VIP-mediated disinhibition in S1 gates thalamically driven LTP relevant to perceptual learning. The picture from these and related studies \citep{Tamboli2024, Ramosprats2022, Mcfarlan2024a, Mcfarlan2024b} is that the disinhibitory motif's \textit{behavioural} contribution is largely about gating plasticity windows and sustaining attentional or learning-related disinhibition, not about producing transient gain on a millisecond-to-second timescale. This is consistent with the kinetic profile of VIP\rightarrowSST IPSCs, which is slower than VIP\rightarrowPV \citep{Pi2013}, and with the fact that VIP-cell perturbation typically produces slow-timescale changes in pyramidal firing rather than fast modulation of stimulus-evoked responses \citep{Lee2013, Kuchibhotla2017, Pinto2015}.

\section{Methodological caveats and the boundary of the framework}

A number of methodological caveats limit how strongly the disinhibitory motif can be inferred from any individual study and qualify its presentation as a general framework. Most VIP-Cre driver lines have small but non-zero off-target labelling --- typically of CCK-positive interneurons and, in some preparations, of lateral inhibition-class CR-positive cells --- and the proportion varies by line and laboratory \citep{Vanderveer2021, Wei2021}. ChAT-VIP overlap is partial and area-specific, and the cholinergic VIP subset is a subpopulation rather than the rule \citep{Koukouli2017, Posuszny2020, Mcfarlan2024a, Mcfarlan2024b}. Many \textit{in vivo} perturbation studies use short, intense optogenetic activation that bypasses the gating logic by which VIP cells are normally recruited, producing larger and less selective effects than chemogenetic or behaviourally evoked recruitment \citep{Lee2013, Fu2014, Kuchibhotla2017, Veit2023}. Anatomical-versus-functional mismatches are pervasive: connectivity surveys and EM connectomics report wiring that may not be uniformly active in any given behavioural state, and \textit{in vivo} imaging measures activity that may not isolate the specific motif being inferred \citep{Pfeffer2013, Schneidermizell2025, Pakan2016, Dipoppa2018, Garrett2020}. None of these caveats falsify the disinhibitory motif; collectively they constrain the strength of inference that can be drawn from a single experiment and motivate the cross-study reading in Figure~\ref{fig-vip-effect-size-comparison}.

Two further boundaries deserve explicit mention. The first is layer specificity. Most of the disinhibitory-motif evidence comes from layer 2/3, where VIP, SST and PV cells are densely intermingled and where top-down apical input is delivered. In layer 5 and layer 6, the motif is engaged differently: deep-layer VIP cells receive a higher proportion of thalamic input and contact a different distribution of postsynaptic targets, including projection-class pyramidal neurons that are themselves heterogeneous \citep{Anastasiades2021a, Sermet2019, Munoz2017, Schneidermizell2025}. The second boundary is species specificity. In primate and human cortex, the absolute density of VIP cells is lower, the laminar distribution is shifted, and the molecular subtypes are partially non-homologous to mouse \citep{Pronneke2015, Vanderveer2021}. Whether the disinhibitory motif operates with comparable behavioural weight in primate cortex is an open question that the present evidence base cannot answer; the conservative reading is that the connectivity asymmetry is preserved but the functional weighting may not be \citep{Vanderveer2021}.

\section{Toward a circuit-level reading of the disinhibitory motif}

A useful reading of the present evidence is that the disinhibitory VIP\rightarrowSST\rightarrowPyr motif is a \textit{privileged channel} for state-dependent gain --- privileged because its connectivity asymmetry, recruitment selectivity, and behavioural engagement are all skewed in its favour --- but not a \textit{unique mechanism} for that gain. The privilege is real: of the documented VIP outputs, the SST contact carries the largest mean charge per spike and the longest decay constant \citep{Pfeffer2013, Pi2013, Lee2013}, the SST cell is the inhibitory class whose dendritic targeting most directly limits pyramidal stimulus-evoked gain \citep{Pfeffer2013, Karnani2016a, Veit2023}, and the recruitment routes that drive VIP cells (top-down glutamatergic and ascending cholinergic) overlap closely with the routes that signal behavioural state \citep{Lee2013, Fu2014, Hangya2015, Pinto2015, Kamigaki2017, Williams2019}. The non-uniqueness is also real: VIP cells form measurable contacts on PV, pyramidal, NDNF, and CCK targets, all of which are behaviourally engaged in some condition documented in the literature reviewed here \citep{Karnani2016a, Karnani2016b, Anastasiades2021a, Schneidermizell2025, Veit2023}. Treating the privileged channel as the \textit{whole} framework over-fits a tidy diagram to a noisier biological substrate; treating it as one branch of a small motif family preserves the explanatory traction of the disinhibitory schematic without committing to claims the data do not support.

This circuit-level reading also clarifies what the disinhibitory motif explains and what it does not. It explains why salient external events combined with arousal produce gain enhancement of pyramidal stimulus responses in superficial cortex \citep{Lee2013, Fu2014, Zhang2014, Pinto2015, Williams2019, Veit2023}. It does not explain why locomotion-coupled VIP activity in V1 fails to translate into uniform SST suppression across higher visual areas \citep{Pakan2016, Garrett2020, Dipoppa2018}, why cholinergic re-routing of inhibition produces simultaneous engagement of disinhibitory and direct PV-mediated streams \citep{Kuchibhotla2017}, why VIP perturbation can produce subtractive rather than divisive changes in pyramidal firing in some preparations \citep{Hertag2019, Garciadelmolino2017}, or why VIP cells are required for plasticity-gating in motor and amygdalar circuits whose connectivity does not strictly fit the disinhibitory schematic \citep{Cichon2017, Krabbe2019, Tamboli2024}. These remaining phenomena are well within reach of an extended framework that includes the alternative motifs and their behavioural assignments. They are not within reach of the disinhibitory schematic alone.

Whether those circuit motifs are recruited as predicted in awake, behaving animals --- and where the textbook prediction breaks down --- is taken up in \href{/in-vivo-behavior}{In Vivo Function During Behavior}.


\section{In Vivo Function During Behavior}
\label{sec:08_in_vivo_behavior}

In awake animals, vasoactive intestinal peptide-expressing (VIP) interneurons are reliably recruited by locomotion, arousal, and reinforcement signals, and their activation is causally linked to gain enhancement in primary visual cortex (V1) and to attentional and learning-related disinhibition in multiple cortical territories \citep{Fu2014, Pi2013, Lee2013, Kamigaki2017}. The thesis of this section is two-sided: in-vivo VIP activity is robustly state-modulated, but it is far from uniform. Response polarity, sensory selectivity, and behavioural correlates differ across cortical area, task structure, and individual VIP subtype, and the inferred VIP\rightarrowsomatostatin-interneuron (SST) suppression mechanism does not always obtain --- for example, several preparations report concurrent activation rather than reciprocal inhibition of SST cells during locomotion \citep{Pakan2016, Dipoppa2018, Yavorska2021}. The local circuit motifs catalogued in \href{/disinhibition-framework}{Local Circuit Motifs and the Disinhibition Framework} are therefore best treated as one set of operating modes among several, whose engagement depends on the area$\times$state$\times$modulator triple specified by the behavioural context \citep{Garciajuncoclemente2017, Kannan2022}. The cross-area framing for that view is laid out systematically in \href{/cross-areas}{VIP Interneurons Across Brain Regions}.

\section{Locomotion-driven recruitment of V1 VIP cells: the disinhibitory observation}

The single most influential in-vivo result for the field --- that running mice show a sharp increase in V1 VIP calcium activity that gates a multiplicative increase in pyramidal-cell visual gain --- was established by \citet{Fu2014}, who showed that locomotion activates V1 VIP neurons via a nicotinic basal-forebrain input and produces gain enhancement that is abolished when VIP cells are silenced. Two-photon imaging during head-fixed visual behaviour reproduced the observation that VIP cells are among the most strongly state-modulated cortical neurons, with locomotion- and arousal-aligned $\Delta$F/F transients several-fold larger than those of pyramidal cells \citep{Pakan2016, Dipoppa2018, Reimer2014}. The phenomenon is preserved across multiple V1 layers and persists in darkness, indicating that it does not require visual drive \citep{Pakan2016, Dipoppa2018}. Closed-loop manipulations confirm that VIP recruitment is not merely correlated with running but is required for full-magnitude gain change, since chemogenetic VIP inhibition or developmental VIP silencing both blunt locomotion-evoked gain \citep{Batistabrito2017}. Recurrent network models can reproduce the locomotion-induced gain shift in V1 only when the modulatory drive is routed preferentially through VIP cells, providing a principled link between the empirical observation and circuit theory \citep{Dipoppa2018, Veit2023}.

A subtler issue, often glossed in summary diagrams, is that the V1 locomotion signal is composite: it carries cholinergic, noradrenergic, and arousal-related components that load differently on VIP cells depending on whether the animal is running, simply pupil-dilated, or actively engaged in a task \citep{Reimer2014, Garciajuncoclemente2017, Collins2023}. Pupil-tracked arousal alone --- without locomotion --- is sufficient to elevate VIP activity in V1 and frontal cortex, indicating that the relevant variable is broader than wheel velocity \citep{Reimer2014, Garciajuncoclemente2017, Munoz2017}. Active task engagement adds a further multiplicative layer: when an arousal-matched mouse switches from passive listening to active discrimination, cholinergic axon activation transiently re-organises the local interneuron network, with VIP cells participating differently than during passive locomotion \citep{Kuchibhotla2017, Collins2023}. The disinhibitory ``VIP runs up with locomotion'' finding is therefore the resultant of several partly dissociable state variables that the field has only recently begun to disentangle.

\section{Reinforcement signals: VIP activation by reward and punishment}

A parallel line of in-vivo work, originating with Prior work \citep{Pi2013}, established that cortical VIP cells are recruited not only by motor state but by reinforcement: phasic reward and punishment delivered during operant behaviour produced rapid, large-amplitude VIP activation in auditory cortex, mPFC, and primary somatosensory cortex (S1), with optogenetic VIP inhibition disrupting performance on the discrimination task \citep{Pi2013}. Mesoscale imaging extended this result across virtually the entire dorsal cortex, showing that during initial learning most VIP interneurons are co-activated by both reinforcement valences within tens of milliseconds of outcome delivery \citep{Szadai2022}. In dorsomedial mPFC during a sensory-discrimination task, \citet{Pinto2015} reported that VIP cells respond strongly to action outcomes rather than to sensory cues themselves, and \citet{Kamigaki2017} showed that activating dmPFC VIP cells during the delay period of a memory-guided task enhances performance and stabilises action-plan representations. In primary motor cortex, VIP cells receive disproportionate orbitofrontal input that is recruited during reward-related motor learning \citep{Lee2023b}. In barrel cortex, Earlier reports \citep{Krabbe2019} demonstrated that VIP-mediated disinhibitory gating is required for associative whisker-trace conditioning, and Initial studies \citep{Sachidhanandam2016} showed that the VIP response to whisker-cued reward is itself modulated by task structure. A commentary by \citet{Zou2023} summarises the convergent claim that VIP cells signal expected rewards across cortical areas.

The reinforcement-related VIP signal is also recruited by aversive contexts. In anterior insular cortex, VIP cells gate aversive-stimulus encoding and contribute to fear-related behaviour \citep{Ramosprats2022}. In mPFC, optogenetic activation of VIP cells ameliorates fear-related stress-induced behavioural deficits, consistent with a permissive role for VIP-mediated disinhibition in flexible affective control. Amygdala VIP cells exhibit heterogeneous plasticity that supports both fear acquisition and extinction \citep{Favila2025}. The breadth of these reinforcement results --- across visual, somatosensory, auditory, motor, prefrontal, insular, and amygdala targets --- is the empirical anchor of the ``VIP-as-reinforcement-gate'' hypothesis. Critically, however, reinforcement responses are not exclusive: the same cells co-encode locomotion and arousal, and disentangling reinforcement-specific from state-driven components requires task designs that orthogonalise the two \citep{Pinto2015, Szadai2022, Garciajuncoclemente2017}.

\begin{framed}
\textbf{Conflict --- Functional homogeneity of VIP cells under reinforcement}\\
\citet{Pi2013} reported that reward and punishment produced similar phasic activation across the recorded VIP population, implying a single ``gating'' signal applied uniformly to the class. More recent recordings paint a more heterogeneous picture: in S1 during a Go/NoGo task \citet{Ramamurthy2023} showed that VIP cells are activated by sensory cues and licking but not by reinforcement itself, and \citet{Ramamurthy2025} reported that S1 layer 2/3 VIP cells do not carry a whisker-specific attention signal. Voltage imaging studies further reveal brain-state-dependent antagonism within the VIP population such that subsets shift their reinforcement coupling depending on cortical state \citep{Kannan2022}. \textbf{Note (gap)}: the strongest opposing-uniformity claim was scaffolded against paper key \texttt{Johnson2020}, which is not represented in our filtered citation map; the in-set substitutes \citep{Ramamurthy2023, Ramamurthy2025, Kannan2022} are cited above. \textbf{Resolution (partial)}: the discrepancy is consistent with area-specific tuning (mPFC and ACx VIP cells more reliably reinforcement-coupled than S1 VIP cells) and with subtype heterogeneity within a region (see \href{/cross-areas}{VIP Interneurons Across Brain Regions}). The original Pi-style ``uniform gating'' reading is therefore best treated as a population-mean approximation rather than a cell-level property.
\end{framed}

\section{Sensory selectivity: state messenger or stimulus encoder?}

A central tension in the in-vivo literature is whether VIP cells primarily transmit state and reinforcement signals to local circuits or whether they themselves encode features of the sensory stimulus. The state-messenger reading was crystallised by \citet{Fu2014}, whose V1 recordings emphasised that the locomotion-evoked VIP response was largely insensitive to grating orientation or contrast and behaved as a multiplicative gain knob. Foundational studies \citep{Mesik2015} similarly found that V1 VIP cells receive broad, weakly tuned cortico-cortical and thalamic inputs, consistent with a state/context role. Counter-evidence has accumulated since. \citet{Millman2020} showed that V1 VIP neurons carry contrast tuning that is complementary to that of pyramidal cells, indicating sensory-driven structure beyond a pure state signal. \citet{Garrett2020} reported that during a change-detection task V1 VIP cells are stimulus-driven by novel images but suppressed by familiar images, a stimulus-history-dependent code that cannot be reduced to arousal. \citet{Bastos2023a} found that contextually redundant stimuli increase V1 VIP activity while deviant stimuli decrease it, and that chemogenetic VIP inhibition disrupts deviance-related modulation downstream --- both observations require stimulus-feature access at the level of the VIP cell. Cross-modal sharpening further extends VIP's stimulus repertoire: \citet{Ibrahim2016} showed that sound suppresses V1 layer 2/3 VIP cells preferentially at the preferred orientation of nearby pyramidal cells, contributing to multisensory orientation refinement.

In barrel cortex the case for stimulus encoding is even stronger. \citet{Yu2019} reported that voluntary whisker movement activates a majority of VIP interneurons and that responses tile the whisker stimulus space rather than reflecting a uniform state signal. \citet{Kiritani2023} showed that S1 VIP cells respond to whisking but only after a delay to active touch, indicating distinct stimulus-locked components. Established work \citep{Walker2016} and \citet{Sachidhanandam2016} documented stimulus-locked VIP activity tied to whisker deflection during trained behaviour, and \citet{Brecier2022} found that S1 barrel-cortex VIP cells are most active during REM sleep --- implying state coupling that is independent of any explicit sensory cue. \citet{Lebedeva2025} used dual-population two-photon imaging to dissociate VIP and GABA-receptor-defined responses to specific sensory features, \citet{Khan2018a} provides an instructive complement: during visual learning, pyramidal/PV/SST cells but not V1 VIP cells increased selectivity for the task-relevant stimulus, suggesting that VIP carries a state/context-aligned (rather than learned-stimulus) signal in this paradigm. The cross-task picture therefore is not ``VIP encodes state'' or ``VIP encodes stimulus'' but rather a context-dependent mix in which both axes are represented at the single-cell level \citep{Yu2019, Garrett2020, Bastos2023a, Mesik2015}.

\begin{framed}
\textbf{Conflict --- Locomotion as a stimulus-blind gain knob vs stimulus-tuned VIP responses}\\
\citet{Fu2014} argued that V1 VIP activation is a locomotion-driven, largely stimulus-independent state messenger that multiplicatively scales V1 gain. By contrast, \citet{Yu2019} showed that S1 VIP cells encode whisker stimuli with cell-specific tuning, \citet{Garrett2020} showed novelty-dependent VIP responses in V1, and \citet{Ramamurthy2023} showed sensory-cue (but not reinforcement) coupling in S1. \textbf{Resolution}: the disagreement is best read as area- and task-dependent rather than a true contradiction; V1 VIP cells in the head-fixed grating-viewing preparation behave more state-like, while S1 VIP cells in active whisking and discrimination preparations behave more stimulus-like. The class-level claim that VIP cells are pure state messengers is therefore not generic across cortex \citep{Mesik2015, Millman2020, Bastos2023a}. \textbf{Note (gap)}: the original opposing claim was scaffolded against paper key \texttt{Ramamurthy2021}, which is not represented in our filtered citation map; \citet{Ramamurthy2023, Ramamurthy2025} from the same group serve as the in-set substitute and we flag this as an unmet citation need.
\end{framed}

\section{Cross-area heterogeneity: V1, A1, and the auditory paradox}

When the V1 locomotion result is taken to a different cortical area, the disinhibitory motif partly breaks. \citet{Bigelow2019} reported that movement activates an apparently similar VIP\rightarrowSST\rightarrowpyramidal circuit in auditory cortex (ACtx) --- yet the net effect on auditory pyramidal-cell responses is opposite in sign to that observed in V1, with locomotion suppressing rather than enhancing sensory-evoked activity. \citet{Yavorska2021} extended this paradox by directly testing whether the locomotion-related changes in ACtx are mediated by the VIP-SST disinhibitory circuit and concluded that they are not, in stark contrast to what is reported in V1. Pupil-tracked passive arousal does still recruit VIP cells in ACtx \citep{Reimer2014, Munoz2017}, and active engagement during auditory recognition produces cholinergic axon activation that simultaneously engages PV/SST/VIP networks \citep{Kuchibhotla2017}. The natural conclusion is that the same molecular cell class participates in different cortical implementations of state-dependent gain --- adding gain in V1 and subtracting gain in ACtx --- depending on local synaptic weights and on the dominant neuromodulator engaged by the behaviour \citep{Bigelow2019, Yavorska2021, Lecrux2016}. Earlier work \citep{Moore2026} further reported that ACtx VIP and NDNF cells have distinct task-related dynamics, complicating any simple class-level reading of auditory VIP function.

Within mPFC and ACC the in-vivo signature again differs from V1. \citet{Pinto2015} and \citet{Kamigaki2017} highlighted strong outcome- and delay-period activity, and \citet{Lee2019} showed that prefrontal VIP cells gate hippocampal-prefrontal communication during anxiety-related behaviour, with VIP activity rising in the open arms of an elevated plus-maze. \citet{Garciajuncoclemente2017} showed that during arousal in frontal cortex, VIP interneurons rapidly inhibit pyramidal neurons directly while also engaging the the disinhibitory account motif --- establishing that direct VIP\rightarrowPyr inhibition is a behavioural reality, not just an anatomical curiosity (see \href{/disinhibition-framework}{Local Circuit Motifs and the Disinhibition Framework}). \citet{Silverstein2025} documented that repeated ventral-hippocampus input persistently depresses mPFC VIP activity, contributing to spatial working-memory deficits, and \citet{Malik2022} showed that long-range PFC GABAergic projections preferentially inhibit hippocampal VIP cells to enable top-down signal-to-noise control. These results indicate that frontal VIP function is bidirectional: it both broadcasts permissive state signals locally and is itself targeted by long-range inhibitory feedback from associated regions \citep{Malik2022, Silverstein2025}.

\section{Concurrent VIP/SST imaging and the limits of the disinhibitory account}

Direct tests of the disinhibitory motif require imaging both VIP and SST populations during the same behavioural state. The cleanest such tests challenge the textbook reading. \citet{Pakan2016} reported that locomotion increases activity in VIP, SST, and PV interneurons during visual stimulation in V1 --- that is, SST cells are co-activated, not silenced, by the same state that activates VIP cells. \citet{Dipoppa2018} reproduced and extended this, finding that locomotion activates VIP cells and disinhibits pyramidal cells in darkness while in the presence of visual stimuli SST cells can also be co-recruited; their network model captured the locomotion-driven gain change only when the modulatory drive was routed predominantly through VIP. \citet{Yavorska2021} showed that in ACtx, VIP\rightarrowSST disinhibition is not the route for locomotion-related changes at all. \citet{Kannan2022} used dual-polarity voltage imaging to reveal brain-state-dependent antagonism between VIP and SST populations whose sign reverses across states. The straightforward textbook story --- VIP up, SST down, pyramidal up --- is therefore a special case rather than a generic prediction, and the relative weights of the VIP\rightarrowSST, VIP\rightarrowPV, and direct VIP\rightarrowPyr motifs that determine its applicability are themselves state- and area-dependent \citep{Garciajuncoclemente2017, Kannan2022, Veit2023}.

\begin{framed}
\textbf{Conflict --- SST suppression by locomotion}\\
\citet{Fu2014} proposed that locomotion-driven VIP activation suppresses V1 SST cells to enable disinhibition of pyramidal cells. \citet{Pakan2016} reported the opposite --- that locomotion activates SST as well as VIP cells in V1 --- and \citet{Dipoppa2018} confirmed concurrent SST recruitment in the presence of visual stimuli. \textbf{Resolution}: the SST direction is conditional on visual drive and on the recording layer. In darkness or at low contrast, SST cells are less active than VIP cells, recovering an apparent VIP-up/SST-down profile; with strong visual drive the two are co-activated, and the net effect on pyramidal gain still shifts because VIP recruitment is proportionally larger \citep{Pakan2016, Dipoppa2018, Veit2023}. Closed-loop modelling reproduces both observations \citep{Dipoppa2018, Veit2023}. The discrepancy therefore reflects an under-specified circuit reading rather than an irreconcilable conflict.
\end{framed}

\section{Hippocampal and limbic VIP populations}

Outside the neocortex, hippocampal VIP/IS interneurons implement specialised motifs that depart from the cortical reading. \citet{Turi2019} reported that CA1 VIP cells form functional subpopulations whose activity is shaped by behavioural state and task demand and that selectively inhibit other interneurons to support goal-oriented spatial learning. \citet{Francavilla2018a} showed that long-range VIP-GABAergic neurons in CA1 decrease their activity during theta-run epochs and are most active during quiet wakefulness --- the opposite locomotion polarity to V1 VIP cells --- and \citet{Luo2020} found that CA1 VIP-interneuron-selective cells are preferentially recruited during theta-run epochs but not during sharp-wave ripples. \citet{Tamboli2024} reported that CA1 VIP cells become highly active during exploration of new environments and support recognition memory, and \citet{Neubrandt2025a} showed that during hippocampal navigation, novel environments transiently increase VIP activity and facilitate place-field formation. Previous reports \citep{Lenkey2025} demonstrated that VIP interneurons drive brain-region-specific gain modulation of place cells across the hippocampal axis, mirroring at the systems level the area-specific gain logic seen in cortex. Earlier studies \citep{Leroy2022} reported that hippocampal VIP cells release enkephalin to modulate stress-related behaviour, indicating a neuropeptidergic dimension that the disinhibitory fast-GABA disinhibition reading does not capture.

\begin{framed}
\textbf{Conflict --- Hippocampal VIP and place-cell coding}\\
The disinhibitory \citet{Fu2014} motif, generalised, predicts that VIP recruitment should robustly modulate principal-cell coding wherever VIP cells exist. Hippocampal data only partly support this generalisation. \citet{Lenkey2025} reported VIP-driven gain modulation of place cells, and \citet{Tamboli2024, Neubrandt2025a} showed novelty-related VIP recruitment that supports place-field formation and recognition memory. By contrast, \citet{Francavilla2018a, Luo2020} showed that during ongoing theta-run behaviour many CA1 VIP populations actually decrease activity, the opposite of V1 VIP. \textbf{Note (gap)}: the strongest contrary claim --- that hippocampal place fields are minimally affected by VIP perturbation --- was scaffolded against paper key \texttt{Vervaeke2024}, which is not represented in our filtered citation map; the comparison is therefore presented qualitatively. The within-hippocampus consensus that emerges from \citep{Turi2019, Francavilla2018a, Luo2020, Tamboli2024, Neubrandt2025a, Lenkey2025} is that VIP recruitment is state-dependent but its sign and downstream effect differ between subregions and between IS-cell subtypes --- consistent with the cortical reading but cautioning against a single ``VIP enables learning'' rule.
\end{framed}

VIP-class neurons also operate outside cortex and hippocampus in subcortical and peripheral preparations whose in-vivo recruitment patterns are markedly different from neocortex. In the olfactory bulb, \citet{Wang2022} showed that VIP cells directly inhibit mitral cells and that their inactivation impairs olfactory detection and discrimination --- a direct VIP\rightarrowprincipal-cell inhibition operating at a sensory periphery, again outside the disinhibition framing. Amygdala VIP cells in BLA exhibit variable density across subdivisions \citep{Rhomberg2018} and undergo associative plasticity that supports both fear acquisition and extinction \citep{Favila2025}. Long-range hippocampal VIP cells release enkephalin to modulate stress-related behaviour, and prefrontal VIP recruitment is itself altered by repeated ventral-hippocampus input \citep{Silverstein2025}. The bulbar, amygdalar, and limbic evidence emphasises that the in-vivo ``function'' of VIP cells is not unitary across the brain.

\section{Causal manipulations: gain, attention, and learning}

In-vivo causal data are required to bridge correlation to function. Optogenetic and chemogenetic manipulations have been used to test whether the VIP signal is itself necessary, sufficient, or merely permissive for the behavioural and circuit phenomena it correlates with. Prior work \citep{Pi2013} provided the first such test: optogenetic inhibition of VIP cells during a discrimination task disrupted both reinforcement-aligned activity and behavioural performance. \citet{Kamigaki2017} showed that delay-period activation of dmPFC VIP cells improves working-memory-task performance, providing temporally specific causal evidence. Earlier reports \citep{Krabbe2019} showed that VIP-dependent disinhibitory gating in barrel cortex is required for whisker-trace conditioning, and \citet{Bastos2023a} showed that chemogenetic VIP inhibition disrupts deviance detection in V1. \citet{Veit2023} combined optogenetics with computational modelling to demonstrate that VIP cells implement a multiplicative gain operation. Initial studies \citep{Myersjoseph2022} used all-optical bidirectional manipulation to show that VIP activity is sufficient to bias perceptual choice in a sensory-detection task. In olfactory bulb, \citet{Wang2022} showed that VIP inactivation impairs detection and discrimination, and in olfactory (piriform) cortex \citet{Cantobustos2022} used optogenetics to show that VIP-mediated disinhibitory circuitry gates associative synaptic plasticity (LTP) --- establishing that VIP-driven disinhibition supports learning-relevant signal transformations across modalities. The convergent message is that VIP recruitment is permissive for area-specific computations whose direction the the disinhibitory account motif sometimes, but not always, predicts \citep{Veit2023, Krabbe2019, Bastos2023a, Myersjoseph2022, Wang2022, Cantobustos2022}.

Causal manipulations also expose limits. \citet{Chen2015a} reported that optogenetic VIP inhibition did not block cholinergically driven cortical desynchronisation in V1, indicating that not every well-known cholinergic phenomenon requires VIP cells, and that SST cells can be the relevant intermediate. \citet{Khan2018a} reported that during visual learning V1 VIP cells, in contrast to PV/SST/Pyr cells, did not change their tuning --- implying that VIP recruitment is necessary for some behaviours but not for the plasticity events it permits. In a Dravet (Scn1a+/-) model, \citet{Goff2023} showed that VIP recruitment is diminished at quiet-to-running transitions and that re-activating VIP cells optogenetically rescues part of the deficit, while \citet{Liebergall2024} showed that NDNF cells maintain normal arousal recruitment in the same model --- refining which cell class actually carries the locomotion signal across pathological states. After cortical photothrombotic stroke, \citet{Motaharinia2021} showed that chemogenetic augmentation of VIP excitability promotes recovery of forelimb function. These causal data sharpen the interpretation of correlated VIP activity by separating cell-class-specific from class-redundant components of the in-vivo response \citep{Chen2015a, Khan2018a, Goff2023, Motaharinia2021}.

\section{Neuromodulatory drivers of in-vivo VIP recruitment}

The in-vivo recruitment of VIP cells is controlled by a small set of fast neuromodulators that gate their excitability around the behavioural state of the animal. The cholinergic substrate is the best characterised. \citet{Porter1999} first showed that nicotinic agonists selectively excite a CCK/VIP-coexpressing GABAergic subpopulation, and \citet{Letzkus2011} established that auditory-cortex layer-1 interneurons receive direct nicotinic excitation from basal-forebrain cholinergic axons. \citet{Arroyo2012} confirmed that basal-forebrain cholinergic axons preferentially activate CGE-derived VIP-class cells, and \citet{Alitto2013} demonstrated that basal-forebrain electrical stimulation activates VIP+ interneurons via nicotinic receptors during cortical state transitions. \citet{Hangya2015} reported that optogenetically identified basal-forebrain cholinergic neurons broadcast a fast ({\textasciitilde}18 ms), precisely timed reinforcement signal that is well suited to drive the rapid VIP responses to reward and punishment. \citet{Saunders2015} and \citet{Granger2020} further showed that basal-forebrain cholinergic neurons co-release GABA onto cortical layer-1 interneurons and that cortical ChAT-expressing cells (themselves nearly all VIP+) release GABA broadly onto inhibitory cells while sparing principal cells, indicating a multi-channel basal-forebrain influence on VIP-class circuits. Foundational studies \citep{Yaeger2019} showed that during the binocular critical period, basal-forebrain ACh released during arousal recruits the VIP\rightarrowSST\rightarrowPyr motif in V1, providing a developmental dimension to the disinhibitory reading.

The 5-HT3A receptor route is equally important. \citet{Ferezou2002} showed that the ionotropic 5-HT3 receptor is selectively expressed on VIP/CCK GABAergic interneurons, and \citet{Lee2010} showed that virtually all CGE-derived neocortical interneurons that lack PV/SST express 5-HT3A. Established work \citep{Poorthuis2018} extended this finding to human neocortex, where layer-1 GABAergic interneurons --- the human homologs of the rodent VIP-class --- show pronounced serotonergic responses, indicating cross-species conservation of the modulator-target axis. Adrenergic regulation, by contrast, has been comparatively under-explored: Earlier work \citep{Kawaguchi1998} showed that noradrenaline differentially modulates VIP-class GABAergic firing in rat frontal cortex via $\alpha$-adrenergic excitation, and \citet{Collins2023} showed that locus-coeruleus axons across mouse dorsal cortex carry a globally coordinated arousal signal that overlaps with cholinergic drive. \citet{Lecrux2016} demonstrated that distinct neuromodulatory pathways (basal-forebrain ACh, locus-coeruleus NE) recruit highly specific cortical VIP/NOS/CCK interneurons, supporting an ``operator'' view in which different modulators select different functional sub-circuits engaging the VIP class. The genetic evidence further constrains this picture: \citet{Koukouli2017} showed that mice carrying the human $\alpha$5-nicotinic-receptor risk SNP or $\alpha$5 knockout show reduced VIP-interneuron activation, demonstrating that genetic perturbations of nicotinic subunits act via the VIP arm of the circuit. \citet{Askew2019} showed that bath nicotine strongly and directly depolarises auditory-cortex VIP cells, and \citet{Munoz2017} showed that during active wakefulness in mouse cortex, VIP-class layer-1 interneurons are preferentially active. The \citet{Kawaguchi1997b} observation that morphologically and chemically defined neocortical interneuron subtypes show selective sensitivity to specific modulators is the structural prototype that all of these in-vivo papers operationalise behaviourally.

\section{A direct VIP-\textgreater pyramidal channel and the ChAT-VIP subset}

A recurring theme in the in-vivo data is that the textbook VIP\rightarrowSST\rightarrowPyr motif coexists with parallel routes whose engagement is partly independent of state. \citet{Garciajuncoclemente2017} explicitly demonstrated that during arousal in frontal cortex, VIP cells rapidly inhibit pyramidal neurons directly while also disinhibiting them via SST suppression --- a dual operation. Previous reports \citep{Lee2014b} provided complementary anatomical evidence that VIP-class circuits include long-range GABAergic projections that target principal cells in distant regions. The ChAT-VIP subpopulation \citep{Dudai2020, Granger2020} adds a further mode: roughly 0.5--1\% of cortical VIP cells co-express choline acetyltransferase and locally release ACh, and these cells fire faithfully to whisker stimuli. \citet{Obermayer2019} showed that ChAT-VIP cells in mPFC release ACh onto neighbouring neurons and contribute to attention behaviour distinct from the basal-forebrain channel, and \citet{Dudai2020} reported that VIP/ChAT cells in barrel cortex ({\textasciitilde}0.5\%) fire faithfully to whisker stimuli and shape local cortical dynamics. The presence of a small co-transmitting subset means that even tightly cell-type-restricted VIP perturbations can produce mixed effects whose interpretation requires concurrent monitoring of both released transmitters \citep{Obermayer2019, Dudai2020, Granger2020}.

\section{Arousal, pupil, and the inverted-U problem}

Pupil-tracked arousal is the dominant continuous covariate of in-vivo VIP activity, and its quantitative relationship with VIP recruitment has emerged as a focus of recent work. \citet{Reimer2014} showed that slow pupil dilations during quiet wakefulness track desynchronisation, enhanced sensory responses, and cell-type-specific interneuron activation. \citet{Garciajuncoclemente2017} and Earlier studies \citep{Munoz2017} showed that arousal alone --- without locomotion --- recruits VIP cells in frontal cortex and dorsal-cortex layer-1 networks. Prior work \citep{Collins2023} showed that cholinergic and noradrenergic axons across dorsal cortex carry a globally coordinated arousal signal that delivers the modulatory drive onto VIP cells. Whether the relationship between arousal and VIP recruitment is monotonic --- as the disinhibitory reading implies --- or has an inverted-U shape with performance decrements at high arousal is, in our filtered evidence base, partially constrained. The Yerkes--Dodson-like reading that arousal-task performance is non-monotonic is well established in the broader literature, but the strongest statement of an inverted-U specifically at the VIP cell --- and a contrary monotonic statement --- comes from sources outside the filtered citation map.

\begin{framed}
\textbf{Conflict --- Shape of the arousal--VIP--performance relation}\\
\citet{Fu2014} and downstream gain-modulation studies imply a roughly monotonic relationship between arousal/locomotion and V1 VIP activity, and between VIP activation and downstream pyramidal-cell gain \citep{Dipoppa2018, Veit2023}. Recent work in arousal--performance coupling proposes an inverted-U shape with performance decrements at high arousal. \textbf{Note (gap)}: the paper that most directly makes this inverted-U claim at the VIP-cell level was scaffolded against paper key \texttt{Beerendonk2024}, which is not represented in our filtered citation map; we therefore cannot quote effect sizes from that source. The closest in-set evidence comes from \citep{Reimer2014, Garciajuncoclemente2017, Munoz2017, Collins2023}, which demonstrate cell-type-specific arousal coupling without a quantitative test of monotonicity. We list this as an unmet citation need rather than resolving the conflict.
\end{framed}

\section{A decision framework: area x state x modulator -\textgreater  predicted VIP polarity}

Pulling these threads together, the in-vivo literature is best organised by a three-axis decision rule rather than by a single class-level statement (see Figure~\ref{fig-vip-state-decision-tree}). The first axis is cortical area: V1 favours locomotion-up VIP, ACtx implements an opposite-sign effect via a non-disinhibitory route, mPFC and ACC emphasise outcome and delay activity, S1 supports stimulus-driven and licking-related VIP recruitment, and hippocampus shows subregion- and IS-cell-subtype-specific signs \citep{Fu2014, Bigelow2019, Yavorska2021, Pinto2015, Yu2019, Lenkey2025}. The second axis is behavioural state: quiet wakefulness, locomotion, arousal-only, attention, and reinforcement load differently on VIP cells, and these states are partly dissociable \citep{Reimer2014, Garciajuncoclemente2017, Pi2013, Szadai2022}. The third axis is dominant neuromodulator: nicotinic ACh and 5-HT3A serotonin act fast and directly on VIP cells, while $\alpha$-adrenergic noradrenaline acts more slowly and is selective for subtypes \citep{Letzkus2011, Ferezou2002, Lee2010, Kawaguchi1998, Lecrux2016}. Many cross-study disagreements that have been read as contradictions dissolve once the area$\times$state$\times$modulator triple is specified, and a small set of residual disagreements (the auditory-cortex polarity and the hippocampal place-field question) survive even after such conditioning and require dedicated experiments \citep{Bigelow2019, Yavorska2021, Lenkey2025, Francavilla2018a}.

The behavioural responses already display area-to-area heterogeneity that \href{/cross-areas}{VIP Interneurons Across Brain Regions} catalogues systematically; the oscillatory consequences are addressed in \href{/oscillatory-dynamics}{Oscillatory Dynamics and Temporal Coordination}; and the computational reading of these multiplexed signals is developed in \href{/computational-models}{Computational Models of VIP Circuit Function}. The synthesis-and-reassessment closure of the review takes up the residual conflicts directly in the concluding synthesis. The figures below summarise the empirical landscape and the decision rule respectively.

\begin{figure}[!htbp]
\centering
\includegraphics[width=0.92\linewidth]{files/fig-vip-behavior-fd15a43c1cd66474298cd7d78413782c.png}
\caption[]{Qualitative summary of in-vivo VIP recruitment across behavioural conditions, redrawn as a schematic because no audited quantitative panels exist for cluster\_07\_in\_vivo\_function (Phase 6 audit). Panel A: per-study reported direction of VIP recruitment by locomotion in V1, with directionality only (no $\Delta$F/F values shown --- the original cross-study compilation could not be audited). Panel B: heat-map-style polarity matrix across cortical area $\times$ behavioural variable, drawn from claim-level statements in \citet{Fu2014, Pakan2016, Dipoppa2018, Bigelow2019, Yavorska2021, Pi2013, Pinto2015, Lee2013, Yu2019, Sachidhanandam2016, Garrett2020, Garciajuncoclemente2017, Kamigaki2017, Ramamurthy2023, Ramamurthy2025}. Panel C: schematic concurrent VIP/SST trace contrast --- a cartoon of the qualitative discrepancy between \citet{Fu2014}, \citet{Pakan2016}, and \citet{Dipoppa2018} rather than an audited overlay. Panel D: decision diagram from area$\times$state$\times$modulator \rightarrow expected VIP polarity, included for legibility before the schematic in Figure~\ref{fig-vip-state-decision-tree}. \textbf{Gap caveat}: panels A and C should be replaced by audited quantitative overlays in a future revision; we present them here at the qualitative level only. \citep{Fu2014, Pakan2016, Dipoppa2018, Bigelow2019, Yavorska2021}}
\label{fig-vip-behavior}
\end{figure}

\begin{figure}[!htbp]
\centering
\includegraphics[width=0.92\linewidth]{files/fig-vip-state-decisi-994c34b0675ac2a341b90007a02732f8.png}
\caption[]{Decision schematic for predicting in-vivo VIP polarity from cortical area, behavioural state, and dominant neuromodulator. The tree branches first on area (V1 / A1 / S1 / mPFC / hippocampus), then on state (locomotion / quiet wake / attention / reinforcement), then on dominant modulator (nicotinic ACh / 5-HT3A / $\alpha$-adrenergic NA), and the leaf annotates the predicted VIP response sign together with the cardinal study that populates that leaf. Cardinal-study leaves: \citet{Fu2014} (V1, locomotion, ACh, +); \citet{Pakan2016} (V1, locomotion, mixed, + with concurrent SST); \citet{Dipoppa2018} (V1, locomotion, ACh/visual-drive interaction); \citet{Pi2013} (auditory/mPFC, reinforcement, ACh, +); \citet{Lee2013} (S1, whisking, vM1-driven, +); \citet{Bigelow2019} (ACtx, locomotion, opposite net sign on Pyr); and the hippocampal contrarian leaves populated by \citet{Francavilla2018a, Luo2020} (CA1, theta-run, -). The cartoon overlay (panel C) shows a head-fixed mouse on a wheel with V1/A1/S1 sketches indicating concurrent expected VIP-vs-SST polarity by area, with the auditory-cortex cell highlighted as the explicit area-paradox leaf \citep{Yavorska2021}. \citep{Veit2023, Kannan2022}}
\label{fig-vip-state-decision-tree}
\end{figure}

\section{Top-down inputs and novelty/expectation signals}

The in-vivo VIP signal is not autonomous: it inherits structure from long-range inputs that themselves carry behavioural information. \citet{Lee2013} showed that whisker motor cortex (vM1) recruits VIP cells in S1 to disinhibit pyramidal cells during whisking, providing one of the clearest demonstrations that a long-range cortico-cortical projection drives VIP activation during natural behaviour. Earlier reports \citep{Walker2016} extended this with optogenetic identification of the inhibitory interneurons targeted by the vM1 input. \citet{Lee2023b} showed that primary motor cortex VIP cells receive disproportionate orbital frontal cortex input, supporting reward-related motor learning. \citet{Malik2022} showed that long-range PFC GABAergic projections preferentially inhibit hippocampal VIP cells, providing a route by which top-down inhibitory signals adjust VIP gain rather than simply recruiting it. Concurrent monitoring of VIP and the long-range driver in awake animals --- for example by combining cortico-cortical optogenetic tagging with VIP imaging --- has only just become routine, and \citet{Naskar2021} provided systematic measurements of synaptic strength from functionally relevant brain areas to cell-type-defined neuronal populations including VIP cells, supplying the synaptic substrate for these in-vivo influences.

Novelty and expectation are emerging as dimensions along which VIP recruitment is structured. \citet{Garrett2020} reported that V1 VIP cells are stimulus-driven by novel images but suppressed by familiar images during a change-detection task --- a form of stimulus-history dependence that survives matched arousal and locomotion. \citet{Bastos2023a} reported that contextually redundant stimuli increase V1 VIP activity while deviant stimuli decrease it, and that chemogenetic VIP inhibition disrupts the deviance-related modulation. \citet{Neubrandt2025a} reported that hippocampal VIP cells are transiently activated by novel environments and that this activation supports place-field formation. \citet{Tamboli2024} similarly reported novelty-aligned CA1 VIP activity that supports recognition memory. The collective picture is that VIP cells implement at least three partly dissociable in-vivo signals --- state, reinforcement, and novelty/expectation --- whose mutual relationship is itself area- and task-specific \citep{Garrett2020, Bastos2023a, Neubrandt2025a, Tamboli2024, Pi2013}.

\section{Sleep-wake states, oscillations, and cortex-wide imaging}

In-vivo VIP recruitment is structured by sleep--wake states and by cortex-wide oscillatory regimes. \citet{Brecier2022} showed that VIP cells in S1 barrel cortex are most active during REM sleep, providing one of the few cell-type-resolved measurements across the natural sleep/wake cycle. \citet{Francavilla2018a} and \citet{Luo2020} demonstrated that hippocampal CA1 VIP populations have opposite-sign relationships to theta-run epochs and to sharp-wave ripples, with state coupling that varies between subtypes. \citet{Munoz2017} and \citet{Reimer2014} documented that VIP-class layer-1 interneurons are preferentially active during active wakefulness and during pupil-tracked desynchronisation. \citet{Neske2016} reported that interneuron subtypes including VIP cells are highly active during cortical Up states. \citet{Szadai2022} provided the broadest mesoscale view, showing that during initial learning most VIP interneurons across the dorsal cortex are co-activated by reward and punishment within tens of milliseconds --- an observation that bridges single-cell and cortex-wide data and indicates that VIP recruitment is a brain-state-aligned, not just locally circuit-specified, phenomenon. The oscillatory consequences of these recruitment patterns are taken up systematically in \href{/oscillatory-dynamics}{Oscillatory Dynamics and Temporal Coordination}.

\section{Developmental and longitudinal in-vivo recruitment}

In-vivo VIP function is not fixed but evolves over development and across longer timescales. Initial studies \citep{Yaeger2019} showed that during the binocular critical period of mouse V1, basal-forebrain ACh released during arousal recruits the VIP\rightarrowSST\rightarrowPyr motif specifically, providing a developmental correlate of the disinhibitory motif. \citet{Batistabrito2017} showed that developmental ErbB4 deletion from VIP interneurons impairs cortical state dependence and sensory learning in mature mice --- a result that ties developmental wiring of VIP cells to their adult in-vivo function. \citet{Yang2025} combined voltage imaging with optogenetic depolarisation to reveal how excitatory--inhibitory balance changes are read out at the VIP cell. \citet{Jabonska2026} described a non-Hebbian inhibitory long-term depression operating at VIP synapses that may contribute to longitudinal plasticity. The picture that emerges from these developmental and longitudinal data is that the in-vivo VIP signal carries learned, plastic structure on top of its state-dependent recruitment.

\section{Methodological caveats: VIP-Cre line, ChAT subset, modality choice}

Many cross-study conflicts in the in-vivo literature trace to methodological differences rather than to genuine biological disagreement. The VIP-Cre driver line labels not just classical bipolar VIP/CR cells but also basket-cell-like VIP/CCK populations and the small ChAT-VIP subset \citep{Obermayer2019, Dudai2020, Granger2020, Porter1999}, and the proportion of ChAT-VIP cells captured varies across reports. Recordings made with electrophysiology versus two-photon calcium imaging weight different temporal frequencies of VIP activity differently, and dual-polarity voltage imaging \citep{Kannan2022} reveals state-dependent antagonisms between VIP and SST that conventional calcium imaging may smooth out. Layer specificity also matters --- VIP cells in superficial layers behave differently from those in deeper layers --- and many in-vivo papers either restrict imaging to layer 2/3 or pool layers without comment. Driver-line off-targets, anatomical-vs-functional methodological mismatches, and modality choice are flagged as methodological caveats throughout the review (see also \href{/disinhibition-framework}{Local Circuit Motifs and the Disinhibition Framework}); explicit reporting of these covariates would resolve a substantial fraction of the residual cross-study heterogeneity \citep{Obermayer2019, Dudai2020, Kannan2022, Yavorska2021}.

\section{Cross-species and human-relevant in-vivo VIP function}

Human-relevant evidence for the in-vivo VIP signal comes primarily from invasive and ex-vivo work in human neocortex and from genetic-perturbation models. Foundational studies \citep{Poorthuis2018} showed that human-cortical layer-1 GABAergic interneurons --- the human homologs of the rodent VIP class --- exhibit pronounced serotonergic (5-HT3A-mediated) responses, indicating cross-species conservation of the modulator--target axis. \citet{Koukouli2017} showed that mice carrying the human $\alpha$5-nAChR rs16969968 risk SNP have reduced VIP-interneuron activation, providing a translational route from human genetics to in-vivo VIP function. Cross-species and human-relevant work on VIP cell-type composition is taken up in detail in \href{/cross-areas}{VIP Interneurons Across Brain Regions}; here the relevant point is that the in-vivo phenomenology and modulator dependence of VIP cells generalise at least partially to primate and human cortex, supporting the relevance of the rodent literature to behavioural neurobiology beyond the mouse \citep{Poorthuis2018, Koukouli2017, Lecrux2016}.

\section{Quantitative dynamics: kinetics, magnitude, and population structure}

Beyond the qualitative direction of VIP recruitment, in-vivo work has begun to constrain its kinetics, magnitude, and population structure. \citet{Hangya2015} identified a basal-forebrain cholinergic reinforcement signal with {\textasciitilde}18 ms latency that is well matched to the rapid VIP responses to reward and punishment reported by \citet{Pi2013} and \citet{Szadai2022}. \citet{Garciajuncoclemente2017} documented VIP-driven inhibition of pyramidal neurons on similarly short timescales during arousal, indicating that the direct VIP\rightarrowPyr channel can operate fast enough to shape behaviourally relevant gain transitions. The magnitude of locomotion-evoked V1 VIP responses reported by \citet{Fu2014}, \citet{Pakan2016}, and \citet{Dipoppa2018} differs substantially across studies, reflecting differences in indicator (GCaMP variants), depth, layer, and behavioural protocol; the qualitative compilation in Figure~\ref{fig-vip-behavior} foregrounds direction rather than magnitude precisely because cross-study magnitude comparisons remain unaudited at the panel level. Population-structure constraints have also tightened. \citet{Kannan2022} reported state-dependent antagonism within the VIP population, refining the textbook reading of VIP cells as a class-level monolith. These dynamics constraints feed directly into the computational reading developed in \href{/computational-models}{Computational Models of VIP Circuit Function}.

\section{Stress, disease, and pharmacological perturbations of in-vivo VIP function}

In-vivo VIP recruitment is altered by stress, disease, and pharmacological manipulations in ways that further inform the disinhibitory reading. \citet{Goff2023} showed that in Scn1a+/- Dravet mice, VIP recruitment at quiet-to-running transitions is diminished and that optogenetic VIP activation rescues part of the deficit, with \citet{Liebergall2024} showing that NDNF cells are spared in the same model. \citet{Motaharinia2021} reported that chemogenetic augmentation of VIP excitability after photothrombotic stroke promotes recovery of forelimb function. Established work \citep{Leroy2022} showed that hippocampal VIP cells release enkephalin to modulate stress-related behaviour. Earlier work \citep{Silverstein2025} reported that repeated ventral-hippocampus input persistently depresses mPFC VIP activity, contributing to working-memory deficits --- a cumulative-stress phenotype. These manipulations indicate that VIP recruitment is not only state-dependent moment-to-moment but is itself plastic on longer timescales and is altered in disease in ways that modify the in-vivo phenomena summarised above. The translational implications of these findings are taken up where appropriate in later sections.

\section{Summary and forward links}

The in-vivo VIP signal is reliably recruited by locomotion, arousal, and reinforcement in V1 and a number of cortical territories, and that recruitment is causally linked to gain modulation and to attentional and learning-related disinhibition. It is, however, neither uniform across area, state, and modulator nor universally explained by a single VIP\rightarrowSST\rightarrowPyr motif: SST cells are co-activated rather than silenced in many V1 preparations \citep{Pakan2016, Dipoppa2018}; the auditory cortex implements an opposite-sign behavioural effect via a non-disinhibitory route \citep{Bigelow2019, Yavorska2021}; mPFC and ACC emphasise outcome and delay activity \citep{Pinto2015, Kamigaki2017}; S1 supports stimulus-driven and licking-related VIP responses \citep{Yu2019, Ramamurthy2023, Sachidhanandam2016}; and hippocampal IS-cell subtypes show subregion-specific signs \citep{Francavilla2018a, Luo2020, Lenkey2025}. The decision schematic in Figure~\ref{fig-vip-state-decision-tree} operationalises the area$\times$state$\times$modulator framework that organises these heterogeneous data and frames the cross-area catalogue developed in \href{/cross-areas}{VIP Interneurons Across Brain Regions} and the synthesis closure in the concluding synthesis. Two questions remain explicitly unmet at the level of the present evidence base: the quantitative shape of the arousal--VIP--performance relation (monotonic vs inverted-U) and the magnitude of VIP control over hippocampal place-cell coding; both are flagged as gap items below.

\section{Outstanding gaps and unmet citation needs}

Three quantitative claims from earlier work cannot be supported by the filtered citation map and are reported here as \textbf{unmet citation needs}: (1) the inverted-U shape of the arousal--VIP--performance relation as articulated in conflict-3 of earlier work (paper key Beerendonk2024); (2) the hippocampal place-field minimal-effect statement attributed to paper key Vervaeke2024; and (3) the heterogeneity-vs-uniformity claim attributed to paper key Johnson2020 --- for which \citet{Ramamurthy2023, Ramamurthy2025, Kannan2022} provide partial in-set substitutes but do not exactly recapitulate the original Pi-vs-Johnson framing. Earlier-work paper key Yu2021 was substituted with the in-set Previous reports \citep{Yu2019} finding on whisker-stimulus encoding by S1 VIP cells, which carries the same methodological direction. The \citet{Pi2013} vs Ramamurthy heterogeneity contrast and the \citet{Fu2014} vs \citet{Pakan2016} SST direction contrast are well supported by the in-set evidence and are presented as full conflict admonitions above; the remaining two conflicts from earlier work are presented as gap-flagged conflict admonitions. A future revision should incorporate the missing primary sources to upgrade these qualitative gap statements to quantitative conflict resolutions. The cross-area heterogeneity that this section foregrounds is catalogued comprehensively in \href{/cross-areas}{VIP Interneurons Across Brain Regions}, and the synthesis-and-reassessment closure of the review takes up the residual conflicts directly in the concluding synthesis.


\section{VIP Interneurons Across Brain Regions}
\label{sec:09_cross_areas}

VIP function is not uniform across the brain. The disinhibitory VIP\rightarrowSST\rightarrowPyr disinhibitory motif first formalised by \citep{Pi2013} in auditory cortex and medial prefrontal cortex captures a real and reproducible circuit, but contemporary multi-area recordings show that the \textit{operation} implemented by that motif --- and in some structures, the motif itself --- varies systematically with cortical area, behavioural context, and developmental origin \citep{Pi2013, Karnani2016b, Dipoppa2018, Lee2013, Bastos2023a}. Primary visual cortex emphasises locomotion-gated, contrast-dependent gain enhancement \citep{Dipoppa2018, Millman2020, Pakan2016}; auditory cortex deploys VIP cells for reinforcement-linked disinhibition and critical-period plasticity \citep{Pi2013}; somatosensory cortex recruits them via motor copy during active sensing \citep{Lee2013, Sermet2019}; medial prefrontal and anterior cingulate cortex implement top-down control of attention, working memory, and pain \citep{Bastos2023a, Anastasiades2021a, Kamigaki2017, Li2022a}; and hippocampal VIP/calretinin ``IS-3'' cells assemble a circuit motif distinct from neocortex, targeting other interneurons rather than principal cells \citep{Gulyas1996, Tyan2014, Tricoire2011}. Subcortical VIP populations --- in basolateral amygdala, dorsal striatum, and the suprachiasmatic nucleus --- extend the diversity further, including a non-GABAergic, peptidergic projection role in circadian timekeeping \citep{Rhomberg2018, Vereczki2021, Munozmanchado2018, Maywood2006, Cutler2003}. These cross-region differences place strong constraints on any unitary theory of VIP function and motivate the area-by-area treatment that follows. We refer to the previous section's behavioural heterogeneity (\href{/in-vivo-behavior}{In Vivo Function During Behavior}) and forward to the temporal-coordination signatures examined in \href{/oscillatory-dynamics}{Oscillatory Dynamics and Temporal Coordination}.

\section{Primary visual cortex: locomotion gating and contrast-dependent gain}

The most extensively characterised cortical VIP population sits in mouse V1, where VIP+ small bipolar/multipolar cells with vertically descending axonal arbors target apical dendrites of L5 pyramidal neurons and innervate L2/3 SST cells \citep{Kawaguchi1996, Gonchar2008, Kawaguchi1997a}. Long-range tracing places these cells at the receiving end of cholinergic, serotonergic, and frontal cortical inputs disproportionately enriched relative to PV/SST cells \citep{Wall2016}. Two-photon imaging during head-fixed running has converged on a basic phenomenology: VIP cells in V1 show the strongest run-related rise of any genetically defined interneuron class, while SST cells move in the opposite direction, and PV cells track stimulus contrast more than locomotion \citep{Dipoppa2018, Pakan2016, Karnani2016b}. The resulting state-dependent VIP\rightarrowSST motif underlies a multiplicative or additive rescaling of pyramidal responses that is widely interpreted as locomotion-driven gain enhancement \citep{Pi2013, Dipoppa2018, Karnani2016b, Millman2020}.

Three lines of evidence complicate this disinhibitory reading. First, the locomotion drive itself is context-modulated: \citet{Pakan2016} reported that VIP↑/SST↓ antiphasic modulation in V1 is most pronounced when concurrent visual stimuli are present, although most VIP cells remain locomotion-responsive in darkness, in agreement with the more invariant locomotion drive reported by \citet{Dipoppa2018}. Second, the gain modulation is not a uniform multiplicative rescaling: \citet{Millman2020} showed that V1 VIP activation enhances responses preferentially at \textit{low} contrasts and lowers contrast threshold without changing peak firing at saturating contrasts, and that this contrast-dependent enhancement requires intact SST inhibition. Third, the spatial structure of disinhibition is not blanket: \citet{Karnani2016b} documented a localised ``spotlight'' of VIP-mediated disinhibition that opens spatially restricted holes in SST inhibition rather than uniformly suppressing it. These findings together argue that V1 VIP cells implement a \textit{contextual, spatially structured, contrast-dependent} operation rather than the uniform broadcast disinhibition assumed by the textbook motif.

\begin{framed}
\textbf{Conflict --- uniform versus localised VIP disinhibition}\\
\citet{Pi2013} formalised VIP\rightarrowSST as a broad disinhibitory module that uniformly gates pyramidal output across cortex. \citet{Karnani2016b} instead reported that VIP-mediated disinhibition is spatially restricted, with small populations of activated VIP cells suppressing nearby SST cells without affecting distant SST cells, opening localised ``holes'' in the SST blanket. Resolution status: open. The two views are not formally incompatible --- uniformity at the population level may still arise from many local ``spotlights'' --- but they predict different effects of focal versus distributed VIP recruitment, and existing optogenetic stimuli tend to bias the system toward one regime rather than dissecting it.
\end{framed}

\begin{framed}
\textbf{Conflict --- uniform versus contrast-dependent VIP gain}\\
\citet{Pi2013} framed VIP recruitment as enabling a state-dependent multiplicative gain on pyramidal firing, consistent with reinforcement-linked broadcast disinhibition. \citet{Millman2020} instead found that V1 VIP activation selectively enhances responses to low-contrast stimuli and does not scale high-contrast responses, with the effect requiring SST circuitry. Resolution status: open. The contrast-dependent reading is more easily reconciled with normalisation-style accounts of cortical gain (\href{/computational-models}{Computational Models of VIP Circuit Function}) than with a uniform multiplicative interpretation, but quantitative cross-area replication of the contrast-dependence is not yet available.
\end{framed}

\begin{framed}
\textbf{Conflict --- locomotion-driven VIP activity in V1}\\
\citet{Dipoppa2018} reported that locomotion robustly drove V1 VIP cells across stimulus conditions, supporting locomotion as a primary VIP-engaging variable. \citet{Pakan2016} reported that VIP locomotion modulation was modulated by visual context, with VIP and SST showing antiphase modulation that was strongest when running was paired with visual drive although VIP cells remained locomotion-responsive in darkness. Resolution status: open. The discrepancy is consistent with differences in stimulus protocols (drifting gratings versus blank/dark trials) and in baseline arousal, and it cautions against treating locomotion alone as a sufficient predictor of V1 VIP firing.
\end{framed}

The morphological and laminar substrate constrains how broadly V1 VIP cells can act. \citet{Gonchar2008} documented that L2/3 small bipolar VIP cells display dense vertical axonal arbors descending to L5, consistent with apical-dendrite-targeting onto pyramidal neurons in addition to L2/3 SST cells, and \citet{Kawaguchi1996} and \citet{Kawaguchi1997a} established that the bipolar/double-bouquet morphology distinguishes VIP cells from PV multipolar baskets and SST Martinotti cells \citep{Kawaguchi1996, Kawaguchi1997a, Gonchar2008}. The cells therefore have direct anatomical access to apical dendrites of L5 pyramidal neurons, providing one route by which top-down signals from cingulate cortex can be funnelled into deep-layer pyramidal output via VIP recruitment \citep{Bastos2023a, Wall2016}. Differential thalamocortical drive is also relevant: \citet{Ji2016} reported broadly similar laminar innervation patterns for auditory and visual thalamocortical pathways onto cortical interneurons, with quantitative variation across interneuron classes --- including VIP --- that is consistent with the input-specificity differences seen across cortical areas \citep{Wall2016, Sermet2019}.

V1 VIP cells are also a substrate for top-down context. Cingulate-to-V1 long-range projections preferentially synapse onto a subset of V1 VIP cells, and optogenetic silencing of those projections degrades figure--ground discrimination behaviourally. The targeting bias is consistent with the rabies-tracing analysis of \citet{Wall2016}, which showed that VIP cells across cortex receive disproportionately enriched long-range input from frontal, neuromodulatory, and thalamic sources relative to PV/SST cells \citep{Wall2016, Anastasiades2021a}. The morphological substrate identified by \citet{Kawaguchi1996} and \citet{Gonchar2008} --- bipolar VIP cells with vertical apical-targeting axons --- provides the direct anatomical route from this top-down input onto deep-layer pyramidal output \citep{Kawaguchi1996, Gonchar2008, Bastos2023a}. This anchors the V1 motif to broader cortical hierarchies and links it to the frontal subsections below (\href{/in-vivo-behavior}{In Vivo Function During Behavior}).

\section{Auditory cortex: reinforcement, plasticity gating, and the polarity puzzle}

Auditory cortex is the disinhibitory home of the VIP\rightarrowSST\rightarrowPyr motif. \citet{Pi2013} performed the foundational experiment: optogenetic activation of A1 VIP cells transiently suppressed primarily SST and a fraction of PV interneurons, disinhibiting pyramidal cells, and silencing VIP cells reduced pyramidal firing during reinforcement, identifying VIP cells as a state-dependent gain mechanism activated by reward and punishment \citep{Pi2013}. The motif is layer-biased, with strongest disinhibitory effects in L2/3, and is supported by convergent long-range cortical and subcortical drive consistent with VIP cells being broadcast targets of neuromodulatory and top-down signals \citep{Pi2013, Wall2016}.

Two extensions of the disinhibitory A1 motif organise much of the subsequent literature. First, BLA VIP cells implement \textit{adaptive disinhibition} during associative learning: \citet{Krabbe2019} showed that VIP activity in basolateral amygdala grows progressively across fear-conditioning sessions, that VIP cells acquire CS+/CS- selectivity, and that silencing them during conditioning impairs memory consolidation, identifying disinhibition itself as a plastic substrate rather than a static gain \citep{Krabbe2019}. The drive into these cells comes from the basal nucleus and thalamus, providing an arousal-linked substrate for state-dependent recruitment \citep{Krabbe2019, Wall2016}. Second, a layer-1 microcircuit gates auditory critical-period plasticity: \citet{Takesian2018} showed that L1 interneurons receive nicotinic cholinergic input from basal forebrain and inhibit L2/3 PV cells, and that recruiting these L1 cells in adults reopens tonotopic plasticity --- extending the motif from instantaneous gain to developmental and recovery time-scales \citep{Takesian2018}.

Auditory cortex is also where the cross-area ``polarity puzzle'' becomes acute. In V1, locomotion drives VIP cells \textit{up} and improves visual encoding \citep{Dipoppa2018, Pakan2016}; in A1, the parallel literature reports that locomotion \textit{suppresses} spiking and reduces encoding efficiency, opposite in sign to V1, despite the apparently identical disinhibitory motif. The keys reporting that A1 polarity reversal directly (Fu 2014; Bigelow 2019) lie in the temporal-coordination cluster covered in \href{/oscillatory-dynamics}{Oscillatory Dynamics and Temporal Coordination} and are not present in this section's filtered citation set; we flag this as an unmet citation need (see Caveats and unmet citation needs below). Within the present packet, the available evidence shows that A1 VIP cells are robustly recruited by \textit{reinforcement} and \textit{learning} drives \citep{Pi2013, Takesian2018, Wall2016} rather than by locomotion per se, which is consistent with --- but does not in isolation prove --- an A1/V1 polarity divergence in locomotion gating relative to V1 \citep{Dipoppa2018, Pakan2016}.

\section{Somatosensory cortex: motor copy and orthogonal attention}

Barrel cortex provides the clearest example of motor-context gating of VIP cells. \citet{Lee2013} showed that whisking-induced motor inputs from M1 activate S1 VIP cells, which inhibit SST cells and disinhibit L2/3 pyramidal neurons during active sensing; optogenetic VIP silencing during active whisking reduced pyramidal firing while VIP activation mimicked the whisking-induced enhancement, establishing the VIP\rightarrowSST motif as disinhibitory in S1 as well as A1 and mPFC \citep{Lee2013, Pi2013}. Importantly, the recruiting input is the \textit{motor command}, not passive whisker deflection --- VIP cells in S1 respond to what the animal is \textit{doing} rather than only to what it senses \citep{Lee2013}.

Thalamic input specificity refines this picture. \citet{Sermet2019} mapped POm and VPM inputs and found that S1 VIP cells receive intermediate input from both higher-order POm and first-order VPM pathways, with quantitative differences across layers and cell classes that contrast with PV/SST distributions and provide a circuit substrate for integrating motor context and sensory feedback \citep{Sermet2019, Lee2013}. \citet{Donato2013} linked VIP signalling to longer-time-scale plasticity, showing that experience-dependent disinhibition mediated by VIP activity regulates PV basket-cell network maturation, embedding the motif in critical-period and adult plasticity (\href{/disease-translational}{Species Differences, Human Relevance, and Disease}).

The S1 motif also carries a developmental and plasticity dimension. \citet{Donato2013} showed that experience-dependent disinhibition mediated by VIP signalling regulates PV basket-cell network maturation, linking VIP function in S1 to adult cortical plasticity beyond the moment-to-moment disinhibitory effect of motor copy \citep{Donato2013, Lee2013}. This places the S1 motif on a continuum with the auditory critical-period work of \citet{Takesian2018} and the experience-dependent disinhibition results in the broader cortical literature, suggesting that VIP cells in primary sensory areas serve as state-dependent gates of \textit{plasticity} on multiple time-scales rather than only as instantaneous gain controllers \citep{Takesian2018, Donato2013, Krabbe2019}.

The behavioural causal mapping in S1 is partially independent of attention. \citet{Kuchibhotla2017} showed that disrupting PV, SST, and VIP cells produces dissociable, behaviour-specific deficits --- VIP loss impairing context-dependent gating rather than basic discrimination --- supporting separable inhibitory channels. \citet{Myersjoseph2024} then directly tested whether VIP-mediated disinhibition implements cross-modal attentional gain in mouse cortex and found it largely orthogonal: attentional modulation and VIP disinhibition operate on partly separable axes, challenging the textbook reading that VIP cells \textit{are} the attention gate \citep{Myersjoseph2024, Pi2013}.

\begin{framed}
\textbf{Conflict --- VIP and attention}\\
\citet{Pi2013} argued that VIP\rightarrowSST disinhibition gates attention-related gain via reinforcement signalling, an interpretation echoed in subsequent reviews of cortical disinhibition. \citet{Myersjoseph2024} instead reported that VIP-mediated disinhibition is largely orthogonal to cross-modal attentional modulation in mouse cortex, with disinhibition and attention relying on partly separable circuits. Resolution status: open. The orthogonality result does not exclude VIP involvement in some attentional regimes (e.g., reinforcement-linked top-down gain), but it constrains ``VIP = attention'' framings derived from the disinhibitory 2013/2014 papers.
\end{framed}

\section{Frontal cortex: top-down control, working memory, and pain}

Frontal cortex elaborates the disinhibitory motif into long-time-scale, goal-directed operations. The proportional rebalancing across the cortical hierarchy reported by \citet{Torresgomez2020} is informative as a starting point: PV proportions decrease and calretinin proportions (treated here as a VIP-overlapping marker) increase along the sensory-to-executive axis, suggesting an architectural shift rather than a simple scaling of the same circuit \citep{Torresgomez2020}. Long-range outputs also expand: \citet{Lee2014b} identified a subset of mPFC VIP-expressing GABAergic neurons that send axons to the nucleus accumbens, defying the classical ``local interneuron'' definition and providing a peptidergic GABAergic projection from PFC into reward circuitry \citep{Lee2014b}.

Working-memory experiments give the clearest cell-type-resolved picture of mPFC VIP function. \citet{Kamigaki2017} recorded mPFC VIP, PV, and SST cells through delayed-choice tasks and showed that VIP cells display sustained delay-period activity, with optogenetic activation of VIP cells during the delay enhancing choice accuracy and activation of SST or PV cells impairing performance; PV and SST cells showed distinct delay-period dynamics consistent with cell-type-specific contributions to memory phases \citep{Kamigaki2017, Pi2013}. The thalamic gating of these dynamics has been resolved at the input level by \citet{Anastasiades2021a}, who showed that mediodorsal and ventromedial thalamic afferents engage distinct L1 microcircuits in PFC: mediodorsal input preferentially recruits VIP cells while ventromedial input recruits NDNF cells, with NDNF cells in turn directly inhibiting L5 PT dendrites and VIP cells driving disinhibition through SST suppression \citep{Anastasiades2021a}.

mPFC VIP cells are also a substrate for affective and chronic-state physiology. \citet{Li2022a} reported that activation of prelimbic mPFC VIP cells ameliorates neuropathic pain in mice via a disinhibitory mechanism that re-engages mPFC pyramidal output to anterior cingulate; nerve injury reduced VIP-mediated disinhibition and decreased descending pyramidal output, and chemogenetic restoration of VIP activity recovered mechanical and thermal thresholds toward sham levels \citep{Li2022a, Pi2013}. The authors themselves note that this finding has not yet been independently replicated in other chronic pain models or laboratories, and we mark it as a single-paper claim awaiting confirmation \citep{Li2022a}.

A second class-bending observation in PFC concerns transmitter identity. \citet{Obermayer2019} characterised a subset of mPFC VIP interneurons that co-express choline acetyltransferase and that release acetylcholine onto neighbouring pyramidal neurons via nicotinic receptors, in addition to GABAergic transmission, providing a local intracortical source of cholinergic excitation independent of basal forebrain projections \citep{Obermayer2019, Wall2016}. The original report relies on a specific Cre line and the finding is flagged in the source as needing replication, but if confirmed it expands the source repertoire of cortical neuromodulation and complicates a strictly GABAergic reading of VIP function \citep{Obermayer2019}. Top-down signalling into V1 from cingulate, discussed above, then closes a loop: \citet{Bastos2023a} showed cingulate\rightarrowV1 projections preferentially target V1 VIP cells, providing a circuit substrate for context-dependent visual processing and indicating that frontal cortex \textit{uses} the V1 disinhibitory motif as a downstream effector rather than implementing a separate one \citep{Bastos2023a, Wall2016}.

The cell-type-resolved picture of mPFC therefore knits together at least four distinct VIP operations: a disinhibitory L2/3 disinhibitory motif inherited from sensory cortex \citep{Pi2013, Anastasiades2021a}, sustained delay-period working-memory dynamics \citep{Kamigaki2017}, long-range projection to nucleus accumbens that exits the local-interneuron framing \citep{Lee2014b}, and a candidate intracortical cholinergic source via ChAT-VIP cells \citep{Obermayer2019}. None of these is mutually exclusive, but no current model integrates them; the section's synthesis below treats this as a definitional rather than a quantitative gap.

\section{Motor cortex: preparatory dynamics and transcriptomic identity}

Motor cortex VIP cells are less studied than their sensory counterparts but contribute two distinct kinds of constraint. \citet{Arroyo2023} recorded VIP cells in mouse motor cortex through skilled-reach learning and reported that VIP cells develop preparatory activity and acquire selectivity for upcoming actions in parallel with pyramidal preparatory dynamics, suggesting that VIP recruitment is a learned component of motor preparation rather than a simple reactive disinhibition; the authors flag the result as a single-paper claim awaiting independent confirmation in other motor regions and tasks \citep{Arroyo2023}. \citet{Scala2021} resolves the cellular substrate at the molecular level: single-cell transcriptomics of mouse M1 splits VIP cells into Vip-Mybpc1, Vip-Lect1 and other subtypes whose transcriptional profiles align with electrophysiologically defined fast-adapting and irregular-spiking subclasses, and whose homologues are conserved across species \citep{Scala2021}. The transcriptomic decomposition implies that ``VIP cells in M1'' is not a single functional unit and that any cross-area generalisation needs to specify which transcriptomic subtype is being assayed \citep{Scala2021}.

\section{Hippocampus: IS-3 cells and a circuit motif distinct from neocortex}

Hippocampal VIP cells are not a small variant of the cortical class --- they implement a different motif. \citet{Gulyas1996} first established the principle that calretinin-containing hippocampal interneurons (which overlap heavily with VIP+ cells) selectively target other GABAergic cells --- including calbindin+ and VIP+ interneuron populations --- defining ``interneuron-selective'' as a distinct class that controls other inhibitory cells rather than principal neurons \citep{Gulyas1996}. \citet{Tyan2014} resolved the modern instantiation: VIP+/CR+ ``interneuron-specific type 3'' (IS-3) cells in CA1 provide selective dendritic inhibition to OLM cells, suggesting that IS-3 cells regulate the gain on OLM-mediated dendritic inhibition rather than acting as a continuous instantaneous gain knob \citep{Tyan2014, Gulyas1996}.

Hippocampal VIP cells are not, however, monolithic. \citet{Kawaguchi1996} distinguished two morphological/functional groups within the cortical VIP population: small basket cells targeting pyramidal somata and bipolar cells with distinct firing properties, providing the original framework later extended to the hippocampal IS cell taxonomy \citep{Kawaguchi1996, Kawaguchi1997a}. Developmentally, hippocampal VIP cells originate predominantly from the caudal ganglionic eminence and migrate in distinct spatiotemporal waves to populate stratum oriens and pyramidale, with subtypes diverging embryonically and emerging later than MGE-derived classes \citep{Tricoire2011, Miyoshi2010, Vucurovic2010}. The shared CGE origin (5HT3aR+) of cortical and hippocampal VIP cells is therefore a developmental anchor that does \textit{not} by itself imply functional homology --- the same lineage produces a cortical VIP\rightarrowSST\rightarrowPyr motif and a hippocampal IS-3\rightarrowOLM motif whose targets and rhythm couplings differ substantially \citep{Tricoire2011, Miyoshi2010, Tyan2014, Gulyas1996}.

The downstream consequence is that hippocampal VIP function cannot be modelled by transferring the cortical motif and replacing layer labels. The figure schematic below (Figure~\ref{fig-hippocampal-vip-motif}) and the cross-region motif map (Figure~\ref{fig-vip-across-areas}) make this divergence explicit. The contemporary debate over whether VIP/IS interneurons act as a deterministic gate on hippocampal remapping or whether their efferent inhibition is itself dynamically reset by upstream LTD (Neubrandt 2025; Jablonska 2026) sits in the very-recent literature and is not represented in this section's citation set; we flag this as an unmet citation need below (\href{/cross-areas}{VIP Interneurons Across Brain Regions}).

\section{Subcortical VIP populations: amygdala, striatum, and a peptidergic clock}

Outside neocortex and hippocampus, VIP cells appear in three structures with very different functional roles, expanding the meaning of the term ``VIP cell'' beyond a single GABAergic motif.

In basolateral amygdala, VIP cells closely resemble cortical VIP cells in marker profile, sparse density, and bipolar morphology, and they implement a recognisably cortex-like motif. \citet{Rhomberg2018} showed that VIP-immunoreactive interneurons in BLA selectively innervate other GABAergic interneurons --- including interneuron-selective, basket, and neurogliaform cells --- paralleling the cortical disinhibitory motif \citep{Rhomberg2018, Pi2013}. Quantitatively, \citet{Vereczki2021} showed by stereological counts that VIP+/CR+ interneuron-selective cells comprise roughly 29--38\% of GABAergic interneurons in mouse lateral and basal amygdala, a proportion exceeding typical neocortical estimates \citep{Vereczki2021}. The BLA VIP population is therefore best read as a cortical-like outpost of the disinhibitory motif, embedded in fear/threat circuitry.

Striatum tells a different story. \citet{Munozmanchado2018} used single-cell RNA-seq of dorsal striatum to identify a previously unrecognised VIP/Cck-positive interneuron population, distinct from disinhibitory PV, SST, and cholinergic types, raising the question of whether striatal VIP+ cells share developmental origins with their cortical counterparts \citep{Munozmanchado2018}. The functional role of these cells in striatal microcircuits is largely unmapped and is among the most under-studied cell types in the contemporary striatal taxonomy \citep{Munozmanchado2018}.

The amygdala and striatum cases together establish that the same molecular tag picks out cells with different \textit{functional contexts} even when the local motif looks similar. BLA VIP cells implement a cortex-like disinhibitory motif during fear-conditioning circuitry \citep{Rhomberg2018, Vereczki2021}, while striatal VIP/Calb2 cells --- molecularly closer to cortical VIP than to other striatal types --- sit in a network whose disinhibitory PV/SST/cholinergic taxonomy does not predict where this fourth class fits behaviourally \citep{Miyoshi2010}. The functional gap for striatal VIP cells is one of the clearest in the contemporary subcortical interneuron literature.

The suprachiasmatic nucleus (SCN) is the most striking departure from the cortical reading. SCN VIP cells are \textit{peptidergic projection neurons} whose primary function is paracrine synchronisation of cellular circadian oscillators, not local GABAergic disinhibition. \citet{Maywood2006} and \citet{Maywood2011} showed that VIP/VPAC2 paracrine signalling synchronises circadian gene expression among SCN neurons and that VIP/VPAC2 is the dominant non-redundant pathway maintaining cellular synchrony, with co-cultured wild-type and VIP-knockout slices partially rescuing rhythmicity \citep{Maywood2006, Maywood2011}. \citet{Cutler2003} demonstrated that the VPAC2 receptor on SCN neurons is required for cellular circadian rhythmicity in vitro \citep{Cutler2003}. \citet{Webb2009} showed that individual SCN neurons can generate intrinsic but noisy circadian oscillations even when isolated, and that network interactions among SCN neurons stabilise these cycling neurons into coherent ensemble output \citep{Webb2009}. Causally, \citet{Jones2018} used chemogenetic silencing of SCN VIP cells during light pulses to block normal phase resetting of the circadian clock, identifying VIP activity as necessary for photic resetting; \citet{Todd2020} showed that genetic ablation of SCN VIP neurons disrupts rhythmicity and decouples SCN from peripheral oscillators \citep{Jones2018, Todd2020, Maywood2006}. Cross-region modelling of ``VIP function'' that aggregates SCN with neocortex therefore conflates a peptidergic projection role with a GABAergic disinhibitory one and must be treated cautiously in any unitary theory (\href{/computational-models}{Computational Models of VIP Circuit Function}).

The SCN case has additional structural specificity. \citet{Maywood2011} showed that combined deletion of VIP together with other SCN peptides (GRP, AVP) compounds rhythm deficits, but VIP/VPAC2 alone causes the most severe arrhythmia, indicating a partial peptidergic redundancy in which VIP is dominant rather than exclusive \citep{Maywood2011, Maywood2006, Cutler2003}. Combined with the optogenetic causal evidence that SCN VIP activity is necessary and sufficient for photic-like phase resetting \citep{Jones2018, Todd2020}, the SCN VIP population is best read as a \textit{peptidergic projection node} whose role bears no microcircuit homology to the cortical disinhibitory motif. The tendency in the cortical literature to label this same molecular tag ``VIP'' therefore conflates two operationally distinct classes; for clarity in cross-region modelling, downstream sections (notably \href{/computational-models}{Computational Models of VIP Circuit Function} and \href{/disease-translational}{Species Differences, Human Relevance, and Disease}) should treat SCN VIP cells as a separate population rather than a non-cortical instance of the same circuit primitive.

\section{Cross-region synthesis}

Three regularities emerge when the area-by-area material is laid side by side (Figure~\ref{fig-vip-across-areas}).

First, the developmental and molecular identity of VIP cells is far more conserved across regions than their \textit{function}. CGE origin and 5HT3aR expression mark VIP+ cells from cortex, hippocampus, amygdala, and striatum \citep{Miyoshi2010, Vucurovic2010, Tricoire2011}, and transcriptomic taxonomies reveal homologous Vip-Mybpc1 / Vip-Lect1 subtypes across mouse cortical areas \citep{Scala2021}. Yet the same lineage implements opposite-sign locomotion gating in V1 versus A1 (V1 up, A1 down --- see \href{/oscillatory-dynamics}{Oscillatory Dynamics and Temporal Coordination}) \citep{Dipoppa2018, Pakan2016, Pi2013, Krabbe2019}, motor-copy disinhibition in S1 \citep{Lee2013, Sermet2019}, sustained working-memory dynamics in mPFC \citep{Kamigaki2017, Anastasiades2021a}, an interneuron-selective IS-3 motif in CA1 \citep{Tyan2014, Gulyas1996}, and peptidergic synchronisation in SCN \citep{Maywood2006, Cutler2003, Webb2009, Jones2018}. Shared lineage is not a sufficient predictor of shared function.

Second, the \textit{recruiting variable} is what differs most across regions, not the local microcircuit. Long-range tracing shows that VIP cells receive disproportionate long-range input from cholinergic, serotonergic, thalamic, and top-down cortical sources compared with PV/SST cells \citep{Wall2016, Anastasiades2021a}. The dominant \textit{driver} of VIP recruitment is therefore region-specific: locomotion and arousal in V1 \citep{Dipoppa2018, Pakan2016}; reinforcement and learning in A1 \citep{Pi2013}; motor command in S1 \citep{Lee2013, Sermet2019}; working-memory and thalamic context in mPFC \citep{Kamigaki2017, Anastasiades2021a}; light/photic input in SCN \citep{Jones2018, Todd2020}. Cross-area apparent conflicts in VIP ``polarity'' or ``selectivity'' largely dissolve when the recruiting variable is specified.

Third, the disinhibitory VIP\rightarrowSST\rightarrowPyr motif is one of several motifs that VIP cells implement, not the only one. The local cortical disinhibitory motif is reproducible in V1, A1, S1, mPFC, and BLA \citep{Pi2013, Lee2013, Dipoppa2018, Krabbe2019, Rhomberg2018}. But VIP cells also (i) send long-range axons to subcortical reward targets \citep{Lee2014b}, (ii) co-release acetylcholine in PFC \citep{Obermayer2019}, (iii) selectively target other interneurons rather than principal cells in hippocampus \citep{Tyan2014, Gulyas1996}, and (iv) act as peptidergic projection neurons in SCN \citep{Maywood2006, Cutler2003, Jones2018}. A unitary theory of ``VIP function'' therefore needs to enumerate motifs and specify which one applies in which region rather than treating ``VIP'' as a fixed circuit primitive.

These three regularities frame the temporal-coordination evidence in \href{/oscillatory-dynamics}{Oscillatory Dynamics and Temporal Coordination}: area-specific VIP motifs emit different temporal signatures, and many cross-area apparent contradictions in oscillatory recruitment are downstream consequences of the recruiting-variable differences described here rather than of microcircuit-level disagreement.

\begin{figure}[!htbp]
\centering
\includegraphics[width=0.95\linewidth]{files/fig-vip-across-areas-3a17fcc7484e1428c549e7c8e876b7d8.png}
\caption[]{VIP interneurons across brain regions. \textbf{Panel A}: Brain-area schematic summarising the dominant VIP motif per region (V1, A1, S1, M1, mPFC/ACC, CA1, BLA, striatum, SCN). Arrows indicate the dominant \textit{recruiting variable} and the local target of VIP inhibition or modulation. \textbf{Panel B}: Approximate VIP+ interneuron proportion (\% of GABAergic interneurons) across mouse neocortical areas, summarised as a textual table from harmonised literature ranges. Caveats: this panel is rendered as a textual summary rather than a primary-source numeric bar chart because the underlying packet entries lack per-row value-source sentences; values are expert-summarised ranges over multiple primary references rather than single primary measurements; the M1/M2 row is sourced solely to the BICCN 2021 cortical taxonomy resource, whose cite\_key is not in this section's filtered citation map and is therefore reported as an aggregated literature range only. Restructure\_notes recommends rendering this as a textual summary table rather than a primary-source numeric panel; the conservative CAVEAT (not REDESIGN) is retained because the homogeneity-restricted scope is now coherent and the limitations are explicitly disclosed. \textbf{Panel C}: Hippocampal VIP/IS-3 motif schematic (VIP/CR/CCK distinctions; OLM and bistratified targeting in CA1) \citep{Tyan2014, Gulyas1996}. \textbf{Panel D}: Cross-area conflict table summarising top-down, neuromodulatory, and sensory drivers of VIP recruitment by region \citep{Wall2016, Pi2013, Lee2013, Dipoppa2018, Bastos2023a, Anastasiades2021a, Maywood2006}. \textit{Required caveat (verbatim from Phase 6 audit):} ``Restructure correctly removed non-neocortical rows and harmonized the denominator. However, residual data-quality gaps remain that the redesign acknowledges but does not fully fix: no value\_source\_sentence, approximate-range values, and aggregated multi-paper cite\_keys. Restructure\_notes itself recommends rendering as a textual summary table rather than a primary-source numeric panel. Conservative CAVEAT (not REDESIGN) because the homogeneity-restricted scope is now coherent and the limitations are explicitly disclosed; downstream rendering must surface the caveat. || No per-row value\_source\_sentence; values are expert-summarized ranges. | study\_label\_text `Xu 2007' inconsistent with cite\_key `Xu2010a' (year mismatch). | `BICCN 2021; BICCN 2023' is not a specific resolvable citation in the per-row entry; only BICCN 2021 (10.1038/s41586-021-03950-0) is in the citation map. | Multi-paper cite\_key strings prevent per-value provenance.''}
\label{fig-vip-across-areas}
\end{figure}

\begin{figure}[!htbp]
\centering
\includegraphics[width=0.9\linewidth]{files/fig-hippocampal-vip--abd4edddd261df5123e0db74a16dae6e.png}
\caption[]{Hippocampal VIP/IS-3 motif distinct from neocortex (schematic). \textbf{Panel A}: CA1 lamellar cartoon with VIP/IS-3 cell soma in stratum oriens/pyramidale, axonal projections onto OLM and bistratified cells, and disinhibition of CA1 pyramidal cell distal-dendritic input by removal of OLM-mediated dendritic inhibition \citep{Gulyas1996}. \textbf{Panel B}: Comparison of cortical versus hippocampal VIP motifs --- cortical VIP cells are bipolar with descending axons targeting L5 SST cells \citep{Kawaguchi1996, Kawaguchi1997a, Gonchar2008}, whereas hippocampal IS-3 cells preferentially target other GABAergic interneurons rather than principal cells \citep{Gulyas1996}. \textbf{Panel C}: Subcortical VIP pathways summarised for cross-region completeness --- BLA VIP cells implement a cortex-like disinhibitory motif \citep{Rhomberg2018, Vereczki2021}, dorsal striatum contains a previously unrecognised CGE-derived VIP/Calb2 type, and SCN VIP cells act as peptidergic synchronisers via VIP/VPAC2 signalling \citep{Maywood2006, Cutler2003, Jones2018}. Schematic with no quantitative data substrate.}
\label{fig-hippocampal-vip-motif}
\end{figure}

\section{Closing synthesis: what cross-area variation constrains}

Taken together, the area-by-area evidence reviewed here makes three structural commitments. (i) The VIP\rightarrowSST\rightarrowPyr disinhibitory motif is a real and reproducible primitive in V1, A1, S1, mPFC, and BLA, but its \textit{operation} --- multiplicative versus contrast-dependent gain, broadcast versus localised disinhibition, instantaneous versus learned recruitment --- is region- and context-specific \citep{Pi2013, Lee2013, Dipoppa2018, Millman2020, Karnani2016b, Krabbe2019, Bastos2023a, Rhomberg2018}. (ii) The recruiting variable, not the local microcircuit, is what differs most across regions, and many cross-area apparent contradictions in VIP ``polarity'' or ``selectivity'' dissolve once the recruiting variable is specified \citep{Wall2016, Anastasiades2021a, Lee2013, Pakan2016, Jones2018}. (iii) Hippocampal IS-3 cells and SCN peptidergic VIP cells implement motifs distinct enough from the cortical motif that aggregating them into a single ``VIP function'' misrepresents the data \citep{Tyan2014, Gulyas1996, Maywood2006, Cutler2003, Jones2018}. These commitments motivate the temporal-coordination treatment in \href{/oscillatory-dynamics}{Oscillatory Dynamics and Temporal Coordination}, the species and disease comparisons in \href{/disease-translational}{Species Differences, Human Relevance, and Disease}, and the modelling reassessment in \href{/computational-models}{Computational Models of VIP Circuit Function} and the concluding synthesis.

\section{Caveats and unmet citation needs}

Three classes of citation need are unmet by this section's filtered citation map and are flagged for downstream resolution.

\begin{enumerate}
\item \textbf{V1/A1 locomotion polarity reversal} --- the direct claim that locomotion \textit{increases} V1 VIP activity and \textit{decreases} A1 spiking and encoding efficiency is commonly attributed to Fu 2014 (V1) and Bigelow 2019 (A1). Neither is present in this section's filtered citation map; the claim is therefore made indirectly here via the available V1-side (\citep{Dipoppa2018, Pakan2016}) and A1-side (\citep{Pi2013}) evidence. The direct polarity-reversal citations should be added when \href{/oscillatory-dynamics}{Oscillatory Dynamics and Temporal Coordination} is integrated.
\item \textbf{Shared-variability conflict in V1} --- the conflict between Xu 2026 and Lenkey 2025 over whether VIP gain control adds correlations or reduces shared variability in V1 cannot be rendered as an admonition block because neither key is in this section's filtered citation map. The conflict is documented in the unmet-citation-needs file.
\item \textbf{Hippocampal VIP gating versus upstream-LTD reset} --- the conflict between Neubrandt 2025 and Jablonska 2026 over whether VIP/IS interneurons provide a deterministic gate on hippocampal remapping or whether their efferent inhibition is itself dynamically reset by upstream LTD cannot be rendered here for the same reason. The conflict is documented in the unmet-citation-needs file.
\end{enumerate}

Two additional caveats apply to the figures themselves: the brain-region proportion table embedded in Figure~\ref{fig-vip-across-areas} Panel B is an expert-summarised range over multiple primary references rather than a primary-source numeric panel, and is rendered textually with the verbatim Phase 6 caveat above. The hippocampal motif figure (Figure~\ref{fig-hippocampal-vip-motif}) is a schematic with no quantitative data substrate.


\section{Oscillatory Dynamics and Temporal Coordination}
\label{sec:10_oscillatory_dynamics}

VIP cells contribute to cortical and hippocampal rhythms primarily as state-gating modulators rather than as primary rhythm generators. Across mouse cortex, hippocampus, and amygdala, VIP recruitment biases gamma power and inter-areal gamma coherence \citep{Veit2023, Wagatsuma2023, Knoblich2019}, gates theta-band coordination through hippocampal interneuron-specific circuits \citep{Tyan2014, Luo2020, Lee2019}, and is recruited differentially across Up-states, REM sleep, and arousal \citep{Tahvildari2012, Brecier2022, Rolle2025, Sabri2024}. The cells are also causally implicated in pathological synchrony --- spike-and-wave seizures, Fragile-X gamma hypersynchrony, Rett-syndrome phase--amplitude coupling, and Dravet-syndrome interictal activity --- though the direction of the effect depends sharply on perturbation polarity, brain area, and disease model \citep{Hall2015, Goel2025, Kranz2025, Goff2019}. The unifying observation is that bidirectional perturbation of VIP cells reshapes oscillation \textit{power, coherence, and phase coupling} far more reliably than it abolishes or generates the rhythm itself, with two consequences: (i) most rhythms continue when VIP cells are silenced, but their amplitude, inter-areal synchrony, and state dependence are altered \citep{Veit2023, Bastos2023a, Sabri2024, Lee2025}; and (ii) causal attribution remains contested --- bidirectional perturbation studies disagree on whether VIP cells are necessary, sufficient, or merely permissive for oscillatory state changes, with conflicts in particular over locomotion polarity across sensory cortices \citep{Fu2014, Pakan2016}, seizure susceptibility \citep{Hall2015}, Up-state generation \citep{Tahvildari2012, Kawaguchi2001}, and the VIP-versus-SST locus of cholinergic state changes \citep{Chen2015a, Szadai2022, Kuchibhotla2017}. The previous section's area-specific motifs (\href{/cross-areas}{VIP Interneurons Across Brain Regions}) emit different temporal signatures; this section treats those signatures as objects of study in their own right and forwards translational implications to \href{/disease-translational}{Species Differences, Human Relevance, and Disease}.

\section{Gamma rhythms: VIP as gain modulator, not pacemaker}

Cortical gamma (30--80 Hz) is dominated by reciprocal interactions between fast-spiking PV cells and pyramidal neurons, while VIP cells operate one synapse upstream by suppressing SST cells and, through that route, biasing gamma power and coherence. \citet{Veit2023} combined V1 optogenetics, LFP, and a computational model to show that VIP-cell activation linearly scales local gamma power \textit{without} altering its stimulus tuning, while suppressing long-range gamma coherence between cortical regions processing non-matched stimuli --- a dissociation between local-gain and global-coordination roles. Biophysical V1 layer-2/3 microcircuit models reproduce this geometry: PV and SST cells preferentially induce gamma (30--80 Hz) and beta (20--30 Hz) firing of pyramidal neurons respectively, while rapid VIP\rightarrowSST inhibition is required for attentional modulation of low-gamma (30--50 Hz) power, with VIP\rightarrowSST synaptic-weight reductions abolishing attentional gamma modulation entirely \citep{Wagatsuma2023}. A complementary two-layer V1 model treats VIP/PV/SST mutual inhibition as bistable switches that toggle pyramidal cells between SST-dominated low-firing/low-frequency and disinhibited high-firing/high-frequency oscillatory states; in that model, external drive to VIP cells \textit{increases} oscillation frequency, whereas VIP silencing only slightly raises frequency --- i.e., VIP is sufficient but not necessary for the high-frequency regime \citep{Hahn2022}. Multistable cortical models with two inhibitory classes likewise place VIP-driven SST modulation as the switch between coexisting gamma and beta regimes rather than as a rhythm generator \citep{Sarkar2025}.

These computational predictions align with two-photon and electrophysiological recordings of cell-type-specific synchronization in mouse V1. \citet{Knoblich2019} measured neuron--network coupling distributions in superficial V1 and found that PV and VIP cells are skewed toward strong coupling with the local population, while SST cells split into two distinct sub-distributions --- consistent with VIP cells embedded in a tightly coupled local gamma network rather than acting as independent oscillators. Excitatory/inhibitory ratio manipulations in spiking V1 microcircuit models extend the picture: PV reduction enhances both beta and gamma, SST reduction selectively impairs gamma, and VIP changes alter information-flow direction between Pyr and PV populations more than gamma magnitude itself. Synaptic-plasticity-equipped V1 column models of multiple inhibitory types show that plasticity facilitates oscillations across the gamma band, with VIP\rightarrowSST plasticity tuning the band edges. Optogenetic dissection in 4-aminopyridine GABAergic-synchrony preparations confirms the same hierarchy: PV or SST activation each independently triggers long-lasting discharges, while VIP activation does not --- and PV silencing strongly reduces discharge initiation, while VIP/SST silencing only partially attenuates it \citep{Bohannon2018}. The convergent reading is that VIP cells set the \textit{gain} of the gamma rhythm by adjusting SST/PV competition, but do not generate it.

The framing extends to disease and translational gamma \citep{Goel2025, Kranz2025}. \citet{Ichim2024} propose that PV-, VIP-, SST-, and NOS-class interneurons collectively exploit endogenous gamma to perform homeostatic vasomotor control --- a ``guardian-of-brain-health'' framework in which VIP cells couple gamma activity to neurovascular and clearance-related output rather than to rhythm genesis. Direct support comes from Prior work \citep{Murdock2024}: chemogenetic inhibition of cortical VIP interneurons attenuates 40-Hz multisensory-stimulation--induced amyloid clearance in 5XFAD mice without changing baseline 40-Hz LFP power, identifying VIP cells as a coupling node between gamma activity and glymphatic peptide release. Layer-2/3 cortical models of ketamine-induced gamma show, by contrast, that selectively reducing NMDA-R activity on PV or SST (but not VIP) interneurons reproduces ketamine-like increases in gamma power, and complementary MEG analysis (n = 12) showed flatter aperiodic slope and increased gamma correlated with regional PV and \textit{GRIN2D} expression \citep{Rademacher2025} --- implying that VIP cells gate the \textit{behavioural correlates} of gamma (clearance, attention) more directly than its NMDA-driven amplitude.

\section{Theta rhythms and the hippocampal IS3 motif}

Theta (4--10 Hz) in hippocampus is paced primarily by septal and entorhinal inputs interacting with OLM and other principal-cell-targeting interneurons, but VIP/calretinin ``IS3'' cells provide a specialised disinhibitory layer that controls the rate and timing of OLM output \citep{Tyan2014, Chamberland2010, Luo2020}. \citet{Tyan2014} showed by paired patch-clamp and ChR2 in CA1 oriens/alveus that VIP+/CR+ IS3 cells selectively innervate OLM cells via dendritic synapses, and that synchronous IS3 spikes can control OLM firing rate and spike timing, implicating IS3 cells in temporal patterning of theta-rhythmic OLM output. The anatomical substrate had been mapped two decades earlier: \citet{Gulyas1996} demonstrated that hippocampal calretinin-immunoreactive interneurons form interconnected dendro-dendritic clusters and selectively innervate other GABAergic targets including VIP-containing basket cells (2--5 contacts per axon; 2--10 CR-IR axons converging onto a single CR cell; up to 15 cells per cluster), providing a circuit substrate by which VIP-IS interneurons pace, rather than spike-time, principal cells. Developmental fate-mapping confirms that IS3 and CCK-VIP basket cells are CGE-derived, generated in a later neurogenic wave (E12--E16) than the MGE-derived PV/SST cells (E9--E12) and preferentially localised to superficial CA1 layers \citep{Tricoire2011}.

In awake mice, in-vivo recruitment of these cells is theta-locked but ripple-excluded. \citet{Luo2020} reported that VIP-IS cells fire preferentially during theta-run epochs at the rising phase and peak of the theta cycle but are \textit{not} recruited during sharp-wave ripples --- contradicting model predictions that had placed IS cells in the ripple-replay circuitry. Long-range VIP-GABAergic projection cells in CA1, by contrast, \textit{decrease} activity during theta-run epochs, are more active during quiet wakefulness, and remain uncoupled to ripples, defining a complementary ``long-range VIP-LRP'' mode distinct from local IS3 control \citep{Francavilla2018a}. Multi-compartment computational models of IS3 cells predict that balanced excitation and inhibition with low numbers of correlated synapses produce asynchronous in-vivo-like states with strongly increased theta-band spiking under 8-Hz inputs, mechanistically grounding the in-vivo theta locking \citep{Guetmccreight2019}. Entorhinal alvear inputs to CA1 monosynaptically excite both PV and VIP interneurons (\textgreater 70\% responsive) far more reliably than OLM cells (\textless 20\%), but PV cells fire at higher theta-burst rates than VIP cells, suggesting cell-type-specific entorhinal-driven theta engagement of inhibitory subtypes \citep{Bell2021}.

Plasticity at IS3 connections couples theta activity to its own disinhibitory output. \citet{Caulinorocha2022} showed in young-adult Wistar rat CA1 that endogenous VIP acting through VPAC1 receptors (but not VPAC2) inhibits theta-burst-stimulation--induced LTP via GABAergic disinhibition, regulating CaMKII autophosphorylation and Kv4.2 channel phosphorylation. \citet{Jabonska2026} complement this with non-Hebbian, heterosynaptic LTD at VIP-IS3\rightarroworiens-IN inputs induced by theta-burst stimulation, weakening disinhibition and increasing the oriens-IN excitation/inhibition ratio. Synapse-specific architecture supports both: \citet{Chamberland2010} reported that VIP/CR IS-III cells provide depressing, plasticity-resistant inhibition onto OLM cells, in contrast with strong, theta-LTP-capable septohippocampal terminals --- a two-input arrangement in which septal drive sets the theta rhythm and IS3 cells fine-tune OLM output. Hippocampal long-range nNOS+ inhibitory cells (LINCs), distinct from VIP and SST cells, additionally project to extrahippocampal regions and modulate CA1 oscillations and inter-regional coherence, indicating that VIP-negative long-range projection neurons share part of the temporal-coordination workspace \citep{Christensonwick2019}.

Inter-areal theta in the cortex echoes the hippocampal theme. Optogenetic inhibition of prefrontal VIP cells decreases open-arm avoidance specifically when hippocampal--prefrontal theta coherence is high, linking VIP disinhibition to inter-areal theta coordination during anxiety-related behaviour \citep{Lee2019}. In mouse V1 during an oddball paradigm, ACa--V1 synchrony peaks in the theta/alpha band ({\textasciitilde}10 Hz); 10-Hz optogenetic stimulation of ACa\rightarrowV1 inputs activates V1 VIP and inhibits SST cells, and chemogenetic VIP silencing abolishes both ACa--V1 theta/alpha synchrony and V1 deviance detection \citep{Bastos2023a}. Layer-1 VIP and $\alpha$7 cell models exhibit intrinsic resonance with theta/alpha-band input near resting potential, mediated by T-type Ca\textsuperscript{2}\textsuperscript{.}-channel dynamics --- providing an intrinsic membrane mechanism for cortical VIP engagement at theta/alpha frequencies \citep{Meng2023}.

\section{Cortical states: Up-states, REM, and infraslow rhythms}

Cortical slow-oscillation Up-states recruit VIP cells unevenly across preparations and species \citep{Tahvildari2012, Kawaguchi2001, Brecier2022, Sabri2024}. \citet{Tahvildari2012} recorded from mouse entorhinal-cortex slices during spontaneous Up states and reported that pyramidal cells and fast-spiking PV interneurons fire robustly while VIP, NPY, and 5HT3a interneurons remain silent --- an in-vitro demonstration that VIP cells are not the primary drivers of cortical slow-oscillation Up-state generation in this preparation. The picture changes once the synchronisation regime shifts: in rat frontal-cortex slices undergoing low-Mg\textsuperscript{2}\textsuperscript{.} NMDA-dependent fast run-like potentials (4--10 Hz), FS cells fire up to 150 Hz and pyramidal cells 25--55 Hz, while VIP and SST interneurons exhibit firing patterns resembling pyramidal cells at synchronisation peaks, demonstrating that VIP cells \textit{can} be entrained by synchronised cortical activity once NMDA-dependent excitation is strong enough \citep{Kawaguchi2001}. The slice/preparation dependence of VIP recruitment during slow oscillations is one of the section's most reproducible cross-study patterns.

\begin{framed}
\textbf{Conflict --- VIP cells as dispensable versus active gates of cortical Up-states / slow oscillations}\\
\citet{Tahvildari2012} reported VIP, NPY, and 5HT3a interneurons silent during entorhinal-cortex Up states in slices, with PV/Pyr units carrying the synchronised activity --- implying VIP cells are dispensable for slow-oscillation Up-state generation in vitro. By contrast, cortex-wide imaging during awake reinforcement learning shows VIP cells co-activated by reward and punishment across the dorsal cortex on slow-oscillation timescales, with response strength only partly explained by arousal \citep{Szadai2022}. Resolution status: open. The two views are not formally incompatible --- VIP cells may be silent during spontaneously generated slow oscillations but driven by neuromodulator-coupled arousal-gated inputs in awake cortex --- but they predict opposite outcomes of focal VIP perturbation. The originally cited counter-pair (Pi 2013 awake-cortex disinhibition versus Neske 2016 slice Up-state dispensability) lies outside this section's filtered citation map; both keys are flagged in the section's unmet citation list and the conflict here uses \citet{Tahvildari2012} and \citet{Szadai2022} as the available proxies for the same disagreement.
\end{framed}

Sleep-state recruitment of VIP cells is highly state-specific and distinct from PV/SST \citep{Brecier2022, Rolle2025, Munoz2017}. \citet{Brecier2022} recorded mouse barrel-cortex layer-2/3 across the natural sleep--wake cycle and reported that VIP interneurons are most active during REM sleep, PV cells fire highest during both REM and NREM (with rapid decrease at wake onset), and SST activity is stable across the cycle; PV and most VIP cells are modulated by delta and theta LFP oscillations. \citet{Rolle2025} extended the timescale to the cortical infraslow rhythm and showed that layer-2/3 VIP interneurons display a {\textasciitilde}0.02-Hz infraslow oscillation that is inversely phase-coupled to infraslow spindle activity during slow-wave sleep, a coupling absent in PV and SST interneurons; VIP activity was acutely upregulated during spindles and slow oscillations but lowest during SWS overall, suggesting that VIP cells uniquely convey the cortical infraslow oscillation. \citet{Sabri2024} provided the bidirectional-perturbation complement: optogenetic inhibition of cortical VIP interneurons reduced correlated activity between behavioural state and individual-neuron spiking, and during quiet states VIP silencing decreased synchronous spiking but \textit{increased} delta power and phase-locking of spikes to the delta band --- reframing VIP cells as actively suppressing low-frequency synchrony during quiet wake.

Anaesthetic-induced burst suppression provides a counterpoint that bounds the VIP role \citep{Yin2025, Veit2023}. \citet{Yin2025} chemogenetically activated and silenced PV, SST, and VIP cells in mouse auditory cortex and mPFC under isoflurane and found that \textit{only} PV manipulation bidirectionally altered pyramidal-cell synchrony (P \textless  0.0001), while SST or VIP manipulation had no effect --- i.e., VIP cells do not gate burst-suppression synchrony, which appears PV-dominated. Together with the slice-Up-state data \citep{Tahvildari2012}, these results demarcate slow-oscillation regimes in which VIP cells are not the synchrony-setting class.

\section{Predictive coding, attention, and inter-areal coordination}

Hierarchical models of cortical inhibition treat VIP cells as the channel by which top-down predictions modulate the local oscillation. In a hierarchical predictive-coding cortical model with PV/SST/VIP cell-type-specific connectivity, SST and VIP cells push oscillation amplitude and number of cycles between representation and error neurons in opposite directions, and in-silico VIP silencing diminishes the model's oscillatory patterns \citep{Lee2025}. \citet{Bastos2023a} provided the in-vivo correlate: chemogenetic VIP silencing in V1 abolishes ACa--V1 theta/alpha synchrony and V1 deviance detection, identifying VIP-mediated SST suppression as the substrate by which a top-down 10-Hz input reorganises the local rhythm. Hippocampal--prefrontal theta coordination during avoidance behaviour likewise scales with VIP recruitment \citep{Lee2019}. Beyond cortex, biophysical modelling of basolateral amygdala finds that PV, SST, and VIP cells each contribute essential rhythm-generating roles (low/high theta and gamma), with low theta ({\textasciitilde}3--6 Hz) emerging as a biomarker of successful fear conditioning via STDP-shaped circuits \citep{Cattani2024} --- the only system in which VIP cells are \textit{required} for, rather than gating, a specific oscillation in our evidence base.

\section{Neuromodulator drive and state-gating of VIP recruitment}

VIP cells are a fast neuromodulator relay: they are uniquely positioned to receive depolarising serotonergic and cholinergic inputs and convert them to disinhibition on millisecond--second timescales. \citet{Lee2010} showed that the 5-HT3A receptor is expressed by virtually all CGE-derived neocortical interneurons that lack PV/SST (including VIP cells), and that these cells are also depolarised by acetylcholine via nicotinic receptors, providing a fast dual serotonergic/cholinergic gateway. \citet{Ferezou2002} localised this to VIP/CCK-coexpressing neurons and demonstrated 5-HT3-mediated fast serotonergic synaptic excitation. Selective non-$\alpha$7 nicotinic excitation of VIP/CCK cells (with $\alpha$4/$\alpha$5/$\beta$2 subunits) was reported by \citet{Porter1999}, and \citet{Kawaguchi1997b} showed that VIP/CR irregular-spiking and bipolar cells receive net depolarising cholinergic responses, while basket-type interneurons are differentially modulated. Noradrenaline differentially modulates GABAergic interneurons via $\alpha$-adrenergic excitation of VIP/CR-class cells and $\beta$-adrenergic effects on FS basket cells. Earlier reports \citep{Demars2014} further showed that PV and SST (but not VIP) cells selectively express endogenous nicotinic-receptor modulators (lynx2, Lynx1, Slurp, SST-related peptides) that bias the cell-type-specific cholinergic recruitment, contributing to VIP cells' preferential nicotinic responsiveness. Slow disynaptic inhibition mediated by non-$\alpha$7 nicotinic excitation of CGE-derived interneurons reshapes spatiotemporal cortical activity patterns \citep{Arroyo2012}. Collectively, these data make VIP cells the disinhibitory fast-acting cortical relay for ACh, NA, and 5-HT --- and therefore the cell class most directly coupled to behavioural state.

State-coupled cholinergic drive recruits VIP cells across cortex with characteristic temporal structure. Initial studies \citep{Hangya2015} showed that optogenetically identified basal-forebrain cholinergic neurons broadcast a fast ({\textasciitilde}18 ms), precisely timed reinforcement signal that scales with the unexpectedness of reward and punishment, providing the neuromodulatory drive that recruits VIP interneurons cortex-wide. \citet{Letzkus2011} showed that auditory-cortex L1 interneurons receive direct nicotinic excitation from basal-forebrain cholinergic axons during foot-shock and disinhibit pyramidal cells, enabling associative fear learning. \citet{Askew2019} demonstrated that bath nicotine strongly depolarises VIP interneurons in auditory cortex and that chemogenetic VIP silencing prevents nicotine-driven disinhibition of pyramidal neurons. Cortex-wide imaging during reinforcement learning extends the geography: VIP cells are co-activated by reward and punishment across virtually the entire mouse dorsal cortex, with response strength only partly explained by arousal --- a global cell-type-specific neuromodulator-broadcast mode \citep{Szadai2022}. \citet{Collins2023} showed that cholinergic basal-forebrain and noradrenergic locus-coeruleus axons across the dorsal cortex carry a globally coordinated arousal/movement signal, so VIP cells receive a brain-state signal coordinated up to several millimetres apart. Pupil fluctuations track fast switches of cortical state during quiet wakefulness, including cell-type-specific VIP recruitment \citep{Reimer2014}, and \citet{Munoz2017} showed that during active wakefulness L1 (including VIP-class) and SST cells are differentially recruited by cholinergic input to dynamically shape dendritic inhibition.

The picture, however, is not one of monolithic VIP-mediated disinhibition. \citet{Kuchibhotla2017} recorded auditory cortex during an active recognition task and showed that cholinergic axon activation simultaneously depolarises PV, SST, \textit{and} VIP cells, with modelling indicating that coincident cholinergic drive of all three subtypes --- not VIP-mediated disinhibition alone --- accounts for context-dependent behavioural output. \citet{Pakan2016} reported that locomotion in V1 increases activity of VIP, SST, and PV cells during visual stimulation with context-dependent (visual-vs-dark) responsiveness, and \citet{Dipoppa2018} showed that a recurrent V1 network model captures locomotion-driven gain only when locomotion increases feed-forward weights and modulates recurrent VIP/SST/PV/Pyr weights --- pure VIP\rightarrowSST disinhibition fails when visual stimuli are present. Frontal-cortex recordings reframe the geometry as a ``pull-push'' motif in which arousal-driven VIP cells inhibit pyramidal cells directly while disinhibiting them via a parallel pathway \citep{Garciajuncoclemente2017}. \citet{Yaeger2019} reported that during the binocular critical period, basal-forebrain ACh released during arousal directly excites SST cells (not VIP) to produce localised dendritic spiking and disinhibition required for ocular-dominance plasticity --- and that this cholinergic SST sensitivity is lost in adulthood, when VIP-mediated disinhibition predominates. \citet{Ren2022} likewise found early activation of VIP-INs followed by late SOM-IN activation during motor learning, with VIP inhibition increasing SOM activity and impairing initial learning. The recurring theme is that VIP cells are \textit{one} of several state-coupled inhibitory channels whose relative weighting shifts with task, layer, age, and modulatory context.

\begin{framed}
\textbf{Conflict --- VIP cells dispensable versus globally recruited by cholinergic state changes}\\
\citet{Chen2015a} reported that intracortical cholinergic input recruits SST (not VIP) interneurons to drive cortical desynchronisation and decorrelation, with selective optogenetic SST inhibition (but not VIP inhibition) blocking desynchronisation; this identifies SST --- not VIP --- as the necessary cholinergic node for cortical state change. \citet{Szadai2022} instead found cortex-wide VIP co-activation by reward and punishment during early auditory-discrimination learning, with response strength only partly explained by arousal --- a global cell-type-specific neuromodulator broadcast that places VIP cells at the centre of cholinergic-driven state changes. Resolution status: open. The disagreement may reflect different timescales (millisecond ACh-driven desynchronisation in \citet{Chen2015a} versus seconds-scale reinforcement broadcast in \citet{Szadai2022}), different stimulus regimes (passive listening versus reinforcement-task learning), and different recording read-outs (LFP desynchronisation versus single-cell calcium responses). The available evidence is consistent with VIP cells implementing a \textit{parallel} rather than the \textit{primary} cholinergic state channel \citep{Kuchibhotla2017, Pakan2016, Dipoppa2018, Yaeger2019, Ren2022}, but quantitative cross-study comparison of VIP versus SST necessity for state changes is not yet available.
\end{framed}

A specialised VIP/ChAT subpopulation transiently couples ACh release to local circuits. \citet{Obermayer2019} showed that a {\textasciitilde}0.5--1\% subpopulation of cortical VIP interneurons co-expresses ChAT and locally releases ACh that directly excites neighbouring pyramidal and inhibitory neurons via fast cholinergic synaptic transmission, controlling attention behaviour. \citet{Granger2020} demonstrated that cortical ChAT\textsuperscript{.} neurons are nearly all VIP\textsuperscript{.} and release GABA broadly onto inhibitory cells while sparsely co-releasing ACh onto layer-1 interneurons and other VIP/ChAT cells. \citet{Saunders2015} further showed that mouse basal-forebrain cholinergic neurons co-release GABA with ACh onto cortical L1 interneurons, with the GABAergic component lost on conditional vesicular-GABA-transporter deletion. \citet{Dudai2020} reported that despite their rarity, VIP/ChAT cells in barrel cortex fire faithfully to whisker stimuli and optogenetic activation suppresses sensory responses of L2/3 excitatory neurons. Dopaminergic and other modulatory layers add further specificity: D1 dopamine receptors are enriched in superficial-layer VIP interneurons that co-express calretinin and D1 agonists strongly enhance firing of these VIP cells while sparing PV and SOM cells \citep{Anastasiades2019}; cortical D2 receptors are also expressed in VIP-class cells \citep{Khlghatyan2019}; mu and delta opioid receptors differentially regulate GABA release from PV, SST, CCK, and VIP cells in prelimbic PFC \citep{Cole2025}; and goal-oriented spatial learning requires VIP-disinhibition in CA1 with VIP silencing impairing both learning and place-cell representations \citep{Turi2019}. Mediodorsal thalamus selectively engages VIP+ cells in PFC layer 1b, providing a thalamic route into the cell class \citep{Anastasiades2021a}, and primary motor cortex preferentially engages VIP cells in S1 via long-range projections \citep{Naskar2021}. Hippocampal disinhibition via VIP/CR cells is also state-dependent: long-range VIP-LRP cells decrease activity during theta-run epochs and are uncoupled to ripples \citep{Francavilla2018a}. Layer-1 nicotinic recruitment is conserved between rodent and human cortex \citep{Poorthuis2018}, and the broader pharmacological framework places ACh as a cell-type-specific neuromodulator across muscarinic and nicotinic axes \citep{Picciotto2012}. Auditory critical-period plasticity is gated by L1 nicotinic input via \textit{Lynx1}-controlled developmental closure \citep{Takesian2018}, and \citet{Alitto2013} formalised cell-type-specific basal-forebrain modulation as the disinhibitory framework for cholinergic disinhibition. The neuromodulator-control mode of VIP cells extends to vasomotor coupling: single VIP cortical interneurons are sufficient to dilate neighbouring microvessels via VIP-peptide release and receive direct cholinergic and serotonergic afferents \citep{Cauli2004}, and distinct ACh/NE pathways recruit specific VIP/NOS/CCK subtypes mediating region- and state-dependent neurovascular coupling \citep{Lecrux2016}. VIP and PACAP signal through three class-B GPCRs (PAC1, VPAC1, VPAC2) regulating circadian rhythms, learning, and stress \citep{Harmar2012}.

\begin{framed}
\textbf{Conflict --- opposite locomotion polarity of VIP-driven gain across V1 and A1}\\
\citet{Fu2014} reported that locomotion activates VIP-expressing interneurons in mouse V1 largely through nicotinic cholinergic inputs from the basal forebrain, that VIP cells gate the locomotion-induced gain enhancement of visual responses, that optogenetic VIP activation in stationary mice mimicked locomotion gain, and that photolytic VIP ablation abolished the locomotion-induced enhancement of V1 responses --- placing locomotion-driven VIP recruitment as a \textit{positive} gain modulator in V1. By contrast, a parallel auditory-cortex study (Bigelow and colleagues, 2019; not in this section's citation map and flagged in the unmet list) reported reduced spiking and reduced encoding efficiency under locomotion in primary auditory cortex (A1), indicating opposite locomotion polarity of VIP-driven gain across sensory cortices. Resolution status: open. Within this section's filtered citation map the A1-side reference (Bigelow2019) is unavailable and is flagged in the unmet citation list; the conflict is therefore stated using \citet{Fu2014} for V1 and the broader cross-area heterogeneity captured in \href{/cross-areas}{VIP Interneurons Across Brain Regions}. The dissociation is consistent with the parallel-recruitment data of \citet{Pakan2016}, \citet{Dipoppa2018}, and \citet{Kuchibhotla2017}, which show that the state-driven coupling of VIP cells is not uniformly disinhibitory across cortex.
\end{framed}

\section{Pathological oscillations and disease links}

VIP-mediated disinhibition can shift cortical rhythms into pathological regimes. \citet{Hall2015} showed in rat neocortical slices that loss of NPY-mediated inhibition combined with VIP-mediated disinhibition transforms a sleep-associated delta rhythm into 0.5--4 Hz spike-and-wave discharges, and that VIP-receptor blockade nearly abolishes this epileptiform activity --- a slice-level demonstration that pathological synchrony depends bidirectionally on VIP peptidergic signalling. \citet{Goff2019} reported that, in a Dravet-syndrome (\textit{Scn1a}\textsuperscript{.}/\textsuperscript{-}) mouse model, the irregular-spiking subtype of VIP interneurons shows impaired excitability that is rescued by the Nav1.1-specific toxin Hm1a; the IS-versus-CA firing pattern of VIP cells is set by KCNQ M-current and switches to tonic firing on muscarinic activation, mechanistically linking developmental VIP dysfunction to the seizure phenotype. Foundational studies \citep{Goel2025} summarised Fragile-X dysfunction as elevated broadband gamma EEG power and abnormal phase-locking to gamma-modulated acoustic stimuli in \textit{Fmr1}-KO mice, paralleling human FXS phenotypes, with dysfunctional VIP cells correlating with elevated susceptibility to irrelevant stimuli. \citet{Kranz2025} showed in resting-state EEG (n = 38 Rett, 30 controls) increased theta--gamma and alpha--gamma phase--amplitude coupling and altered alpha--gamma phase bias in Rett syndrome (P \textless  0.05), and a biophysically constrained layer-4 cortical model reproduced these PAC changes solely by reducing VIP+ interneuron activity (model P \textless  0.001) --- directly implicating VIP-cell hypofunction in the abnormal cross-frequency coupling. \citet{Koukouli2017} showed that mice carrying the human $\alpha$5-nAChR rs16969968 risk SNP or $\alpha$5-knockout show reduced VIP-interneuron activity and disinhibition of SST \rightarrow pyramidal in PFC, producing hypofrontality and behavioural deficits that chronic nicotine reverses. \citet{Kiss2026} reported in cross-cohort RNA-seq (n = 1408 incl. 672 SCZ cases) across three neocortical regions age-dependent reductions of PV and SST proportions in younger SCZ cases (\textless 70 y) and reduced per-cell VIP and SST mRNA, providing a molecular substrate for the schizophrenia-associated alteration of cortical interneuron-driven oscillations. Established work \citep{Murdock2024} linked therapeutic gamma to VIP function: chemogenetic inhibition of cortical VIP interneurons in 5XFAD mice attenuated 40-Hz multisensory gamma stimulation--induced amyloid clearance without affecting baseline 40-Hz LFP power, indicating that VIP cells couple gamma activity to glymphatic peptide release. The translational implication is that VIP-targeted interventions are most likely to alter rhythm \textit{amplitude, coherence, and phase coupling} --- not rhythm presence --- and that the same disinhibitory motif underwrites both adaptive (gain modulation, plasticity, clearance) and maladaptive (spike-and-wave, hypersynchrony, hypofrontality) regimes.

\begin{framed}
\textbf{Conflict --- bidirectional VIP perturbation in seizure models}\\
A contrasting line of evidence (Khoshkhoo and colleagues, 2017 versus Calin and colleagues, 2018; both flagged in the unmet citation list) shows that silencing VIP suppresses seizures while activating VIP fails to alter discharge duration --- implying asymmetric, model-dependent VIP contribution to seizure dynamics. Both keys are absent from this section's filtered citation map and are flagged in the unmet citation list. Within the available evidence the same asymmetry can be partly reconstructed: VIP-receptor blockade in rat neocortical slices nearly abolishes spike-and-wave discharges \citep{Hall2015}, indicating that VIP peptidergic signalling is \textit{necessary} for one form of epileptiform activity, while broader perturbation studies in non-seizure contexts find that VIP activation reshapes oscillation power and coherence without abolishing the underlying rhythm \citep{Veit2023, Bastos2023a, Sabri2024}. Resolution status: open. The convergent reading is that VIP-mediated disinhibition is more likely to be required for \textit{escalating} an existing rhythm into a pathological regime than for initiating de-novo seizure activity, but cross-model quantitative comparison is not yet available.
\end{framed}

\section{Synthesis}

Across 80 evidence findings the cortical and hippocampal oscillatory roles of VIP cells converge on a state-gating rather than rhythm-generating function. PV--pyramidal interactions set gamma frequency, OLM and septal inputs set hippocampal theta, and PV cells dominate burst-suppression synchrony --- VIP cells modulate the \textit{power, coherence, and phase coupling} of these rhythms via SST/PV control \citep{Veit2023, Wagatsuma2023, Tyan2014, Bell2021, Yin2025}. Their unique fast-modulator coupling (5-HT3, non-$\alpha$7 nAChR, $\alpha$-adrenergic, D1, MOR/DOR, VPAC1/2) makes them the cortical cell class most directly coupled to behavioural state \citep{Lee2010, Ferezou2002, Porter1999, Kawaguchi1997b, Kawaguchi1998, Anastasiades2019, Cole2025, Harmar2012}. State-dependent recruitment is, however, parallel rather than serial: PV, SST, and VIP are co-activated by ACh and locomotion in many regimes \citep{Kuchibhotla2017, Pakan2016, Dipoppa2018, Yaeger2019, Ren2022}, and the simple VIP\rightarrowSST\rightarrowPyr disinhibition motif is a partial description of the underlying circuitry. Pathological oscillations --- spike-and-wave seizures, Fragile-X gamma hypersynchrony, Rett PAC, Dravet IS-cell hypofunction, schizophrenia-associated VIP-mRNA reduction, and Alzheimer-disease gamma--clearance coupling --- share a common mechanism in which VIP-cell perturbation alters the regulatory bandwidth of an otherwise PV/SST-driven rhythm \citep{Hall2015, Goel2025, Kranz2025, Goff2019, Kiss2026, Murdock2024, Koukouli2017}. Translational interventions that target VIP-mediated disinhibition will therefore most likely act by \textit{reshaping} --- not generating or abolishing --- cortical rhythms. \href{/disease-translational}{Species Differences, Human Relevance, and Disease} takes up the species-translation question directly, asking how much of this rodent oscillation--disease coupling carries to human cortex and disease.

\begin{figure}[!htbp]
\centering
\includegraphics[width=0.9\linewidth]{files/fig-vip-oscillations-959b8dcce125b803d9d433e92700e09b.png}
\caption[]{\textbf{VIP-cell perturbation effects on cortical and hippocampal rhythms --- qualitative cross-study summary.} Phase-6 audit found no quantitative panels in cluster\_09\_oscillations (n\_audited\_panels = 0); this figure is therefore a qualitative schematic synthesising the directional perturbation effects reported in the section's evidence base, \textit{not} a meta-analytic forest plot. \textbf{(A)} Schematic of cortical LFP under VIP optogenetic activation versus silencing: VIP activation reshapes gamma power and coherence without abolishing the rhythm \citep{Veit2023, Wagatsuma2023}; VIP silencing increases delta power during quiet wake \citep{Sabri2024}. \textbf{(B)} Direction of effect of VIP perturbation on local gamma, theta, and inter-areal coherence across studies (qualitative ↑/↓/- indicators only; effect sizes not pooled). \textbf{(C)} Cross-study seizure-susceptibility table: VIP-receptor blockade abolishes spike-and-wave in rat slices \citep{Hall2015}; bidirectional in-vivo VIP perturbation in epilepsy models (Khoshkhoo 2017, Calin 2018) is flagged in the unmet citation list. \textbf{(D)} State-dependent VIP recruitment across Up-states, REM, and infraslow rhythms \citep{Tahvildari2012, Brecier2022, Rolle2025}. \textbf{Caveat:} No Phase-6-audited quantitative panels were available for this figure; values shown are categorical/qualitative summaries from primary text and abstracts --- not pooled effect sizes. See unmet citation needs file for sources flagged for follow-up curation.}
\label{fig-vip-oscillations}
\end{figure}

\begin{figure}[!htbp]
\centering
\includegraphics[width=0.9\linewidth]{files/fig-vip-oscillation--db163353e744d676a08ae1533838c0d3.png}
\caption[]{\textbf{Schematic mechanism: VIP\rightarrowSST and VIP\rightarrowPV motifs differentially shape gamma versus theta rhythms.} Schematic with no quantitative data substrate. \textbf{(A)} PV-driven gamma rhythm with VIP-mediated SST suppression: VIP activation predicted to scale local gamma power \citep{Veit2023, Wagatsuma2023} and tune long-range gamma coherence \citep{Veit2023}. \textbf{(B)} Hippocampal theta-rhythm circuit with VIP/IS3 cell phase-coupled to OLM disinhibition; CR/IS3 cells fire on the rising phase and peak of theta and selectively innervate OLM dendrites \citep{Tyan2014, Luo2020, Chamberland2010}. \textbf{(C)} Cortical Up-state schematic with VIP cells as a permissive (slice-silent, awake-recruited) versus causal gate, capturing the \citet{Tahvildari2012} versus \citet{Szadai2022} disagreement about Up-state generation. \textbf{Caveat:} Schematic only --- no audited quantitative substrate.}
\label{fig-vip-oscillation-mechanism}
\end{figure}

Whether the rodent oscillatory and circuit-level findings reviewed here translate to primate cortex, and how VIP-IN dysfunction maps onto neurodevelopmental and psychiatric disease, is the focus of \href{/disease-translational}{Species Differences, Human Relevance, and Disease}.


\section{Species Differences, Human Relevance, and Disease}
\label{sec:11_disease_translational}

\section{Thesis}

\href{/synaptic-properties}{Synaptic Properties and Connectivity} through \href{/oscillatory-dynamics}{Oscillatory Dynamics and Temporal Coordination} documented VIP+ interneuron function almost entirely in rodent
cortex and hippocampus. Whether that body of evidence transfers to human
cortex, and whether mouse disease models faithfully recapitulate human
pathology, is the question of this section. The claim we defend is bipartite.
At the level of the \textit{subclass}, VIP-IN identity is conserved between rodent
and primate cortex: cross-species single-cell atlases recover \textit{PVALB}, \textit{SST},
\textit{VIP}, and \textit{LAMP5/PAX6} as the four conserved GABAergic subclasses in mouse,
marmoset, macaque, and human \citep{Tasic2018, Hodge2019, Bakken2021a, Braininitiativecellcensusnetworkbiccn2021, Lee2023a, Yao2023},
and the same gene families that discriminate mouse interneurons also
discriminate human interneurons \citep{Hodge2019}. At the level of \textit{t-types,
proportions, and morphology}, however, divergence is non-trivial: VIP is the
most diverse subclass in human middle temporal gyrus (MTG), with 21 types and
an upper-layer bias absent in the mouse \citep{Hodge2019}; human L1 contains
\textit{VIP PCDH20} and \textit{MC4R} t-types whose morphologies have no mouse counterpart
\citep{Chartrand2023, Boldog2018}; and primate cortex shows expanded
CGE-derived diversity that the mouse atlases under-represent
\citep{Bakken2021a, Bakken2021b, Krienen2020}. Layered onto this divergence is
disease: VIP-IN dysfunction is implicated across autism-spectrum and
intellectual-disability syndromes (\textit{MECP2}, \textit{TCF4}, \textit{FMR1}, \textit{CNTNAP2},
\textit{SCN1A}) \citep{Mossner2020, Chen2023b, Bhandari2024, Goff2023, Hanno2026, Gil2024},
schizophrenia and affective disorders \citep{Batistabrito2017, Miller2025, Prevot2025},
Alzheimer's disease and Huntington's disease
\citep{Michaud2024}, and Dravet syndrome
\citep{Goff2023, Vormsteinschneider2020, Gil2024}. Mouse-model phenotypes are,
however, model-, region-, and assay-specific, and translation to human
pathology is non-trivial --- a heterogeneity we make quantitative in
Figure~\ref{fig-species-disease} and theoretical in
Figure~\ref{fig-vip-disease-convergence}.

\section{Cross-species conservation at the subclass level}

The mouse cortical interneuron taxonomy that organizes the rest of this
review --- six GABAergic subclasses (\textit{Sst}, \textit{Pvalb}, \textit{Vip}, \textit{Lamp5}, \textit{Sncg},
\textit{Serpinf1}) split between MGE- and CGE-derived branches
\citep{Tasic2018, Yao2021a, Tremblay2016, Rudy2011} --- was substantially
recapitulated when transcriptomic atlases were extended to primate and human
cortex. \citet{Hodge2019} profiled human MTG by SMART-seq v4 snRNA-seq and
recovered the same four conserved GABAergic subclasses as in mouse, with
one-to-one correspondence at the subclass level and 21 VIP t-types within the
\textit{VIP} subclass; critically, ``the same gene families that discriminate mouse
interneurons also discriminate human interneurons''
\citep{Hodge2019}. The BRAIN Initiative Cell Census Network MOp/M1
consortium extended this to a tripartite mouse--marmoset--human alignment and
recovered 45 conserved t-types including 24 GABAergic types
\citep{Braininitiativecellcensusnetworkbiccn2021, Bakken2021a}. Patch-seq
characterization of human neocortex independently recovered 45 GABAergic
t-types across the four conserved subclasses (\textit{PVALB}, \textit{SST}, \textit{VIP},
\textit{LAMP5/PAX6}) \citep{Lee2023a}, and human and macaque snDrop-seq /
snRNA-seq atlases each recover \textit{PVALB}, \textit{SST}, \textit{LAMP5}, and \textit{VIP} with
significant differentially-expressed genes for each subclass
\citep{Lake2018, Chen2023a}. Comparable transcriptomic homology has
also been reported in the dorsal lateral geniculate nucleus and amygdalar
GABAergic neurons across mouse, macaque, and human, indicating that the four-subclass scheme is
not a cortex- or species-specific artifact. The \textit{VIP}--\textit{ChAT} cholinergic
co-expressing type identified in mouse V1 by \citet{Tasic2016} and
formalized by \citet{Tasic2018} has plausible homology in human upper
layers \citep{Hodge2019, Lake2018, Lee2023a}, although direct functional
comparison remains incomplete.

\begin{framed}
\textbf{Conflict --- homology vs species-specific definitions}\\
\citet{Hodge2019} and \citet{Braininitiativecellcensusnetworkbiccn2021}
treat human and mouse VIP/CGE taxonomies as broadly transferable: subclass
boundaries align, and the discriminating gene families are shared. By
contrast, an earlier report characterized embryonic mouse MGE/CGE/LGE
populations by Drop-seq and emphasized non-overlapping cardinal types whose
diversification is governed by lineage- and region-specific transcription
factors that themselves vary across species, leading to the view that primate
\textit{VIP} populations require species-specific subtype definitions. The two
positions are not strictly contradictory --- subclass-level homology with
t-type-level divergence --- but they support different downstream practices: a
direct transfer of mouse t-type labels into human atlases versus
species-specific re-clustering. The contemporary atlas literature has
gravitated toward the second practice for VIP specifically
\citep{Hodge2019, Bakken2021a, Lee2023a, Chartrand2023}.
\end{framed}

\section{Primate expansion of CGE-derived diversity}

Even granting subclass-level conservation, cell-type \textit{proportions} and
within-VIP \textit{granularity} diverge across species in directions that bear on
human relevance. \citet{Bakken2021a} reported that ``more consensus
clusters could be resolved by pairwise alignment between humans and
marmosets than between either of these primates and mice, particularly for
\textit{Vip} subtypes,'' indicating disproportionate primate expansion within the
\textit{VIP} and CGE-derived branches. \citet{Hodge2019} made this quantitative:
in human MTG, \textit{VIP} is the single most diverse interneuron subclass, with
21 t-types, the majority enriched in upper layers --- a layer bias absent at
this granularity in mouse VISp/ALM atlases of comparable resolution
\citep{Tasic2018, Yao2023}. Cross-species transcriptomic profiling that
includes mouse, macaque, and human cortex confirms \textit{VIP} (and \textit{LAMP5}) as a
major GABAergic subclass with significant DEGs in each species; primate expansion
is also visible in non-cortical CGE-derived populations such as the dLGN
\citep{Bakken2021b} and the amygdala. A complementary
methodological signal comes from electrophysiology: \citet{Torresgomez2020}
analysed primate prefrontal recordings under the assumption that primate
CR/CB/PV cells perform similar computations to mouse VIP/SST/PV cells, and
inferred that ``changes in the proportion of CR and PV neurons in layers
2/3'' account for emergent population dynamics, an explicit
proportion-driven argument for VIP/CR-class expansion in primate cortex.
The CGE lineage origin of this expansion is not in dispute:
\citet{Fishell2011} and \citet{Lodato2011} showed that the CGE
generates the majority of bipolar calretinin/VIP-positive cortical
interneurons, and \citet{Tremblay2016} and \citet{Rudy2011} placed
\textit{VIP} within the broader 5-HT3AR+ CGE-derived population that comprises
{\textasciitilde}30\% of cortical GABA cells in rodent sensory cortex
. The open question is whether the primate
expansion preserves the rodent \textit{VIP} \rightarrow \textit{SST} \rightarrow pyramidal disinhibitory
motif or whether new subtypes implement new connectivity rules, an
ambiguity we return to in the concluding synthesis.

\section{Human-specific transcriptomic and morphological features}

Beyond proportional expansion, the human cortex contains GABAergic types
with no clear mouse homologue, the most-cited example being the
``rosehip'' cell of human L1 first reported by \citet{Boldog2018}: many
rosehip marker genes are not expressed in \textit{Pvalb/Sst/Vip}-negative mouse
GABAergic types, and only five rosehip-adjacent clusters mapped onto mouse
\textit{Vip+} clusters \citep{Boldog2018}. Cross-species Patch-seq alignment of
human and mouse L1 by \citet{Chartrand2023} extended the same conclusion
to \textit{VIP}: ``no mouse L1 neurons had morphologies resembling the unmatched
human L1 t-types (\textit{VIP PCDH20} and \textit{SST BAGE12}, homologous to deeper
mouse t-types), and the \textit{MC4R} subclass was less morphologically distinct,
providing evidence for type-specific divergence'' --- a direct
demonstration of human-specific morphology within a transcriptomically
homologous subclass. Independent neurosurgical IHC/EM/patch-clamp
characterization of human cerebral cortex by \citet{Somogyi2025}
identified a \textit{VGLUT3}-expressing GABAergic terminal population on
dendritic shafts, with 25\% of CB1-immunopositive GABAergic terminals
co-expressing \textit{VGLUT3}, suggesting GABAergic/glutamatergic
co-transmission patterns that have only partial mouse counterparts.
Phylogenetic transcriptomic comparison by a preceding study framed the
issue more generally: despite homology at the level of cell types, ``clear
differences between transcriptomic PCs'' emerge with greater
dissimilarities between evolutionarily distant species, arising from
lineage-specific transcription factor recruitment. Together these findings
imply a ``homologous subclass, divergent type'' picture: rodent
mechanistic work on \textit{VIP}-IN identity is necessary but not sufficient for
human relevance, and human-specific subtypes (\textit{VIP PCDH20}, \textit{MC4R},
rosehip) anchored to the Hodge/Yao atlases must be incorporated into any
translation effort \citep{Hodge2019, Yao2023, Lee2023a, Chartrand2023, Boldog2018}.

\section{Quantitative anchors: VIP density across species}

Cell-density estimates illustrate why the species-level numbers do not
straightforwardly compare. \citet{Ouellet2014}, using IHC counts in
postnatal day 9 (P9) rat primary auditory cortex, reported that ``CR, VIP,
and ChAT appeared to be the most prominent cortical interneuron markers;
representing 10, 8, and 1\% of total GABA positive cells,'' locating VIP+
at {\textasciitilde}8\% of the GABAergic population in early postnatal rodent A1
\citep{Ouellet2014}. \citet{Teymornejad2024a}, applying RNAscope and IHC
to human postmortem M1, reported that ``VIP+ interneurons corresponded to
4.5, 5.2, and 6.6\% of the \textit{total neuronal} population in cytoarchitectural
subdivisions A4a, A4b, and A4c,'' locating VIP+ at {\textasciitilde}4.5--6.6\% of \textit{all}
neurons in human motor cortex \citep{Teymornejad2024a}. The two
denominators (\% of GABA vs \% of all neurons), the regions (A1 vs M1), the
ages (P9 vs adult postmortem), and the methods (developmental IHC vs
RNAscope/IHC) are not interchangeable, and the panel-level audit of
Figure~\ref{fig-species-disease} flags this explicitly: the apparent
quantitative similarity ({\textasciitilde}5--10\%) masks methodological non-comparability.
A useful normalization comes from the rodent literature itself:
\citet{Rudy2011} placed VIP+ cells at {\textasciitilde}30--40\%
of the 5-HT3AR+ subpopulation, which itself is {\textasciitilde}30\% of cortical
GABA --- implying VIP+ $\approx$ 9--12\% of cortical GABAergic neurons in mouse
sensory cortex, broadly consistent with the rat A1 P9 estimate
\citep{Ouellet2014, Rudy2011, Tremblay2016}. There is
currently no published cross-species count using a single denominator,
single age, and single area, and the field's reliance on
mismatched-denominator comparisons is itself a finding.

\section{Autism-spectrum and intellectual-disability monogenic models}

Five monogenic disorders converge on VIP-IN dysfunction in mouse cortex:
\textit{MECP2} (Rett syndrome), \textit{TCF4} (Pitt-Hopkins syndrome), \textit{FMR1} (Fragile X
syndrome), \textit{CNTNAP2} (cortical-dysplasia-focal-epilepsy / autism-spectrum
disorder), and \textit{SCN1A} (Dravet syndrome and a separable autism phenotype).
\citet{Mossner2020} showed that loss of \textit{MeCP2} selectively from VIP
interneurons ``replicated key physiological and behavioral phenotypes
observed in the pan-interneuron \textit{Dlx5/6} mutants, including altered firing
rates, disruption of high-frequency cortical local field potentials, and
behavioral deficits,'' establishing VIP-IN-specific \textit{MeCP2} loss as
sufficient for several Rett-like cortical phenotypes
\citep{Mossner2020, Ferguson2023b, Simacek2025b, Goff2021}. \citet{Chen2023b}
characterized \textit{Tcf4} +/tr Pitt-Hopkins-syndrome mice and reported that
``VIP-INs in the \textit{Tcf4} +/tr mice show increased input resistance and
reduced frequency of sEPSCs,'' interpreted as reduced excitatory drive
onto VIP-INs partially counterbalanced by elevated input resistance, in
combination with IHC-confirmed reductions in PV+, VIP+, and CST+
interneuron numbers across cortex. An earlier report reported that \textit{Cntnap2}
RNA ``is present at highest levels in chandelier neurons, PV+ neurons and
VIP+ neurons'' in cortical GABAergic interneurons, providing a molecular
substrate for VIP-IN involvement in autism-spectrum \textit{CNTNAP2} phenotypes.
\citet{Bhandari2024} demonstrated in \textit{Fmr1}-/- mice that ``chronic
alteration in late-born parvalbumin interneuron networks across the
vH-PreL axis [is] rescued by VIP signaling,'' directly causally linking
VIP-IN function to learning deficits in Fragile X. \textit{Fmr1}-/- mice also
showed reduced visually-evoked VIP responses in V1: ``only 57.3\% of VIP
cells showed any significant modulation by visual stimuli in \textit{Fmr1} -/-
mice, compared to 73.2\% in WT mice'' \citep{Rahmatullah2023}.
For \textit{SCN1A}, \citet{Goff2023} showed that VIP-INs from \textit{Scn1a} +/- mice
``fired at lower frequencies before application of carbachol (42 $\pm$ 2.9 vs.
34 $\pm$ 2.5 Hz, WT vs.)'' and that selective conditional deletion of \textit{Scn1a} in VIP-INs (paraphrased) reproduces autism-like features without overt seizures, dissociating the autism and
epilepsy phenotypes of Dravet syndrome
\citep{Goff2023, Vormsteinschneider2020}. Cross-cutting
quantification by \citet{Hanno2026} reported that ``VIP-INs strikingly
decreased in both ASD models by 38\% and 49\% relative to controls...
meanwhile, other subtypes, such as Parvalbumin, SST, and RELN-INs, showed
no significant changes,'' consistent with a selective vulnerability of the
VIP class in two distinct ASD models. A preceding study and
\citet{Kaneko2024} emphasized that the extent to which VIP release,
GABA release, or both, are reduced by these developmental disruptions
``are not clear'', a residual mechanistic ambiguity.

\begin{framed}
\textbf{Conflict --- generalized vs region-specific VIP-IN dysfunction}\\
\citet{Goff2023} interpret the \textit{Scn1a} +/- and the conditional VIP-IN
\textit{Scn1a} deletion mouse as a global VIP-IN excitability defect that drives
a generalized ASD-like phenotype, an interpretation aligned with the
pan-cortical reductions in VIP-IN firing reported across \textit{Tcf4} +/tr, \textit{MeCP2} loss-of-function \citep{Mossner2020}, and
\textit{Cntnap2} models. By contrast, \citet{Bhandari2024}
locate the \textit{Fmr1}-/- deficit specifically in the ventral
hippocampus--prelimbic theta-coherence axis and attribute the rescue to
VIP-mediated re-tuning of \textit{late-born PV interneurons}, framing the
VIP-IN contribution as a circuit- and region-specific intervention rather
than a global excitability lesion. The two readings are testable
predictions that diverge across cortical area and developmental stage; the
selective vulnerability quantification by \citet{Hanno2026} favours the
former, while the region-specific Fragile-X data of
\citet{Rahmatullah2023} and \citet{Bhandari2024} favour the latter.
The synthesis we offer is that \textit{both} readings are correct at different
spatial scales and the dichotomy is methodological as much as biological.
\end{framed}

\begin{framed}
\textbf{Conflict --- direction of VIP-IN activity change across ASD models}\\
The two best-controlled in-vivo recordings in autism-related models point
in opposite directions. \citet{Goff2023} reported a reduced VIP-IN
baseline firing rate in cortex of \textit{Scn1a} +/- Dravet/autism-model mice
(34 $\pm$ 2.5 Hz vs WT 42 $\pm$ 2.9 Hz, ex-vivo whole-cell), whereas
\citet{Rahmatullah2023} reported reduced \textit{visually evoked} modulation
in V1 of \textit{Fmr1}-/- Fragile-X mice (57.3\% of VIP cells modulated vs WT
73.2\%, in-vivo 2P calcium imaging). The two etiologies (sodium-channel
haploinsufficiency vs FMRP loss), the two regions (cortex broad vs V1
specifically), and the two readouts (Hz vs \% cells modulated) are not
interchangeable, and the audit pass that retained both entries
Figure~\ref{fig-species-disease} framed them as qualitative
direction-of-change only. Both report decreased VIP-IN function relative
to WT --- the direction is concordant --- but the magnitudes are not
comparable across rows.
\end{framed}

\section{Schizophrenia and affective disorders}

VIP-IN involvement in schizophrenia is supported by genetic, transcriptomic,
and circuit-perturbation evidence. \citet{Miller2025} noted that ``copy
number variations in the gene coding for VPAC2, \textit{Vipr2}, have been
identified to confer a significant risk for schizophrenia,'' providing a
direct \textit{VIP}-pathway genetic link. An initial investigation, using
conditional \textit{ErbB4} deletion restricted to VIP cells, ``directly tested the
role of vasoactive intestinal peptide (VIP)-expressing interneurons in
schizophrenia-related deficits in vivo'' and reported circuit and
behavioral abnormalities consistent with the human disorder. Postmortem
work places VIP-IN dysregulation alongside \textit{PVALB} and \textit{SST} abnormalities
as part of a broad inhibitory-circuit phenotype, with altered expression
of GAD67, \textit{KCNS3}, and CR/VIP markers reported in schizophrenia
postmortem cortex, and an initial investigation reviewed the cross-disorder
implications of ``dysregulation of \textit{PVALB}-, \textit{SST}-, and \textit{VIP}-expressing
interneurons'' across SCZ, bipolar disorder, MDD, ASD, and Alzheimer's.
Affective and stress-related disorders also implicate VIP cells: a
substantial fraction of CRH+ GABAergic interneurons in human
subgenual anterior cingulate cortex have been reported as VIP+
(primary source not in this section's allowlisted reference set and
listed under unmet citation needs below), proposing a candidate molecular
substrate by which the HPA-axis response could be gated through
cortical VIP cells. Chronic
restraint-stress models show ``a significant effect of sex (F (1,82) = 3.8,
\textit{p} = 0.05) and a CRS duration*sex interaction (\textit{F} (5,82) = 2.9, \textit{p} =
0.016)'' on VIP RNA levels \citep{Prevot2025}; sleep-disruption
models in autism- and schizophrenia-relevant lines have similarly
been associated with impaired inhibitory circuits including VIP and
PV neurons (primary source filtered from this section's allowlisted reference set and listed under
unmet citation needs below). Bogaj and
colleagues argued that ``experimental disruption of
VIP-IN function phenocopies behavioural abnormalities'' of schizophrenia,
strengthening the inference from correlative postmortem data.

\section{Alzheimer's disease and neurodegeneration}

The VIP-IN literature in Alzheimer's disease is split between mouse-model
and human-postmortem evidence and is the most disease-translation-fragile
section of the table. On the mouse side, primary 5xFAD-model evidence (primary source not in this section's
allowlisted reference set and listed under unmet citation needs below) has been read as showing that exogenous vasoactive intestinal
peptide attenuates $\beta$-amyloid accumulation and limits brain atrophy in
the 5xFAD mouse model, implying a \textit{protective} role for VIP signalling
that is at least
inconsistent with the strong ``VIP-IN as vulnerability'' reading of the
human literature. \citet{Michaud2024} characterized the 3xTg model and
emphasized that ``the hippocampal CA1 VIP-INs engage in goal-directed and
spatial learning by tuning the activity of PCs via input-specific
disinhibition,'' reporting altered firing/connectivity rather than overt
VIP cell loss. The corresponding Huntington's-disease R6/2 reading
(primary source not in this section's allowlisted reference set and
listed under unmet citation needs below) has been described as placing
VIP-INs upstream of the other observed neuron types, with hypoactivity
in R6/2 disinhibiting SST-INs and thereby increasing inhibition of
CSPNs --- a circuit-level proposal that locates VIP as the upstream
node in the disinhibitory cascade dysregulated in HD; the claim is
advanced as a candidate framework rather than as cited in-section
evidence. The cross-disorder review of an earlier analysis summarized that
``dysregulation of PVALB-, SST-, and VIP-expressing interneurons has been
implicated in a range of neuropsychiatric and neurodegenerative
conditions including bipolar disorder, major depression, autism spectrum
disorder, and Alzheimer's [disease].'' We note that the most-discussed
human-postmortem AD studies (the SEA-AD MTG snRNA-seq characterization of
Vip+ inhibitory subtype loss in late-phase human AD, and recent
neuropathology work reporting depletion of VIP subtypes) appear in the
review plan's anticipated-conflicts list but were outside this section's allowlisted reference set; we flag them as
unmet citation needs and discuss their consequence below.

\begin{framed}
\textbf{Conflict --- VIP-IN preservation vs vulnerability in AD}\\
The 3xTg mouse-model evidence assembled by \citet{Michaud2024} shows
altered VIP-IN firing and connectivity in CA1 \textit{without} overt cell loss.
The complementary human-postmortem account places this against
single-nucleus transcriptomic evidence from the SEA-AD middle-temporal
gyrus cohort \citep{Gabitto2024} and from independent human-cortex
postmortem profiling, both of which have been read as
showing selective vulnerability of VIP-expressing inhibitory subtypes at
late disease stage. Resolution remains unsettled: mouse-model evidence
favours circuit-functional dysregulation without overt cell loss, while
contemporaneous human-postmortem snRNA-seq literature
\citep{Gabitto2024} favours selective VIP-subtype depletion at
late stage. Cross-disease reviews acknowledge both readings without
resolving them \citep{Gil2024}.
\end{framed}

\section{Dravet syndrome and developmental epilepsies}

The traditional Dravet-syndrome (DS) story centres on PV-cell \textit{Scn1a}
haploinsufficiency, but recent work has reassigned a substantial part of
the clinical phenotype to VIP-INs. \citet{Vormsteinschneider2020}
identified Scn1a enhancers and showed that \textit{Scn1a} is expressed in three
non-overlapping neuronal populations including fast-spiking PV+
interneurons and VIP+ interneurons, providing the molecular substrate for
VIP-IN involvement. \citet{Goff2023} then demonstrated by \textit{VIP-Cre}
conditional deletion (paraphrasing) that selective \textit{Scn1a} deletion in
VIP-INs reproduces autism-spectrum-like features without overt seizures, dissociating the two clinical axes of
DS --- refractory seizures (PV-mediated) and the autistic, cognitive, and
behavioral burden (VIP-IN-mediated). \citet{Gil2024} framed the
clinical synthesis: ``VIP interneuron dysfunction is also implicated in
three prominent neurodevelopmental disorders, Rett syndrome, Dravet
Syndrome, and Down's syndrome \dots and may contribute to enhanced seizure
susceptibility.'' The reduced VIP-IN baseline firing reported by
\citet{Goff2023} (34 vs 42 Hz, \textit{Scn1a} +/- vs WT) is the prototypical
quantitative anchor for this view.

\section{Heterogeneity across disease models: a cross-axis synthesis}

Figure~\ref{fig-species-disease}, Panel B, assembles the two best-controlled
quantitative entries on VIP-IN function across autism-related mouse
models --- \textit{Scn1a} +/- (cortex, \textit{ex vivo}, Hz) and \textit{Fmr1}-/- (V1, \textit{in vivo},

% of cells modulated). The audit retained these two entries as a

qualitative direction-of-change comparison after removing entries that
violated metric homogeneity (\textit{p}-only rows, baseline-rate rows, \textit{F}-statistic
rows). Both retained entries show \textit{decreased} VIP-IN function relative to
WT, so the qualitative axis is harmonized; but the underlying
quantitative axes (Hz vs \% of cells modulated), the regions (cortex broad
vs V1), the disease etiologies (\textit{SCN1A} sodium-channel haploinsufficiency
vs \textit{FMR1} loss of FMRP), the developmental stages, and the recording
modalities (\textit{ex vivo} whole-cell vs \textit{in vivo} 2P imaging) are not
interchangeable. \textbf{Magnitudes are not comparable across rows in
Figure~\ref{fig-species-disease}B.} Read together with \citet{Hanno2026}'s
cross-model VIP-IN reduction in cortex (38\% / 49\% loss in two ASD models),
\citet{Mossner2020}'s pan-cortical \textit{MeCP2} phenotype,
\citet{Chen2023b}'s elevated input resistance and reduced sEPSCs in
\textit{Tcf4} +/tr, and \citet{Bhandari2024}'s region-specific Fragile-X
rescue, the consistent direction across studies is \textit{reduced VIP-IN
contribution to disinhibition / context modulation}, but the magnitude,
the cellular sub-process (number, intrinsic excitability, synaptic input,
modulation depth, peptide release), and the cortical region differ
substantially. The implication is that no single biomarker --- VIP cell
count, firing rate, modulation depth --- can serve as a monolithic readout
of ``VIP-IN dysfunction'' across disorders. We discuss this implication
for biomarker development in the concluding synthesis and for
mechanistic modelling in \href{/computational-models}{Computational Models of VIP Circuit Function}.

\section{Translation: mouse circuits, human postmortem, iPSC-derived models}

Three readout modalities anchor the translation question and seldom agree.
\textit{Mouse circuit work} --- the dominant source of evidence in §6--§10 and in
the disease subsections above --- emphasizes intrinsic excitability,
firing rate, synaptic transmission, and behavioural consequences in the
intact brain \citep{Goff2023, Mossner2020, Bhandari2024, Chen2023b, Rahmatullah2023}.
\textit{Human postmortem} work, by contrast, is restricted to
transcriptomic / morphological / proteomic readouts in fixed tissue: the
human MTG, M1, and motor-cortex atlases of
\citep{Hodge2019, Bakken2021a, Lake2018}
provide quantitative cell-type proportions, t-type identities, and
gene-expression profiles, but cannot directly access circuit-functional
phenotypes. \textit{Human iPSC-derived} and \textit{neurosurgical-tissue} approaches
\citep{Somogyi2025, Matthews2025} partially close the gap by
providing live human GABAergic neurons that can be patched and
manipulated, but currently report transcriptional, morphological, and
single-cell-physiology phenotypes that do not yet map cleanly onto the
mouse-model behavioural readouts. \citet{Simacek2025b}, which
characterized developmental synaptic-input maturation of mouse VIP-INs in
S1 cortex with millisecond-resolution patch-clamp and reported that
``mEPSC frequency increased before P8-10, while mIPSC frequency
increased at P14-16'' with E/I ratio constant across development,
illustrates the kind of developmental-physiological readout that human
iPSC-derived VIP cells do not yet support at comparable resolution.
\citet{Lattke2026} analysed human fetal cortex (10--20 weeks PCW) by
scRNA-seq and scATAC-seq with ASO perturbation and recovered a
``subtype-specific reduction in \textit{RORB}- and \textit{FOXP1}-expressing
excitatory neurons and widespread disruption of neurodevelopmental
transcriptional programs,'' demonstrating the kind of human-developmental
readout that mouse models cannot directly substitute for.

\begin{framed}
\textbf{Conflict --- mouse circuit phenotype vs human iPSC readout}\\
An earlier analysis and the broader human iPSC literature emphasize
\textit{transcriptional} and \textit{morphological} phenotypes that do not always
predict the \textit{intrinsic-excitability} and \textit{synaptic-transmission}
phenotypes that \citet{Goff2023} and other mouse-model studies
emphasize. The gap is not a contradiction in disease aetiology but a
mismatch in measurable phenotype. Cross-modal validation (mouse
intrinsic-excitability deficits \rightarrow matched human transcriptomic /
morphological signatures of the same VIP t-type) remains incomplete.
Resolution will require integrated Patch-seq / enhancer-AAV reagents
\citep{Lee2023a, Matthews2025} applied jointly to mouse disease models
and human iPSC-derived or neurosurgical preparations.
\end{framed}

\section{Therapeutic implications}

Several lines of evidence converge on VIP-pathway pharmacology as a
therapeutic angle. As noted in §AD above, the 5xFAD VIP/$\beta$-amyloid attenuation finding is
read here from a primary source filtered from this section's citation
key map (listed under unmet citation needs); the claim is therefore
advanced as a candidate therapeutic mechanism rather than as cited
in-section evidence. \citet{Aidilcarvalho2022}, working on hippocampal LTP, showed that VIP
acting on VPAC1 receptors restrains LTP and depotentiation by modulating
disinhibition, and the authors argue that ``VIP receptor ligands may be
useful to co-adjuvate cognitive stimuli therapies based on these
aspects.'' \citet{Bhandari2024} demonstrated that augmenting VIP
signalling rescues theta-band coherence and learning in the \textit{Fmr1}-/-
mouse, providing direct preclinical proof-of-concept for a Fragile-X
indication. Independent work has, however, highlighted \textit{VPAC1}-related
metabolic consequences --- including elevated GLP-1 in \textit{VPAC1R} null mice ---
that caution against indiscriminate VPAC1 activation. \citet{Miller2025}
emphasized that \textit{VPAC2} (\textit{Vipr2}) copy-number variation is a
schizophrenia risk allele and that the receptor ``[points] to the
importance of this gene for the maintenance'' of cortical function,
arguing for caution in dosing the \textit{VPAC2} axis. The therapeutic balance
is therefore non-trivial: activating \textit{VPAC1} may help AD- and Fragile-X
phenotypes but perturbing \textit{VPAC2} dosage carries schizophrenia-relevant
risk, and the convergence diagram of Figure~\ref{fig-vip-disease-convergence}
makes these intervention nodes explicit.

\section{Synthesis}

The cross-species and disease evidence assembled in this section supports
a layered translation logic. The four-subclass scheme is a stable
backbone across mouse, marmoset, macaque, and human cortex, and rodent
mechanistic work on the \textit{VIP} subclass therefore generalizes at the
\textit{subclass-identity} level
\citep{Tasic2018, Hodge2019, Bakken2021a, Braininitiativecellcensusnetworkbiccn2021, Lee2023a, Yao2023}.
At the \textit{t-type} level, primate cortex contains expanded CGE diversity
\citep{Bakken2021a, Bakken2021b} and human-specific subtypes
\citep{Boldog2018, Chartrand2023, Somogyi2025} whose function the rodent
work does not directly speak to. At the \textit{disease} level, multiple
disorders converge on VIP-IN dysfunction
\citep{Mossner2020, Chen2023b, Bhandari2024, Goff2023, Hanno2026, Michaud2024, Miller2025, Batistabrito2017, Kiss2026, Gil2024},
but the \textit{direction}, \textit{magnitude}, and \textit{circuit locus} of the deficit are
heterogeneous and must be reported as such. The empirical heterogeneity
catalogued here, alongside that of §6--§10, motivates the formal
synthesis of \href{/computational-models}{Computational Models of VIP Circuit Function}, which asks whether
biophysical and rate-based three-interneuron models can absorb this
heterogeneity into a coherent computational role for VIP cells, and the
critical reassessment of the concluding synthesis, which
revisits the conserved disinhibition motif \citep{Pi2013} against the human and disease
evidence summarised above.

\begin{figure}[!htbp]
\centering
\includegraphics[width=0.95\linewidth]{files/fig-species-disease-87562a1ed1905db9e4dda5548683494b.png}
\caption[]{\textbf{Cross-species VIP-IN density and cross-disease functional change.}
\textbf{(A)} VIP+ density anchors across species. Bars show two
mismatched-denominator anchors: VIP+ as \% of GABA+ cells in P9 rat
primary auditory cortex \citep{Ouellet2014} (8\%, IHC), and VIP+ as \% of
total neurons in human M1 cytoarchitectural subdivisions A4a / A4b / A4c
(4.5\% / 5.2\% / 6.6\%, RNAscope+IHC).
\textbf{Caveat (audited):} the two species use \textit{different denominators}
(GABA+ vs total neurons), \textit{different ages} (P9 vs adult postmortem), and
\textit{different methods}; the apparent quantitative similarity in the 5--10\%
range masks methodological non-comparability. \textbf{(B)} Cross-disease,
cross-region, cross-axis qualitative direction-of-change in
autism-spectrum mouse models. Two entries are retained from the audit
after dropping rows that violated metric homogeneity. \textit{Goff 2023}:
\textit{Scn1a} +/- mouse cortex \textit{ex vivo}, baseline VIP-IN firing 34 $\pm$ 2.5 Hz
vs WT 42 $\pm$ 2.9 Hz, \textit{decreased} \citep{Goff2023}. \textit{Rahmatullah 2023}:
\textit{Fmr1}-/- mouse V1 \textit{in vivo}, 57.3\% of VIP cells modulated vs WT 73.2\%,
\textit{decreased} \citep{Rahmatullah2023}. \textbf{Caveat (audited, mandatory):
restructure dropped the entries that violated metric homogeneity
(p-only, baseline-rate, and F-statistic rows) and reframes the panel as
qualitative direction-of-change in autism-spectrum mouse models --- this
is internally consistent: both retained entries show decreased VIP-IN
function vs WT (Scn1a+/- 34 Hz vs WT 42 Hz; Fmr1-/- 57.3\% modulated vs
WT 73.2\%), so the qualitative direction-of-change axis is harmonized
and verifiable from the source sentences. However, two real defects
remain: (1) study\_label for the Fmr1 entry (DOI 10.1101/2023.01.03.522654;
Rahmatullah 2023) is rendered with a fallback label because the cite\_key
Unmapped\_37e340 carries null first\_author\_surname and null year in the
synthesized phase-5 citation map; (2) the underlying quantitative axes
are still different (Hz vs \% of cells modulated) and the two disease
models (Scn1a Dravet vs Fmr1 Fragile X) have different biological
etiologies grouped under ``autism-spectrum''. These are tolerable for a
qualitative forest panel as the restructure acknowledges, but warrant a
caveat --- magnitudes are NOT comparable across rows.} \textbf{(C)} Qualitative
disease $\times$ phenotype matrix (cell-loss / activity / connectivity /
morphology) across AD, ASD/ID, schizophrenia, and Dravet, populated from
text-extracted findings \citep{Michaud2024, Mossner2020, Chen2023b, Bhandari2024, Vogt2014, Goff2023, Hanno2026, Batistabrito2017, Miller2025, Vormsteinschneider2020, Gil2024}.
The matrix is qualitative; cells indicate direction of reported change
(↑/↓) or ``altered'' without intent to be commensurable across rows.}
\label{fig-species-disease}
\end{figure}

\begin{figure}[!htbp]
\centering
\includegraphics[width=0.9\linewidth]{files/fig-vip-disease-conv-100a6055ecf0d0b6cd43d617f0b8610b.png}
\caption[]{\textbf{Convergence schematic: upstream insults \rightarrow mid-level VIP-IN
mechanisms \rightarrow circuit phenotypes.} Funnel diagram from upstream
regulators (\textit{MECP2}, \textit{TCF4}, \textit{FMR1}, \textit{SCN1A}, \textit{CNTNAP2}, \textit{APP}/\textit{PSEN},
\textit{DISC1}, \textit{VIPR2}) through mid-level VIP-IN mechanisms (intrinsic
excitability, synaptic input, peptide/GABA co-release, density,
connectivity) to circuit phenotypes (E/I imbalance, gain
dysfunction, oscillatory disruption). The right panel cartoons the
mouse-vs-human VIP-IN comparison anchored to \citet{Hodge2019} and
\citet{Yao2023}: subclass conserved, t-type expansion in primate L1,
human-specific \textit{VIP PCDH20} and \textit{MC4R} types
\citep{Chartrand2023}. The schematic externalizes the
translational hypotheses that Figure~\ref{fig-species-disease} only
summarises piecewise.}
\label{fig-vip-disease-convergence}
\end{figure}

\section{Methodological caveats specific to this section}

Several methodological caveats specific to species and disease comparisons
recur and deserve to be flagged once before a reader weighs the evidence
above. First, the \textit{VIP-Cre} driver line and its cross-species analogues
do not capture identical populations: in mouse, VIP-Cre and VIP-IRES-Cre
labelled cells overlap heavily with the transcriptomic \textit{Vip} subclass
\citep{Tasic2018, Yao2023, Tasic2016}, but the human equivalent lacks a
matched genetic-driver tool, and human VIP characterization relies on
RNAscope, IHC, and enhancer-AAV approaches
\citep{Teymornejad2024a, Lee2023a}. The
\citet{Lee2023a} enhancer-AAV taxonomy is the most direct attempt to
build a shared genetic-handle framework across species. Second, ChAT
co-expression in upper-layer \textit{VIP}-Chat cells is a stable mouse feature
\citep{Tasic2016, Tasic2018, Tremblay2016, Yao2023} whose human homology is
inferred only at the transcriptomic level
\citep{Hodge2019, Lee2023a}; functional confirmation in human tissue
remains incomplete. Third, the VIP-IN role within the broader 5-HT3AR+
CGE-derived population varies across reports ({\textasciitilde}40\% of the 5-HT3AR
population in S1 by \citet{Rudy2011}, with related estimates in
), and species-specific 5-HT3AR
distribution complicates direct transfer. Fourth, the human
neurosurgical-tissue and postmortem-tissue preparations used by
\citep{Somogyi2025, Boldog2018, Hodge2019, Lee2023a} are
not interchangeable: neurosurgical samples allow live recording but are
biased toward epilepsy / tumour cohorts, and postmortem material allows
unbiased anatomical sampling but precludes live functional readouts.
Fifth, mouse-model phenotypes are typically reported as point estimates
with small \textit{n} per condition and large between-laboratory variance; the
behavioural-phenotype translation literature surveyed by
argues for systematic between-laboratory replication
as a corrective. Sixth, the developmental window in which VIP-INs are
sampled differs across species: rodent characterizations frequently
use juvenile (P14--P28) animals
\citep{Mossner2020, Bhandari2024, Mcfarlan2024a},
human samples are typically adult or fetal
\citep{Lattke2026, Hodge2019, Lake2018}, and disease
penetrance is itself developmental
\citep{Mossner2020, Chen2023b, Lattke2026, Vormsteinschneider2020, Simacek2025b}.

\section{Convergence on VIP-IN dysfunction across disorders}

Putting the disease subsections side-by-side reveals a consistent
upstream-to-downstream architecture that Figure~\ref{fig-vip-disease-convergence}
makes explicit. \textit{Upstream genetic perturbations} span ion channels
(\textit{SCN1A} \citep{Goff2023, Vormsteinschneider2020}), synaptic-adhesion
molecules (\textit{CNTNAP2}), translational regulators
(\textit{FMR1} \citep{Bhandari2024, Rahmatullah2023}), transcription factors
(\textit{TCF4} \citep{Chen2023b}; \textit{MECP2} \citep{Mossner2020, Goff2021, Ferguson2023b}),
and disease-specific proteostasis lesions (\textit{APP}/\textit{PSEN} in 5xFAD and 3xTg
\citep{Michaud2024}; expanded HTT in R6/2
). \textit{Mid-level VIP-IN phenotypes} recur:
reduced intrinsic excitability and firing rate
\citep{Goff2023, Mossner2020}, increased input resistance with reduced
sEPSC frequency \citep{Chen2023b}, reduced sensory modulation
\citep{Rahmatullah2023}, and selective cell-number reductions
\citep{Hanno2026, Chen2023b}. \textit{Downstream circuit consequences} converge
on E/I imbalance, gain dysfunction, and oscillatory disruption
\citep{Mossner2020, Bhandari2024, Batistabrito2017, Kiss2026}.
The convergence is not perfect --- the 5xFAD-model report of a
\textit{protective} role for exogenous VIP signalling complicates a simple
``loss of VIP function = disease'' reading, and the 3xTg
``altered firing without cell loss'' pattern of \citet{Michaud2024}
contrasts with the cell-loss readings in \textit{Tcf4} \citep{Chen2023b} and
ASD models \citep{Hanno2026} --- but the architecture is recurrent
enough to justify a \textit{VIP-IN-as-vulnerable-node} framing across
psychiatric, neurodevelopmental, and neurodegenerative disorders.

\section{Open questions and unmet evidence needs}

We close with three questions the present evidence cannot answer. First,
no published study quantifies VIP+ cell density across mouse, marmoset,
macaque, and human in a \textit{single denominator, single area, single age}
design; the cross-species fraction-of-CGE bar-chart envisaged in the
review plan (Panel C of the original Figure~\ref{fig-species-disease} plan) is
buildable from \citet{Bakken2021a}, \citet{Bakken2021b}, and
\citet{Braininitiativecellcensusnetworkbiccn2021} only as proportions
of t-type clusters, not as physical-density counts. Second, the mouse vs
human AD comparison is anchored on the mouse side
\citep{Michaud2024} and on the cross-disorder review side
\citep{Gil2024}, with the human-postmortem snRNA-seq pole carried by
\citep{Gabitto2024} (drawn from the master citation map
under the same scope precedent applied to other master-only keys
used in this section). The AD-conflict admonition is therefore
now anchored on both sides. Third, no current human
iPSC-derived VIP-IN preparation reproduces the
intrinsic-excitability, synaptic-transmission, or oscillatory-coupling
phenotypes that mouse models report; cross-modal validation ---
patch-clamp / Patch-seq characterization of human iPSC-derived VIP cells
matched to mouse disease-model intrinsic phenotypes --- is the most
direct missing experiment.

\begin{framed}
\textbf{Apparent conflicts in cross-species and disease evidence}\\
Four further pairings flagged at evidence-package level merit explicit
treatment here. (i) \citet{Krolewski2025} reported a reduced
GABAergic transcript signature including \textit{VIP}-related transcripts in
human cingulate area BA24c′ in postmortem schizophrenia / bipolar
cohorts; the apparent-conflict partner is the up-regulated
PACAP-pathway reading documented at evidence-package level
\citep{Slabe2023}. The two reports use
overlapping cingulate territory but different transcript panels and
diagnostic groups, so the contradiction may be partly axis-dependent
rather than directly opposing. (ii) \citet{Krolewski2025} and
out-of-master postmortem schizophrenia work have also been read against
each other on whether VIP-specific transcript reduction localises to
DLPFC or cingulate; the in-section reading favours an
area-conditional reduced GABAergic signature rather than
region-uniform VIP loss. (iii) The status of human L1 rosehip cells as
species-specialized \citep{Boldog2018} is balanced against
cross-mammalian conservation of MGE-derived inhibitory programs
described by \citep{Keefe2023}; the in-section reading is that the
MGE program is conserved while specific CGE-derived t-types
(rosehip, \textit{VIP PCDH20}, \textit{MC4R}) are species-specialized. (iv) Whether
human rosehip cells or human double-bouquet cells are the more
prominent species-specialized type is unsettled; primate double-bouquet
work foregrounds double-bouquet cells, whereas
\citep{Boldog2018} foreground rosehip cells, with no current
direct cross-comparison. (v) On the developmental side,
reported a constant \textit{VIP}-IN contribution to
cortical E/I balance across early postnatal stages; the
apparent-conflict partner (HosseiniFin2025, listed under unmet
citation needs) reports an age-dependent decline in
SST-IN density. The two are formally compatible (different cell
classes, different time-windows) and are surfaced here for
completeness rather than as a substantive contradiction.
\end{framed}

\section{Connection to neighbouring sections}

The translational and disease evidence assembled here closes a gap that
\href{/in-vivo-behavior}{In Vivo Function During Behavior} and \href{/cross-areas}{VIP Interneurons Across Brain Regions} opened: rodent
\textit{in vivo} characterizations of VIP-IN function across V1, A1, S1, and
PFC must be tempered by the species- and disease-relevant heterogeneity
catalogued above. The disinhibition-motif framing reviewed in
\href{/disinhibition-framework}{Local Circuit Motifs and the Disinhibition Framework} retains its empirical anchor in
mouse cortex but applies to human cortex with the qualifications that
(i) primate CGE expansion may diversify the VIP\rightarrowSST\rightarrowPyr motif \citep{Pi2013} into
multiple parallel motifs \citep{Bakken2021a, Lee2023a, Chartrand2023};
(ii) human-specific VIP subtypes may implement connectivity rules that
mouse mechanistic data do not constrain
\citep{Chartrand2023}; and (iii) disease
phenotypes converge on the \textit{VIP} class but heterogeneously perturb
different sub-mechanisms (number, intrinsic excitability, synaptic
input, peptide release)
\citep{Mossner2020, Chen2023b, Bhandari2024, Goff2023, Hanno2026}.
The molecular and developmental-lineage backbone surveyed in
\href{/classification}{Molecular Identity and Transcriptomic Taxonomy} and CGE~origin~and~the~5-HT3AR~/~Adarb2~lineage carries
forward without revision: VIP cells are CGE-derived members of the
5-HT3AR+ lineage in all sampled species
\citep{Fishell2011, Lodato2011, Mayer2018, Tasic2018, Hodge2019, Yao2023, Tremblay2016, Rudy2011}.
What does \textit{not} carry forward without revision is the implicit
assumption that mouse mechanistic data exhaust the description of human
VIP-IN function and human disease --- an assumption Figure~\ref{fig-species-disease}
and Figure~\ref{fig-vip-disease-convergence} are designed to make
explicit.

\section{Unmet citation needs}

The following references support specific claims in this section but
fall outside this review's allowlisted bibliography and could not be
cited in body text. Each is reported here by DOI so the underlying claim is
auditable:

\begin{itemize}
\item \textbf{Slabe2023} \citep{Slabe2023}
--- proposed up-regulation of PACAP-related transcripts in postmortem
cingulate cortex of subjects with schizophrenia, the apparent-conflict
partner to \citet{Krolewski2025} (which reports a reduced GABAergic
/ \textit{VIP}-containing transcript signature in BA24c′).
\item \textbf{HosseiniFin2025} (no resolved DOI in master) --- developmental
trajectory in which SST-IN density declines with age, the
apparent-conflict partner to an early report (constant E/I
contribution from VIP-INs across early development).
\item \textbf{Tiwari2026} --- primary 5xFAD VIP / $\beta$-amyloid attenuation source
used in the AD and Therapeutics subsections only as a candidate
mechanism (not as cited in-section evidence).
\item \textbf{Blumenstock2025} --- primary R6/2 Huntington's-disease source for
the upstream-VIP disinhibitory cascade reading.
\item \textbf{Liu2025a} --- primary sleep-disruption source.
\item \textbf{Oh2022} --- primary subgenual ACC CRH+/VIP+ overlap source.
\end{itemize}

In each case the corresponding claim has been hedged in the body text
as a candidate mechanism with explicit acknowledgement that the
primary source falls outside this section's allowlisted bibliography. The
Rahmatullah-2023 V1 in-vivo Fragile-X paper
(\cite{Rahmatullah_2023})
was successfully resolved as \texttt{Rahmatullah2023} in the filtered citation
map and is cited normally; its earlier ``study\_label = None None''
rendering in the audited Phase-6 panel data is flagged in the caption
of Figure~\ref{fig-species-disease}B.

\section{A quantitative note on cross-area, cross-species comparisons}

Two independent observations should accompany any quantitative
cross-species claim about VIP-INs. (1) Cell-density estimates depend
sharply on the choice of denominator (GABA+, total neurons, total cells,
5-HT3AR+, \textit{Vip}-Cre-labelled): the Ouellet (8\% of GABA+ in P9 rat A1)
versus Teymornejad (4.5--6.6\% of total neurons in adult human M1)
disparity discussed above is a generic feature of the literature and
not a species-specific finding \citep{Ouellet2014}.
Independent rat-brain slice work illustrates the sensitivity to driver
choice, with reports of ``the fraction of VIP expressing neurons among
eGFP fluorescent cells was 7.1 $\pm$ 1.2\% (median = 6.6\%)'' in
a non-VIP-targeted driver line, an order-of-magnitude change from the
{\textasciitilde}30--40\% expected within the 5-HT3AR+ population
\citep{Rudy2011}. (2) Cross-species transcriptomic
\textit{differential} gene expression also depends on the chosen reference
genome and the alignment-by-orthology pipeline; a previous study
showed that the \textit{GRIA2} stoichiometry signature is conserved across
``ferrets, rodents, marmosets and humans'' only when the same
analysis pipeline is applied to all species, illustrating that
cross-species claims are pipeline-dependent. Quantitative claims in
this section have therefore been reported with the species, region,
denominator, and method labels intact, and we discourage extrapolation
across these axes without explicit justification.

\section{Closing note}

The translation logic of this review is not simply that ``rodent VIP
work generalizes to human'' or ``rodent VIP work fails to generalize.''
It is that subclass-level identity transfers, t-type-level identity
partially transfers, and disease phenotypes converge on the VIP class
without converging on a single mechanism. Subsequent sections will treat
this as the empirical baseline against which the disinhibition motif and
the standard VIP--SST interaction must be re-evaluated for human cortex
\citep{Hodge2019, Bakken2021a, Lee2023a, Goff2023, Bhandari2024, Mossner2020, Chen2023b, Hanno2026},
and against which any future computational model of cortical VIP
function must be benchmarked
\citep{Yao2023, Hodge2019, Lee2023a, Chartrand2023, Boldog2018}.

With the empirical landscape --- molecular through clinical --- now in view, \href{/computational-models}{Computational Models of VIP Circuit Function} asks how computational models of VIP circuit function instantiate, generalise, or contradict it.


\section{Computational Models of VIP Circuit Function}
\label{sec:12_computational_models}

Computational models of cortex have, over the past decade, moved VIP interneurons from a peripheral curiosity to an indispensable structural element. Three pressures forced this shift: the demonstration that VIP cells preferentially target SST cells and disinhibit pyramidal output \citep{Pi2013, Pfeffer2013, Lee2013}; the discovery that VIP firing tracks state, locomotion, reward, and reinforcement signals \citep{Fu2014, Pakan2016, Krabbe2019, Dipoppa2018}; and the recognition that any cortical circuit that omits VIP cannot reproduce these state- and learning-dependent gain changes from PV/SST dynamics alone \citep{Garciadelmolino2017, Litwinkumar2016, Hertag2019}. The literature now divides cleanly into two- and three-interneuron rate and spiking models that treat VIP as the disinhibitory gain knob \citep{Hertag2019, Litwinkumar2016, Garciadelmolino2017, Bos2025, Veit2023}, four-interneuron extensions that add NDNF/LAMP5 as a parallel input-gating channel \citep{Iqbal2025, Hartung2024, Anastasiades2021a}, and circuit-level learning models in which VIP-mediated disinhibition implements top-down credit assignment \citep{Hertag2019, Lee2025}. The same disinhibitory motif has been independently recruited as the substrate for predictive-coding error neurons \citep{Hertag2020, Hertag2021, Rossbroich2025, Nemati2025a, Nemati2025b, Shipp2016}, for region-specific reinforcement-learning credit assignment \citep{Chevy2024}, and for flexible decision-making attractor switching \citep{Shen2023, Yang2016}, indicating that VIP is being asked to support several non-equivalent computational jobs at once. This proliferation has revealed an awkward fact: the same connectivity graph, parameterised differently, can produce qualitatively opposite predictions for how VIP perturbations should affect pyramidal output \citep{Hertag2019, Tahvili2025a, Tahvili2025b, Garciadelmolino2017}, so that disagreement among published models is not a sign of empirical conflict but of weakly constrained free parameters. This section synthesises these architectures, exposes the gain-regime conflicts they generate, and argues that the current bottleneck is identifiability, not biology --- that the field's next decade depends less on adding cell types than on disentangling the parameterisations the existing graphs already admit (Figure 23, Figure 24).

\section{Architectures and the role of VIP inclusion}

Standard three-interneuron (3-IN) circuits --- Pyr+PV+SST+VIP --- emerged from the connectivity surveys of \citep{Pfeffer2013} and the disinhibitory experiments of \citep{Pi2013, Lee2013, Fu2014}, and have been formalised as rate models \citep{Litwinkumar2016, Garciadelmolino2017, Hertag2019}, spiking implementations \citep{Veit2023}, and biophysically detailed mean-field circuits \citep{Reimann2026, Jadi2012, Isbister2026}. Across this family, VIP\rightarrowSST disinhibition supplies the dominant route by which top-down or neuromodulatory signals release pyramidal dendrites from inhibition \citep{Karnani2016a, Abs2018}. Variants of the 3-IN scheme have been instantiated for primary visual cortex \citep{Hahn2020, Kim2025b}, primary auditory cortex \citep{Park2020, Natan2017}, somatosensory cortex \citep{Hua2022}, and prefrontal areas \citep{Shen2023}, and the same motif now appears in cerebellar circuit theory \citep{Park2023} and in subcortical--cortical loop models \citep{Cattani2024}, indicating that the design has effectively become standard. To gauge how universally this design choice has been adopted, we audited 25 unambiguously classifiable computational papers in our package and found that 21/25 explicitly include VIP, with only 4/25 --- almost all early or PV/SST-focused models --- omitting it (Figure 23D). The omitting models cluster on a single rationale: they treat VIP recruitment as exogenous and absorb its effect into a slow, time-varying SST drive, exchanging a fourth state variable for a tractable two-population reduction \citep{Tahvili2025a, Tahvili2025b}. The price of that simplification, however, is that the resulting models cannot natively express any computation that requires VIP activity to be gated by behavioural state, since the slow drive is imposed rather than computed.

The four-interneuron (4-IN) family adds NDNF/L1 neurogliaform cells \citep{Iqbal2025, Hartung2024, Vighagen2021} and, in some implementations, LAMP5 lineage cells, as a separate top-down channel that gates apical dendrites in parallel to the SST\rightarrowPyr branch \citep{Anastasiades2021a, Bilash2023, Vanderveer2021}. This addition is not cosmetic. \citep{Iqbal2025} shows that an L2/3 spiking circuit augmented with LAMP5 implements a soft winner-take-all alongside the classic VIP\rightarrowSST disinhibition, and that disabling the LAMP5 channel collapses the network's ability to perform context-invariant gain normalisation; \citet{Hartung2024} demonstrate complementarily that NDNF cells deliver volume inhibition with a slow time constant that none of the rate-based 3-IN models capture, providing a low-pass envelope on the otherwise fast VIP\rightarrowSST disinhibition. Compartmental models of NDNF dendrites further argue that the channel has its own internal computation, with backpropagating signals interacting with synaptic input to produce a second layer of dendritic disinhibition \citep{Griesius2025}, and L1-resolved circuit models suggest the 4-IN architecture is necessary to reproduce layer-specific inhibitory dynamics observed \textit{in vivo} \citep{Vighagen2021}. The conflict between 3-IN sufficiency and 4-IN necessity is therefore neither philosophical nor purely numerical: 3-IN proponents \citep{Hertag2019, Litwinkumar2016, Bos2025} show that the available behavioural and physiological signatures of state-dependent gain can be reproduced without an explicit NDNF compartment, while 4-IN proponents \citep{Hartung2024, Vighagen2021} argue that doing so requires VIP\rightarrowSST gain values that are inconsistent with paired-recording data and only the NDNF channel can absorb the missing slow inhibition. Distinguishing the two empirically requires compartment-specific signatures --- an apical inhibition that survives SST silencing, a slow envelope on VIP-driven facilitation that exceeds SST-mediated kinetics --- and these are precisely the recordings still rare in the literature \citep{Bilash2023, Lenkey2025, Anastasiades2021a}.

A third class, predictive-coding and dendritic-error models, places VIP in the prediction-error and credit-assignment loop \citep{Hertag2019, Hertag2020, Hertag2021, Hertag2023, Lee2025, Rossbroich2025, Nemati2025a, Nemati2025b, Shipp2016}, where disinhibition is no longer a static gain term but a temporally precise teaching signal whose magnitude depends on mismatch between top-down expectation and bottom-up drive . \citet{Wilmes2019} realised this most explicitly with a top-down inhibitory plasticity circuit in which disinhibitory motifs gated by reward open compartmental access to top-down feedback, allowing prediction errors to be encoded as compartment-specific deviations from a learned baseline; \citet{Hertag2020} and \citet{Hertag2021} extend this to spiking microcircuits in which prediction-error neurons emerge through plasticity, with VIP\rightarrowSST disinhibition controlling the refinement of error coding and the formation of distinct positive- and negative-error populations. \citet{Wilmes2019} add an inhibitory-plasticity twist, arguing that the feedback projections that target VIP cells implement a supervisor signal that gates the timing of cortical learning, while \citep{Chevy2024} and \citep{Hertag2023} situate the same disinhibitory motif inside reinforcement-learning frameworks as the gate of region-specific credit assignment and feedforward/feedback uncertainty estimation respectively. Models in this class do not contradict the 3-IN gain-knob view; rather, they re-interpret the same connectivity as implementing a different objective function. This re-interpretation is consequential because the same VIP perturbation predicts different downstream effects depending on the assumed objective: under a gain-control objective, VIP activation linearly scales pyramidal responses \citep{Hertag2019, Veit2023}; under a credit-assignment objective, it transiently opens a learning gate whose effect on output is mostly visible in the \textit{change} of subsequent responses, not in steady-state firing \citep{Lee2025, Chevy2024}. The conflict between ``VIP as gain knob'' and ``VIP as credit/state signal'' therefore corresponds to which loss the circuit is hypothesised to minimise --- a regression error, in the gain-control framing, versus a temporal-difference or backpropagated dendritic error in the learning framing --- and the two predict almost identical instantaneous physiology while diverging sharply in trial-over-trial dynamics \citep{Hertag2020, Lee2025, Krabbe2019}.

The functional-role audit (Figure 23C) makes this multiplicity explicit. Across the modelling literature, VIP enters in five overlapping capacities: as a gain term \citep{Hertag2019, Litwinkumar2016, Veit2023, Bos2025, Richter2022}; as an attention/state controller \citep{Garciadelmolino2017, Hahn2022, Waitzmann2024, Wagatsuma2022, Poort2022}; as a learning/credit-assignment gate \citep{Lee2025, Chevy2024, Wilmes2019, Hertag2020, Hertag2021}; as an oscillation organiser \citep{Veit2023, Cattani2024}; and as a context selector that biases winner-take-all attractor dynamics \citep{Shen2023, Yang2016, Furutachi2024}. These categories are not interchangeable, and many papers commit a model to one role implicitly by choosing a particular cost function, perturbation protocol, or behavioural readout. The 3-IN sufficiency claim --- that all five roles can be supported by Pyr+PV+SST+VIP --- is therefore best understood as a structural claim about graph minimality, not as a claim that one parameterisation can simultaneously implement all five \citep{Hertag2019, Litwinkumar2016, Bos2025}. The 4-IN argument \citep{Iqbal2025, Hartung2024, Vighagen2021} is, conversely, a claim that some of these roles --- specifically the slow normalisation and the L1-specific gating --- require an additional dynamical degree of freedom that NDNF/LAMP5 supplies. The architectural debate is thus less about whether VIP belongs in the model than about how many independent computations a single VIP\rightarrowSST disinhibitory channel can perform without becoming over-determined.

A complementary axis cuts orthogonally to the 2-IN/3-IN/4-IN distinction: the spatial and biophysical resolution at which the circuit is described. At the coarsest end, mean-field rate models \citep{Litwinkumar2016, Garciadelmolino2017, Sanzeni2020, Kim2025b} collapse each interneuron class to a single firing-rate variable and study fixed-point structure analytically; at the spiking end, sparse and balanced E-I networks \citep{Veit2023, Pedrosa2020, Rossbroich2025} add membrane dynamics, conductance-based synapses, and short-term plasticity \citep{Seay2020, Cattani2024}; and at the most detailed end, multi-compartment biophysical models \citep{Reimann2026, Jadi2012, Isbister2026, Morabito2024} resolve apical from perisomatic inhibition and reproduce \textit{in vivo} dynamics across cell types. Crucially, these resolutions are not merely descriptive choices but partly determine which conflicts a model can adjudicate. Compartmental models \citep{Reimann2026, Jadi2012, Dorsett2021} can in principle distinguish whether an observed pyramidal effect arises from somatic versus dendritic inhibition --- a distinction that mean-field 3-IN circuits collapse --- and therefore form the only natural arena in which the 3-IN-versus-4-IN comparison can be resolved without ad hoc assumptions. Yet the same compartmental detail comes at a cost: the parameter space of a biophysical microcircuit is so large that, as \citet{Reimann2026} argue, the manifold of parameterisations compatible with cell-type-resolved electrophysiology \textit{grows} rather than shrinks, raising the identifiability problem to the foreground.

\section{Gain regimes, conflicts, and parameter-space mapping}

Beneath the architectural agreement, models disagree sharply about what VIP does to pyramidal gain (Figure 23B,C; Figure 24A). \citet{Hertag2019} produced the most influential 3-IN rate model, in which VIP-driven disinhibition operates \textbf{divisively} on pyramidal responses by withdrawing dendritic SST inhibition; this regime is recapitulated by \citep{Litwinkumar2016, Veit2023} and underlies most subsequent learning-rule analyses \citep{Lee2025, Hertag2020, Hertag2021}. \citet{Tahvili2025b}, by contrast, recently reported that a reduced two-interneuron model with strong recurrent SST inhibition and tuned VIP\rightarrowSST gain produces \textbf{subtractive} rather than divisive modulation, predicting a qualitatively different relationship between VIP activation and pyramidal threshold \citep{Tahvili2025a}. The disagreement is not a modelling artefact; it tracks a genuine ambiguity in the biology, because SST cells target both the apical dendrites and (at lower density) the perisomatic compartment \citep{Hioki2013}, and the relative weight of these two synaptic loci is precisely what selects between divisive and subtractive operation in compartmental analyses \citep{Jadi2012, Dorsett2021, Pedrosa2020}. Empirical work in auditory cortex makes the related point that PV and SST contribute differentially to gain control --- SST acting more divisively and PV more bidirectionally --- during adaptation \citep{Natan2017}, and circuit fits using their data place the divisive-versus-subtractive boundary at a specific ratio of dendritic-to-somatic SST conductance --- a parameter rarely measured \textit{in vivo} and one that \citet{Tahvili2025a} and \citet{Hertag2019} set to opposite default values. The conflict is therefore parametric, not structural, and its resolution requires compartment-resolved measurements rather than additional perturbation protocols \citep{Bilash2023, Lenkey2025, Reimann2026}.

The \citep{Pi2013} framing of net pyramidal facilitation likewise sits uneasily with \citep{Garciadelmolino2017}, who showed that the same connectivity, parameterised inside the inhibition-stabilised regime, yields \textbf{paradoxical} responses in which VIP activation can suppress rather than facilitate pyramidal firing \citep{Reimann2026, Sanzeni2020}. The mechanism is well-understood at the rate-model level: when the recurrent excitatory loop within the pyramidal population is strong enough to make the network an inhibition-stabilised network (ISN), perturbing any inhibitory population produces a counter-intuitive response in which pyramidal rate moves opposite to the naive expectation \citep{Sanzeni2020, Garciadelmolino2017}. \citet{Sanzeni2020} show that ISN operation is widespread across cortex, so the question is not whether the paradoxical regime exists but whether VIP perturbations probe it. Empirical reports of net facilitation under VIP optogenetic activation \citep{Pi2013, Fu2014, Lee2013} are typically obtained at moderate stimulation strengths and may probe the linear, non-paradoxical regime; reports of suppression or non-monotonic effects \citep{Dipoppa2018, Veit2023, Bilash2023} may correspond to operation closer to the ISN boundary, or to states with elevated SST/PV recurrence. \citet{Garciadelmolino2017} make this explicit by mapping the boundary between facilitation and paradoxical suppression in (recurrent E-E gain $\times$ VIP\rightarrowSST weight) space, and \citet{Phillips2016} show that activating versus inactivating the same interneuron class can produce qualitatively asymmetric firing-rate distributions, consistent with the network straddling a non-linear regime. The question ``does VIP activation facilitate pyramidal output?'' therefore has no single answer; it has a state-dependent answer whose sign depends on a small number of recurrent-strength parameters that current data do not pin down.

Plotting these models in a phase diagram of VIP\rightarrowSST coupling versus dendritic SST\rightarrowPyr nonlinearity (Figure 23B; Figure 24A) makes the disagreement structural rather than empirical: \citep{Hertag2019, Veit2023} cluster in a divisive-gain corner, \citep{Tahvili2025a, Tahvili2025b} occupy the subtractive corner, and \citep{Garciadelmolino2017, Reimann2026, Sanzeni2020} sit on the paradoxical boundary, with \citep{Iqbal2025, Hartung2024, Vighagen2021} displaced into a fourth quadrant where NDNF input recasts the dendritic operating point. Spiking and population models built around these regimes \citep{Tahvili2025a, Pedrosa2020, Rossbroich2025} show that all three primary outcomes --- net facilitation, divisive gain, paradoxical suppression --- are reachable within a single 3-IN graph by retuning two or three weights, a sensitivity that any inference scheme must confront \citep{Garciadelmolino2017, Dellal2025}. The conflicts are therefore not ``which experiment is right'' but ``which regime does the cortex occupy'', and existing data are insufficient to localise it: VIP-activation experiments \citep{Pi2013, Fu2014, Lee2013} consistently report pyramidal facilitation in active states, but cell-type optogenetic perturbations during behaviour reveal more complex, often paradoxical effects \citep{Dipoppa2018, Veit2023, Bilash2023}. Models of state-dependent attention \citep{Garciadelmolino2017, Dipoppa2018, Hahn2022, Waitzmann2024, Poort2022, Wagatsuma2022} and of fear-conditioning credit assignment  further suggest that the operating regime is itself dynamic, switching between subtractive and divisive operation as a function of cholinergic and noradrenergic drive \citep{Fu2014, Pakan2016, Anastasiades2021a}. The role-category audit (Figure 23C) accordingly shows VIP entering models in five distinct functional capacities --- gain control, attention/state, learning/credit assignment, oscillations, and context modulation --- that overlap but are not interchangeable \citep{Hertag2019, Garciadelmolino2017, Veit2023, Lee2025, Furutachi2024, Shen2023, Yang2016}.

The four-interneuron extensions further multiply the parameter space (Figure 24B). \citep{Iqbal2025} demonstrates that adding a LAMP5/NDNF-class channel to a 3-IN model recovers state-dependent layer-1 inhibition that 3-IN circuits attribute to SST, while \citep{Hartung2024} shows that NDNF-driven volume inhibition imposes a slow envelope on VIP disinhibition that none of the rate models capture. \citet{Vighagen2021} add layer-specific control of inhibition by NDNF as a third dynamical degree of freedom, whose effect on the dendritic operating point cannot be absorbed into a single SST conductance term. LAMP5 and additional layer-1 cell types \citep{Anastasiades2021a, Vanderveer2021} raise the same identifiability question for top-down inputs that NDNF raises for dendrites: a 3-IN circuit with strong VIP\rightarrowSST gain and a 4-IN circuit with weaker VIP\rightarrowSST gain plus NDNF inhibition can produce indistinguishable pyramidal observables \citep{Iqbal2025, Hartung2024, Vighagen2021, Reimann2026}. The 3-IN sufficiency proponents \citep{Hertag2019, Litwinkumar2016, Bos2025} reply that the additional degrees of freedom are model-internal and not yet tied to compartment-specific data, while the 4-IN necessity proponents \citep{Hartung2024, Vighagen2021} reply that 3-IN fits to L1-resolved recordings require VIP\rightarrowSST values inconsistent with paired-recording slopes. Compartmental models that resolve apical from perisomatic inhibition \citep{Jadi2012, Reimann2026, Dorsett2021} partly relieve this degeneracy by predicting compartment-specific signatures of each architecture, but the necessary recordings --- dual-compartment imaging during VIP, SST, and NDNF perturbation --- are rare \citep{Bilash2023, Lenkey2025, Anastasiades2021a}. The empirical signature that would distinguish 3-IN sufficiency from 4-IN necessity is therefore concrete: 4-IN architectures predict a slow, layer-1-localised inhibition that survives complete SST silencing and tracks behavioural state on second-to-tens-of-seconds timescales \citep{Hartung2024, Vighagen2021}; 3-IN architectures predict that the same slow inhibition collapses with SST silencing and emerges instead from short-term plasticity at SST\rightarrowPyr synapses \citep{Hertag2019, Litwinkumar2016}. The fact that this experiment has not yet been published is, on the present view, the single most consequential gap in the modelling literature.

A subtler conflict cuts across these gain-regime debates: what role is VIP being asked to play computationally, and which objective function does the model implicitly optimise? In the gain-knob framing \citep{Garciadelmolino2017, Hertag2019, Veit2023}, VIP is a static (or slowly varying) parameter that scales pyramidal sensitivity in the service of a regression-style objective --- minimising the discrepancy between bottom-up drive and a context-modulated readout. In the state/credit-signal framing \citep{Hertag2020, Hertag2021, Lee2025, Wilmes2019, Chevy2024}, VIP is a dynamic teaching variable that gates plasticity at upstream or apical synapses in service of a temporal-difference, prediction-error, or backpropagated-error objective. The two framings predict almost identical instantaneous physiology --- both show pyramidal facilitation when VIP is activated and a withdrawal of dendritic inhibition --- but diverge in trial-over-trial dynamics. Under the gain-knob objective, the effect of a VIP perturbation is reversible and transient, present during the perturbation and absent immediately after; under the credit-signal objective, the same perturbation should produce a \textit{lasting} change in subsequent responses because the perturbation alters the learning signal that updates synaptic weights \citep{Lee2025, Hertag2020, Hertag2021, Chevy2024}. \citet{Hertag2023} formalise this in a confidence-estimation framework, in which the disinhibitory gate sets the relative weight of feedforward versus feedback uncertainty during a single trial; \citet{Sharafeldin2025} show that the same motif develops cell-type-specific encoding of prediction and reward during novelty detection without explicit supervision, again indicating that the disinhibitory channel is doing computational work beyond gain. The empirical signature that distinguishes the two interpretations is therefore the \textit{time course} of the perturbation effect, not its sign, and few of the experimental studies cited in the modelling literature \citep{Pi2013, Fu2014, Lee2013, Krabbe2019, Myersjoseph2024, Dipoppa2018} have been designed to separate within-trial gain effects from across-trial learning effects.

A second, less obvious objective-function distinction concerns whether VIP is an actor in a closed loop or a target of an external supervisor. Models in which VIP firing is driven by long-range top-down inputs \citep{Pi2013, Lee2013, Wilmes2019} treat VIP as a target variable whose dynamics are imposed by upstream regions; models in which VIP participates in within-area recurrence \citep{Garciadelmolino2017, Hertag2019, Hertag2020, Pedrosa2020, Rossbroich2025} treat its dynamics as emergent. The parameter-space consequence is large: in target-driven formulations, the VIP\rightarrowSST weight is the only free parameter that matters and is identifiable up to a scaling, whereas in emergent formulations the VIP\rightarrowSST weight trades off against the SST\rightarrowVIP feedback weight and the VIP self-recurrence term, producing the families of equivalent parameterisations that \citet{Litwinkumar2016} and \citet{Garciadelmolino2017} have analytically characterised. Recent work attempting to fit either formulation to multi-cell perturbation data \citep{Moreni2025a, Reimann2026} recovers narrow posteriors only when one of the two formulations is enforced \textit{a priori}, suggesting that the field's apparent disagreement about VIP gain may reduce, in part, to a disagreement about which causal arrows in the connectivity graph should be treated as exogenous.

\section{Identifiability and empirical validation}

The unifying problem (Figure 24C) is that distinct circuits --- strong VIP\rightarrowSST with weak SST\rightarrowPyr, weak VIP\rightarrowSST with strong SST\rightarrowPyr, or 3-IN+NDNF --- can produce nearly identical pyramidal firing-rate traces under VIP optogenetic activation \citep{Garciadelmolino2017, Reimann2026, Vighagen2021}.  Identifiability analyses \citep{Dellal2025} make this point explicitly by fitting alternative parameterisations to the same VIP-perturbation dataset and showing that several recover the data with comparable likelihood, while \citet{Reimann2026} argue from a large-scale biophysical model that the addition of dendritic compartments and conductance-based synapses \textit{increases} rather than reduces the parameter manifold compatible with cell-type-resolved electrophysiology. Conceptually, this is a structural identifiability problem: the observable (pyramidal firing rate response) is a low-dimensional projection of a high-dimensional parameter vector, and the projection has a non-trivial null space. \citet{Litwinkumar2016} and \citet{Moreni2025a} argue that the VIP\rightarrowSST weight is identifiable from population data once SST\rightarrowPyr is independently constrained, but the §6 audit of paired-recording slopes shows that the available SST\rightarrowPyr connectivity values have wide confidence intervals --- wide enough that the VIP\rightarrowSST estimate inherits the uncertainty and inflates substantially, which the \citep{Litwinkumar2016} proof only partly mitigates. The identifiability claim is therefore conditional, not absolute, and rests on prior precision that paired-recording data do not yet supply.

The conflict has practical consequences for empirical design. Single-cell-type perturbations --- activating VIP alone, silencing SST alone --- populate exactly the directions in parameter space where the null space is largest, because they probe a single row of the connectivity matrix at a time. Constraint thus has to come from designs that pry the regimes apart: simultaneous PV, SST, and VIP perturbation with all-optical readout \citep{Veit2023, Lenkey2025}; compartment-resolved imaging during VIP manipulation \citep{Bilash2023, Pakan2016, Anastasiades2021a}; combined imaging+modeling that exploits learning-induced changes as additional constraints \citep{Poort2022, Lee2025}; and joint fitting against connectivity priors from molecular taxonomies \citep{Yao2023, Rudy2011}. Recent attempts in this direction \citep{Moreni2025a, Jiang2024b, Terwal2021, Huang2025a, Chou2021} couple inference frameworks to multi-cell perturbation data and recover narrower posterior regions, but the conflicts of Figure 24A persist within each posterior, and no single dataset has yet adjudicated between the \citep{Hertag2019} divisive and \citep{Tahvili2025b} subtractive accounts. Recent identifiability analyses extend this by arguing that even with multi-perturbation data, the structural identifiability of cell-type-resolved models requires explicit reporting of the full posterior \citep{Rademacher2025} over parameter combinations, not a single best fit; a recommendation that, if adopted, would force the modelling community to publish posteriors over conflicting regimes rather than choose one.

A second axis of identifiability concerns the validity of the model's behavioural readouts. Many of the most-cited perturbation results \citep{Pi2013, Fu2014, Lee2013, Krabbe2019, Myersjoseph2024} use behaviour-aggregated firing-rate changes as the primary observable, which collapses temporally rich responses into a single scalar and discards exactly the information that distinguishes gain-knob from credit-signal interpretations. \citet{Hertag2020}, \citet{Hertag2021}, and \citet{Wilmes2019} show that prediction-error and credit-assignment models produce their distinctive signatures in the \textit{temporal} structure of pyramidal responses --- phase-locked to mismatch events, ramping across trials with learning, or selectively gated by behavioural context --- and these are signatures that scalar readouts cannot resolve. Conversely, modelling and perturbation studies indicate that the gain-knob versus paradoxical-suppression distinction is most cleanly tested by perturbation-amplitude curves, in which the sign of the response to small versus large VIP perturbations diverges. Few empirical studies provide such curves, so the existing evidence underdetermines the model class even before parameter inference begins. The identifiability problem is therefore not only a problem of high-dimensional fitting; it is also a problem of designing experiments whose outputs occupy the directions in parameter space along which competing models actually differ \citep{Veit2023, Lenkey2025, Dellal2025, Reimann2026}.

The path forward is therefore not ``the next better model'' but coordinated empirical--computational design: pre-registered architectures whose distinguishing predictions are tested with cell-type-, layer-, and compartment-specific perturbations, and inference pipelines that report the full set of compatible parameterisations rather than a single best fit \citep{Dellal2025, Reimann2026, Sancheztodo2026, Moreni2025a}. The candidate roles for VIP --- gain knob, attentional gate, top-down teaching signal, oscillation organiser, context selector --- are not mutually exclusive \citep{Hertag2019, Garciadelmolino2017, Wagatsuma2023, Lee2025, Furutachi2024, Millman2020, Shen2023, Yang2016, Hertag2020, Hertag2021}, and may correspond to different operating regimes the same circuit visits during different behavioural states. The empirical signatures that would adjudicate the regimes are concrete: compartment-specific imaging during VIP perturbation to separate divisive from subtractive operation \citep{Bilash2023, Lenkey2025, Reimann2026}; SST-silenced VIP perturbations to test 4-IN necessity \citep{Iqbal2025, Hartung2024, Vighagen2021}; perturbation-amplitude curves to localise the ISN versus non-ISN regime \citep{Garciadelmolino2017, Sanzeni2020}; and trial-resolved, learning-coupled readouts to separate gain-knob from credit-signal interpretations \citep{Lee2025, Hertag2020, Hertag2021, Chevy2024}. Each of these designs probes a direction in parameter space that current single-perturbation datasets leave unconstrained, and each is now technically feasible. The next decade's progress depends less on adding cell types than on solving this identifiability problem with the cell types we already have, and on a discipline of model reporting that publishes the full set of parameterisations consistent with the data rather than a single, fragile best fit .

A direct consequence of these structural problems is that ``validation'' in the modelling literature has been operationalised heterogeneously, and several of the most prominent claims of empirical support are weaker than they first appear. \citet{Veit2023} presents spiking-circuit fits that recover the established VIP-perturbation phenomenology, but the fits use the same VIP-activation dataset that the model was conceptually motivated by, so that the agreement does not constrain the parameterisation against the alternative subtractive \citep{Tahvili2025a, Tahvili2025b} or paradoxical \citep{Garciadelmolino2017, Sanzeni2020} regimes. \citet{Krabbe2019} provide independent behavioural validation of VIP recruitment during fear conditioning, but their physiology is consistent with both the gain-knob \citep{Garciadelmolino2017, Hertag2019} and the credit-signal \citep{Wagatsuma2023, Lee2025, Hertag2020, Hertag2021} framings, because the experiment lacks the trial-resolved, learning-locked readouts that would separate them. \citet{Dipoppa2018} and \citet{Pakan2016} are similarly compatible with multiple architectures: their state-dependent VIP modulation could be implemented by 3-IN parameter retuning, by 4-IN NDNF channel recruitment \citep{Iqbal2025, Hartung2024, Vighagen2021}, or by an upstream cholinergic gating signal \citep{Fu2014, Anastasiades2021a}, and the data do not arbitrate. The pattern is that experimental confirmations of ``the disinhibitory motif'' are abundant \citep{Pi2013, Lee2013, Fu2014, Pfeffer2013, Karnani2016a, Pakan2016, Krabbe2019, Dipoppa2018, Myersjoseph2024, Veit2023}, while experimental confirmations of any \textit{specific parameterisation} of that motif are rare, and the claims that population-level fits identify the VIP\rightarrowSST weight \citep{Litwinkumar2016, Moreni2025a} rest on prior assumptions about SST\rightarrowPyr that paired-recording studies do not yet justify.

The picture is improved, but not closed, by the most recent generation of inference-coupled experimental designs. \citet{Moreni2025a} demonstrates that simulation-based inference applied to multi-cell perturbation recordings can recover narrower posterior regions in the VIP\rightarrowSST $\times$ SST\rightarrowPyr plane than single-perturbation fits, but the recovered posteriors still straddle the divisive/subtractive boundary in several datasets. \citet{Jiang2024b} and \citet{Huang2025a} couple perturbation experiments to spiking models and recover parameter ranges that are consistent across replicates within a behavioural paradigm but inconsistent across paradigms --- a pattern that itself argues for state-dependent operating regimes. \citet{Chou2021} similarly recover cross-cell-type parameter posteriors that exclude the most extreme regions of parameter space without selecting a single regime. Read together, these inference studies suggest that identifiability is improving but conditional: the field is converging on the \textit{boundaries} of the compatible parameter manifold faster than it is converging on the cortex's actual operating point within that manifold. The \citep{Litwinkumar2016} claim that VIP\rightarrowSST is identifiable from population data is therefore best interpreted as a local-identifiability claim under strong priors, not as a global-identifiability claim from data alone, and is consistent with the §6 audit observation that paired-recording weights have wide confidence intervals.

A final, often-overlooked aspect of validation concerns the integration of computational models with the molecular and connectomic data now becoming available. Cell-type-resolved transcriptomic atlases \citep{Yao2023, Rudy2011} provide independent constraints on which interneuron subclasses exist, their relative abundances, and their input/output specificities --- constraints that were not available when most of the rate models discussed here were first published. Connectomic reconstructions of cortical microcircuits \citep{Hua2022, Talapka2025} provide weight-distribution priors that can in principle pin down the SST\rightarrowPyr connectivity that \citep{Litwinkumar2016} argues is the limiting factor in VIP\rightarrowSST identifiability.  and \citet{Isbister2026} integrate these data sources into anatomically constrained large-scale models that exhibit \textit{in vivo}-like dynamics, but precisely because the resulting models inherit hundreds of free parameters, their predictions span the full divisive/subtractive/paradoxical phase space depending on which subset of parameters is fitted to data. The hope that anatomical detail will resolve the gain-regime conflicts is therefore only partly justified: more biophysical realism increases the \textit{capacity} of the model to fit any given dataset, which compounds the identifiability problem unless the additional parameters are independently constrained. Recent work in this direction \citep{Sancheztodo2026, Moreni2025a} is beginning to publish posterior-resolved fits, suggesting that the discipline is correctable; but the publication norm of reporting a single best-fit parameter set still dominates and continues to obscure the conflicts the field is attempting to resolve.

The implication for the §6 paired-recording audit and for any subsequent meta-analysis is straightforward. If the VIP\rightarrowSST weight is identifiable from population data only under strong priors on SST\rightarrowPyr \citep{Litwinkumar2016}, then the wide confidence intervals reported by the §6 audit are not noise to be averaged away but are propagated directly into the posterior over VIP\rightarrowSST. Any modelling claim that depends on a specific VIP\rightarrowSST value --- and most claims about gain regime, paradoxical operation, and credit-signal magnitude do --- should accordingly be re-stated as conditional on a stipulated SST\rightarrowPyr distribution. Until paired-recording datasets supply tighter priors, or until inference frameworks adopt the full-posterior reporting standard , the field's consensus on VIP function will remain a consensus on connectivity topology, not on circuit operation. Resolving the five conflicts catalogued in this section therefore reduces, in the end, to two empirical asks: tighter compartment-resolved data on SST\rightarrowPyr, and trial-resolved, perturbation-amplitude, multi-cell experimental designs whose outputs sample the directions of parameter space along which the competing models actually differ \citep{Veit2023, Lenkey2025, Bilash2023, Iqbal2025, Hartung2024, Vighagen2021, Sanzeni2020, Phillips2016, Hertag2020, Hertag2021, Wagatsuma2023, Lee2025, Chevy2024}.

\begin{figure}[!htbp]
\centering
\includegraphics[width=1\linewidth]{files/fig-computational-mo-fee8b88e3a65cedb87c7c1bf2f0f498c.png}
\caption[]{\textbf{Computational models of VIP circuit function.} \textbf{(A)} Three established architectures --- 3-IN rate/spiking circuits \citep{Hertag2019, Litwinkumar2016}, 4-IN models including NDNF \citep{Iqbal2025, Hartung2024}, and predictive-coding/credit-assignment circuits with VIP-mediated disinhibition in the error pathway \citep{Wagatsuma2023, Lee2025}. \textbf{(B)} Phase diagram of 11 representative models in (VIP\rightarrowSST weight) $\times$ (dendritic SST\rightarrowPyr nonlinearity), showing clustering into divisive, subtractive, paradoxical, and 4-IN/NDNF regimes. \textbf{(C)} Functional roles assigned to VIP across the modelling literature, partitioned into five overlapping categories. \textbf{(D)} VIP-inclusion audit. \textbf{Caveat (verbatim):} of 40 modelling rows in the evidence base, 25 were unambiguously classifiable as VIP-included or VIP-excluded; 15 rows were unclassifiable (review/theoretical articles, or insufficient detail) and are excluded from the count. Of the 25 classifiable models, 21 include VIP and 4 do not.}
\label{fig-computational-models}
\end{figure}

\begin{figure}[!htbp]
\centering
\includegraphics[width=1\linewidth]{files/fig-vip-model-empiri-f6cedb67351c2919bd990af78edcf2d0.png}
\caption[]{\textbf{Empirical conflicts mapped to model parameter regimes.} \textbf{(A)} Conflict pairs Hertäg 2019 vs Tahvili 2025b (3-IN divisive vs 2-IN subtractive) and Pi 2013 vs García del Molino 2017 (net disinhibition vs paradoxical inhibition) sit in distinct regions of the same parameter plane; NDNF-channel models \citep{Hartung2024} displace the operating point further. \textbf{(B)} Architectural-choice tree: which conflicts each architecture (2-IN, 3-IN, 4-IN; rate, spiking, predictive-coding) can in principle address. \textbf{(C)} Identifiability problem --- three different parameterisations of the same observable: pyramidal firing-rate response to VIP optogenetic activation can be reproduced by strong VIP\rightarrowSST + weak SST\rightarrowPyr, weak VIP\rightarrowSST + strong SST\rightarrowPyr, or a 3-IN+NDNF circuit, illustrating why VIP-perturbation experiments alone cannot constrain architecture \citep{Dellal2025, Reimann2026}.}
\label{fig-vip-model-empirical-mapping}
\end{figure}


\section{Conclusion}
\label{sec:13_conclusion}

Reading the eleven body sections together, the review establishes a position on VIP-expressing cortical interneurons that is more restrictive than the textbook reading and more structured than the catalogue of exceptions. At the subclass level, VIP-INs are a CGE-derived, 5-HT3AR+/Adarb2+ GABAergic population whose identity is reproduced across rodent, primate, and human single-cell atlases, recovered as one of four conserved cortical GABAergic subclasses \citep{Rudy2011, Tasic2018, Hodge2019, Bakken2021a, Lee2023a, Yao2023}. Within the subclass, the lineage decomposes into multiple t-types whose discreteness depends on dataset, taxonomic resolution, and integration method, and whose intrinsic, morphological, and connectivity profiles are graded along the bipolar/bitufted/multipolar and irregular-spiking/continuous-adapting axes that \href{/morphology}{Morphological Diversity} and \href{/electrophysiology}{Intrinsic Electrophysiology} chart \citep{Tasic2016, Tasic2018, Gouwens2020a, Pronneke2015, Cauli1997, Cauli2014, Apicella2022, Schneidermizell2025}. The disinhibitory reading of VIP function --- VIP cells, recruited by behavioural state, suppress SST cells and so release pyramidal dendrites from inhibition --- is supported as a population-average operating mode in primary visual, auditory, somatosensory, prefrontal, and cingulate cortex \citep{Pi2013, Lee2013, Pfeffer2013, Fu2014, Cichon2017, Williams2019, Veit2023}, but it is one operating mode among several, and the literature provides equally well-documented preparations in which VIP recruitment activates SST cells, inverts pyramidal modulation, or fails to modulate pyramidal output at all \citep{Pakan2016, Dipoppa2018, Garrett2020, Yavorska2021, Bigelow2019, Anastasiades2021a, Bastos2023a}.

The five cross-cutting tensions previewed in \href{/}{Introduction} admit, on this reading, partial resolutions rather than dissolutions. The discrete-versus-continuous question for VIP t-types is best read as resolution-dependent: supervised clustering on transcriptomic data recovers nominally discrete groups, but cross-modal Patch-seq integration, morphological metrics, and intrinsic-property continua reveal graded boundaries between many of those groups, with continuity predominant within the \textit{VIP Lamp5}, \textit{VIP Sncg}, and irregular-spiking branches \citep{Tasic2018, Gouwens2020a, Pronneke2015, Apicella2022, Emmenegger2018, Schneidermizell2025}. The subtractive-versus-divisive gain ambiguity reduces to operating regime: VIP perturbation is divisive when SST recruitment is concurrent and dendritic inhibition is dominant, subtractive when the disinhibitory release falls on perisomatic compartments, and approximately mixed when both are engaged, with the rate-model literature reproducing each of these outcomes from the same wiring under different parameter draws \citep{Kuchibhotla2017, Hertag2019, Litwinkumar2016, Garciadelmolino2017, Veit2023, Tahvili2025a, Tahvili2025b}. The net-disinhibition versus paradoxical-effect conflict reduces to whether the local circuit is in the inhibition-stabilised regime: outside it, VIP activation produces the textbook disinhibitory gain; inside it, the same wiring produces paradoxical pyramidal suppression on VIP activation, and the same connectomic prior is consistent with both behaviours \citep{Garciadelmolino2017, Sanzeni2020, Reimann2026, Hertag2019, Tahvili2025a, Tahvili2025b}.

The two remaining tensions are less tractable from the present evidence. The identifiability of the VIP\rightarrowSST weight in four-population rate equations is, on the present data, an unresolved bottleneck: VIP\rightarrowSST and VIP\rightarrowPV cross-couplings trade off against SST\rightarrowPyr and PV\rightarrowPyr terms in a way that connectomic priors do not constrain, and several published models with incompatible VIP\rightarrowSST weights fit the same in-vivo locomotion, attention, and reward data \citep{Hertag2019, Litwinkumar2016, Garciadelmolino2017, Veit2023, Reimann2026, Sabri2024}. \href{/computational-models}{Computational Models of VIP Circuit Function} argues that this is the field's principal next-decade problem and that simultaneous SST/PV/VIP/Pyr recordings under structured perturbations --- paired with electron-microscopy connectomes that report not only contacts but also synapse weights --- are the minimum required to break the degeneracy \citep{Reimann2026, Schneidermizell2025, Hertag2019, Bos2025}. Cross-species translation runs along an analogous axis: VIP-IN identity is conserved at the subclass level between rodent and primate cortex, but the human VIP subclass expands to twenty-one t-types in middle temporal gyrus, includes upper-layer enrichment and human-specific \textit{VIP PCDH20} and \textit{MC4R} morphologies absent in mouse, and over-represents in disease-associated programmes that the rodent atlases do not capture \citep{Hodge2019, Boldog2018, Chartrand2023, Bakken2021a, Bakken2021b, Krienen2020, Lee2023a}.

These conflicts converge on a single integrated reading. VIP function is best treated not as a fixed circuit operation but as the output of an area$\times$state$\times$modulator triple that \href{/in-vivo-behavior}{In Vivo Function During Behavior} and \href{/cross-areas}{VIP Interneurons Across Brain Regions} describe in cortical terms and that \href{/oscillatory-dynamics}{Oscillatory Dynamics and Temporal Coordination} extends into the temporal domain via gamma gain modulation, theta-locked IS-3 firing, Up-state biasing, and predictive-coding/attention frames \citep{Veit2023, Cichon2017, Hertag2020, Hertag2021, Shipp2016, Francavilla2018a, Tyan2014, Tricoire2011, Gulyas1996}. Within that triple, primary visual cortex VIP cells implement contrast-dependent, locomotion-gated gain enhancement \citep{Fu2014, Dipoppa2018, Pakan2016, Millman2020}; auditory cortex deploys them for reinforcement-linked disinhibition and critical-period plasticity, with a polarity puzzle around movement-driven SST activation \citep{Pi2013, Yavorska2021, Bigelow2019}; somatosensory cortex recruits them via motor copy during active sensing \citep{Lee2013, Sermet2019}; medial prefrontal and anterior cingulate cortex implement top-down control of attention, working memory, and pain \citep{Bastos2023a, Anastasiades2021a, Kamigaki2017, Li2022a}; and hippocampal IS-3 cells assemble a circuit motif that targets other interneurons rather than principal cells, breaking the cortical disinhibitory assumption at the wiring level \citep{Gulyas1996, Tyan2014, Tricoire2011, Kawaguchi1996, Francavilla2018a}. The same triple maps onto the disease literature reviewed in \href{/disease-translational}{Species Differences, Human Relevance, and Disease}: VIP-IN dysfunction has been documented in autism, Rett, Dravet, schizophrenia, and Alzheimer-related models, with the sign of the effect depending on perturbation polarity, brain area, and disease model rather than on a uniform vulnerability of the subclass \citep{Mossner2020, Goff2023, Mcfarlan2024a, Bhandari2024, Hall2015, Goel2025, Kranz2025, Goff2019}.

What the corpus does not yet support is also identifiable. It does not support a single rate or spiking model that reproduces all in-vivo VIP recordings with one set of cross-coupling weights \citep{Hertag2019, Garciadelmolino2017, Veit2023, Bos2025, Reimann2026}. It does not support an unambiguous mapping from t-type to behavioural function: most in-vivo recordings target VIP cells via Cre-driver labelling that pools multiple t-types, so the t-type composition of the imaged or photo-stimulated population is rarely specified in the studies that motivate the disinhibitory reading \citep{Pi2013, Fu2014, Pakan2016, Dipoppa2018, Pronneke2015, Schneidermizell2025}. It does not yet support direct transfer of mouse circuit conclusions to human cortex, given the primate expansion of CGE-derived diversity and the human-specific morphologies in upper layers \citep{Hodge2019, Boldog2018, Chartrand2023, Bakken2021a, Krienen2020}. And it does not yet support a settled account of how VIP-mediated disinhibition implements credit assignment or predictive coding: the disinhibitory motif has been recruited as the substrate for both, but the empirical tests that would discriminate the two readings have not been performed \citep{Hertag2019, Hertag2020, Hertag2021, Wagatsuma2023, Lee2025, Rossbroich2025, Nemati2025a, Nemati2025b, Sabri2024, Shipp2016}.

The forward agenda the review supports follows from these limits rather than from the strengths. Three lines of work would discriminate the unresolved cases. First, simultaneous multi-class recordings (VIP, SST, PV, Pyr) under structured perturbations, paired with synapse-weight-resolved connectomes, would constrain the VIP\rightarrowSST weight and break the rate-equation degeneracy that \href{/computational-models}{Computational Models of VIP Circuit Function} exposes \citep{Schneidermizell2025, Reimann2026, Hertag2019, Bos2025, Sabri2024}. Second, t-type-resolved in-vivo recordings --- Patch-seq-tagged or driver-line-restricted to single VIP t-types rather than to the pooled subclass --- would allow the area$\times$state$\times$modulator triple to be decomposed by t-type and would test whether the gain regime, polarity, and target compartment of VIP modulation are t-type properties or only area properties \citep{Gouwens2020a, Pronneke2015, Apicella2022, Schneidermizell2025, Tasic2018}. Third, comparative recordings in primate and human ex-vivo and in-vivo preparations --- leveraging the cross-species t-type alignments already in hand --- would test whether the rodent disinhibitory reading transfers and how the human-specific VIP types behave functionally \citep{Hodge2019, Boldog2018, Chartrand2023, Bakken2021a, Krienen2020, Lee2023a}. Until those discriminating experiments are performed, the calibrated map this review provides is the safest intermediate position: the disinhibitory motif holds as a population-average operating mode under specifiable area, state, and modulator conditions; it fails or runs in reverse under specifiable other conditions; the wiring is reproducible at the subclass average but ambiguous at the parameter level; and the cross-species transfer is partial, with disease-relevant divergences clustered in upper-layer human-specific VIP types \citep{Pi2013, Pakan2016, Dipoppa2018, Garciadelmolino2017, Reimann2026, Hodge2019, Chartrand2023, Mcfarlan2024a}.


\section{Methods}
\label{sec:Methods}

\begin{figure}[h]
\centering
\begin{verbatim}
{
  "section": "methods -overview",
  "body": "Curated literature corpus for the computational review pipeline."
}
\end{verbatim}
\caption*{methods-overview}
\end{figure}

This Methods section documents how the comprehensive critical review of vasoactive-intestinal-peptide-expressing (VIP) cortical interneurons was produced. It is a transparency artifact, written from pipeline metadata and gate artifacts only; no post-hoc characterizations have been added. Eight subsections cover the original review request (M.1), the scope and table of contents derived from it (M.2), the evidence corpus and curation (M.3), the citation infrastructure (M.4), the drafting and critic protocol (M.5), a phase-by-phase pipeline ledger (M.6), the bibliography and integration record (M.7), and finally limitations, reproducibility, and data availability (M.8). The pipeline executed under coordinator skill \texttt{comprev-coordinator-v27.md} with two delegated agent roles: LITREVIEW (scientific judgment) and DATAML (mechanical work).

\begin{quote}
\textbf{Snapshot.} This Methods file is a Phase-13 draft snapshot covering Phases 1 through 13. The Pipeline Ledger (M.6) ends with a \texttt{Phases 14-20 (pending refresh)} placeholder row that will be replaced at Phase 20a (Methods Ledger Refresh) once the document-assembly, citation-triples, verification, fix-preparation, fix-application, final-build, and repository-push phases have executed.
\end{quote}

\section{Review Request}\label{sec-methods-review-request}

This review was initiated by a single user prompt that fixed the topic, scope, two-agent architecture, evidence-saturation parameters, and high-level table of contents. Phase 1 (Scoping) translated this prompt into the structured \texttt{gate\_scope.json} (title, audience, target paper count, cluster definitions, and 13-section table of contents) that drove all subsequent phases. The full text of the request is preserved verbatim in \texttt{provenance/review\_request.md} and reproduced below.

\begin{verbatim}
Start a comprehensive critical literature review titled:
"VIP Interneurons in Cortical Computation: From Molecular Identity to Circuit Function"

This is a v2 of an existing review (https://github.com/AllenNeuralDynamics/ComputationalReviewVIP).
The v1 cited 405 papers across 12 body sections. The v2 goal is to saturate the VIP interneuron
literature - - cover every relevant primary research paper, plus the adjacent SST, PV, disinhibition,
and CGE -lineage literature needed for context. Database sizing estimates the total relevant
envelope at 2,500 -3,000 unique papers.

Initialize this repo fresh from the template:
https://github.com/AllenNeuralDynamics/ComputationalReviewTemplate
The template ships all pipeline skills (v27), plugins, GitHub Actions, and site infrastructure.

The pipeline uses two agents - - LITREVIEW (scientific judgment) and DATAML (mechanical work).

The three key skill files:
 - skills/comprev -coordinator -v27.md - - The coordinator protocol. Read this FIRST. It defines
  all 20 phases, the coordinator protocol, gate artifacts, evidence parameters, and the plan
  structure. Follow it phase by phase.
 - skills/comprev -reviewer -agent.md - - The worker skill for LITREVIEW agents. Pass this to every
  LITREVIEW delegation so the agent can load it.
 - skills/comprev -figure -construction.md - - Already published as a skill on LITREVIEW agents.
  Section writers load it for figure production.

GitHub Repository: https://github.com/AllenNeuralDynamics/ComputationalReviewVIP
Push all outputs to this repo in Phase 20.

Evidence parameters:
 - Saturate the literature: search until <2% new unique papers in the last 100 records, confirmed
  across two consecutive database passes
 - Snowball 2 rounds (forward + backward citation chasing on top -50 most -cited per cluster)
 - No fixed per -cluster floor - - let the saturation criterion drive depth per cluster (some clusters
  are naturally smaller than others; forcing a uniform floor would pad thin clusters with marginal
  papers)
 - Total bibliography target: >=2,000 unique papers

Table of Contents:
1.  Introduction
2.  Molecular Identity and Transcriptomic Taxonomy
3.  Developmental Origins and Postnatal Maturation
4.  Morphological Diversity
5.  Intrinsic Electrophysiology
6.  Synaptic Properties and Connectivity
7.  Local Circuit Motifs and the Disinhibition Framework
8.  In Vivo Function During Behavior
9.  VIP Interneurons Across Brain Regions
10. Oscillatory Dynamics and Temporal Coordination
11. Species Differences, Human Relevance, and Disease
12. Computational Models of VIP Circuit Function
13. Synthesis and Conclusion: Reassessing the Canonical VIP Disinhibitor
\end{verbatim}

\section{Scope and Table of Contents}\label{sec-methods-scope}

Phase 1 produced the \texttt{gate\_scope.json} artifact (vid \texttt{gates/gate_scope.json}), which froze the editorial scope before any literature work began. The thesis question asks how VIP-expressing GABAergic interneurons -- defined by caudal-ganglionic-eminence (CGE) lineage and a distinctive transcriptomic, morphological, and physiological profile -- implement disinhibitory and gain-modulatory operations that shape cortical computation, and to what extent the canonical ``VIP disinhibitor'' framing survives scrutiny under contemporary multi-modal evidence. The intended audience is systems and circuits neuroscientists, computational neuroscientists, and graduate-level researchers working on cortical interneurons.

The review is organized into 13 sections (1 introduction + 11 body sections + 1 synthesis/conclusion) with a total target of approximately 63,000 words across all sections. Per-section targets and topic clusters are listed below.

\bigskip\noindent
\begin{tabular}{p{\dimexpr 0.250\linewidth-2\tabcolsep}p{\dimexpr 0.250\linewidth-2\tabcolsep}p{\dimexpr 0.250\linewidth-2\tabcolsep}p{\dimexpr 0.250\linewidth-2\tabcolsep}}
\toprule
Section & Title & Target words & Source clusters \\
\hline
01 & Introduction & 2,500 & cluster\_06\_disinhibition\_circuit, cluster\_01\_molecular\_taxonomy, cluster\_12\_sst\_pv\_context \\
02 & Molecular Identity and Transcriptomic Taxonomy & 5,500 & cluster\_01\_molecular\_taxonomy, cluster\_12\_sst\_pv\_context \\
03 & Developmental Origins and Postnatal Maturation & 4,500 & cluster\_02\_development\_lineage \\
04 & Morphological Diversity & 4,000 & cluster\_03\_morphology \\
05 & Intrinsic Electrophysiology & 4,000 & cluster\_04\_intrinsic\_electrophysiology \\
06 & Synaptic Properties and Connectivity & 5,500 & cluster\_05\_synaptic\_connectivity, cluster\_12\_sst\_pv\_context \\
07 & Local Circuit Motifs and the Disinhibition Framework & 6,000 & cluster\_06\_disinhibition\_circuit, cluster\_12\_sst\_pv\_context \\
08 & In Vivo Function During Behavior & 6,500 & cluster\_07\_in\_vivo\_function, cluster\_13\_neuromodulation \\
09 & VIP Interneurons Across Brain Regions & 5,000 & cluster\_08\_brain\_regions \\
10 & Oscillatory Dynamics and Temporal Coordination & 4,500 & cluster\_09\_oscillations, cluster\_13\_neuromodulation \\
11 & Species Differences, Human Relevance, and Disease & 5,500 & cluster\_10\_species\_disease, cluster\_01\_molecular\_taxonomy \\
12 & Computational Models of VIP Circuit Function & 5,000 & cluster\_11\_computational\_models, cluster\_06\_disinhibition\_circuit \\
13 & Synthesis and Conclusion: Reassessing the Canonical VIP Disinhibitor & 4,500 & cluster\_06\_disinhibition\_circuit, cluster\_07\_in\_vivo\_function, cluster\_11\_computational\_models \\
\bottomrule
\end{tabular}

\bigskip

Evidence parameters specified in the scope: search until \textless 2\% new unique papers in the last 100 records confirmed across two consecutive database passes; 2 rounds of forward+backward snowball citation chasing on the top-50 most-cited papers per cluster; no fixed per-cluster paper floor (saturation drives depth); total bibliography target \textgreater =2,000 unique papers. Fourteen evidence clusters (with one cluster -- \texttt{cluster\_13\_cge\_lineage\_context} -- added by the Phase-2 LITREVIEW frame as an adjacent-context cluster) covered the topic envelope, mapping onto the 13 review sections as shown above.

\section{Evidence Corpus and Curation}\label{sec-methods-evidence}

Phase 2 (LITREVIEW evidence-gathering) ran 14 parallel cluster delegations across the databases authorized in the scope (PubMed/NCBI E-utilities, Europe PMC, OpenAlex, bioRxiv) plus open-access full-text resolvers (Unpaywall, Semantic Scholar, PubMed Central) and the Elsevier and Springer Nature publisher APIs. Search queries, deduplication, and snowball passes are recorded per cluster inside the evidence JSONs referenced from \texttt{gate\_evidence\_compliance.json} (vid \texttt{gates/gate_evidence_compliance.json}). Aggregate Phase-2 outcome:

\bigskip\noindent
\begin{tabular}{p{\dimexpr 0.500\linewidth-2\tabcolsep}p{\dimexpr 0.500\linewidth-2\tabcolsep}}
\toprule
Metric & Value \\
\hline
Unique DOIs across all clusters & 1,037 \\
Total findings extracted & 1,779 \\
Total conflicts catalogued & 84 \\
Overall full-text rate & 51.83\% \\
Overall finding-level compliance rate & 99.94\% \\
Per-cluster gates passing & 14 / 14 \\
\bottomrule
\end{tabular}

\bigskip

The Phase-2V validator (\texttt{comprev-evidence-validator}, DATAML) re-checked compliance over 6 iterations against ten mechanical schemas (DOI resolves, source not equal to title, source length, claim differs from source, has cite-key, source provenance, full-text honesty, no fabricated keys, deduplication, framework-extract size). The strict bibliography-target check (\textgreater =2,000 papers per scope) failed at iteration 6 with 1,037 unique DOIs returned, but the finding-level pipeline (1,779 findings; 0.9994 compliance) passed all integrity checks. The coordinator accepted this saturation outcome as the final corpus on the basis that the Phase-2 saturation criterion (\textless 2\% new unique in last 100 records, two consecutive passes) had triggered before the 2,000-paper target was reached -- i.e. the addressable VIP literature is empirically smaller than the database-sizing pre-estimate.

Per-cluster curation totals (Phase 5 ``Evidence Curation'' actor; gate vid \texttt{gates/gate_evidence_curated.json}; validator gate vid \texttt{gates/gate_evidence_curated_approved.json}):

\bigskip\noindent
\begin{tabular}{p{\dimexpr 0.333\linewidth-2\tabcolsep}p{\dimexpr 0.333\linewidth-2\tabcolsep}p{\dimexpr 0.333\linewidth-2\tabcolsep}}
\toprule
Cluster & Papers & Findings (Phase 2) \\
\hline
cluster\_01\_molecular\_taxonomy & 208 & 218 \\
cluster\_02\_development\_lineage & 128 & 160 \\
cluster\_03\_morphology & 33 & 58 \\
cluster\_04\_intrinsic\_electrophysiology & 51 & 93 \\
cluster\_05\_synaptic\_properties\_connectivity & 335 & 335 \\
cluster\_06\_disinhibition\_circuit & 91 & 91 \\
cluster\_07\_in\_vivo\_function & 130 & 130 \\
cluster\_08\_brain\_regions & 123 & 153 \\
cluster\_09\_oscillations & 80 & 80 \\
cluster\_10\_species\_disease & 100 & 100 \\
cluster\_11\_computational\_models & 154 & 154 \\
cluster\_12\_sst\_pv\_context & 44 & 44 \\
cluster\_13\_neuromodulation & 91 & 91 \\
cluster\_13\_cge\_lineage\_context & 72 & 72 \\
\bottomrule
\end{tabular}

\bigskip

Phase 5 expanded the 1,779 Phase-2 findings into 2,726 section-curated findings while preserving zero loss against the Phase-2 baseline (every Phase-2 finding is reachable in at least one section curated literature corpus). Coverage of unique Phase-2 DOIs reached 99.8\% (1,047 / 1,049). The validator confirmed:

\begin{itemize}
\item \texttt{NO\_INTRA\_SECTION\_DUPLICATES} (PASS): no duplicate finding/source pair within any section.
\item \texttt{CROSS\_SECTION\_DIFFERENTIATION} (PASS): 0 avoidable duplications; 272 unavoidable single-citation-source / multi-cluster cases accepted as documented exit.
\item \texttt{ZERO\_LOSS} (PASS): 2,726 curated \textgreater = 1,779 input.
\item \texttt{ALL\_CONFLICTS\_ASSIGNED} (PASS): 51 / 51 normalized conflicts assigned to sections.
\item \texttt{COVERAGE\_75} (PASS): 99.8\% of Phase-2 unique DOIs cited in at least one section.
\item \texttt{CITE\_KEY\_ASSIGNED} (PASS): 0 missing keys; 192 of 194 originally synthesized keys resolved via CrossRef and arXiv; 2 fallback keys retained.
\item Two documented exits: PAPER\_TIER\_FLOOR for sections 4 and 9 (driven by Phase-2 saturation in those clusters), and CROSS\_SECTION\_DIFFERENTIATION (272 single-CSS multi-cluster findings).
\end{itemize}

Anti-compression metrics: per the Phase-5V audit, every section retains \textgreater =75\% of its source-cluster findings rather than collapsing them to one-claim-per-paper summaries; this is the operational safeguard against narrative compression that would conceal heterogeneity within the literature.

\section{Citation Infrastructure}\label{sec-methods-citation-infra}

Phase 3 (DATAML actor) constructed the project-wide citation key map and author name table that all downstream phases consume as the single source of authorial truth. Inputs were the 1,037 unique DOIs surviving Phase 2; outputs were \texttt{citation\_key\_map.json} (vid \texttt{provenance/citation_key_map.json}) and \texttt{author\_name\_table.json} (vid \texttt{provenance/author_name_table.json}).

\bigskip\noindent
\begin{tabular}{p{\dimexpr 0.500\linewidth-2\tabcolsep}p{\dimexpr 0.500\linewidth-2\tabcolsep}}
\toprule
Metric & Value \\
\hline
DOIs input & 1,037 \\
DOIs mapped to a cite-key & 1,037 \\
CrossRef resolutions (HTTP 200) & 1,037 / 1,037 \\
Europe PMC fallbacks needed & 0 \\
CrossRef failures & 0 \\
Cite-key collisions disambiguated & 106 \\
ASCII-folding events on author surnames & 51 \\
\bottomrule
\end{tabular}

\bigskip

Cite-key format is \texttt{\^[A-Z][A-Za-z]+{\textbackslash}d\{4\}[a-z]?\$} (first-author surname, ASCII-folded; four-digit year; optional disambiguator suffix). Phase 3V (DATAML validator) re-resolved every entry against CrossRef and ran five mechanical checks. The validation gate (vid \texttt{gates/gate_citation_infrastructure_validation_v2.json}) reports a per-entry pass rate of 99.81\% with the following breakdown: 1,037 / 1,037 CrossRef-resolves, 1,037 / 1,037 key-format pass, 1,037 / 1,037 key-uniqueness pass, 1,035 / 1,037 author-match pass (the two failures were both resolved-by-name-match consortium-author entries -- the Petilla Interneuron Nomenclature Group \texttt{Thepetillainterneuronnomenclaturegroupping2008} and the BRAIN Initiative Cell Census Network \texttt{Braininitiativecellcensusnetworkbiccn2021} -- which carry no CrossRef family field), and 1,037 / 1,037 year-match pass. Cardinality and framework-bidirectional checks passed; zero blocking failures.

The keymap is the contract that Phase 5 (evidence curation), Phase 7 (drafting), Phase 9 (bibliography), Phase 10 (integration), and Phase 11 (introduction/conclusion) all consume read-only. Section drafts are constrained by a per-section \texttt{cite\_keys\_allowlist} derived from this map.

\section{Drafting and Critic Protocol}\label{sec-methods-drafting}

Phase 4 (LITREVIEW) approved a 13-section framework (vid \texttt{provenance/scaffold.json}; gate vid \texttt{gates/gate_scaffold_approved.json}) covering 26 figure slots, 6 recurring themes, and 60 conflicts surfaced for explicit treatment. Every section carries a thesis sentence, a ``connection-prev'' / ``connection-next'' pair, and a unique cross-reference label.

Phase 6 (LITREVIEW critic, blinded ``Figure Audit'') converged in 3 iterations from 36 panels to 34 final panels (8 PASS, 26 CAVEAT, 0 REDESIGN, 0 FABRICATION). Twenty-two fabrication flags were raised across iterations 1-3 and all were resolved before convergence. Gate vid \texttt{gates/gate_figure_audit.json}.

Phase 7 (LITREVIEW section-writing) ran 11 parallel section actors. The DATAML send-back validator enforced eight mechanical checks per section (anchor present, every cite-key in the project allowlist, figure-fence syntax, figure has image path, no bare author names in prose, no entries from the project forbidden-lexicon list, no \texttt{\{label\}} directive duplicating the heading anchor, no duplicate display name across sections). The send-back loop ran for 3 iterations:

\bigskip\noindent
\begin{tabular}{p{\dimexpr 0.333\linewidth-2\tabcolsep}p{\dimexpr 0.333\linewidth-2\tabcolsep}p{\dimexpr 0.333\linewidth-2\tabcolsep}}
\toprule
Iteration & Total mechanical failures & Sections passing \\
\hline
1 & 211 & 0 / 11 \\
2 & 4 & 9 / 11 \\
3 & 0 & 11 / 11 \\
\bottomrule
\end{tabular}

\bigskip

Final Phase-7V gate vid: \texttt{gates/gate_sections_drafted_iter3.json}.

Phase 8 (LITREVIEW critic, ``Section Critic'') ran a separate blinded critic pass against each Phase-7 section draft, raising MUST\_FIX, SHOULD\_CAVEAT, and NIT severities under an enforced rule: \textbf{MUST\_FIX must reach zero before the critic loop can close}. Iteration 1 produced 22 MUST\_FIX items distributed across 8 sections (sec\_02, sec\_03, sec\_05, sec\_06, sec\_07, sec\_08, sec\_11, sec\_12); iteration 2 returned 0 MUST\_FIX across all 11 sections, with 49 SHOULD\_CAVEAT items carried forward as non-blocking polish. Sec\_08 invoked the ``PRESENT-AS-GAP'' rule for 4 framework-required cite-keys whose papers turned out to lack VIP-specific findings in the Phase-2 corpus: rather than fabricate evidence, those claims were rewritten as explicit research gaps. Final Phase-8 gate vid: \texttt{gates/gate_critic_complete.json}.

The actor / critic separation enforced throughout Phases 6, 7, and 8 is an information-barrier protocol: critic frames are never given the actor's drafting context, and actors are never shown raw critic deliberations -- only structured MUST\_FIX / SHOULD\_CAVEAT items routed through the coordinator. This actor-critic separation is the principal mechanism by which the pipeline avoids the rubber-stamping failure mode of single-pass LLM drafting.

\section{Pipeline Ledger}\label{sec-methods-pipeline}

The table below records the actor / critic / validator role for each phase, the gate artifact version that closed the phase, the validation result, and a one-line outcome note. All version IDs are artifact identifiers and resolve via the project artifact store. This table is a Phase-13 draft snapshot; the final row is a placeholder for Phases 14 to 20, which Phase 20a (Methods Ledger Refresh) will replace with executed-phase outcomes.

\bigskip\noindent
\begin{tabular}{p{\dimexpr 0.200\linewidth-2\tabcolsep}p{\dimexpr 0.200\linewidth-2\tabcolsep}p{\dimexpr 0.200\linewidth-2\tabcolsep}p{\dimexpr 0.200\linewidth-2\tabcolsep}p{\dimexpr 0.200\linewidth-2\tabcolsep}}
\toprule
Phase & Role / Agent & Gate VID & Status & Outcome notes \\
\hline
1 & actor (LITREVIEW + DATAML) & \texttt{gates/gate_scope.json} & PASS & scope.json + provenance/review\_request.\{md,txt\}; 13 body sections, 14 evidence clusters, target \textgreater =2000 papers, frame 4ee0a0e5 (LITREVIEW scoping) -\textgreater  266fb089 (DATAML materialisation) \\
1V & validator (DATAML) & \texttt{gates/gate_scope.json} & PASS & Phase-1 validation rolled into gate\_scope.json; frame e79b4741 \\
2 & actor (LITREVIEW x14) & \texttt{gates/gate_evidence_compliance.json} & PASS & Per-cluster evidence JSONs; 1037 unique DOIs, 1779 findings, 84 conflicts; saturation logged across 14 clusters; per-cluster gate 14/14 \\
2V & validator (DATAML) & \texttt{gates/gate_evidence_compliance.json} & PASS & Iter6 strict-fail on bibliography target only (1037\textless 2000); coordinator accepted saturation outcome at 1037 unique DOIs; 14/14 per-cluster gates PASS, compliance 0.9994; frames 3a462e11, ba58bb1b \\
3 & actor (DATAML) & \texttt{provenance/citation_key_map.json} & PASS & citation\_key\_map (1037 keys; 0 collisions; 106 disambiguations; 0 CrossRef failures) + author\_name\_table.json (cdd3608c) \\
3V & validator (DATAML) & \texttt{gates/gate_citation_infrastructure_validation_v2.json} & PASS & Per-entry pass rate 0.998; 1037/1037 CrossRef-resolved; 0 collisions; 2 consortium-author no-family entries reconciled by name match \\
4 & actor (LITREVIEW) & \texttt{gates/gate_scaffold_approved.json} & PASS & Framework (vid 01bea2ad) approved: 13 sections, 26 figure slots, 60 conflicts\_to\_present, 6 recurring themes, total target 63000 words \\
5 & actor (DATAML) & \texttt{gates/gate_evidence_curated.json} & PASS & 13 curated literature corpora + framework extracts + section citation maps + section author tables; 2726 findings curated (zero loss vs 1779 input); 51 conflicts assigned \\
5V & validator (DATAML) & \texttt{gates/gate_evidence_curated_approved.json} & PASS & PASS\_WITH\_DOCUMENTED\_EXITS: coverage 99.8\% of Phase-2 unique DOIs; 0 intra-section duplicates; cross-section differentiation 0 avoidable; 2 documented exits (sec\_04/sec\_09 below tier floor due to upstream saturation; 272 unavoidable single-CSS multi-cluster cases) \\
6 & critic (LITREVIEW) & \texttt{gates/gate_figure_audit.json} & PASS & Blinded figure audit, 3 iterations: 36 panels -\textgreater  34 final (8 PASS, 26 CAVEAT, 0 REDESIGN, 0 FABRICATION); 22 fabrication flags raised then resolved before convergence \\
7 & actor (LITREVIEW x11) & \texttt{gates/gate_sections_drafted_iter3.json} & PASS & Iter3 of section drafting: 11/11 sections PASS across 8 mechanical checks (anchor, cite-key allowlist, figure fences, no bare authors, forbidden lexicon, label directive, duplicate name) \\
7V & validator (DATAML) & \texttt{gates/gate_sections_drafted_iter3.json} & PASS & Send-back loop: iter1=211 failures -\textgreater  iter2=4 -\textgreater  iter3=0; 1037 valid cite-keys loaded for allowlist check \\
8 & critic (LITREVIEW x11) & \texttt{gates/gate_critic_complete.json} & PASS & Iter1=22 MUST\_FIX across 8 sections; iter2=0 MUST\_FIX across all 11 sections; 49 SHOULD\_CAVEAT items carried forward as non-blocking polish; sec\_08 PRESENT-AS-GAP rule applied for 4 corpus-missing framework cite-keys \\
9 & actor (DATAML) & \texttt{content/references.bib} & PASS & references.bib: 1037 entries; 413 used keys; 0 syntax errors; 0 duplicate cite-keys; 2 cosmetic ASCII-fold encoding patches (Argunsah2024, Calin2023) \\
9V & validator (DATAML) & \texttt{gates/gate_bibliography.json} & PASS & 12 checks: balanced braces, key uniqueness, key-keymap bijection, DOI-keymap match, 1037/1037 CrossRef HTTP-200, author/title/year cross-check, ASCII-clean, used-keys-in-bib for all 11 sections, no template contamination; 2 minor diacritic encoding notes (non-blocking) \\
10 & actor (LITREVIEW) & \texttt{gates/gate_integration.json} & PASS & 6 cross-section integration passes (6a transitions, 6b crossref, 6c terminology, 6d continuity audit, 6e figure inventory, 6f hygiene); 10/11 sections modified, 54 edits; cite-key set (413 unique) preserved; 0 broken refs; 22 figures inventoried \\
11 & actor (LITREVIEW) & \texttt{gates/gate_intro_conclusion.json} & PASS & Introduction (1018 words, 76 cite-keys) + Conclusion (1079 words, 71 cite-keys); 0 new cite-keys (all keys drawn from 413-key body cite-set); all paragraphs \textgreater =4 cite-keys; 21 \{numref\}/\{ref\} cross-references resolved \\
12 & critic (LITREVIEW) & \texttt{gates/gate_bookend_critic.json} & PASS & Bookend critic; iter1=1 MUST\_FIX (forbidden lexicon `canonical' in intro) + 5 SHOULD\_CAVEAT; iter2=0 MUST\_FIX after single-word patch; intro v2 vid 8c2df2b9; conclusion vid d2ae9ef1; 5 SHOULD\_CAVEAT carried forward \\
13 & actor (DATAML) & \texttt{\textless this gate\textgreater } & PASS & Methods.md transparency document with M.1 Review Request, M.2 Scope and TOC, M.3 Evidence corpus and curation, M.4 Citation infrastructure, M.5 Drafting and critic protocol, M.6 Pipeline ledger, M.7 Bibliography and integration, M.8 Limitations, reproducibility and data availability \\
14-20 & Pending - to be populated at Phase 20a & \texttt{-} & Pending & Phase 14 (document assembly), Phase 15 (citation triples), Phase 16 (verification), Phase 17 (fix preparation), Phase 18 (LITREVIEW fix application), Phase 19 (fix application + final myst build), Phase 20 (repository push). Phase 20a (Methods Ledger Refresh) replaces this row with executed-phase outcomes and refreshes the Phase-13 snapshot. \\
\bottomrule
\end{tabular}

\bigskip

\textbf{Phases 14-20 (pending refresh).} The Phase-13 numbers above are provisional placeholders for the downstream phases; Phase 20a re-renders this ledger after Phases 14 (document assembly), 15 (citation triples), 16 (citation verification: VERIFIED / MINOR / CHIMERIC / HALLUCINATED / MISATTRIBUTED counts), 17 (fix preparation), 18 (LITREVIEW fix application), 19 (final fix application + myst build), and 20 (repository push) have executed and produced their gate artifacts. Until that refresh runs, treat the Phases 14-20 row as \texttt{Pending - to be populated at Phase 20a}. Forbidden stale phrasings (``are scheduled'', ``had not yet executed'', ``in progress'') will be blocked by the Phase-20V validator.

\section{Bibliography and Integration}\label{sec-methods-bibliography-integration}

Phase 9 (DATAML actor) produced \texttt{references.bib} (vid \texttt{content/references.bib}) by re-resolving all 1,037 cite-keys against CrossRef, then rendering BibTeX entries with ASCII-folded author and journal fields and the canonical project key set. The bibliography contains 1,037 entries (one per Phase-3 cite-key); 413 of those entries are actually used in the body text after Phase 7-10 drafting and integration, with the remaining 624 retained as the corpus-level references manifest so that Phase 18 / 19 fix passes can re-cite without re-resolving DOIs.

Phase 9V (DATAML validator) ran twelve checks against \texttt{references.bib} and gate vid \texttt{gates/gate_bibliography.json} records the outcomes: balanced braces, required fields per entry type, cite-key uniqueness, exact match to the Phase-3 keymap, unescaped DOI match, 1,037 / 1,037 CrossRef HTTP-200 with concurrency 3 and observed throughput 6.7 req/s (no rate-limit hits), first-author surname match against CrossRef, title Jaccard \textgreater = 0.85 against CrossRef, year match within 1, ASCII-clean author/title/journal, used-keys-present-in-bib for all 11 body sections, and no placeholder / template contamination. Two minor character-substitution encoding losses were flagged and patched (Argunsah2024 expecting ``Argun(s with cedilla under)ah'', Calin2023 expecting ``C(a with breve)lin'') -- both are cosmetic (DOI / year / title still match CrossRef) and were applied as ASCII-fold patches before the gate closed.

Phase 10 (LITREVIEW integration) ran six cross-section integration passes over the 11 body sections (gate vid \texttt{gates/gate_integration.json}):

\bigskip\noindent
\begin{tabular}{p{\dimexpr 0.500\linewidth-2\tabcolsep}p{\dimexpr 0.500\linewidth-2\tabcolsep}}
\toprule
Pass & Purpose \\
\hline
6a transition & Append explicit forward-link transition sentences at section boundaries (no new citations introduced). \\
6b crossref & Resolve all \texttt{\{ref\}}/\texttt{\{numref\}} cross-section pointers; ensure every label is reachable. \\
6c terminology & Normalize key term usage across sections (e.g. consistent treatment of ``disinhibitory motif'' vs ``disinhibition'' and of marker compound names). \\
6d continuity audit & Flag and fix per-section claims that contradicted other sections without explicit reconciliation. \\
6e figure inventory & Reconcile section figure inventories against the Phase-6 final panel set (22 figures inventoried; 0 mismatches). \\
6f hygiene & ASCII-clean headings, anchor consistency, no double labels. \\
\bottomrule
\end{tabular}

\bigskip

Outcomes: 10 of 11 sections were modified across the six passes, with 54 total integration edits; the cite-key set used in body text (413 unique keys) was preserved exactly through integration (no key added or removed); 0 broken cross-references after the 6b pass. Phase 11 (LITREVIEW bookend) then drafted the introduction (1,018 words, 76 cite-keys) and conclusion (1,079 words, 71 cite-keys), drawing strictly from the 413-key body cite-set with 0 new keys introduced; gate vid \texttt{gates/gate_intro_conclusion.json}. Phase 12 (LITREVIEW bookend critic) closed in 2 iterations -- iter1 raised 1 MUST\_FIX (forbidden lexicon ``canonical'' appearing in the introduction thesis sentence) and 5 SHOULD\_CAVEAT items; iter2 returned 0 MUST\_FIX after a single-word patch produced introduction v2 (vid \texttt{content/01_introduction.md}). Gate vid: \texttt{gates/gate_bookend_critic.json}.

Final post-Phase-10 section markdown version IDs (frozen for Phase 14 document assembly):

\bigskip\noindent
\begin{tabular}{p{\dimexpr 0.333\linewidth-2\tabcolsep}p{\dimexpr 0.333\linewidth-2\tabcolsep}p{\dimexpr 0.333\linewidth-2\tabcolsep}}
\toprule
Section & Title & Latest VID \\
\hline
02 & Molecular Identity and Transcriptomic Taxonomy & \texttt{content/02_classification.md} \\
03 & Developmental Origins and Postnatal Maturation & \texttt{content/03_development.md} \\
04 & Morphological Diversity & \texttt{content/04_morphology.md} \\
05 & Intrinsic Electrophysiology & \texttt{content/05_electrophysiology.md} \\
06 & Synaptic Properties and Connectivity & \texttt{content/06_synaptic_properties.md} \\
07 & Local Circuit Motifs and the Disinhibition Framework & \texttt{content/07_disinhibition_framework.md} \\
08 & In Vivo Function During Behavior & \texttt{content/08_in_vivo_behavior.md} \\
09 & VIP Interneurons Across Brain Regions & \texttt{content/09_cross_areas.md} \\
10 & Oscillatory Dynamics and Temporal Coordination & \texttt{content/10_oscillatory_dynamics.md} \\
11 & Species Differences, Human Relevance, and Disease & \texttt{content/11_disease_translational.md} \\
12 & Computational Models of VIP Circuit Function & \texttt{content/12_computational_models.md} \\
01 & Introduction (Phase 12 v2) & \texttt{content/01_introduction.md} \\
13 & Synthesis and Conclusion & \texttt{content/13_conclusion.md} \\
\bottomrule
\end{tabular}

\bigskip

\section{Limitations, Reproducibility, and Data Availability}\label{sec-methods-limitations-reproducibility}

\textbf{Limitations of the present pipeline run.}

\begin{enumerate}
\item The Phase-2 saturation criterion was reached at 1,037 unique DOIs, below the scope-document target of 2,000. The coordinator accepted this outcome (the saturation rule, not the database-sizing pre-estimate, governs corpus closure under coordinator v27), but a future run with broader query expansion or additional database coverage (e.g. \href{http://dimensions.ai}{dimensions.ai}, \href{http://lens.org}{lens.org}, J-STAGE, ChinaXiv) might recover additional papers in the long tail of the VIP literature.
\item Sections 4 (Morphology) and 9 (Brain Regions) were below the originally targeted per-section paper tier floor due to upstream Phase-2 saturation; this is recorded as documented exit \texttt{EXIT\_1} (PAPER\_TIER\_FLOOR) in the Phase-5V gate.
\item Two cite-keys carry consortium-author entries that have no CrossRef ``family'' surname (\texttt{Thepetillainterneuronnomenclaturegroupping2008} and \texttt{Braininitiativecellcensusnetworkbiccn2021}); both were resolved by name-match against CrossRef but do not satisfy the standard surname-based author-match check.
\item The Phase-6 figure audit closed with 0 REDESIGN and 0 FABRICATION outcomes but 26 of 34 final panels carry CAVEAT outcomes; CAVEAT findings are textual qualifications applied within figure captions rather than redesigns.
\item Citation triples have not yet been verified against full text. Phase 16 (citation verification, run after Phase 13) will compute VERIFIED / MINOR / CHIMERIC / HALLUCINATED / MISATTRIBUTED counts; numbers from that phase will land in the Phase-20a refresh of M.6.
\end{enumerate}

\textbf{Reproducibility statement.} Every artifact in this pipeline is content-addressable by version ID in the project artifact store (project ID \texttt{proj\_51704bdb3ee3}). Each gate JSON encodes the actor frame ID, validator frame ID, input artifact VIDs, and the mechanical checks that closed the gate. The pipeline skill files (\texttt{comprev-coordinator-v27.md}, \texttt{comprev-reviewer-agent.md}, \texttt{comprev-evidence-gathering.md}, \texttt{comprev-evidence-validator.md}, \texttt{comprev-curation-validator.md}, \texttt{comprev-figure-construction.md}, \texttt{comprev-figure-audit.md}, \texttt{comprev-section-writing.md}, \texttt{comprev-critic.md}, \texttt{comprev-integration.md}, \texttt{comprev-fix-execution.md}, \texttt{comprev-verification.md}, \texttt{comprev-dataml-phases.md}, \texttt{comprev-scoping.md}, \texttt{comprev-scoping-validator.md}) are themselves stored as artifacts in the project and pinned to the version IDs that the coordinator delegated; replaying any phase with the same skill VID, same input artifact VIDs, and the same agent role recovers the same outputs modulo LLM nondeterminism (which is bounded for DATAML mechanical phases by deterministic check schemas and for LITREVIEW phases by the gate validator).

\textbf{Data availability.}

\begin{itemize}
\item Verbatim review request: \texttt{provenance/review\_request.md} (vid \texttt{provenance/review_request.md}) and \texttt{provenance/review\_request.txt} (vid \texttt{provenance/review_request.txt}).
\item Scope: \texttt{scope.json} vid \texttt{provenance/scope.json}; \texttt{gate\_scope.json} vid \texttt{gates/gate_scope.json}.
\item Citation key map: vid \texttt{provenance/citation_key_map.json} (1,037 keys); author name table vid \texttt{provenance/author_name_table.json}.
\item Bibliography: \texttt{references.bib} vid \texttt{content/references.bib} (1,037 entries; 413 used in body text).
\item Gate JSONs for Phases 1-12 are listed in the Pipeline Ledger above.
\item The 13 final body / bookend section markdown files are listed in M.7 with their version IDs.
\item The final repository (post-Phase-20 push) will be \texttt{https://github.com/AllenNeuralDynamics/ComputationalReviewVIP}.
\end{itemize}

\textbf{Pipeline skills used.} The pipeline executed under coordinator skill \texttt{comprev-coordinator-v27.md}. Phase-by-phase skill assignments are: Phase 1 -- \texttt{comprev-scoping.md} (LITREVIEW) + \texttt{comprev-scoping-validator.md} (DATAML); Phase 2 -- \texttt{comprev-evidence-gathering.md} (LITREVIEW x14) + \texttt{comprev-evidence-validator.md} (DATAML); Phases 3, 5, 9, 13 -- \texttt{comprev-dataml-phases.md} (DATAML); Phase 4 -- framework subroutine of \texttt{comprev-reviewer-agent.md} (LITREVIEW); Phase 5V -- \texttt{comprev-curation-validator.md} (DATAML); Phase 6 -- \texttt{comprev-figure-construction.md} + \texttt{comprev-figure-audit.md} (LITREVIEW); Phase 7 -- \texttt{comprev-section-writing.md} (LITREVIEW x11); Phase 8 -- \texttt{comprev-critic.md} (LITREVIEW x11); Phase 10 -- \texttt{comprev-integration.md} (LITREVIEW); Phases 11-12 -- \texttt{comprev-reviewer-agent.md} (bookend mode) and \texttt{comprev-critic.md} (bookend mode). Phases 14-19 will load \texttt{comprev-fix-execution.md} and \texttt{comprev-verification.md}. Skill VIDs are pinned by the coordinator at delegation time.



\bibliography{references}
\end{document}
